% Paper V-A v3.218 / Framework v3.058: native rooted two-contact positive return and restricted Hopf no-go; actual Cert_hier OPEN.
% Paper V-A v3.216: exact C3/channel-only source exclusions; Cert_hier and pole obligations OPEN.
% Framework v3.029: finite Clifford/color/family operator envelope, guarded GP3 copy branch, relative/absolute trace and process descent.
% Framework source integration v3.024 (audit v1.7--1.8): primitive-defect and fixed-chart separation; physical ownership unchanged.
% Paper-series v3.008: source-matched finite Hurwitz charged hierarchy jet tests (no physical origin promotion).
% Paper-series v2.382: audit-identified standalone repair integration; mathematical scopes and open physical owners preserved.
% Paper-series v2.381: complete algebraic exceptional torsion atlas and all-cocycle twist-stability; original claims retained.
% Paper-series v2.380: unified bibliography, exact exceptional double/length/gauge mathematics in I-B; previous theorem labels retained.
% Paper-series v2.379: full equivariant Dirac/Hodge/eta/torsion atlas in Paper I-B; paper-root references synchronized.
% Paper-series v2.378: 2026 exceptional literature consolidation and reference synchronization. The mathematical additions are located in Paper I-B.
% Paper-series v2.377: reference-only rebind to current Paper I-B v1.1; original mathematics retained except new Paper I-B comparison chapter.
% Framework v2.370: derived exceptional-root/Hurwitz comparison in I-A/I-C; other papers reference-only synchronized.
% Framework v2.368: preserve prior II-B convention/unital fixes and verified II-D corrections; current references synchronized.
% Framework v2.358: current 30-root source reference synchronization; selectively updated III-E, III-F, VI-B, VI-C.
% Framework v2.356: bibliography/source-registry synchronization; mathematical body unchanged.
% Framework v2.352: integrated mathematical/operational closure; physical occurrence boundaries preserved.
% Framework v2.348: Paper V audit, terminal observable scope, canonical Gram eigenvalues and exact weak-memory variance.
% Framework v2.345: F40 native three-readout and source-visibility integration.
% Framework v2.338: exact Petersen 120-carrier, oriented two-source and source-ownership results. No hadron mass prediction.
% Framework v2.321: Paper V Petersen rooted 3+12 induction, exact sector compression, conjugation orbit refinement and source-incidence firewall. Original source theorem bodies preserved.
% Framework v2.275: exact positive dense-twin word cubature, finite repeatability and eighth-moment classification, nested sampling frames, joint action and POVM-jet exports; preserves v2.274.
% Framework v2.274: two-scalar positive convolution certificates, first-leaked exceptional coset detection, branch-symmetric twin mixing, and finite endpoint moments; all v2.273 branches preserved.
% Framework v2.273: restored missing Problem 3 theorems from v2.268; exact exceptional frames, twin geometry, convolution repeatability and projector/occurrence firewalls; merged on v2.272 without overwriting independent branches.
% Framework v2.272: merged v2.270 covariant native-readout mathematics with v2.271 Problem 4 finite-action Hopf hypercharge/holonomy corrections; no universal bulk-gauge uniqueness claim.
% Independent branches: v2.270 readout/guard exports and v2.270 K16 action-consumed charge closure; all prior scoped assumptions retained.
% Framework v2.267: exact Paper II-A converse/observable certificates, III-F action tests, V-B right orbit and V-C native gauge scope; prior v2.266 preserved.
% Framework v2.266: exact Paper II-C outcome-count/readout converse results and scoped downstream imports; 30-root references synchronized.
% Framework v2.265: II-C coordinate-stage constitution and exact first-symbol host inheritance; prior mathematical domains retained.
% Framework v2.264: II-B reverse and action-bridge integration; preserves prior mathematical domains; active 30-root references synchronized.
% Framework v2.263: Paper I-D adds bridge-controlled root-carry parity and source-relative old-cap cancellation; other root mathematical bodies retained; active 30-root references synchronized.
% Framework v2.262: merges Paper II-F three-gap closure from v2.260 into v2.261, preserving newer II-B/II-D repairs; 30-root references synchronized.
% Framework v2.261: merged Paper II-B GP2-level/geometric repairs; preserves v2.260 II-D repairs; active references synchronized.
% Framework v2.260: Paper II-D signed RS, finite-symbol typing, and conditional curvature/Jordan repairs; references synchronized.
% Framework v2.259: conditional native SU(2) atlas and explicit later discharge; v2.258 mathematical bodies retained elsewhere.
% Framework v2.258: preserves parallel II-A/dynamics update and adds complete readout/frequency closures.
% Framework v2.257: exact minimal complete position readout and arbitrary-preparation frequency law.
% Framework v2.256: exact native inertia and sharp bounded ordered-source closure; preserves concurrent v2.255 work.
% Framework v2.255: preserves parallel detector integration and adds exact finite readout closures.
% Framework v2.254: shared bibliography/reference sync; substantive Pasterski detector-Ward additions only IV-B/IV-C/VII-C.
% Framework v2.253: typed raw coordinate defect in I-A and synchronized VI-A import; active references synchronized.
% Framework v2.252: optional shell-frame source-return calculation in III-G; active references synchronized.
% Framework v2.251: forward common-action execution, joint source jets and scoped coupled control.
% Framework v2.251: exact full-heat variations, fermion source jets and explicit physical completion frontier.
% Framework v2.251: certified interacting source jets, pole stability and scoped completion proofs.
% Framework v2.251: interacting finite-source evaluation and coupled local vacuum classification.
% Framework v2.251: exact finite source execution and native causal common-parent bridge.
% Framework v2.251: constructive parent and scoped source-incidence audit; see release report.
% Framework v2.251: common-action two-jet, corrected source/response closure and synchronized references.
% Framework v2.244: scoped rooted source-order Schur diagnostics and synchronized paper references.
% Framework v2.242: signed rooted phase lift, complete transverse return and exact family cubic.
% Framework v2.241: proved rooted conormal geometry, conditional cap realization and scoped action tests.
% Framework v2.240: observable-minimal action and response certificates integrated; inherited v2.239 geometry retained.
% Framework v2.234: restructured main physics owner, technical proofs in V-D.
\documentclass[12pt,a4paper]{article}

\usepackage{mathtools}
\usepackage{float}
\usepackage{caption}
\usepackage{microtype}

% Shared series preamble: packages, bibliography, colors, theorem/proof setup, and notation.
%% Macros-v2_201: downstream full-carrier geometry update, five current root bindings, no mathematical macro change.
%% Macros-v2_200: Paper III-F v1.280 current-source binding for framework v3.074; mathematical macro definitions unchanged.
%% Macros-v2_198: current project v3.072 III-D and bibliography binding; mathematical macros unchanged.
%% Current project v3.070 source-reference bindings; previous macro bodies unchanged.
% Current II-D reference root for PapersMinimal v3.063; mathematical macros otherwise unchanged.
% Macros-v2_188: re-based current paper identities and A--G family; no change to geometry.
% Macros-v2_184: active Paper III-E v1.200/III-F v1.267 binding for framework v3.031; mathematical macros unchanged.
% Macros-v2_183: updated shared cross-reference bindings for Paper1B/III-E/III-F v3.030; mathematical macros unchanged.
% Macros-v2_182: active Paper II-F/V-A cross-references for framework v3.029; definitions unchanged.
% Paper-series v2.381: new I-B torsion/modular results, bibliography and current source roots.
% Paper-series v2.380: new exceptional bibliography and synchronized paper references.
% Paper-series v2.379 current 31-paper macro/aux registry.
% Paper-series v2.378: current 31-paper AUX/bibliography registry updated for Paper I-B v1.2 literature additions.
% Framework v2.371: current-source reference registry refreshed; old theorem labels retained.
% Framework v2.370: current 30-paper source/AUX root rebinding; native mathematical macros unchanged.
% Framework v2.368: synchronized merged 30-root registry; mathematical notation unchanged.
% Framework v2.358: active 30-paper source registry synchronized; existing mathematical macros unchanged.
% Framework v2.357: selective new source-owned certificates; active reference roots updated for five revised papers.
% Framework v2.356: active 30-root registry synchronized to the 30 papers in this release.
% Framework v2.275: exact positive dense-twin word cubature, finite repeatability and eighth-moment classification, nested sampling frames, joint action and POVM-jet exports; preserves v2.274.
% Framework v2.274: filename and root-source synchronization.
% Framework v2.273: restored missing Problem 3 theorems from v2.268; exact exceptional frames, twin geometry, convolution repeatability and projector/occurrence firewalls; merged on v2.272 without overwriting independent branches.
% Framework v2.272: merged v2.270 covariant native-readout mathematics with v2.271 Problem 4 finite-action Hopf hypercharge/holonomy corrections; no universal bulk-gauge uniqueness claim.
% Independent branches: v2.270 readout/guard exports and v2.270 K16 action-consumed charge closure; all prior scoped assumptions retained.
% Framework v2.269: synchronized 30 active AUX roots; mathematics and bibliography preserved from Macros-v2_161.
% Framework v2.266: Paper II-C converse-derived outcome resources; II-E, III-C, IV-A/B, V-A and VI-B downstream mathematical closures; 30-root registry synchronized.
% Framework v2.265: II-C coordinate-stage constitution and exact first-symbol host inheritance; prior mathematical domains retained.
% Framework v2.264: paper II-B reverse/audit and IV-A observable-minimal exchange; synchronized external references to 30 active roots.
% Framework v2.263: Paper I-C adds bridge-controlled root-carry parity and source-relative old-cap cancellation; other root mathematical bodies retained; active 30-root references synchronized.
% Framework v2.262: integrates Paper II-F radius, Casimir and readout results; active 30-root references synchronized.
% Framework v2.261: merged Paper II-B GP2-level/geometric repairs; preserves v2.260 II-D repairs; active references synchronized.
% Framework v2.260: Paper II-D signed RS, finite-symbol typing, and conditional curvature/Jordan repairs; references synchronized.
% Framework v2.259: conditional native SU(2) atlas and explicit later discharge; v2.258 mathematical bodies retained elsewhere.
% Framework v2.258: preserves parallel II-A/dynamics update and adds complete readout/frequency closures.
% Framework v2.257: exact minimal complete position readout and arbitrary-preparation frequency law.
% Framework v2.256: exact native inertia and sharp bounded ordered-source closure; preserves concurrent v2.255 work.
% Framework v2.255: preserves parallel detector integration and adds exact finite readout closures.
% Framework v2.254: Pasterski bibliography and synchronized 30-paper external-reference registry.
% Framework v2.253: typed raw coordinate defect in I-A and synchronized VI-A import; active references synchronized.
% Framework v2.252: synchronized source registry; optional III-G shell-frame calculation.
% Framework v2.251: forward common-action execution, joint source jets and scoped coupled control.
% Framework v2.251: exact full-heat variations, fermion source jets and explicit physical completion frontier.
% Framework v2.251: certified interacting source jets, pole stability and scoped completion proofs.
% Framework v2.251: synchronized active source registry; mathematical notation unchanged.
% Framework v2.251: synchronized active source registry; mathematical notation unchanged.
% Framework v2.251: synchronized active source registry; mathematical notation unchanged.
% Framework v2.251: synchronized source registry; no mathematical notation changes.
% Framework v2.244: synchronized 30-paper AUX registry and rooted phase-source action-order tests.
% Framework v2.242: signed rooted phase lift, complete transverse return and exact family cubic.
% Framework v2.241: proved rooted conormal geometry, conditional cap realization and scoped action tests.
% Framework v2.239: all 30 active source roots synchronized; no mathematical macro changes.
% Framework v2.238: Paper IV A/B/C split and 30-root reference registry.
% Framework v2.237: Paper VII-D technical split and 28-root reference registry.
% Framework v2.236: Paper VI-C technical-volume reference added; all current root bindings advanced.
% Framework v2.235: 26-root AUX and bibliography bindings synchronized; mathematical macros unchanged.
% BEGIN INLINE SOURCE: Paper00-TeXBuildCompatibility-v1_0.tex
% Rendering-only repairs for inherited long appendices and bitmap text fonts.
% Integer counters through 26 retain the original labels; later values use AA,
% AB, ... . This file changes no action coefficient, state or theorem premise.
\RequirePackage{alphalph}
\makeatletter
\let\BHPrimitiveAlph\@Alph
\newalphalph{\BHExtendedAlph}{\BHPrimitiveAlph}{26}
\let\@Alph\BHExtendedAlph
\AtBeginDocument{%
  \@ifpackageloaded{microtype}{\microtypesetup{expansion=false}}{}%
}
\makeatother

% END INLINE SOURCE: Paper00-TeXBuildCompatibility-v1_0.tex

% Framework v2.164: current reference/bibliography profile; mathematical definitions retained.
% Framework v2.162: finite metric, functional Hilbert and scoped Ward integration.
% Framework v2.161: inherits Macros-v2_57.tex exactly except the current external-reference binding.
% Macros v2.57 / framework v2.156: current reference registry only; mathematical definitions unchanged.
% Framework v2.143 / Macros v2.55: current external-reference bindings; mathematical notation unchanged.
% Framework v2.110 / Macros v2.54: current external-reference registry and bibliography; mathematical notation inherited.
% Framework v2.67 / Macros-v2_53: adds Born-action-cell and joint-EH readout symbols; absolute unit realization is metrological; MATCH_{J/scal}^{EH} remains the sole G24 normalization bridge.
% Framework v2.35 / Macros-v2_52: grouped paper numbering, active AUX registry, and Born-v6_20 bibliography binding synchronized; notation unchanged.
% Framework v2.26: active-source AUX and bibliography bindings synchronized; notation unchanged.
% Framework v2.09: reviewed early-reciprocity integration; scoped physical imports.
% Framework v2.08: Born--Euler mathematics and conditional complete-action potential criterion.
% Framework v2.07: active AUX registry includes Paper I-B v1.2; shared bibliography Born-v6_18.
% Macros-v2_46 / framework v2.04: mathematical atlas and downstream realization correction.
% Macros-v2_46 / framework v2.03: scoped structure integration and current dependency bindings.
% Macros-v2_46 / framework v2.01: shared source bindings; prior revision comments are historical.
% v2.42 / framework v1.99: Synchronized RG/AS owner and bibliography bindings.
% v2.41 / framework v1.99: active manuscript AUX bindings and geometric-closure bibliography synchronized.
% Release-v1.80 centralized source synchronization for scalar-circle/comparison-band integration.
% Release v1.79: Hurwitz hierarchy and kinetic-normalized action response integration; physical mass matching remains conditional.
% Release v1.78: approved M8 closure-batch integration; physical and analytic domains retained.
% Current v1.78 registry: all active TeX roots are synchronized below.
% Earlier release comments are historical scope records, not current status.
% Release-v1.77 synchronization: native same-owner Noether closure, generator--flow/history typing, analytic-certification firewalls, scalaron/K16, LF/QCD, quantum-Ward, and RAR boundary-owner transport are registered; shared mathematical notation is unchanged.
% Release-v1.76 synchronization: registry points to the v1.76 active paper roots.
% Release-v1.76 synchronization: registers the local A+Phi+g Noether/constraint closure and dynamics-frontier propagation; shared mathematical notation is unchanged.
% Release-v1.75 synchronization: registry points to the v1.75 active paper roots.
% Release-v1.75 synchronization: integrates the joint-Dirac Einstein bulk real form and the local odd-scalar heat-coefficient action on explicit analytic first-jet/interior domains; v1.74 Yang--Mills closure is retained. Fermionic reality, pure-metric curvature-squared/exact-nonlocal terms, noncanonical boundaries, physical vacuum selection, and actual solution-history tokens remain open.
% Release-v1.74 synchronization: registry points to the v1.74 active paper roots and retains the v1.73 theorem labels.
% Release-v1.74 synchronization: canonical Macros-v2_34.tex; generated local nonlinear Yang--Mills action, constrained connection-field phase space, source-free Gauss momentum map, and Legendre closure are integrated on their proved domains. The remaining non-gauge ACT_cl sectors and actual solution-history selection remain open.
% Release-v1.71 synchronization: closes the former JOINT_x/JOINT_5
% connection-occurrence gap by canonical structure-group extension of the
% nonflat Born connection, faithful associated representation, and the local
% cotangent--Clifford operator constructor.  Actual fluctuation members,
% actions, vacua, solutions, calibrations, and analytic QFT domains remain
% separately typed.  The active-AUX registry is advanced once.
% Release-v1.71 canonical macro and active-AUX registry.
% ======================================================================
% Macros-v2_39.tex -- canonical publication macro file
% Release v1.71 retains the v1.70 null-ray/QCD boundaries and additionally
% removes the duplicate structural connection actuality premise.  It does not
% claim global bundle topology, select a gauge/Higgs member, or promote the
% finite construction to an interacting continuum theory.
% Release v1.68 constructs the finite light-front total-fraction operators and
% parton number/momentum operator measures, proves their simplex and Hausdorff-limit
% theorems, and separates positive cut measures from scheme-renormalized
% light-ray distributions.  It also separates finite real Schur poles, stable
% norm-resolvent/Riesz pole-residue limits, and continued resonance
% pole/residue limits.  The Paper-VII-C scattering package is claim-indexed:
% local QFT and spectral claims import no S-matrix or infrared premise, while
% asymptotic amplitudes, LSZ, S-matrix poles, completeness, and cross sections
% import only their named massive or massless/long-range clauses.
% Release v1.67 factorizes BND48_H into a same-action Hessian/Schur
% descendant and still-open typed boundary-trace, critical-lift, and
% refinement-limit conditions.  It proves the finite-collar Dirichlet-to-
% Neumann law sqrt(y)tanh(L sqrt(y)), its nonuniform critical limit, and the
% finite-return ultra-infrared turnover.  The Paper-I-A Hardy constant mode is
% an exact geometric near-hit, not yet the physical metric boundary mode.
% The current local f(R) Hessian remains massive and therefore still refutes
% a RAR prediction on the minimal action branch.
% Release v1.66 derives the transverse depth-six normalizer bridge from the
% already generated C6 Reynolds image and Paper-I-A binary-dihedral normalizer.
% The bridge is the unique least 2D3-stable scalar-containing realization of
% the faithful quotient 2D3/C6=C2, has trace defect six and hence K=48, but is
% transverse to the nested active carrier.  Exact K=44 and K=52 contraction,
% semigroup, and energy-midpoint identities are registered as structural
% readouts rather than projector events.  The active E48 carrier gap remains
% unchanged.  A same-action boundary-symbol/Calderon intertwiner is isolated
% as the necessary and still unproved owner bridge to the static DtN sort;
% representation equivalence alone neither derives RAR nor creates a K=48
% action junction, pole, occupation, or cosmological transition.
% Release v1.65 derives the native E_48-centered carrier mass gap from the
% GP1--GP3-generated active spectrum, distinguishes its carrier-free interval
% from a particle located at E_48, and withdraws the unsupported promotion to
% a positive continuum Yang--Mills mass gap.  The latter remains a valid
% conditional transfer theorem on an enlarged continuum branch, but it is not
% an active framework prediction or an unfinished native-gap proof.
% Release v1.64 synchronizes Derivation Batches 14--22.  It registers the
% GP1--GP3 heat-profile orbit, same-Hessian electroweak response and pole
% ownership, K16/K48 action-response corrections, perturbation state and
% non-Gaussianity typing, finite-vacuum/gap/QCD limit repairs, factorized
% position-area ownership, the actual static metric response, and the ordered
% Wilson/string-tension limits.  Strict continuum convergence, positive
% physical gap, surviving higher cumulants, and positive tension remain open.
% Release v1.61 closes TIMEWIDTH16 as an independent direct-production datum.
% JUNC_{2|7} continues to own the finite intrinsic microscopic profile/width,
% while the actual same-action background/production history determines its
% physical-time pullback, support, and duration.  The current Paper-I-A--VII-C AUX
% registry is advanced once; shared mathematical notation is unchanged.
% Release v1.60 derives the final-octet primitive-remnant depletion, canonical
% Hopf-current representative, exact Hopf descent to dOmega/(4pi), and
% closed-FLRW a^{-3} neutral-current shape.  It retypes K=48 as the onset
% of the final one-unit geometric transfer and records cDM/RAR as distinct
% physical projections of that common neutral parent while preserving
% occurrence, Hilbert-owner, stress-matching, and static-kernel premises.
% Release v1.59 transports the v1.58 exceptional-shell/K16 closure through
% Papers III-F and VI-B--VII-C.  It replaces Paper VI-B's obsolete mixed-Hessian
% output typing by the third-variation pair vertex and Bogoliubov occupation
% support, derives the post-seed Higgs-pair feedback vertex and its exact
% alpha--delta kernel factors, propagates the unique K16/K48 action-junction
% split, identifies the QCD delta count with the upstream 2O shell atom only
% as a common-scalar numerical compatibility, and narrows THERM16 to dynamical
% approach/beta selection because the finite Gibbs/KMS target is already
% unique on a supplied simple observable algebra once its Hamiltonian and
% inverse temperature are supplied; direct-sum central weights remain typed.
% Release v1.58 registers the exceptional nearest-shell closure, the derived
% transverse K=16 action event, the unique interior action junction, the
% type-correct ordered pair vertex, the thin-layer squeeze law, and the
% radiation-like/matter-like kinematic split.  REG_0, transfer magnitude,
% physical layer width, and thermalization remain explicitly conditional.
% Release v1.57 registers the active-depth-six obstruction, the native
% K=48 carrier desert, the typed 7|9 owner reduction, and the dynamical
% reheating handoff across Papers III-B--VI-B.
% Release v1.56 registers the least-support n_2 inflation dressing of the
% persistent n_7 Laurent response, the analytic global-convexity and positive
% endpoint-measure proofs, the derived inflation core and K=16 coefficient
% freeze, and the no-particle dark-sector split into static RAR plus the
% conditional homogeneous effective-curvature response at reflected K=48.
% There is no n_6 action, scalaron-DM mass, or second Einstein--Hilbert term.
% Release v1.55 registers the protected-octet resolved/residual reflection
% q_sp+Z_W=8, its j<->8-j and K<->64-K scale covariance, and the derived
% rank-one-exit quadratic grade C_1^2/U=C_2 in Papers III-C, III-E, and VIII.
% The reflection is a typed trace/scale theorem, not a time reversal or a
% physical event-occurrence theorem; reheating transfer remains conditional.
% Release v1.54 registers the exact positive-part Laurent indices, the
% active-domain/terminal-endpoint action split, the no-second-EH firewall, and
% the q_DM=6/K_DM=48 scalaron readout typing across Papers III-D/E/F, VI,
% VIII, and IX.
% Release v1.53 registers the common nonnegative Laurent obstacle projection
% for n_1^+ and n_7^+, discharges the three former exit/freeze assumptions,
% and advances the Paper-III-D/E/F, VIII, and IX consumers.  Shared notation
% and bibliography ownership are unchanged.
% Release v1.51 registers the PQTI derivation batch and its updated
% Paper-III-A/C/D/E/F, VIII, and IX consumers.  Shared notation and
% bibliography ownership are unchanged.
% Release v1.50 advances the Paper-III-B AUX target to v1.40 after the
% rigorous audit repair; shared notation and bibliography ownership are unchanged.
% Release v1.49 advances the Paper-II-E/F AUX target to v3.14 after the
% minimal audit repair; shared notation and bibliography ownership are unchanged.
% Release v1.46 synchronizes Paper II-A's minimal frozen constructor-origin
% theorem and its five dependent AUX targets.  Shared notation is unchanged.
% Release v1.45 updates the Paper-II-B AUX target to v2.23 after the
% GP3-rigidity and appendix-placement repair; shared notation is unchanged.
% Release v1.44 advances the centralized Paper-I-A--VII-C AUX registry after the
% Paper-I-A v5.16 protected-completion scope repair.  Every current
% paper imports this one file; no paper-level opt-in or \externaldocument list
% is permitted.  The bibliography remains centralized here through Born-v6_13.
% Shared bibliography, theorem infrastructure, colors, notation alignment,
% and current Paper-I-A--VII-C external-reference registration.
% This file centralizes common imports, bibliography sorting, colors, theorem
% infrastructure, and the notation shared by the complete series.
% v2.9 (release v1.43) made the Paper-I-A--VII-C AUX registry automatic and
% centralized the bibliography and all external-reference ownership.
% v2.8 release-registry refresh (v1.33) registers the rooted generated-parent
% cascade and the POL/REG/free-law/11+4/Born-group closure revisions.  No
% shared mathematical notation or Paper-II-A fundamental principle is changed.
% v2.8 release-registry refresh (v1.32) registers the universal generated-
% continuation revisions in Papers II-A, III-A, III-C, IV, and VIII.  No
% shared mathematical notation or Paper-II-A fundamental principle is changed.
% v2.8 release-registry refresh (v1.30) registers the CRP4/projective-state
% source revisions and their direct registry-only consumers.  No shared
% mathematical notation or Paper-II-A axiom is changed.
% v2.7 release-registry refresh (v1.27) updates only the current external
% Paper-I-A--VII-C basenames after the placement/assumption synchronization; no
% shared mathematical notation or Paper-II-A axiom is changed.
% v2.7 synchronizes the geometric defect-repair revision: Paper-II-A v3.42,
% Paper-III-C v1.26, Paper-III-D v1.27, Paper-III-E v1.32, and Paper-III-F
% v1.34.  It registers the common trace-defect/spectral-filtration theorem,
% its owner-typed response pullback, and the reduced GRED and TERM interfaces.
% No shared mathematical notation is changed.
% v2.6 synchronizes Paper-III-E v1.31's cone-type repair, minimal-response
% classification, and explicit classical action-role boundary, together with
% Paper-III-F v1.32's downstream response-atlas synchronization.  No shared
% mathematical notation is changed.
% v2.5 synchronizes the Paper-III-D v1.26 faithful-trace-duality proof,
% exact ordered-owner joint locus, and typed alternative-presentation
% bookkeeping repair, together with its Paper-III-E and closure-frontier
% references.  No shared mathematical notation is changed.
% v2.4 synchronizes the Paper-III-C v1.25 orbit/process/owner repair, its
% downstream Papers III-B--F, V, VI, VIII, and IX revisions, and Born-v6_13.
% v2.3 synchronizes the Paper-II-A v3.41 witness/support/principal inheritance
% layer, the Paper-III-A v1.24 joint-realizability/readout firewall, all dependent
% host-rate revisions, and the current Paper-I-A--VII-B source registry.
% v2.2 synchronizes the generated active-realization, determinant-depth,
% dimension-weighted hierarchy, and typed cosmological-event interfaces.
% v2.1 synchronizes the exceptional-support interface revisions: Paper I-A
% v5.9, Paper II-A v3.30, Paper II-E/F v3.7, and Paper IV v4.16.
% v2.0 introduces the stage-safe G0 occurrence notation, synchronizes the
% active Paper-I-A--IV source registry (current rigorous Paper-III closure:
% III-A v1.17, III-B v1.22, III-C--E v1.17, III-F v1.14, IV v4.14), and makes
% external-reference self-skip depend on logical paper role rather than a
% versioned job name.
% ======================================================================

% ======================================================================
% macros.tex — Shared macros for the Born Geometry series
% ======================================================================

% ======================================================================
% 1. PACKAGES, BIBLIOGRAPHY, COLORS, AND UNIT DEFINITIONS
% ======================================================================
% Core packages shared across the complete Paper-I-A--VII-C series.  Loading them here makes
% the macro file a genuine series preamble while remaining compatible with existing
% paper preambles that load the same packages before \input.
\usepackage{amsmath}
\usepackage{amssymb}
\usepackage{amsthm}
\usepackage{mathrsfs}
\usepackage{array}
\usepackage{booktabs}
\usepackage{xcolor}
\makeatletter
\@ifpackageloaded{geometry}{}{\usepackage{geometry}\geometry{margin=1in}}
\makeatother
\usepackage{mdframed}

% Shared support packages required by the macros below and by most papers.
\usepackage{ragged2e}
\usepackage{longtable}
\usepackage{tabularx}
\usepackage{enumitem}
\usepackage[most]{tcolorbox}

% Series-wide numeric bibliography convention using the portable BibTeX path.
% All Papers I-A--VII-C use the ordinary \cite interface.
\usepackage[numbers,sort&compress]{natbib}
\providecommand{\printbibliography}{%
  \bibliographystyle{plainnat}%
  \bibliography{Born-v6_60}%
}

% Shared colors.  \providecolor avoids clobbering a paper-specific override.
\providecolor{postulatecolor}{RGB}{230,242,255}
\providecolor{compare}{RGB}{230,245,230}
\providecolor{definitioncolor}{RGB}{230,242,255}

\providecolor{openriskcolor}{RGB}{255,235,235}
\providecolor{bridgecolor}{RGB}{240,240,255}
\providecolor{gold}{RGB}{255,215,0}

% ======================================================================
% 2. CUSTOM ENVIRONMENTS
% ======================================================================
% Common table column types used in all papers.  Guarded so that a paper may
% define its own versions before loading this file.
\makeatletter
\@ifundefined{NC@find@L}{%
  \newcolumntype{L}[1]{>{\RaggedRight\arraybackslash}p{#1}}%
}{}
\@ifundefined{NC@find@C}{%
  \newcolumntype{C}[1]{>{\Centering\arraybackslash}p{#1}}%
}{}
\@ifundefined{NC@find@Y}{%
  \newcolumntype{Y}{>{\RaggedRight\arraybackslash}X}%
}{}
\makeatother

% ======================================================================
% 3. THEOREM STYLES AND ENVIRONMENTS
% ======================================================================
% Custom underlined theorem style (useful for remarks or notes)
\newtheoremstyle{underlined}
  {3pt}   % Space above
  {3pt}   % Space below
  {}      % Body font (empty means normal/roman text, typical for remarks)
  {}      % Indent amount
  {}      % Theorem head font (can be \itshape, \bfseries, etc.)
  {.}     % Punctuation after theorem head
  { }     % Space after theorem head
  {\underline{\thmname{#1}\thmnumber{ #2}}\thmnote{ (#3)}} % Theorem head formatting

% Standard plain theorem style environments.  Each environment is guarded so
% the shared macro file can be loaded after a paper-specific theorem setup.
\theoremstyle{plain}
\makeatletter
\@ifundefined{theorem}{\newtheorem{theorem}{Theorem}[section]}{}
\@ifundefined{lemma}{\newtheorem{lemma}{Lemma}[section]}{}
\@ifundefined{proposition}{\newtheorem{proposition}{Proposition}[section]}{}
\@ifundefined{corollary}{\newtheorem{corollary}{Corollary}[section]}{}
\makeatother

% Definition style environments (typically non-italicized body)
\theoremstyle{definition}
\makeatletter
\@ifundefined{definition}{\newtheorem{definition}{Definition}[section]}{}
\@ifundefined{postulate}{\newtheorem{postulate}{Postulate}[section]}{}
\@ifundefined{closure}{\newtheorem{closure}{Closure}[section]}{}
\makeatother

% Applying the custom underlined style
\theoremstyle{underlined}
\makeatletter
\@ifundefined{remark}{\newtheorem{remark}{Remark}[section]}{}
\@ifundefined{convention}{\newtheorem{convention}{Convention}[section]}{}
\@ifundefined{principle}{\newtheorem{principle}{Principle}[section]}{}
\makeatother

% Framed proof environment.
% Every \begin{proof} ... \end{proof} is automatically framed/shaded.
% This avoids mdframed/trivlist nesting, which can drop the top border after
% theorem-like environments.
\makeatletter
\renewenvironment{proof}[1][\proofname]{%
  \par\addvspace{6pt}%
  \noindent
  \begin{tcolorbox}[%
    enhanced,
    breakable,
    colback=black!4,
    colframe=black!35,
    boxrule=0.35pt,
    arc=1pt,
    left=6pt,
    right=6pt,
    top=4pt,
    bottom=4pt,
    boxsep=0pt,
    before skip=0pt,
    after skip=6pt,
    parbox=false
  ]%
  \pushQED{\qed}%
  \normalfont\itshape #1\@addpunct{.}\enspace
  \normalfont\ignorespaces
}{%
  \unskip\nobreak\hfill\popQED
  \end{tcolorbox}%
  \par\addvspace{2pt}%
}
\makeatother

\mdfdefinestyle{postulatebox}{
  backgroundcolor=yellow!10,
  linecolor=black,
  linewidth=0.4pt,
  skipabove=\baselineskip,
  skipbelow=\baselineskip,
  nobreak=true
}

% ======================================================================
% PROOF-READABILITY HELPERS
% ======================================================================

\newcommand{\ProofPlan}[1]{%
  \par\smallskip
  \noindent\textbf{Proof plan. }#1%
  \par\smallskip
}

\newlist{proofsteps}{enumerate}{2}
\setlist[proofsteps,1]{
  label=\textbf{(\alph*)},
  ref=\alph*,
  leftmargin=2.2em,
  itemsep=0.45\baselineskip,
  topsep=0.35\baselineskip
}
\setlist[proofsteps,2]{
  label=\textbf{(\roman*)},
  ref=\roman*,
  leftmargin=2.0em,
  itemsep=0.25\baselineskip,
  topsep=0.2\baselineskip
}


% ======================================================================
% 4. NUMBER SYSTEMS AND ALPHABETS
% ======================================================================
\providecommand{\RR}{\mathbb{R}} % Real numbers
\providecommand{\PP}{\mathbb{P}} % Projective space / Probability
\providecommand{\CC}{\mathbb{C}} % Complex numbers
\providecommand{\ZZ}{\mathbb{Z}} % Integers
\providecommand{\NN}{\mathbb{N}} % Natural numbers
\providecommand{\HH}{\mathbb{H}} % Quaternions / Hyperbolic space
\providecommand{\LFRayGeom}{\mathsf{LF\!\mbox{-}RAY\!\mbox{-}GEOM}}
\providecommand{\LFRep}{\mathsf{LF\!\mbox{-}REP}}
\providecommand{\LFState}{\mathsf{LF\!\mbox{-}STATE}}
\providecommand{\LFLim}{\mathsf{LF\!\mbox{-}LIM}}
\providecommand{\LFScat}{\mathsf{LF\!\mbox{-}SCAT}}
% Number systems / projective spaces
\providecommand{\CP}{\mathbb{CP}}
\providecommand{\RP}{\mathbb{RP}}

% ======================================================================
% 5. GROUPS AND LIE THEORY
% ======================================================================
\providecommand{\SO}{\ensuremath{\mathrm{SO}}}     % Special Orthogonal group
\providecommand{\SU}{\ensuremath{\mathrm{SU}}}     % Special Unitary group
\providecommand{\Spin}{\ensuremath{\mathrm{Spin}}} % Spin group
\providecommand{\CS}{\ensuremath{\mathrm{CS}}}     % Chern-Simons (capitalized)
\providecommand{\cs}{\ensuremath{\mathrm{cs}}}     % Chern-Simons (lowercase)

% ======================================================================
% 6. COMMON MANIFOLDS AND TOPOLOGY
% ======================================================================
\providecommand{\Sone}{S^1}                 % Standard circle

% ======================================================================
% 7. MATH OPERATORS (Declared via \DeclareMathOperator)
% ======================================================================
\DeclareMathOperator{\End}{End}   % Endomorphisms
\DeclareMathOperator{\Aut}{Aut}   % Automorphisms
\DeclareMathOperator{\spec}{spec} % Lowercase spectrum
\DeclareMathOperator{\tr}{tr}     % Lowercase trace (if you already have \Tr)
\DeclareMathOperator{\supp}{supp} % Support
\DeclareMathOperator{\TrA}{Tr_{A}} % Partial trace over subsystem A
\DeclareMathOperator{\TrB}{Tr_{B}} % Partial trace over subsystem B
% Lie/algebra operators
\DeclareMathOperator{\Ad}{Ad}
\DeclareMathOperator{\Sym}{Sym}
\DeclareMathOperator{\su}{su}
% ======================================================================
% 8. STANDARD OPERATOR-LIKE SYMBOLS (Safe definitions via \providecommand)
% ======================================================================
\providecommand{\vol}{\operatorname{vol}}   % Volume form
\providecommand{\Vol}{\operatorname{Vol}}   % Total volume
\providecommand{\Ric}{\operatorname{Ric}}   % Ricci tensor
\providecommand{\Rm}{\operatorname{Rm}}     % Riemann tensor
\providecommand{\Tr}{\operatorname{Tr}}     % Uppercase trace
\providecommand{\Spec}{\operatorname{Spec}} % Uppercase spectrum
\providecommand{\diag}{\operatorname{diag}} % Diagonal matrix
\providecommand{\rank}{\operatorname{rank}} % Rank
\providecommand{\scal}{\mathcal R} % Scalar curvature
\providecommand{\Cl}{\operatorname{Cl}}     % Clifford algebra
\providecommand{\diam}{\operatorname{diam}} % Diameter

% ======================================================================
% 9. FUNCTIONS, TRIGONOMETRY, AND CALCULUS
% ======================================================================

% Codifferential naming conventions
% (Avoid \ddag: commonly defined by physics.sty and symbol packages)
\providecommand{\codiff}{\mathrm{d}^{\dagger}} % Codifferential operator

% Quantum mechanics bra-ket notation (needed for basis tables)
\providecommand{\ket}[1]{\lvert #1 \rangle}
\providecommand{\bra}[1]{\langle #1 \rvert}
\providecommand{\braket}[2]{\langle #1 \mid #2 \rangle}

% ======================================================================
% 10. PHYSICAL SCALES AND CONSTANTS (Paper 1c / Born Geometry)
% ======================================================================
\providecommand{\Rp}{\ensuremath{R_{\mathrm P}}}     % matched common EH curvature scale (not defined from observed G)
\providecommand{\hBorn}{\ensuremath{h_{\mathrm B}}} % physical representative 2\pi\hbar of one primitive Born scalar action cell
\providecommand{\UEHjoint}{\ensuremath{U_{J}}} % joint-Dirac EH curvature-coordinate scale 8 kappa_J
\providecommand{\Gread}{\ensuremath{G_{\rm scal}}} % scalar-action coordinate c_op^3/(h_B U); physical matching explicit
\providecommand{\GreadJoint}{\ensuremath{G_{J,\mathrm{raw}}}} % raw ordinary-trace coordinate = G_scal/96
\providecommand{\Rds}{R_{\mathrm{dS}}}    % de Sitter Curvature
\providecommand{\Mpl}{M_{\mathrm{Pl}}}    % Planck mass
\providecommand{\lp}{\ell_{\mathrm{Pl}}}  % Planck length
\providecommand{\Pmax}{P_{\mathrm{max}}}  % Maximum momentum
\providecommand{\Mplbar}{\bar{M}_{\mathrm{Pl}}} % Reduced Planck mass
\providecommand{\eff}{\mathrm{eff}}       % Subscript for effective quantities
\providecommand{\Cosmo}{\ensuremath{\Lambda}}                % Cosmological constant (use \Cosmo in source)
\providecommand{\LCDM}{\ensuremath{\Cosmo\mathrm{CDM}}}  % \LambdaCDM (typesets with \Cosmo)
\providecommand{\Lcut}{\ensuremath{\Lambda_{\mathrm{cut}}}}  % Cutoff scale (brace when used as script)
\providecommand{\Rcs}{\ensuremath{R_{\mathrm{cs}}}}      % Configuration-space radius
\providecommand{\CosmoEff}{\ensuremath{\Cosmo_{\eff}}}   % Effective cosmological constant

% Disambiguation: \Rmom is the momentum-space (S^3) radius. 
% Reserve plain R for the 4D curvature scalar in f(R).
\providecommand{\Rmom}{\ensuremath{R_{\mathrm{mom}}}} 
\providecommand{\Punit}{P_{0}} % Momentum unit set by \Rmom (= \hbar/\Rmom)
\providecommand{\Rvec}{\ensuremath{R_{\mathrm{vec}}}} % invariant vector-space radius in the reduced round sector
% Specific numeric values and expressions
\providecommand{\Vlogcap}{L_{\mathrm{Hopf}}^{(1)}}
\providecommand{\Hopf}{\operatorname{Hopf}}
% Scalar phase (S^1) sector + universal cover notation (Paper I-A)
\providecommand{\tper}{T_\star}      % Time period/scale
\providecommand{\Eper}{E_\star}      % Energy period/scale
\providecommand{\ttilde}{\tilde t}   % Universal cover time
\providecommand{\Etilde}{\tilde E}   % Universal cover energy

% ======================================================================
% 10A. SERIES-WIDE LOGICAL / NOTATIONAL DISAMBIGUATIONS
% ======================================================================
% Use these symbols to avoid collisions among: (i) the Paper--III
% no-insertion/geometric-closure class, (ii) Paper--I-A trimer diagonals,
% (iii) Paper--V-A/V-C family-space diagonals, (iv) scalar-index running,
% and (v) protected shell-clock labels.
\providecommand{\Rprim}{\ensuremath{\mathsf R_{\mathrm{prim}}^{\Sigma}}} % primitive finite structural readouts
\providecommand{\Rloop}{\ensuremath{\mathsf R_{\circ}^{\Sigma}}} % closed-loop readout role
\providecommand{\Rbdry}{\ensuremath{\mathsf R_{\partial}^{\Sigma}}} % boundary-conditioned readout role
\providecommand{\Rgen}{\ensuremath{\mathsf R_{\mathrm{fin}}^{\Sigma}}} % finite generated closure of primitive roles
\providecommand{\PhysOcc}{\operatorname{PhysOcc}}
\providecommand{\Actual}{\operatorname{Actual}}
\providecommand{\GeomLic}{\operatorname{GeomLic}}
\providecommand{\InvDist}{\operatorname{InvDist}}
\providecommand{\Dcl}{\ensuremath{\mathsf{D}_1}}      % unmarked positive Dirac-even closed spectral-readout form class
\providecommand{\Dheat}{\ensuremath{\mathsf D_{\mathrm{heat}}}} % one-parameter heat-semigroup subclass
\providecommand{\DtriI}{\ensuremath{D^{\mathrm{tri}}_1}} % Paper--I-A cyclic-trimer diagonal
\providecommand{\DtriII}{\ensuremath{D^{\mathrm{tri}}_2}}
\providecommand{\DfamI}{\ensuremath{D^{\mathrm{fam}}_1}} % Family-space diagonal selector
\providecommand{\DfamII}{\ensuremath{D^{\mathrm{fam}}_2}}
\providecommand{\alpharun}{\ensuremath{\alpha_{\mathrm{run}}}} % scalar spectral-index running
\providecommand{\Ksh}{\ensuremath{K_{\mathrm{sh}}}} % protected shell-clock coordinate
\providecommand{\ddef}{\ensuremath{\delta_{\mathrm{def}}}} % universal defect scalar after Paper--III realization

% ======================================================================
% 10B. SERIES-WIDE LOGICAL STATUS TAGS
% ======================================================================
% Upright small-caps wrapper for audit/status labels.  It prevents LaTeX from
% trying to use unavailable small-caps italic font shapes inside theorem bodies
% such as TU/lmr/m/scit.
\providecommand{\AuditLabel}[1]{}
\renewcommand{\AuditLabel}[1]{{\normalfont\scshape #1}}

% These tags are bookkeeping labels for consistency audits. They do not
% create new mathematical hypotheses; they record how a datum is being used
% in the proof chain.
\providecommand{\StatusPrimitive}{}
\renewcommand{\StatusPrimitive}{\AuditLabel{Primitive}}
\providecommand{\StatusDerived}{}
\renewcommand{\StatusDerived}{\AuditLabel{Derived}}
\providecommand{\StatusImported}{}
\renewcommand{\StatusImported}{\AuditLabel{Imported}}
\providecommand{\StatusConditional}{}
\renewcommand{\StatusConditional}{\AuditLabel{Conditional}}
\providecommand{\StatusBridge}{}
\renewcommand{\StatusBridge}{\AuditLabel{Bridge-derived}}
\providecommand{\StatusComparison}{}
\renewcommand{\StatusComparison}{\AuditLabel{Comparison-only}}

% ======================================================================
% 10C. SERIES-WIDE CLAIM-WORDING AUDIT LABELS
% ======================================================================
% These are textual audit labels used to keep modal language consistent
% across the series.  They introduce no new mathematical assumptions.
\providecommand{\ClaimForced}{}
\renewcommand{\ClaimForced}{\AuditLabel{forced}}
\providecommand{\ClaimUnique}{}
\renewcommand{\ClaimUnique}{\AuditLabel{unique}}
\providecommand{\ClaimCanonical}{}
\renewcommand{\ClaimCanonical}{\AuditLabel{canonical}}
\providecommand{\ClaimMinimal}{}
\renewcommand{\ClaimMinimal}{\AuditLabel{minimal}}
\providecommand{\ClaimClosed}{}
\renewcommand{\ClaimClosed}{\AuditLabel{closed}}
\providecommand{\ClaimPrediction}{}
\renewcommand{\ClaimPrediction}{\AuditLabel{prediction}}

% ======================================================================
% 10D. SERIES-WIDE THEOREM-CHAIN AUDIT LABELS
% ======================================================================
% These labels mark the direction of load-bearing theorem chains. They are
% audit annotations only; they introduce no new mathematical assumptions.
\providecommand{\ChainInput}{}
\renewcommand{\ChainInput}{\AuditLabel{Chain input}}
\providecommand{\ChainOutput}{}
\renewcommand{\ChainOutput}{\AuditLabel{Chain output}}
\providecommand{\ChainGuard}{}
\renewcommand{\ChainGuard}{\AuditLabel{Chain guard}}
\providecommand{\ChainNoBackflow}{}
\renewcommand{\ChainNoBackflow}{\AuditLabel{No backflow}}

% ======================================================================
% 10E. SERIES-WIDE LOCAL THEOREM-AUDIT LABELS
% ======================================================================
% These labels mark a statement-level dependency audit around individual
% theorem/definition nodes. They are bookkeeping labels only.
\providecommand{\LocalAudit}{}
\renewcommand{\LocalAudit}{\AuditLabel{Local audit}}
\providecommand{\UsesOnly}{}
\renewcommand{\UsesOnly}{\AuditLabel{Uses only}}
\providecommand{\Exports}{}
\renewcommand{\Exports}{\AuditLabel{Exports}}
\providecommand{\DoesNotUse}{}
\renewcommand{\DoesNotUse}{\AuditLabel{Does not use}}

% ======================================================================
% 10F. SERIES-WIDE NUMERICAL-ROW AUDIT LABELS
% ======================================================================
% These labels classify individual rows in numerical, benchmark, status,
% and comparison tables.  They are audit annotations only and do not change
% any theorem, value, formula, or comparison number.
\providecommand{\RowAudit}{}
\renewcommand{\RowAudit}{\AuditLabel{Numerical-row audit}}
\providecommand{\RowTheorem}{}
\renewcommand{\RowTheorem}{\AuditLabel{Theorem-fixed row}}
\providecommand{\RowBridge}{}
\renewcommand{\RowBridge}{\AuditLabel{Bridge-fixed row}}
\providecommand{\RowRepresentative}{}
\renewcommand{\RowRepresentative}{\AuditLabel{Representative row}}
\providecommand{\RowComparison}{}
\renewcommand{\RowComparison}{\AuditLabel{Comparison row}}
\providecommand{\RowCalibrated}{}
\renewcommand{\RowCalibrated}{\AuditLabel{Calibrated row}}
\providecommand{\RowEstimate}{}
\renewcommand{\RowEstimate}{\AuditLabel{Estimate row}}
\providecommand{\RowOpen}{}
\renewcommand{\RowOpen}{\AuditLabel{Open row}}
\providecommand{\RowNoSelector}{}
\renewcommand{\RowNoSelector}{\AuditLabel{No selector}}

% ======================================================================
% 10G. TABLE HELPERS
% ======================================================================
% Table helpers used by Paper V-C.  The unused readout-label preset has been
% removed; only active table helpers remain.
\providecommand{\gpsection}[1]{%
  \addlinespace[0.55em]
  \multicolumn{3}{@{}l}{\textsc{#1}}\\[-0.25em]
  \midrule
}
\providecommand{\mathcell}[1]{%
  \(\begin{gathered}#1\end{gathered}\)%
}

\newcommand{\DerBy}[1]{%
  \begin{tabular}[t]{@{}l@{}}#1\end{tabular}%
}

\newcommand{\TableNotes}[1]{%
  \par\medskip
  \begin{minipage}{0.94\linewidth}
  \footnotesize\emph{Notes.} #1
  \end{minipage}\par
}
% ======================================================================
% 10H. PAPER VI-A / QCD NOTATION
% ======================================================================
% Shared symbols used by Paper VI-A and downstream QCD/Yang--Mills handoffs.
\providecommand{\Pqcd}{\mathfrak P_{\rm VIII}^{\rm QCD}}
\providecommand{\Iqcd}{\mathfrak I_{\rm QCD}^{(7)}}
\providecommand{\muQCD}{\mu_{\rm QCD}^{\rm int}}
\providecommand{\LambdaQCD}{\Lambda_{\rm QCD}}
\providecommand{\LambdaBH}{\Lambda_{\rm BH}^{\rm QCD}}
\providecommand{\HQCD}{\mathcal H_{\rm QCD}^{\rm BH}}
\providecommand{\HQCDfin}{H_{\rm QCD}^{\rm fin}}
\providecommand{\Hraw}{\mathcal H_{\rm raw}}
\providecommand{\HphysQCD}{\mathcal H_{\rm phys}^{\rm QCD}}
\providecommand{\Vcol}{V_C}
\providecommand{\Ps}{P_{\rm sing}}
\providecommand{\PiOct}{\Pi_8}
\providecommand{\MSbar}{\overline{\rm MS}}
\providecommand{\YM}{\mathrm{YM}}
\providecommand{\Ktri}{K_{\triangle}}
\providecommand{\KYM}{K_{\YM}^{\rm BH}}
\providecommand{\aBH}{a_{\rm BH}}
\providecommand{\epsBH}{\varepsilon_{\rm BH}}
\newcommand{\RowConstructive}{\textnormal{\textsc{Constructive}}}
\newcommand{\RowObligation}{\textnormal{\textsc{Bridge obligation}}}
\newcommand{\RowTarget}{\textnormal{\textsc{Target row}}}
\newcommand{\RowValidation}{\textnormal{\textsc{Validation row}}}
\providecommand{\slashed}[1]{\not{#1}}

% ======================================================================
% 10I. PAPER VI-B / CONTINUOUS-QFT NOTATION
% ======================================================================
% Shared symbols used by Paper VI-B and any downstream continuous-QFT or
% Yang--Mills handoff paper.  All definitions are guarded so individual
% papers may override them locally before loading this file if needed.
\providecommand{\A}{\mathfrak A}
\providecommand{\Hphys}{\mathcal H_{\rm phys}}
\providecommand{\QCD}{\mathrm{QCD}}
\providecommand{\BH}{\mathrm{BH}}
\providecommand{\loc}{\mathrm{loc}}

\providecommand{\RBH}{\mathcal R_{\rm BH}}
\providecommand{\DG}{\mathfrak D_G}
\providecommand{\PW}{\mathcal{P\!W}} % Peter--Weyl spectral-symbol subspace
\providecommand{\Toe}{\operatorname{Toe}} % Toeplitz compression/representative
\providecommand{\Hol}{\operatorname{Hol}} % Holonomy

% ======================================================================
% 11. LATE-LOADED PACKAGES (Must be at the end of the preamble)
% ======================================================================
% Load late packages only if a paper has not already loaded them, avoiding
% option clashes in paper-specific preambles.
\makeatletter
\@ifpackageloaded{hyperref}{}{\usepackage{hyperref}}
\makeatother
% --- Series-wide disambiguations (Pass 3) ---
\providecommand{\diracdelta}{\delta_{\rm D}} % Dirac delta distribution (delta reserved for hierarchy)
\providecommand{\Hg}{H_{\mathrm g}} % Hodge harmonic projector (keep H for Hubble)
\providecommand{\Gr}{G_{\mathrm r}} % Green operator (keep G for Newton)
\providecommand{\FD}{\mathrm{D}} % functional/variational derivative symbol (delta reserved)
\providecommand{\Id}{\mathbf{1}} % identity/Kronecker tensor (avoid \delta_{ij})


% BEGIN INLINE SOURCE: Paper00-ActiveExternalReferences-v2_130.tex
% Framework v2.191: current reference binding; existing definitions retained.
% Framework v2.190: current22-root registry with native clock-evolution proofs.
% Framework v2.189: exact current22-root external registry, including retained v2.188 hardening owners.
% Framework v2.187: current registry synchronization only; source release v2.186 preserved.
% Framework v2.186: authoritative current 22-paper registry.
% Framework v2.182: all current 22 logical paper AUX bindings.
% Framework v2.170: current I-B v1.43 external AUX binding.
% Framework v2.167: exact current 22-root registry after CP--BL2 integration.
% Framework v2.166: exact current 22-root compact-release AUX registry.
% Framework v2.165: authoritative current root bindings.
% Framework v2.164: exact current 22 manuscript AUX bindings.
% Framework v2.162: finite metric, functional Hilbert and scoped Ward integration.
% Framework v2.161: current 22 manuscript AUX bindings; source only.
% Framework v2.156 / registry v2.113: synchronized current 22 manuscript roots; source-only registry.
% Framework v2.143 / registry v2.111: active manuscript AUX basenames; no generated AUX or PDF included.
% ======================================================================
% Central external-reference imports, numbers only, no PDF jumping
% ======================================================================

% xr-hyper is loaded centrally so that papers may retain their existing
% \usepackage{xr-hyper} lines without redefining \externaldocument after this
% registry is installed.  If plain xr was loaded earlier, retain it.
\makeatletter
\@ifpackageloaded{xr-hyper}{}{%
  \@ifpackageloaded{xr}{}{\usepackage{xr-hyper}}%
}
\makeatother

% Make external-document registration idempotent.  The registry identifies a
% logical paper by its Paper<number><optional letter> prefix.  Prefix-qualified
% imports remain independent.  Current paper roots contain no local
% \externaldocument calls; this guard also makes accidental duplicates harmless.
\ExplSyntaxOn
\prop_new:N \g_bh_external_refs_prop
\cs_new_eq:NN \BHOriginalExternalDocument \externaldocument

\cs_new_protected:Npn \bh_external_document:nnnn #1#2#3#4
 {
  \tl_set:Nn \l_tmpa_tl {#3}
  \regex_replace_once:nnN
   { \A (?: xref/ )? (Paper[0-9]+[A-G]?) .* \Z }
   { \1 }
   \l_tmpa_tl
  \tl_set:Nx \l_tmpb_tl { #1 | \tl_use:N \l_tmpa_tl }
  \prop_if_in:NVTF \g_bh_external_refs_prop \l_tmpb_tl
   {
    \PackageInfo{BornHopf}
     {Skipping~duplicate~external~refs~for~#3}
   }
   {
    \prop_gput:NVn \g_bh_external_refs_prop \l_tmpb_tl {#3}
    \IfNoValueTF{#4}
     { \BHOriginalExternalDocument[#1][#2]{#3} }
     { \BHOriginalExternalDocument[#1][#2]{#3}[#4] }
   }
 }

\RenewDocumentCommand{\externaldocument}{O{} O{} m o}
 { \bh_external_document:nnnn {#1}{#2}{#3}{#4} }

\NewDocumentCommand{\BHExternalRefs}{m}
 {
  % Do not import the aux file of the logical paper currently being compiled,
  % even when the job and registry use different version suffixes.
  \tl_set:Nx \l_tmpa_tl { \c_sys_jobname_str }
  \regex_replace_once:nnN
   { \A (Paper[0-9]+[A-G]?) .* \Z }
   { \1 }
   \l_tmpa_tl
  \tl_set:Nn \l_tmpb_tl { #1 }
  \regex_replace_once:nnN
   { \A (Paper[0-9]+[A-G]?) .* \Z }
   { \1 }
   \l_tmpb_tl
  \str_if_eq:eeTF { \tl_use:N \l_tmpa_tl } { \tl_use:N \l_tmpb_tl }
   {
    \PackageInfo{BornHopf}
     {Skipping~external~refs~for~#1~because~this~is~the~current~document}
   }
   {
    \IfFileExists{xref/#1.aux}
     { \externaldocument[][nocite]{xref/#1}[#1.pdf] }
     {
      \IfFileExists{#1.aux}
       { \externaldocument[][nocite]{#1} }
       {
        \PackageInfo{BornHopf}
         {External~aux~file~#1.aux~not~found;~no~labels~imported}
       }
     }
   }
 }
\ExplSyntaxOff

% External references are automatic.  This file is the sole owner of both the
% bibliography and the current Paper-I-A--VII-D AUX registry (30 roots including I-C).  Missing auxiliaries
% are informational on the first pass and are populated by the source-only
% two-pass series build.  Basenames deliberately omit temporary download/copy
% suffixes such as ``(1)'' or ``-1''.
  % Paper I-A: round S3 compendium.
  \BHExternalRefs{Paper1A-Geometry-v5_181}

  % Paper I-B: complete exceptional geometry and finite spectral structures.
  \BHExternalRefs{Paper1B-ExceptionalGeometry-v1_26}
  % Paper I-D: technical companion of I-A through I-C.
  \BHExternalRefs{Paper1D-TechnicalMaterial-v1_69}
  % Paper I-C: reciprocal M8 compendium.
  \BHExternalRefs{Paper1C-M8Geometry-v1_146}

\ifdefined\BHMathOnly
  % Papers I-A and I-B import only the mathematical atlas labels.
\else
  % Paper II: five physical-mathematical main papers plus one technical proof volume.
  \BHExternalRefs{Paper2A-GeometricConstitution-v3_227}
  \BHExternalRefs{Paper2B-MomentumCarrier-v2_147}
  \BHExternalRefs{Paper2C-ReciprocalLocalization-v2_176}
  \BHExternalRefs{Paper2D-SymmetrySpectra-v2_145}
  % The former II-F export labels remain in II-E; current II-F is the technical volume.
  \BHExternalRefs{Paper2E-KinematicExports-v3_140}
  % Paper II-F: complete technical proof and calculation owner for II-A--II-E.
  \BHExternalRefs{Paper2F-TechnicalMaterial-v1_71}

  % Paper III: Dynamics and Intrinsic Geometric Constitution (six components).
  \BHExternalRefs{Paper3A-GeneratedDynamics-v1_181}
  \BHExternalRefs{Paper3B-SpectralAction-v1_192}
  \BHExternalRefs{Paper3C-ProtectedHierarchy-v1_196}
  \BHExternalRefs{Paper3D-RunningResponse-v1_168}
  \BHExternalRefs{Paper3E-GravityAction-v1_204}
  \BHExternalRefs{Paper3F-CompleteDynamics-v1_280}
  % Paper III-G: complete preserved technical proof owner.
  \BHExternalRefs{Paper3G-TechnicalMaterial-v1_57}

  % Paper IV-A: generated finite quantum mechanics.
  \BHExternalRefs{Paper4A-FiniteQuantumMechanics-v4_217}
  % Paper IV-B: complete finite quantum processes and physical response.
  \BHExternalRefs{Paper4B-QuantumResponse-v1_55}
  % Paper IV-C: extended proofs and finite technical calculations.
  \BHExternalRefs{Paper4C-TechnicalMaterial-v1_50}

  % Paper V-A: The Standard Model.
  \BHExternalRefs{Paper5A-FiniteMatterGeometry-v3_219}

  % Paper V-B: Finite Spectrum and Flavor Closure.
  \BHExternalRefs{Paper5B-TerminalSpectrum-v5_202}

  % Paper V-C: Finite Interaction Boundaries of the Standard Model.
  \BHExternalRefs{Paper5C-InteractionResponse-v2_185}
  % Paper V-D: preserved technical proof and calculation owner.
  \BHExternalRefs{Paper5D-TechnicalMaterial-v1_58}

  % Paper VI-A: Gravity and Closed Cosmology.
  \BHExternalRefs{Paper6A-GravityCosmology-v2_268}

  % Paper VI-B: Geometric Dark Sector and Cosmological Readouts.
  \BHExternalRefs{Paper6B-CosmologyReadouts-v1_215}

  % Paper VI-C: complete technical proof and comparison owner.
  \BHExternalRefs{Paper6C-TechnicalMaterial-v1_53}

  % Paper VII-A: Finite QCD Kernel and Nonperturbative Closure.
\BHExternalRefs{Paper7A-QCDKernel-v1_184}

  % Paper VII-B: Hadronic, Thermal, and Nuclear QCD Readouts.
  \BHExternalRefs{Paper7B-QCDReadouts-v1_168}

  % Paper VII-C: Continuous Quantum Field Theory.
  \BHExternalRefs{Paper7C-ContinuousQFT-v2_208}

  % Paper VII-D: complete technical material for VII-A through VII-C.
  \BHExternalRefs{Paper7D-TechnicalMaterial-v1_60}

\fi

% END INLINE SOURCE: Paper00-ActiveExternalReferences-v2_130.tex





\setlength{\emergencystretch}{3em}
\raggedbottom

\title{\textbf{Born--Hopf Geometry in a Closed Universe V-A: Finite Matter Geometry, Charges and Source Constitution}\\[0.45em]
\large Faithful representation, rooted matter cells and charged source incidence}

\author{Peter Kugelmann\thanks{\texttt{pkugelmann@proton.me}}}
\date{Version 3.219 --- 20 November 2026 \\ Framework v3.075}

\begin{document}
% Publication line-breaking flexibility for technical prose (no math scaling).
\emergencystretch=2em
\maketitle
\begin{abstract}
Paper V-A develops the constituted finite matter geometry over the
Born--Hopf carrier: its distinct charge and chiral roles, faithful global
quotient, rooted family transport, finite Dirac and charged-source operations,
and conditional gauge/fermion exports. Exact rooted source constructions are
separated from the actual same-action incidence and pole certificates required
for a measured particle. Extended character, anomaly and source tests are
preserved in Paper V-D.
\end{abstract}
\tableofcontents
\clearpage

\section{Geometric inputs, theorem domains and occurrence status}
\label{p5a:sec:organized-scope}
\subsection{Framework status of Paper V-A}
\label{p5:sec:framework-status}

\subsubsection{Framework prefix}
Paper~V-A is the Standard--Model carrier layer of the framework architecture.  Its ambient input is the prefix
\begin{equation}
\label{p5:eq:framework-prefix}
\mathfrak F_{\le4}
=
\mathfrak F_{\le3}
+
\text{Paper-IV finite quantum closure},
\end{equation}
where \(\mathfrak F_{\le3}\) contains the Paper-I-A reference geometry, the Paper-II compact kinematic carrier, and the Paper-III dynamical constitution.  The present paper works entirely inside this prefix.
The source theorems used below are from
Paper--II \cite{Kugelmann2026b} and Paper--III \cite{Kugelmann2026c}.

The prefix and the geometric bridge theorems proved below determine the finite
carrier and its Standard--Model representation roles.  The generated
Paper-III nonflat connection and the native structure-group/operator
constructors also determine one full local joint gauge family; finite carriers
still do not establish the analytic Lorentzian hypotheses of spin--statistics
and renormalized QFT.  Accordingly Paper~V-A retains only
\begin{equation}
\label{p5:eq:theorem-scope-convention}
\mathfrak F_{\le4}+\mathsf G_5^{\rm fin}+\mathsf R_{5,\mathcal T}
\ \vdash\ \text{Paper-V-A conclusion}.
\end{equation}
Here \(\mathsf G_5^{\rm fin}\) is shorthand for the theorem outputs of the
finite geometric construction, not an assumption, and
\(\mathsf R_{5,\mathcal T}\subseteq
\{\mathsf R_5^{\rm SSQFT},\mathsf R_5^{\rm PCT}\}\) is exactly the residual
subset named by theorem \(\mathcal T\), together with any further analytic or
numerical interface hypotheses displayed in its statement.  Here
\(\mathsf R_5^{\rm rep}=\varnothing\): the former structural hypercharge
residual is discharged by parent-relation PCAD, while relative chirality is
supplied by the generated diagonal-fibre ancestry theorem.  If a later claim
requires a particular nonzero field, vacuum, source, or background member,
that occurrence must be stated in its own target sort and is not a
representation-role premise.  The legacy symbol
\(\mathsf R_5^{\rm conn}\) denotes the derived theorem certificate
\(\mathsf J_5^{\rm conn,der}=({\rm JOINT}_5^{\rm der})\), not a residual.
The same-connection, same-curvature, and same-profile relations are derived
with it and are not another selector.  The domains
\(\mathsf R_5^{\rm SSQFT}\) and
\(\mathsf R_5^{\rm PCT}\subseteq\mathsf R_5^{\rm SSQFT}\) are analytic theorem
domains.  The words ``derived'', ``minimal'', ``unique'', and ``canonical''
below are therefore always relative to the named domain.  No GP is used to
infer a missing cross-role bridge or analytic QFT domain from the mere
existence of a matrix algebra; the connection certificate instead follows
from the explicit nonflat geometric constructor word proved below.

The literature-verified ordinary spin--\(Z_6\) anomaly obstruction is
narrowed in \ref{p5:sec:v2210-lit-global-domain}; its actual regulated
background/measure descent remains claim-scoped.  This refinement does not
reopen finite representation or charged source closure.

\subsubsection{Layer ledger}
The logical roles of the papers used here are:
\begin{table}[H]
\centering
\caption[Paper-V-A framework-prefix ledger]{Paper-V-A framework-prefix ledger.}
\label{p5:tab:framework-prefix-ledger}
\footnotesize
\renewcommand{\arraystretch}{1.15}
\begin{tabularx}{\textwidth}{@{}p{0.18\textwidth}X p{0.24\textwidth}@{}}
\toprule
\textbf{Layer} & \textbf{Role in Paper~V-A} & \textbf{Status} \\
\midrule
Paper~I-A & Round \(S^3\), Hopf normalization, \(\Spin(4)\) representation data, and finite carrier reference facts. & Reference layer \\
Paper~II & Compact momentum carrier and finite kinematic truncation family, including the generated and co-realized physical \(m_p=1\) momentum role. & Kinematic prefix \\
Paper~III & Scalar Born repair, reciprocal cell, dynamical closure class, hierarchy and effective action data used as dynamical background. & Dynamical prefix \\
Paper~IV & Finite operator bridge, finite-particle closure, and finite quantum interface used by the visible gauge--matter carrier. & Quantum prefix \\
Paper~V-A & Geometric bridge theorems; finite carrier, structurally generated hypercharge and derived chirality, global-group, anomaly, Yukawa, trace-weight, joint-operator, and residual-domain verification. & Derived finite layer with generated hypercharge, generated relative chirality and joint connection family, plus explicitly named member/analytic theorem domains where required \\
Paper~V-B and later & Paper~V-B supplies the elementary finite spectrum, masses, scalar comparison/readout charts, flavor and fermion/Majorana readouts; Paper~V-C proves the native electroweak solder reduction and a finite interaction boundary while keeping scalar stationary/effective-action claims, QCD/strong-CP, and thermal/QFT interfaces separately typed; later papers supply conditional cosmology, finite QCD execution, and continuum-QFT completion. & Export layers \\
\bottomrule
\end{tabularx}
\end{table}

\subsubsection{Task and export boundary}
The Paper--V-A task is to close every finite constitutive bridge without using an
observed charge, mass, mixing angle, or threshold value and then to expose the
irreducible remainder.  Tensor composition and hidden-multiplicity removal
follow from the generated relation-free supports and GP3 leastness.  The
hypercharge stabilizer and both relative-grading projectors are computed
structurally; physical structural hypercharge follows from the generated
parent relation by PCAD, while visible chirality follows from the generated
even ancestry relation.  Common kinetic normalization follows from the
uninserted trace on \(\mathsf J_5^{\rm conn,der}\), whose factor relations are
same-datum consequences rather than premises.  When spin--statistics or local
renormalized QFT is invoked, \(\mathsf R_5^{\rm SSQFT}\) records the analytic
theorem domain.  Physical particle/antiparticle completion uses the narrower
\(\mathsf R_5^{\rm PCT}\) domain.  Neither analytic domain is a selector.

\subsection{Introduction}
\label{p5:sec:introduction}

Paper~V-A reconstructs the minimal Standard--Model representation on its named
geometric bridges from the
oriented finite geometry.  The construction closes the axial quaternionic
frame, the component Lie algebras, relation-free tensor composition, the
coordinate scalar type, the distinguished-line Abelian character kernel, the
two relative gradings, and the common trace theorem.  It then verifies the charge, anomaly, Yukawa,
global-group, trace-weight, and one-loop outputs on their stated residual
domains.
Numerical flavor, mass, threshold, and comparison laws are not part of this
layer.

\label{p5:part:closed-sm}

\paragraph{Orientation.}
This opening part separates generated finite constitution from cross-role
bridges, physical occurrence, and analytic domains.  Oriented incidence, the
minimal tensor carrier, scalar type, Abelian candidate plane and character
kernel, both relative-grading projectors, conjugate-representation
availability, the generated even external/internal ancestry, the full local
joint-connection family, the PCAD hypercharge stabilizer, and the typed
common-trace theorem are finite structural outputs.  Physical antiparticle
completion and the Lorentzian QFT domains are not inferred from those outputs
alone; particular nonzero field/vacuum/background members remain G0/member
occurrences in their own target sorts.

\subsection{family-correspondence firewall}

\begin{remark}[Higgs coordinate versus finite label-exclusion ownership]
\label{p5:rem:v181-higgs-label-owner-firewall}
The coordinate theorem
\(\mathcal V_H\otimes1_Q\otimes1_G\) is family blind at the point of its
generation.  It therefore does not by itself activate a selected
$25$-label projector, the uniform $25$-label state, or any particular
incidence-derived Yukawa coefficient.  If a later physical weighted
Yukawa/source word is proved to consume the Paper--V-B self-exclusion owner,
Paper~V-B's owner-retention theorem fixes its exclusion readout; the present
coordinate theorem cannot supply that source-word occurrence.
\end{remark}

\label{p5:sec:post-v125-family-correspondence}

The generated Yukawa/Majorana support quiver and the family carrier
$G\cong\mathbb C^3$ do not yet determine physical flavor coefficients.  This
revision sharpens the correct missing object.

A metric-compatible family connection is insufficient.  If a quiver arrow is
unitary, $Y_a=U_a$, then $Y_aY_a^*=I$; even allowing one scalar edge weight,
$Y_a=y_aU_a$, gives $Y_aY_a^*=|y_a|^2I$.  Such data cannot produce three
nondegenerate singular values.  A physical flavor constructor must therefore
be a \emph{weighted} dagger-quiver/Dirac correspondence, with nontrivial
positive polar parts as well as sector-relative unitary data.

This does not promote the full Weyl algebra $\operatorname{End}(G)$ into a
selector.  The Weyl closure remains an operator grammar.  Likewise, the
Paper-I-A half-step transports between the two alternating Hurwitz trimers
cannot be reused as family Yukawa transports without a new cross-owner
relation: in the present Paper-V-A typing those trimers feed distinct
component/color-side and family-replication roles, while all generated
core/scalar operations extend relation-free as the identity on $G$.

The optional finite spectral datum $(\mathcal H_F,J_F,\Gamma_F,D_F,\rho)$
remains optional comparison structure.  Writing the same unrestricted Yukawa
coefficients inside a block $D_F$ does not reduce the physical flavor orbit or
supply the scalar vacuum.


\label{p5:part:conclusion-export-ledger}

\subsection{Selector-status partition}
\label{p5:sec:selector-status-partition}
For Paper~V-A the finite constitutive outputs and residual interfaces have the
following distinct logical status.
\begin{enumerate}[label=(\roman*),leftmargin=2.2em]
\item \textbf{Generated/derived/co-realized finite base:} color/family
      incidence, the multiplicity-one carrier \(F\), the scalar doublet type,
      \(B-L\), the Abelian candidate plane and distinguished-line character
      kernel, both projectors \(P_\pm\), functorial availability of \(\bar F\),
      the generated even relative-chirality ancestry, and the
      representation-theoretic trace calculation.
\item \textbf{Structural hypercharge and occurrence boundary:}
      Proposition~\ref{p5:prop:v187-pcad-yfix} generates physical structural
      hypercharge, so \(\mathsf R_5^{\rm rep}=\varnothing\).  Particular
      nonzero fields, sources, vacua, or backgrounds remain separately typed
      member/G0 occurrences.  \((\mathrm{CHI}_5)\) is only a legacy label for
      the derived relation \eqref{p5:eq:chi5-derived}.
\item \textbf{Derived connection/operator ancestry:}
      \(({\rm JOINT}_5^{\rm der})=\mathsf J_5^{\rm conn,der}\) contains one
      canonical nonflat connection on the full candidate bundle and the
      complete unmarked gauge/odd Dirac family.  Paper~III-F derives the
      common factor connection, curvature, and profile.  Neither the anchor
      nor those relations are a further selector or occurrence premise.
\item \textbf{Analytic theorem domains:}
      \(\mathsf R_5^{\rm SSQFT}\) supplies the hypotheses under which
      spin--statistics and the local background-gauge/heat-kernel calculation apply, while
      \(\mathsf R_5^{\rm PCT}\subseteq\mathsf R_5^{\rm SSQFT}\) supplies the
      refinement used for physical conjugate-sector completion.  Neither
      domain selects a finite member.
\item \textbf{Downstream realization families:} Yukawa/Majorana matrices,
      Higgs-vacuum data, numerical couplings, thresholds, and optional finite
      spectral-triple realizations.
\end{enumerate}
The chronology is one-way:
\[
\boxed{\begin{aligned}
&\mathcal H_p\cong\mathcal H_1,\quad m_p=1,\\
&({\rm JOINT}_5^{\rm der})\text{ is generated from the prefix and is
independent of }\mathsf R_5^{\rm SSQFT},\\
&\text{Yukawa matrices, the Higgs vacuum, and }J_F,D_F.
\end{aligned}}
\]
No later Paper--V-A representation or interaction datum may be used to reselect
the already generated Paper--II momentum carrier.


\subsection{Exceptional-parent indifference of the finite Standard-Model core}
\label{p5:sec:v236-parent-indifference}
\begin{theorem}[Local tetrahedral-core factorization]
\label{p5:thm:v236-local-core-parent-indifference}
Every finite Paper--V-A construction in this paper which consumes Hurwitz
incidence, the alternating trimers or the active tetrahedral branch factors
through one same-owner tetrahedral core
\[
 H\cong2T.
\]
The theorem does not require that \(H\) be the maximal exceptional parent.
The same local construction applies when \(H\) occurs directly, when
\(H=[2O,2O]\) is generated from an octahedral parent, or when \(H\) is one
member of the five tetrahedral cores generated inside an icosahedral parent.
Parent-specific normalizer or five-core incidence data are retained only by
claims that actually consume them.

On each such branch, Paper~III-C now generates the complete unmarked active
axis family from \(H\); a separate body-diagonal occurrence is not required.
A marked oriented Hurwitz link may still generate the matching particular
active member, and that later member relation does not select the global
exceptional parent.
\end{theorem}

\begin{remark}[Preserved v2.35 closures]
\label{p5:rem:v236-preserved-closures}
The structural hypercharge theorem \(Y=T_R^3+\frac12(B-L)\), the derived joint
connection/operator family, and the generated relative chirality are retained
unchanged.  This section retypes only their exceptional ancestry; it does not
reopen YFIX, JOINT, or chirality as residual selectors.
\end{remark}
% END v2.36 LOCAL CORE PARENT-INDIFFERENCE


% BEGIN ROOTED-GEOMETRIC-SRCVAL-CLOSURE v3.92
\subsection{Framework handoff after the Standard-Model carrier}
\label{p5:sec:framework-handoff}

\subsubsection{Handoff ledger}
\begin{table}[H]
\centering
\caption[Standard-Model carrier handoff]{Standard-Model carrier handoff.}
\label{p5:tab:standard-model-carrier-handoff}
\footnotesize
\renewcommand{\arraystretch}{1.15}
\begin{tabularx}{\textwidth}{>{\raggedright\arraybackslash}p{0.28\textwidth}
                                >{\raggedright\arraybackslash}X
                                >{\raggedright\arraybackslash}p{0.23\textwidth}}
\toprule
\textbf{Export} & \textbf{Status after Paper~V-A} & \textbf{Downstream use} \\
\midrule
Structural visible carrier & Finite carrier, incidence, composition, both
relative gradings, generated visible-chirality ancestry, and
conjugate-representation availability derived.  Local gauge fields use
\(\mathsf J_5^{\rm conn,der}=({\rm JOINT}_5^{\rm der})\); Fermi statistics is
derived on \(\mathsf R_5^{\rm SSQFT}\), and physical conjugate-sector
completion on \(\mathsf R_5^{\rm PCT}\). & Papers~V-B--VI-A \\
Support data & The \(q=5\) rank-three carrier, its \(C_6\)-action, and the
\(6=5+1\), \(5=3+2\) grammar are derived upstream as an equivariant orbit.
The oriented link mark fixes the matching active representative through the
joint zero-defect graph.  Color/family incidence and the one-field
coordinate scalar are derived; Majorana support is optional and its compatible
location is calculated. & Paper~V-B spectrum/flavor readout and
Paper~V-C interaction/EWSB completion \\
Charge and anomaly structure & The candidate plane and character kernel are
finite outputs.  Hypercharge, electric charge, the faithful global group, and
anomaly coefficients follow from the structurally generated hypercharge line
without fitting observed charges; the chiral census uses the derived even ancestry relation, and physical local
anomalies use \(\mathsf J_5^{\rm conn,der}+\mathsf R_5^{\rm SSQFT}\). & Papers~V-B--VI-A \\
Trace and one-loop diagnostics & Relative trace weights are exact; coupling
relations use the structurally generated hypercharge representation on the
typed uninserted kinetic domain, and the one-loop vector is a derived graded boundary response on
\(\mathsf R_5^{\rm SSQFT}\). &
Papers~V-C--VI-B \\
Fixed terminal mass/incidence and native current packets; physical
\(W/Z/H\) and fermion pole/rest comparisons & The rooted charged
zero-order source and comparison owner are generated. A scalar stationary
history and actual physical pole/rest identification remain separately
consumed domains. & Paper~V-B terminal/current readouts and typed physical
comparison \\
Physical coupling/pole comparisons, separately claimed scalar equilibrium/history,
rare-process rates and thermal history & The native solder/stabilizer and gauge-kernel
reduction are generated on their stated branch. The additional state/history and
detector comparison data remain claim-scoped. & Paper~V-C native interaction
outputs and their explicit downstream interfaces \\
Phenomenological fits & Not used as carrier-selection data. & Validation layer only \\
\bottomrule
\end{tabularx}
\TableNotes{The handoff table separates finite theorem outputs, generated hypercharge and
chirality ancestry, separately typed member occurrences, the derived
joint-connection certificate, the analytic
SSQFT/PCT domains, and later interaction/comparison data.}
\end{table}

\subsubsection{Downstream division of labor}
Paper~V-B owns the generated terminal energy/mass and incidence packets;
its later operators and current maps consume those packets without
back-generating masses from a spectrum. Actual source-visible roots,
particle/rest identification, requested reporting coordinates and limiting
comparisons retain the precise downstream domains that their claims consume.  Paper~V-C owns the finite interaction boundary:
its representation/color factors and hot-SM degree count are theorem outputs
on their stated domains. Native electroweak reduction is already generated
by the occurring chiral solder and its positive gauge Gram
(Theorem~\ref{p7v2:thm:v258-native-ew-solder-reduction} and
Corollary~\ref{p7v2:cor:v2109-rooted-ew-occurrence}). Physical pole/coupling
comparisons, separately claimed scalar equilibrium/history,
QCD/strong-CP completion, thermal state and kinetics, and
baryogenesis remain explicit downstream interfaces.  Later papers own
nonperturbative QCD, conditional cosmological propagation, and continuum-QFT
completion.  These layers may consume the Paper--V-A finite base as a theorem
output while retaining the analytic SSQFT/PCT domains and any explicitly named
member/calibration data their claims actually use; the joint-connection certificate is
imported as a theorem output.


% BEGIN ROOTED-FERMION-DESCENDANTS v3.93
\subsection{Geometric particle carriers and the complete source readout}
\label{p5:sec:v224-geometric-source-comparison}
The existing 48-dimensional internal carrier, relative-Weyl sector,
rooted odd solder and hierarchy--family operator maps retain their
generated status. Their analytic 192-dimensional Clifford completion
does not introduce new observed species. The new geometric comparison
in Proposition~\ref{p1:prop:v224-icosian-complement} organizes the
48 antipodal pairs of the icosian complement as four regular $A_4$
modules. Under its specified extended actions,
\[
 F=W\otimes Q\otimes G\cong4\operatorname{Reg}(A_4)
                              \cong\mathbb C[(2I\setminus2T)/\{\pm1\}].
\]
This character equality is stronger than dimension matching. It is still
a comparison representation: its nonunique intertwiner has not been
identified with the complete physical gauge, Dirac, grading and source
transport. It neither rescales the ordinary trace nor selects a particle
source amplitude.

The existing relative grading has 24 internal states of each sign.
With $P_+=(I+\gamma_x\otimes\Gamma_W)/2$, the primitive-owner
retraction gives the insertion $P_+/48$, of rank 96 and trace two.
Proposition~\ref{p1b:prop:v224-retained-primitive} therefore derives the
pure-metric coefficient $\kappa_J/96=U/8$ from existing operations.
This uses the already generated matter carrier downstream; it is not a
new backward selection of Paper II's fixed-radius momentum coordinate.
If this readout is consumed by a complete action, its graded, gauge,
scalar, fermionic and boundary terms follow their actual transport, as in
Proposition~\ref{p3f:prop:v224-normalization-transport}. Their physical
normalizations cannot be left independent of that action comparison.

The rooted full-rank source of
Theorem~\ref{p5:thm:v392-rooted-srcval-incidence} and
Proposition~\ref{p5:prop:v392-rooted-member-full-rank} supplies a
candidate unit section $M_0=K_F/\sqrt{48}$ for a later boundary map.
Its already proved occurrence needs no new existence token. Identifying
that section as the gravity boundary datum and fixing its norm are still
different, substantive action maps. A unit fermionic source is not a
derived bosonic background or scalar VEV. Full support also means a
nonzero positive companion perturbation on this same internal fibre
cannot be invisible to its normal root return; the exact test and second
contacts are in Proposition~\ref{p3f:prop:v224-companion-descendant}.

The charged source synthesis and rooted hierarchy arrows must be composed
on their native owners before scalar readout. The existing $24/25$
tight-frame cross overlap, exterior powers and Hodge compensation remain
exact comparison algebra. Assigning those exterior owners to nine
physical charged grades still requires their actual owner and relative
chiral-frame contractions. Identical frames on both legs can cancel the
apparent correction after kinetic normalization. No new physical mass,
mixing or gauge prediction is asserted by the representation comparison.
% END FRAMEWORK v2.224: particle geometry and source comparison


% BEGIN FRAMEWORK v2.228: CHARGED ORIGIN AND CHARACTER TESTS
\subsection*{charged role-incidence scope correction}
\label{p5:sec:v3106-charged-role-incidence-scope}
\noindent\textbf{Current scope.} Paper~V-B v5.96 and the main-text v3.107 section close the tau singlet endpoint and the discrete mu/tau $E_{36}$ direction choice.  The v3.106 theorem remains the controlling firewall against treating block membership or structural phase availability as a complete in-plane current evaluation.

\begin{theorem}[Generated packet synthesis does not evaluate an absent role arrow]
\label{p5:thm:v3106-symbolic-incidence-not-evaluation}
Let a charged source word be written formally as
\[
 R_s^{\rm ch}=\sum_{r\in\mathcal I_s}m^{\rm fin}_{s,r}\,\iota_{s,r}.
\]
Paper--II-A co-consumption preserves every independently generated
$\iota_{s,r}$ and forbids deleting or rescaling it inside the frozen word.
It does not, however, construct or numerically evaluate a role arrow that has
not already been supplied by the source constructor.  In particular, block
membership in a hierarchy/family operator space is weaker than an explicit
role-to-left-family incidence.
\end{theorem}

\begin{corollary}[Retyped charged-lepton incidence frontier]
\label{p5:cor:v3106-charged-lepton-incidence-frontier}
The charged terminal masses and their mathematical terminal supports retain
all previously derived statuses.  The additional physical readout
\[
 J_e:\mathcal O_e\longrightarrow G_L
\]
or an equivalent role-resolved projector triple must still be evaluated by
an actual source/current constructor before a numerical lepton-current CP
invariant is claimed.  Accordingly the stronger reading of the v3.97
preterminal closure as an evaluation of the charged-lepton left-family frame
is superseded.  No new source-value selector is introduced by this
retyping.
\end{corollary}

\begin{definition}[Charged-lepton orientation target]
\label{p5:def:v3106-charged-lepton-orientation-target}
Let $P_e$ be the rank-one physical electron-role projector in $G_L$ when the
source-owned lepton incidence has been evaluated, and retain the generated
rooted odd family operator $C_{2,G}=-i(X_G-X_G^\dagger)$.  Define
\[
 \boxed{\zeta_e:=\frac1{\sqrt3}\operatorname{Tr}(P_eC_{2,G}).}
\]
The current framework supplies the host and the structural operator
$C_{2,G}$; the value of $\zeta_e$ is an owner-incidence readout and is not a
new free parameter.
\end{definition}
% END CHARGED-ROLE-INCIDENCE-SCOPE-CORRECTION v3.106

% BEGIN NATIVE-MATTER-VACUUM-LEDGER v2.216

\paragraph{Technical proof owner and claim-status guard.}
Detailed auxiliary matrix, representation, finite-spectral and comparison
calculations formerly embedded in this manuscript now belong to
Paper~V-D. The original named theorem/proof labels remain
cross-referenced, and no optional finite comparison is promoted
to physical source occurrence, interacting mass pole or measured response.

\section{Generated finite matter carrier and joint connection}
\label{p5a:sec:organized-carrier}
\subsection{Finite geometric constitution bridges}
\label{p5:sec:finite-geometric-constitution-bridges}

\begin{theorem}[Generated \(1+3\) incidence and complementary multiplicity]
\label{p5:thm:generated-color-family-incidence}
On the fixed oriented Hurwitz-link branch of Paper~I-A, let its positive axial unit be
\(u_+\).  Paper~I-A,
Theorem~\ref{p1:thm:axial-opposite-trimer-quaternionic-frame}, gives
\[
 \{h_{u_+}\}\sqcup\{h_\eta:\eta\in T_-\}
\]
as an orthonormal quaternionic frame and maps it by
\(L_{\bar h_{u_+}}\) to
\(\{1\}\sqcup\operatorname{Im}\HH\).  Therefore
\[
 Q=\ell\oplus C
 \cong\CC1\oplus\operatorname{Im}(\HH)_\CC,
 \qquad
 G=\CC^{T_+}\cong\CC^3 .
\]
The first trimer is characterized as the unique link trimer orthogonal to the
actual axial unit; the second is its link complement.  Hence the \(1+3\)
component and rank-three multiplicity roles are incidence outputs, not a
choice between two indistinguishable trimers.  Orientation reversal produces
the conjugate/unitarily equivalent representative.
\end{theorem}

\begin{proof}
The cited inner-product identity is
\(\langle h_\varepsilon,h_\delta\rangle
=(1+\varepsilon\cdot\delta)/4\).  The axial sign and every member of \(T_-\),
and every two distinct members of \(T_-\), have sign dot product \(-1\).
This proves the orthonormal frame and the \(1+3\) split.  The diagonal-link
partition then leaves exactly \(T_+\).  The equivalence claim is the
orientation statement of the cited theorem.
\end{proof}

\begin{theorem}[Coordinate generation of one blind Higgs doublet]
\label{p5:thm:generated-coordinate-higgs}
On the protected half-spin block, the round coordinate representation is
\[
V_{\rm coord}=(\tfrac12,\tfrac12),
\]
with its canonical four-dimensional real form \(V_{\rm coord}^{\RR}\).
The two retained Clebsch--Gordan channels
\(V_{\rm coord}\otimes W_{R\pm}\to W_L\) have multiplicity one.  Their
scale-free self-adjoint real intertwiner type can therefore be written
\[
 \mathcal V_H
 =
 \left\{
 \begin{pmatrix}0&\Phi(H)\\ \Phi(H)^\dagger&0\end{pmatrix}
 :H\in\CC^2
 \right\},
 \qquad
 \Phi(H)=(\widetilde H,H),\quad
 \widetilde H=i\sigma_2\bar H .
\]
Thus \(\mathcal V_H\cong V_{\rm coord}^{\RR}\cong\RR^4\cong\CC^2\) contains one complex
\(SU(2)_L\) doublet and its pseudoreal conjugate, equivalently
\[
 H:W_{R-}\longrightarrow W_L,\qquad
 \widetilde H:W_{R+}\longrightarrow W_L .
\]
The coordinate representation and its protected intertwiner type are
generated before \(Q\) and \(G\) are composed.
Relation-free GP2 composition therefore transports it uniquely as
\[
 \mathcal V_H\otimes1_Q\otimes1_G .
\]
Color and family blindness and minimal multiplicity are consequently
theorems, not scalar-role assumptions.  This statement uses only the abstract
representation and multiplicity-one intertwiner spaces.  A concrete Toeplitz
coordinate quadruple, its reduced matrix elements, adjoint phases, and its
scale \(\rho_{\rm prot}\) are not used.  Nonzero Yukawa matrices, a scalar
potential, and a vacuum expectation value are not conclusions of this theorem.
\end{theorem}

\begin{proof}
The Clebsch--Gordan decompositions contain exactly one retained copy of the
opposite half-spin module in each direction.  The quaternionic \(2\times2\)
realization identifies their self-adjoint real form with \(H\in\CC^2\);
pseudoreality of the weak fundamental fixes the other column to
\(i\sigma_2\bar H\).  This leaves four real, rather than eight independent,
scalar coordinates, independently of a common nonzero normalization.  The
generated type has no \(Q\)- or \(G\)-label.  Paper~IV's relation-free
composition theorem therefore extends it by the identities.  GP2 no-omission
requires the generated intertwiner type to occur, and Paper~II-B,
Proposition~\ref{p2:prop:gp3-multiplicity-bound}, removes any additional
unmarked multiplicity factor.
\end{proof}

% BEGIN FRAMEWORK v3.029: GP3 MATTER OWNER BRIDGE
\begin{remark}[Common-parent scalar type versus independent copy naturality]
\label{p5:rem:v329-commonparent-higgs}
The one generated coordinate-doublet \emph{type} follows from the
common $V_{\rm coord}^{\RR}$ parent and its proven $1_Q\otimes1_G$
identity extension. Equality $\dim M_\ell=\dim M_C=k$ on an
amplified candidate does not canonically identify those two spaces.
Under independent unitary presentation changes on $M_\ell$ and $M_C$,
separate scalar coefficients on the $\ell$ and $C$ central blocks
are natural. A common scalar coefficient therefore comes from the
actual generated common-parent constructor (or an independently
certified joint multiplicity owner), not from rank equality or
unmarked unitary naturality alone. The physical non-repetition of
an \emph{occurring} scalar field retains the complete GP3 comparison
condition of Proposition~\ref{p2:prop:gp3-multiplicity-bound}.
\end{remark}
% END FRAMEWORK v3.029

% BEGIN C01 V2174 FIELD DICTIONARY

% BEGIN INLINE SOURCE: Paper5A-FieldCoefficientDictionary-Main-v1_0.tex
% v2.174: evaluated coordinates of the already generated field type.
\begin{corollary}[Explicit pair-to-field coefficient dictionary]
\label{p5:cor:v2174-pair-field-dictionary}
Under the orientation-preserving ancestry of
Theorem~\ref{p3b:thm:diagonal-fibre-spin-ancestry}, write
\[
 \operatorname{Sym}^2(S_L\oplus S_R)
  =\operatorname{Sym}^2S_L\oplus(S_L\otimes S_R)
                   \oplus\operatorname{Sym}^2S_R.
\]
With $\epsilon=i\sigma_2$, the maps $K\mapsto K\epsilon$ on the
symmetric left/right pairs and $K\mapsto-K\epsilon$ on mixed pairs
give their adjoint and scalar coefficient presentations. For real $h_a$,
\begin{equation}
 A(h)=h_4I_2-i\sum_{j=1}^3h_j\sigma_j
       =(\epsilon\bar H,H),\qquad
 H=(-h_2-ih_1,h_4+ih_3)^T.
 \label{p5:eq:v2174-pair-field-dictionary}
\end{equation}
This is the explicit real quadruple of
Theorem~\ref{p5:thm:generated-coordinate-higgs}, extended by the already
licensed $1_Q\otimes1_G$.
\end{corollary}
\begin{proof}
The invariant identity $g^T\epsilon=\epsilon g^{-1}$ gives the adjoint
and mixed intertwining rules. Substitution of the Pauli matrices gives the
two columns of $A(h)$. The already generated scalar type is color/family
blind, so its established identity extension applies.
\end{proof}
This is a coefficient dictionary, not a new source occurrence or quantum
ladder constructor. Internal left/right spin labels retain their distinction
from Lorentz chirality. Normalized basis images, charges and the active gauge
scope are evaluated in Appendix~\ref{p5:app:v2174-field-coefficient-dictionary}.

% END INLINE SOURCE: Paper5A-FieldCoefficientDictionary-Main-v1_0.tex

% END C01 V2174 FIELD DICTIONARY


\begin{theorem}[Unit Clifford solder as the normalized generated scalar direction]
\label{p5:thm:v377-unit-solder-normalized-scalar}
For \(H=(h_1,h_2)^T\in\CC^2\), the generated coordinate intertwiner
\[
 \Phi(H)=(\widetilde H,H),\qquad
 \widetilde H=i\sigma_2\bar H,
\]
obeys
\[
 \boxed{\Phi(H)^\dagger\Phi(H)=(H^\dagger H)I_2},
 \qquad
 \boxed{\det\Phi(H)=H^\dagger H}.
\]
Consequently \(\Phi\) restricts on \(H^\dagger H=1\) to the standard
quaternionic \(SU(2)\) realization.  In the compatible quaternionic spin
frame of Paper~I-A, and modulo the already removable relative chiral phase,
for every generated unit quaternion \(\hat h\) there is a unique normalized
\(\widehat H\in\CC^2\) such that
\[
 \boxed{\Phi(\widehat H)=c(\hat h).}
\]
Thus the unit Clifford solder and the normalized direction of the generated
scalar module are the same geometric datum; no scalar magnitude or vacuum
expectation value is introduced.
\end{theorem}

\begin{proof}
Writing \(H=(h_1,h_2)^T\) gives
\[
 \Phi(H)=
 \begin{pmatrix}
  \bar h_2 & h_1\\
  -\bar h_1 & h_2
 \end{pmatrix}.
\]
Direct multiplication yields
\(\Phi(H)^\dagger\Phi(H)=(|h_1|^2+|h_2|^2)I_2\), while direct expansion of
the determinant gives \(\det\Phi(H)=|h_1|^2+|h_2|^2\).  Hence the unit sphere
maps into \(SU(2)\).  Conversely every \(U\in SU(2)\) has the unique standard
form
\(
U=\left(\begin{smallmatrix}\bar h_2&h_1\\-\bar h_1&h_2\end{smallmatrix}\right)
\)
for its unique unit second column \(H=(h_1,h_2)^T\).  Paper~I-A's compatible
quaternionic spin frame represents the generated unit Clifford element by the
same \(SU(2)\) matrix, up to the already removable relative chiral phase.
\end{proof}

\begin{theorem}[Distinguished-line character kernel]
\label{p5:thm:exact-oriented-vector-hypercharge}
The oriented Hopf Cartan and axial quaternionic component canonically fix the
one-dimensional representation module
\[
 L_0:=W_{R+}\otimes\ell .
\]
No choice of a nonzero vector, phase, norm, or Hopf section in \(L_0\) is
made.  Let
\[
 B-L=\frac13P_C-P_\ell,\qquad
 X=aT_R^3+b(B-L).
\]
Then the kernel of the infinitesimal character of \(L_0\), in the
projective/traceless two-plane, is the one-dimensional line
\[
 X|_{L_0}=0
 \quad\Longleftrightarrow\quad
 b=\frac a2,
 \qquad
 Y=T_R^3+\frac12(B-L)
\]
after normalizing the inherited \(T_R^3\) coefficient to one.  This derives
the unique kernel ray in the candidate plane without a Majorana-neutrality
premise.  It does not, by itself, identify that ray as the physically
surviving Abelian subgroup.
\end{theorem}

\begin{proof}
The oriented Hopf Cartan has weight \(+\tfrac12\) on \(W_{R+}\), while
\(B-L\) has weight \(-1\) on \(\ell\).  Therefore the character of \(X\) on
\(L_0\) is
\[
 \chi_0(X)=\frac a2-b,
\]
and its kernel gives \(b=a/2\).  The common identity direction
has already been removed by projective/traceless representation, so this is
the unique nonzero connected line.  Because a character kernel depends only
on the one-dimensional module, it is invariant under every nonzero rescaling
or phase choice and requires no identification of vectors in its two tensor
factors.  Reversing the oriented token conjugates the representation but
preserves the unmarked gauge line.
\end{proof}

\begin{corollary}[Parent-relation stabilizer hypercharge]
\label{p5:cor:yfix-hypercharge}
On \(\mathsf G_5^{\rm fin}\), the nonzero projected parent relation of
Proposition~\ref{p5:prop:v187-pcad-yfix} has connected stabilizer
\[
 \operatorname{Lie}(K_Y)
 =\mathbb R\!\left(T_R^3+\frac12(B-L)\right).
\]
Thus the physical structural Abelian generator is
\(Y=T_R^3+\tfrac12(B-L)\), up to the inherited overall generator/coupling
convention; no independently selected nonzero pointing is required.
\end{corollary}
\begin{proof}
The projected relation has torus weight \((-1,+2)\), so its connected
stabilizer is the kernel \(-a+2b=0\).  This is exactly the PCAD descendant of
the generated parent relation.
\end{proof}

\begin{theorem}[Relative-grading dichotomy]
\label{p5:thm:generated-orientation-coherent-chirality}
Let \(\gamma_x\) and \(\Gamma_W\) be the external and internal grading
operators on the oriented local Clifford branch, with
\[
 \gamma_x|_{S_L}=-1,\quad\gamma_x|_{S_R}=+1,\qquad
 \Gamma_W|_{W_L}=-1,\quad\Gamma_W|_{W_R}=+1 .
\]
The representation data admit the two projectors
\[
 P_\pm=\frac12(1\pm\gamma_x\otimes\Gamma_W),
\qquad
\operatorname{rank}P_+=\operatorname{rank}P_-.
\]
They commute with every even external and internal action, and simultaneous
orientation reversal leaves both unchanged.  Their ranges are
\[
 \operatorname{Ran}P_+
 =(S_L\otimes W_L)\oplus(S_R\otimes W_R)
\]
and
\[
\operatorname{Ran}P_-
=(S_L\otimes W_R)\oplus(S_R\otimes W_L).
\]
Consequently the separate gradings do not select the visible relative parity;
GP3 cannot delete one equal-rank invariant realization merely because the
other is phenomenologically desired.  This theorem does not derive
statistics.
\end{theorem}

\begin{proof}
Since \(\gamma_x^2=\Gamma_W^2=1\), the displayed operators are complementary
spectral projectors.  The \(\pm1\) external and internal eigenspaces have the
same dimensions, which gives equal ranks.  Every even action commutes with
both gradings, hence with \(P_\pm\).  Replacing
\((\gamma_x,\Gamma_W)\) by \((-\gamma_x,-\Gamma_W)\) leaves their product
unchanged.  Expanding the two product eigenspaces gives the stated ranges.
\end{proof}

\begin{corollary}[Generated diagonal ancestry selects visible chirality]
\label{p5:cor:chi5-visible-chirality}
On the generated Paper~III-B/III-E ancestry chain, the even map
\(J_\Delta\gamma_x=\Gamma_WJ_\Delta\) selects the relative-even fibre product,
where here $W=W_{\mathrm{spin}}=W_L\oplus W_R$ is the complex internal
spin carrier used by $\Gamma_W$, not Paper~III-E's real $W_7$ response block;
hence
\[
P_{\rm phys}=P_+
=\frac12(1+\gamma_x\otimes\Gamma_W),
\qquad
\operatorname{Ran}P_{\rm phys}
=(S_L\otimes W_L)\oplus(S_R\otimes W_R).
\]
No independent \(( {\rm CHI}_5 )\) premise is used.
\end{corollary}

\begin{proof}
The intertwining identity maps the external \(\varepsilon\)-eigenspace
isomorphically to the internal \(\varepsilon\)-eigenspace.  Thus the
co-realization relation fixed by the bridge is equality of grading labels,
\(\varepsilon_x=\varepsilon_W\).  On \(S_x\otimes W\) this is exactly the
\(+1\) eigenspace of \(\gamma_x\otimes\Gamma_W\), whose projector is \(P_+\).
An odd bridge would instead satisfy
\(J_{\rm or}\gamma_x=-\Gamma_WJ_{\rm or}\) and select \(P_-\); the sign is
therefore genuine bridge data.
\end{proof}

\begin{corollary}[Functorial core conjugation; optional Majorana data]
\label{p5:cor:functorial-conjugate-completion}
Orientation reversal and complex conjugation functorially supply
\(F\mapsto\bar F\), canonically up to unitary equivalence.  The rooted
transport-cell theorem below independently supplies the finite family
involution \(I_G\mapsto I_G\), \(X_G\leftrightarrow X_G^{-1}\).  These two
facts must not be conflated: family inversion alone does not reverse gauge
charge.  No finite real spectral-triple operator \(J_F\) is asserted here.
A symmetric Majorana operator may be added later on a fully self-conjugate
neutral core; neither its occurrence nor its nonzero eigenvalues enter the
constitution or the character-kernel calculation.
\end{corollary}

\begin{corollary}[PCT continuum antiparticle identification]
\label{p5:cor:pct-conjugate-completion}
On \(\mathsf R_5^{\rm PCT}\), the antiunitary PCT theorem identifies the local
positive-energy field sector transforming on \(F\) with the conjugate physical
antiparticle sector transforming on \(\bar F\) \cite{Haag1996}.  The finite
internal conjugate representation and rooted inverse transport are generated
structurally upstream; PCT supplies their continuum local-QFT
particle--antiparticle interpretation and does not create an additional finite
source branch.
\end{corollary}

\begin{theorem}[Minimized boundary of the Standard--Model constitution]
\label{p5:thm:two-residual-sm-constitution}
The Papers~I-A--IV prefix and the preceding finite bridge theorems derive the
carrier, \(1+3\) incidence, threefold multiplicity, scalar doublet,
nonabelian gauge algebra, Abelian candidate plane and its character kernel,
both relative-grading projectors, generated even chirality ancestry, and
conjugate-representation availability.  Physical structural hypercharge is generated by
Proposition~\ref{p5:prop:v187-pcad-yfix}; visible chirality is derived.
Theorem~\ref{p5:thm:canonical-full-joint-connection} realizes the generated
derivations in one joint local connection family and supplies a nonflat
canonical member; Paper~III-F then derives the factor relations and, on its
typed kinetic domain, the common
uninserted-trace functional.  If \(\mathsf R_5^{\rm SSQFT}\) is verified,
Paper~IV's spin--statistics interface yields exterior powers and exact CAR and
the standard local-QFT anomaly and loop theorems are licensed.  Physical
conjugate-sector completion uses the refinement \(\mathsf R_5^{\rm PCT}\).
\end{theorem}


\subsubsection{Level of the Paper--V-A operators}
\label{p5:sec:operator-level}
Paper~II classifies bare kinematic profiles, Paper~III supplies their
dynamical realization and the diagonal-fibre spin-isometry carrier, and
Paper~IV closes the finite quantum calculus on generated typed roles.
Paper~V-A applies those theorems on \(\mathsf G_5^{\rm fin}\), with physical
structural hypercharge generated by Proposition~\ref{p5:prop:v187-pcad-yfix},
using \(\mathsf J_5^{\rm conn,der}\) for the derived local connection, and
\(\mathsf R_5^{\rm SSQFT}\) only for statistics or local-QFT statements.

\begin{equation}
\label{p5:eq:src-op-subset-envelope}
\operatorname{Op}_{\rm src}(\mathfrak F_{\le4})
\subsetneq
\End(\mathcal H_{m_p=1})
\end{equation}
remains a strict inclusion: the event envelope contains many operator modules
that are not selected by the generated incidence, coordinate, stabilizer, and
grading tests.

\begin{proposition}[Componentwise closure form]
\label{p5:prop:componentwise-closure-form}
The structural theorem form used in Paper~V-A is
\begin{equation}
\label{p5:eq:componentwise-theorem-form}
\mathfrak F_{\le4}+\mathsf G_5^{\rm fin}
\vdash
\left(
\mathfrak{su}(2)_L,
\mathfrak{su}(3)_c,
\CC[Z_3]_{\rm rep}
\right)
\xRightarrow[\substack{\text{GP2 no omission}\\
\text{GP3 no hidden multiplicity}}]{\text{relation-free product}}
F.
\end{equation}
The coordinate and distinguished-line character-kernel theorems then
determine the Abelian candidate ray.  The physical hypercharge direction and
charge table follow from Proposition~\ref{p5:prop:v187-pcad-yfix}.  This formula records
theorem outputs, not hidden role clauses.
\end{proposition}

\begin{theorem}[Relation-free product of the generated supports]
\label{p5:thm:componentwise-no-deletion}
On the declared relation-free physical candidate envelope, the \(W,Q,G\) supports carry the generated commuting incidence and Clifford tests without an additional licensed cross-action relation. The upstream joint geometric quadratic is retained in its geometric sort; this candidate-envelope statement does not assert its absence or identify its action with the physical target action.  Paper~IV's GP2 composition theorem
therefore gives their tensor product.  GP2 no-omission supplies at least one
copy of each marked component, while Paper~II-B,
Proposition~\ref{p2:prop:gp3-multiplicity-bound}, supplies at most one
unmarked copy.
\end{theorem}

\begin{proof}
Direct sums encode alternatives or superselection.  Here the simultaneously
generated commuting test families are independently coexistent full event
domains.  Relation-free composition gives the product module.  On every
generated simple block its multiplicity \(m_\alpha\) satisfies
\(m_\alpha\ge1\) by GP2 faithfulness/no-omission and \(m_\alpha\le1\) by the
cited GP3 multiplicity boundary.  Hence \(m_\alpha=1\).
\end{proof}

\begin{lemma}[Presentation covariance is not gauge promotion]
\label{p5:lem:unresolved-support-full-unitary-closure}
Let \(V\) be a finite carrier with maximal event envelope \(\End(V)\).
Simultaneous unitary conjugation is a covariance of presentations.  With a
constituted tuple \(\mathcal T\) held fixed, the actual unitary symmetry is only
its stabilizer
\[
\operatorname{Stab}_{U(V)}(\mathcal T)
=\{u\in U(V):uTu^{-1}=T\ \text{for every }T\in\mathcal T\}.
\]
Neither \(\End(V)\) nor its presentation covariance promotes
\(\mathfrak{su}(V)\) to a physical local gauge algebra.
\end{lemma}

\begin{proof}
Paper~IV proves
\(u\End(V)u^{-1}=\End(V)\) only when every represented datum is conjugated
together.  Holding the projectors, grading, finite grammar, or geometric
operators fixed changes the question from presentation covariance to
stabilizer symmetry.  A local gauge interpretation additionally uses the
derived connection \(\mathsf J_5^{\rm conn,der}\).  The finite
derivation algebra and its stabilizer are already outputs of
\(\mathsf G_5^{\rm fin}\).
\end{proof}

\begin{lemma}[Generated status of generation blindness]
\label{p5:lem:replication-not-gauge-support}
The opposite trimer
\[
G=\mathbb C^{T_g}\cong\mathbb C[Z_3]
\]
carries a cyclic replication grammar.  Its full matrix algebra alone does not
forbid continuous unitary actions.  In the present construction, however, the
weak, trimer-derivation, and distinguished-line character-kernel operations are generated
on \(W\otimes Q\) before the complementary support is composed.  Paper~IV's
relation-free extension is therefore
\[
\rho_{\rm gauge}(a)=\rho_{\rm core}(a)\otimes1_G
\]
and generation blindness follows from GP2.  Downstream Yukawa coefficients
may lie in \(\End(G)\) without changing the gauge representation.
\end{lemma}

\begin{proof}
Paper~II-F explicitly distinguishes the trimer's finite grammar from any
later physical use.  Since the clock and shift generate \(M_3(\CC)\), bare
algebra alone would not suffice.  What suffices is provenance: the core
derivations are already generated when \(G\) is added relation-free.  GP2
preserves them by identity.  Yukawa endomorphisms are later coefficients in
intertwiners and are not gauge generators.
\end{proof}

\begin{lemma}[Diagonal-fibre left--right carrier and Hopf split]
\label{p5:lem:protected-chirality-right-hopf-split}
The vector part of Paper~III's diagonal round fibre has connected spin-isometry
carrier
\[
\Spin(4)_{\rm fib}\cong SU(2)^a_L\times SU(2)^a_R
\quad\text{on}\quad
W_L\oplus W_R.
\]
Under the generated equivariant identification, the Hopf
operator resolves
\[
W_R=W_{R+}\oplus W_{R-}
\]
and the split stabilizer retains
\(SU(2)_L\times U(1)_{T_R^3}\).  The Born-odd tangent module is not this gauge
Lie algebra.  The internal half-spin labels are also distinct from physical
Lorentz chirality.
\end{lemma}

\begin{proof}
Paper~III's Born-parity theorem gives a symmetric-pair decomposition with
\([\mathfrak k_-,\mathfrak k_-]\subseteq\mathfrak k_+\) and explicitly states
that \(\mathfrak k_-\) is not a gauge Lie algebra.  The same theorem identifies
the diagonal round-fibre isometry and its spin lift
\(\Spin(4)_{\rm fib}\).  Stabilizing an oriented right Cartan preserves the two
right eigenlines and leaves \(U(1)_{T_R^3}\).  The Yang--Mills connection
family is realized by the derived certificate
\(\mathsf J_5^{\rm conn,der}\).
\end{proof}

\begin{theorem}[Generated independent component composition]
\label{p5:thm:component-independence-composition}
The following composition is on the declared physical candidate envelope. A later licensed joint relation must be retained in its stabilizer calculation; shared geometric ancestry alone does not impose a tensor quotient or erase such a relation.

On \(\mathsf G_5^{\rm fin}\), iterating Paper~IV's
independent event-envelope theorem gives, uniquely up to unitary equivalence,
\begin{equation}
\label{p5:eq:componentwise-joint-carrier}
F
=
(W_L\oplus W_{R+}\oplus W_{R-})
\otimes(\ell\oplus C)
\otimes G,
\end{equation}
with no unobserved multiplicity and \(\dim_\mathbb C F=48\).  Using the
derived even ancestry relation \eqref{p5:eq:chi5-derived}, set
\[
F_L=W_L\otimes(\ell\oplus C)\otimes G,
\qquad
F_R=(W_{R+}\oplus W_{R-})\otimes(\ell\oplus C)\otimes G.
\]
The matter one-particle Weyl carrier is
\begin{equation}
\label{p5:eq:physical-one-particle-factorization}
\mathcal H_{\rm 1p}^{\rm matt}
=(S_L\otimes F_L)\oplus(S_R\otimes F_R).
\end{equation}
The separate unbridged reduct admits both \(P_+\) and \(P_-\), but the
generated diagonal-fibre ancestry selects \(P_+\).  Complex conjugation supplies the representation
\(\bar F\) functorially but not its simultaneous physical occurrence.
Conjugate-sector physical completion follows only on
\(\mathsf R_5^{\rm PCT}\), or from an explicitly supplied optional finite real
structure.
\end{theorem}

\begin{proof}
Let \(\mathcal A_{\rm gen}\) be the actual finite-dimensional algebra
generated by the component incidence/Clifford operations together with the
family Weyl grammar.  It is a finite semisimple algebra, so every represented
carrier decomposes as
\[
\mathcal K\cong
\bigoplus_{\alpha}V_\alpha\otimes M_\alpha,
\qquad
\mathcal A_{\rm gen}\cong
\bigoplus_{\alpha}\End(V_\alpha)\otimes1_{M_\alpha}.
\]
Only the generated summands occur in this expression.  GP2
faithfulness/no-omission gives \(\dim M_\alpha\ge1\) for each marked
summand.  The host signature contains no mark on \(M_\alpha\); Paper~II-B,
Proposition~\ref{p2:prop:gp3-multiplicity-bound}, therefore gives
\(\dim M_\alpha\le1\).  Thus \(\dim M_\alpha=1\) blockwise.  The six component
dimensions are
\[
(2\cdot1\cdot3,\;2\cdot3\cdot3,\;
1\cdot1\cdot3,\;1\cdot1\cdot3,\;
1\cdot3\cdot3,\;1\cdot3\cdot3)
=(6,18,3,3,9,9),
\]
whose sum is \(48\).  This proves the displayed \(F\) using the generated
algebra rather than the ambient
\(\End(W)\otimes\End(Q)\otimes\End(G)\).
Corollary~\ref{p5:cor:chi5-visible-chirality} gives the derived Weyl-carrier formula;
Corollary~\ref{p5:cor:pct-conjugate-completion} gives
the analytic conjugate-sector statement.
\end{proof}

\begin{corollary}[Visible one-particle carrier on the graded bridge]
\label{p5:cor:one-particle-derived-later}
Relative to Paper~IV, which retains a generic one-particle assignment family,
the generated incidence and relation-free tensor structure of Paper~V-A,
together with the upstream even ancestry theorem, give
\[
\mathcal H_{\rm 1p}^{\rm matt}
=(S_L\otimes F_L)\oplus(S_R\otimes F_R)
\]
as a \emph{DERIVED-LATER} specialization once the bridge is constructed.
This
does not assign particle meaning to every finite Paper--IV spectral block.
Moreover
\[
\operatorname{CoReal}(W_L,W_{R+},W_{R-},\ell,C,G\mid F)
\]
in the generated carrier: these are simultaneous constituents, not competing
members from which Paper~V-A selects one.
\end{corollary}

\begin{lemma}[Generated compact candidate plane and character kernel]
\label{p5:lem:compact-phase-line-componentwise}
On \(\mathsf G_5^{\rm fin}\), projective/traceless representation reduces the
generation-blind diagonal centralizer, modulo the common identity, to
\begin{equation}
\label{p5:eq:componentwise-abelian-plane}
Y=aT_R^3+bX_{1|3},
\qquad
X_{1|3}=\frac13P_C-P_\ell.
\end{equation}
The character kernel of \(L_0=W_{R+}\otimes\ell\) gives \(b=a/2\).  Choosing the
representative whose coefficient in the convention-fixed \(T_R^3\) coordinate
is one sets \(a=1\), hence
\begin{equation}
\label{p5:eq:componentwise-hypercharge-line}
Y=T_R^3+\frac12X_{1|3}.
\end{equation}
This is the character-kernel ray.  It is the physical structural compact phase line by
Corollary~\ref{p5:cor:yfix-hypercharge}.
\end{lemma}

\begin{proof}
The complete six-block centralizer and its coordinate-scalar reduction are
proved in Lemmas~\ref{p5:lem:one-abelian-slot-derived}
and~\ref{p5:lem:intrinsic-abelian-centralizer-plane}.  On
\(L_0=W_{R+}\otimes\ell\), the two weights are \(T_R^3=1/2\) and
\(X_{1|3}=-1\); the character kernel gives \(a/2-b=0\), as in
Theorem~\ref{p5:thm:exact-oriented-vector-hypercharge}.
\end{proof}

\begin{lemma}[Generated coordinate scalar intertwiner]
\label{p5:lem:threshold-intertwiner-before-higgs}
Inside the scale-free protected coordinate-intertwiner type, the generated family- and color-blind
off-diagonal module connecting the derived negative right line to the left
block is
\begin{equation}
\label{p5:eq:primitive-threshold-intertwiner}
\mathsf I\in\operatorname{Hom}(W_{R-},W_L)\otimes1_{\ell\oplus C}\otimes1_G.
\end{equation}
After hypercharge is fixed, it transforms as one left weak doublet, color
singlet and family blind, with \(Y=1/2\).  Its conjugate supplies the
\(W_{R+}\to W_L\) map.  This is the unique self-adjoint real coordinate
quadruple.  Other singlet, triplet, colored, or family-dependent modules may
exist in the full event envelope, but they are not generated coordinate
descendants.
\end{lemma}

\begin{proof}
The Hom module is the fundamental \(SU(2)_L\) representation because the
source is a weak singlet and the target is \(W_L\).  Identity action on \(Q,G\)
gives color and family blindness.  Its abelian weight is the target charge
minus the source charge, \(1/2\).  Complex conjugation gives the opposite
weight.  Theorem~\ref{p5:thm:generated-coordinate-higgs} proves the real
self-adjoint coordinate and multiplicity statements.
\end{proof}


\subsubsection{Microscopic-carrier and reciprocal-position readout convention}
\label{p5:sec:mp-mx-carrier-convention}
All finite gauge, charge, color, family, anomaly, and coefficient statements use the compact microscopic support
\begin{equation}
\label{p5:eq:mp-mx-convention}
\boxed{m_p=1.}
\end{equation}
Spacetime-local fields, running scales, thresholds, and comparison charts are reciprocal-position readouts at \(m_x\gg1\), where \(m_x\) labels the finite harmonic window \(\mathcal P_x(m_x)\) on the already smooth Born-derived position role.  The parameter \(m_x\) is not an additional gauge, family, internal carrier factor, microscopic site count, or second compact carrier.

\begin{remark}[Finite-carrier status of uniqueness claims]
\label{p5:rem:finite-carrier-uniqueness}
Unless explicitly stated otherwise, uniqueness claims in this paper are
relative to the generated carrier \(\mathsf G_5^{\rm fin}\).
GP3 removes kernels, spectators, or unmarked multiplicity only when doing so
preserves the generated incidence and grading tests.  It does not assert that
every operator in the full event envelope occurs physically.
\end{remark}

\begin{remark}[Upstream finite imports and logical order]
\label{p5:rem:upstream-finite-import-order}
On the actualized core-support role \(\mathsf{Supp}_{2T}\), Paper~I-A supplies
the finite-cell and trimer geometry and Paper~II records the neutral trimer
grammar.  Paper~II-B supplies the generated and co-realized \(m_p=1\)
momentum carrier; Paper~III supplies the diagonal-fibre isometry; Paper~IV
supplies quantum representations of generated typed roles.  Paper~V-A closes
the finite incidence and coordinate composition, computes the stabilizer and
grading alternatives, derives the even relative-chirality bridge, and closes structural hypercharge by parent-relation descent while exposing
only the separately typed field/member-occurrence boundary.
Equation~\eqref{p5:eq:omega5-envelope} records the analytic SSQFT/PCT
domains, while
\(\mathsf J_5^{\rm conn,der}\) is a distinct theorem output.
\end{remark}

\begin{theorem}[Upstream derived active-orbit handoff]
\label{p5:thm:derived-active-realization-handoff}
On an occurring Hurwitz-even mass-shell/contact branch, Papers~III-B--III-C
derive the equivariant orbit of nested active supports
\[
 \operatorname{rank}_{\mathbb R}P_{\rm act}:8\longrightarrow4
 \longrightarrow3\longrightarrow1,
 \qquad
 q=0,4,5,7,
 \qquad
 K_{\rm sh}=0,32,40,56.
\]
For each orbit index \(u\), their carrier ranges and additional finite action
data are, respectively,
\[
\begin{aligned}
 q=4:&\ A^{\rm even}\text{ with Hurwitz units }2T,\\
 q=5:&\ E_3(u)\text{ with }C_6(u)\curvearrowright E_3(u),\\
 q=7:&\ \mathbb R\mathbf1\text{ with the trivial action}.
\end{aligned}
\]
Thus \(C_6(u)\) is an action group, not the rank-three carrier.  More
precisely, Paper~III-C constructs the equivariant family of finite
self-adjoint active-depth operators \(\mathcal Q_{\rm act}(u)\) whose upper
spectral projectors have these ranks.  Their common terminal projector is
\[
 \mathbf1_{\{7\}}(\mathcal Q_{\rm act}(u))
 =P_0=\mathbb E_{\rm IR},
\]
the unique unital parity-even \(O(3)\)-equivariant scalar projector.
Consequently \((\mathrm{IR}_{O(3)})\) is inherited as a theorem-status
certification from the same G0-occurring branch, not supplied as a second
occurrence.
At \(q=5\), the six-point link permutation module for \(C_6(u)\) has
canonical projectors of ranks \(1,3,2\), hence the exact split grammar
\[
 \boxed{6=5+1,\qquad5=3+2}.
\]
These statements require neither a free \((\mathrm{ACT}_{\rm Cl})\) family
nor a seven-step \((\mathrm{FLAG})\).  All body-diagonal representatives are
conjugate.  The unmarked orbit, ranks, group types, and split grammar are
derived; a particular \(u\), cyclic embedding, or
\(\mathcal Q_{\rm act}(u)\) requires a mark and is not selected by G0.
\end{theorem}

\begin{corollary}[No native generated depth-six carrier]
\label{p5:cor:no-native-depth-six-carrier}
For every active-orbit member imported above,
\begin{equation}
 \label{p5:eq:no-native-depth-six-carrier}
 \boxed{\mathbf1_{\{6\}}\bigl(\mathcal Q_{\rm act}(u)\bigr)=0.}
\end{equation}
Indeed, Paper~III-C classifies all \(C_3\)-equivariant subprojectors of the
rank-three carrier.  Its rank-two rotating plane is orthogonal to the scalar
endpoint and therefore is not a nested upper-support stage.  Hence the
minimal generated Standard-Model carrier cannot acquire a native depth-six
sector by refining the active chain.  This is a relative generation theorem:
it does not prohibit adding a new external field/carrier with independent
ancestry and occurrence data.
\end{corollary}

\begin{proof}
Apply Theorem~\ref{p3c:thm:active-depth-six-obstruction} to the active orbit
co-realized in
Theorem~\ref{p5:thm:derived-active-realization-handoff}.  The scope clause is
the finite-carrier convention of
Remark~\ref{p5:rem:finite-carrier-uniqueness}.
\end{proof}

\begin{theorem}[Joint zero-defect matching of link and active orbits]
\label{p5:thm:joint-link-active-orbit-matching}
Let \(\mathscr L_{\rm link}\) be Paper~I-A's equivariant family of oriented
cube-link diagonals \(\delta\), cyclic stabilizer actions, and alternating
trimer pairs generated from \(\mathsf{Supp}_{2T}\).  Let
\(\mathscr C_{\rm act}\) be Paper~III-C's equivariant active-chain orbit.
The two families use the same canonical \(2O\)-normalizer diagonal orbit.  In
an oriented link frame write
\(\delta_\varepsilon=(\varepsilon,-\varepsilon)\),
\(\varepsilon=(\varepsilon_1,\varepsilon_2,\varepsilon_3)\in\{\pm1\}^3\),
and define
\[
 \jmath:\mathcal O_{\rm link}\longrightarrow\mathcal O_{\rm diag},
 \qquad
 \jmath(\delta_\varepsilon)
 :=\frac{\varepsilon_1e_{23}+\varepsilon_2e_{31}
               +\varepsilon_3e_{12}}{\sqrt3}.
\]
This is the orientation-induced Hodge identification; it is an equivariant
bijection and sends the link stabilizer lift to
\(C_6(\jmath(\delta))\).  In the formal product atlas
\(\mathcal O_{\rm link}\times\mathcal O_{\rm diag}\), the joint GP2
zero-realization-defect locus is precisely
\[
 \boxed{
 \mathfrak Z_{5}^{\rm diag}
 =\{(\delta,u):u=\jmath(\delta)\}.}
\]
Consequently the physical joint family is the graph of \(\jmath\), not the
Cartesian product of two independently markable diagonal families.  The graph
does not select an unmarked diagonal.  If a later oriented-link occurrence
marks \(\delta\), however, it generates the matching active member
\(\mathcal Q_{\rm act}(\jmath(\delta))\), its
\(C_6(\jmath(\delta))\)-action, and its rank-three support without a second
active-diagonal mark.
\end{theorem}

\begin{proof}
The bivectors \(e_{23},e_{31},e_{12}\) are an oriented orthonormal basis of
the imaginary quaternionic/even Clifford three-space.  Hence the displayed
formula has norm one and is a bijection on the eight oriented sign diagonals.
For every orientation-preserving cube rotation \(g\), Hodge equivariance gives
\(\jmath(g\delta)=\operatorname{Ad}_g\jmath(\delta)\).  Moreover
\[
 c_{\jmath(\delta_\varepsilon)}
 =\frac12\bigl(1+\varepsilon_1e_{23}+\varepsilon_2e_{31}
                    +\varepsilon_3e_{12}\bigr),
\]
so Paper~I-A's binary lift of the oriented link stabilizer is exactly the cyclic
group used by Paper~III-C.

Paper~I-A's generated diagonal-stabilizer/six-cycle theorem constructs, for
each common root \(\sigma=(H,v,\delta)\), the oriented cyclic stabilizer and
alternating trimer pair equivariantly from the Hurwitz core.  Paper~III-C's
active-orbit theorem constructs, from that same root through
\(u=\jmath(\delta)\), the cyclic action
\(C_6(u)\curvearrowright E_3(u)\), the projector \(\Pi_3(u)\), and the
active-depth operator.  Thus the typed common-root constructor
\[
 \mathfrak F_5(H,v,\delta)
 :=\bigl(\delta,\jmath(\delta)\bigr)
\]
has the displayed graph as its exact essential image.  A mismatched tuple is
outside that image: it either inserts a second diagonal
and cyclic action or fails that generated relation, giving a nonzero insertion
or relation component of the joint GP2 realization defect.  A graph tuple
retains exactly one diagonal, both licensed constructor images, and all their
action relations, so its joint defect is zero.

The constructor is host-natural, contains no parameter or representative
choice, conservatively retains the oriented-diagonal comparison, and its graph
with both licensed outputs is GP1-admissible with zero joint defect.  Hence
Paper~II-A's witness-transport and joint-unresolved-family exhaustion theorems
restrict the realized Paper--V-A branch to the displayed graph.  Projection of the graph to
\(\mathcal O_{\rm link}\) is an equivariant bijection with inverse
\(\delta\mapsto(\delta,\jmath(\delta))\), so the restriction propagates the
unmarked orbit rather than selecting a member.  Once \(\delta\) is marked by
an independently occurring oriented link, generated-role inheritance carries
that existing mark through the graph to the unique matching active member.
\end{proof}

\begin{remark}[What the active-realization theorem supplies to the incidence bridge]
Theorem~\ref{p5:thm:derived-active-realization-handoff} derives the
equivariant hexagonal action/support family and its two alternating rank-three
supports.  It does not alone distinguish the roles or select a member.
Theorem~\ref{p5:thm:joint-link-active-orbit-matching} proves that the already
oriented link token supplies the matching member without a second mark.
Combined with the actual axial unit,
Paper~I-A's orthogonality theorem identifies the opposite trimer with the
\(\operatorname{Im}\HH\) component and the complement with multiplicity.
The local joint-connection family and its same-process factor relations are
derived by Papers~III-F and V-A.  Spin--statistics is
derived on the analytic SSQFT domain.
\end{remark}

\paragraph{Finite-support incidence.}
On the actualized core-support role \(\mathsf{Supp}_{2T}\), Paper~I-A supplies
the Hurwitz-link
decomposition
\[
N(v)=\{u_+,u_-\}\sqcup T_+\sqcup T_-,
\qquad |T_+|=|T_-|=3,
\]
together with its Walsh--Clifford octet and binary-dihedral/Weyl data, while
Theorems~\ref{p5:thm:derived-active-realization-handoff} and
\ref{p5:thm:joint-link-active-orbit-matching} supply the \(q=5\) cyclic
action/support stage up to conjugacy and synchronize it with the same link
diagonal.
Paper~I-A's axial-frame theorem identifies the trimer orthogonal to the actual
axial unit with \(C\cong\operatorname{Im}\HH_\CC\); the complementary trimer
is \(G\).  The labels \(T_c,T_g\) record this generated incidence and may swap
under orientation reversal together with conjugation of the representative.

\begin{theorem}[Generated Hurwitz-link support and roles]
\label{p5:thm:hurwitz-link-support-import}
Assume the actualized core-support role \(\mathsf{Supp}_{2T}\), an occurring
\(q=5\) active orbit from
Theorem~\ref{p5:thm:derived-active-realization-handoff}, and the already fixed
oriented Hurwitz-link token used in this Paper--V-A branch.  If its diagonal is
\(\delta\), Theorem~\ref{p5:thm:joint-link-active-orbit-matching} supplies the
matching active member \(u=\jmath(\delta)\) without a second mark.  Paper~I-A then supplies
\[
N(v)=\{u_+,u_-\}\sqcup T_+\sqcup T_-,
\qquad |T_+|=|T_-|=3.
\]
The two trimers are distinct supports of equal rank.  Orthogonality to the
actual axial unit identifies one as \(T_c\), while link complement identifies
the other as \(T_g\), and
\begin{equation}
\label{p5:eq:two-trimer-supports}
C:=\mathbb C^{T_c},
\qquad
G:=\mathbb C^{T_g}.
\end{equation}
No preferred sitewise map \(T_c\to T_g\) is needed.  The distinction is
orientation-relative and reversal gives the conjugate equivalent
representative.
\end{theorem}

\begin{proof}
The Hurwitz-link theorem supplies the two alternating transverse three-cycles
inside every \(q=5\) representative.  The upstream active-realization theorem
fixes the conjugacy class and cyclic action type, while the joint-orbit theorem
propagates the already marked link diagonal to the active member.  G0 types
the occurrence but performs no representative selection.
Theorem~\ref{p5:thm:generated-color-family-incidence} then
uses quaternionic orthogonality and complement to prove the unequal roles.
Orientation reversal exchanges the presentation and conjugates the resulting
representation.
\end{proof}

\begin{theorem}[Binary-dihedral Weyl skeleton of the incidence trimer]
\label{p5:thm:binary-dihedral-color-skeleton}
Let \(T_c\) be the trimer generated as \(C=\CC^{T_c}\).  Paper~I-A proves
that the stabilizer of the oriented cube-link diagonal is \(2C_3\), while the
stabilizer of the same unoriented diagonal is its binary-dihedral completion
\(2D_3\), with
\[
1\longrightarrow2C_3\longrightarrow2D_3\longrightarrow C_2\longrightarrow1,
\qquad
2D_3/\{\pm1\}\cong D_3\cong S_3\cong W(A_2).
\]
On this trimer, the projective Weyl group permutes the three basis projectors
and the six ordered off-diagonal transitions.  Consequently
\[
\End_0(C)
=
\mathfrak h_{A_2}
\oplus
\bigoplus_{\alpha\in\Phi(A_2)}\mathfrak g_\alpha,
\qquad
\dim\mathfrak h_{A_2}+|\Phi(A_2)|=2+6=8.
\]
The group \(2D_3\) is a finite binary Weyl-normalizer skeleton.  It is not the
continuous gauge group.
\end{theorem}

\begin{proof}
The stabilizer and quotient statements are the binary-dihedral diagonal-lift
and \(A_2\)-Weyl theorems of Paper~I-A.  The three diagonal projectors have a
two-dimensional traceless span, identified with the \(A_2\) Cartan plane,
while the six matrix units \(E_{ij}\), \(i\ne j\), are the six root
directions.  The projective quotient \(D_3\cong S_3\) acts by permuting the
trimer labels and hence the root system.  This establishes the finite skeleton
before the continuous Lie closure below.
\end{proof}

\begin{theorem}[Exact Lie closure of the Paper--II trimer grammar]
\label{p5:thm:trimer-su3-lie-closure}
In the cyclic basis of the Paper--II-F trimer theorem, let
\[
D_1=\diag\!\left(\frac23,-\frac13,-\frac13\right),
\qquad
D_2=\diag(0,1,-1),
\]
\[
C_1=X+X^\dagger,\qquad C_2=-i(X-X^\dagger).
\]
Then
\[
\operatorname{Lie}_{\RR}
\langle iD_1,iD_2,iC_1,iC_2\rangle
=\mathfrak{su}(3).
\]
Indeed \(D_2\) is redundant:
\(\operatorname{Lie}_{\RR}\langle iD_1,iC_1,iC_2\rangle
=\mathfrak{su}(3)\).
This is the algebraic color closure because \(C\) is the generated
quaternionic imaginary component.  Its realization by a local connection
uses \(\mathsf J_5^{\rm conn,der}\).
\end{theorem}

\begin{proof}
Let \(\mathfrak l\) be the real Lie algebra generated by
\(iD_1,iC_1,iC_2\).  Every generator is traceless and skew-Hermitian, so
\(\mathfrak l\subseteq\mathfrak{su}(3)\).  Complexify.  Since
\[
C_1+iC_2=2X,\qquad C_1-iC_2=2X^\dagger,
\]
\(\mathfrak l_{\CC}\) contains \(D_1,X,X^\dagger\).  With indices modulo
three and \(Xe_j=e_{j+1}\),
\[
X=E_{10}+E_{21}+E_{02}.
\]
The three summands are eigenvectors of \(\operatorname{ad}_{D_1}\) with
distinct eigenvalues \(-1,0,1\), respectively.  Lagrange polynomials in
\(\operatorname{ad}_{D_1}\), implemented by finite linear combinations of
iterated commutators, therefore isolate \(E_{10},E_{21},E_{02}\).
Applying the same argument to \(X^\dagger\) gives the three opposite matrix
units.  Their commutators give
\[
[E_{ij},E_{ji}]=E_{ii}-E_{jj},
\]
so \(\mathfrak l_{\CC}=\mathfrak{sl}_3(\CC)\).  Hence
\(\dim_{\RR}\mathfrak l=8=\dim_{\RR}\mathfrak{su}(3)\), and the inclusion is
an equality.
\end{proof}

\begin{theorem}[Generated quark--lepton component and split stabilizer]
\label{p5:thm:color-selection}
Let
\[
Q=\ell\oplus C,
\qquad \dim_{\mathbb C}\ell=1,
\qquad \dim_{\mathbb C}C=3,
\qquad C=\mathbb C^{T_c}.
\]
On the generated incidence component, the trimer grammar closes to
\(\mathfrak{su}(3)\) on \(C\) by
Theorem~\ref{p5:thm:trimer-su3-lie-closure}.  The connected special-unitary
stabilizer of the generated split \(Q=\ell\oplus C\) is
\[
S(U(1)_\ell\times U(3)_C),
\]
with Lie algebra
\[
\mathfrak{su}(3)_c\oplus\mathfrak u(1)_{B-L}.
\]
The unique traceless block generator, with positive triplet normalization, is
\[
B-L=\frac13P_C-P_\ell.
\]
No \(SU(4)\) parent or off-diagonal quark--lepton gauge action follows from
this stabilizer theorem.
\end{theorem}

\begin{proof}
The Lie-closure theorem supplies the traceless nonabelian action on \(C\).
Preserving \(P_\ell,P_C\) gives block-diagonal unitaries
\(u_\ell\oplus u_C\); imposing unit total determinant gives
\(S(U(1)\times U(3))\).  Its Lie algebra is
\(\mathfrak{su}(3)\oplus\mathfrak u(1)\).  A block-scalar traceless generator
has the form \(aP_C+bP_\ell\) with \(3a+b=0\); choosing \(a=1/3\) gives the
displayed direction.
\end{proof}

\begin{remark}[Convergent color-triplet closures]
\label{p5:rem:color-forward-delta-crosslink}
The continuous \(A_2\) Lie algebra is now derived directly from the displayed
Paper--II trimer operations, not from full-matrix presentation covariance.
Its designation as color follows from the generated
\(\operatorname{Im}\HH\) incidence and \(A_2\) action.  Actual local gauge
realization is exactly the connection residual; no downstream running or
numerical fit is used.
\end{remark}


\begin{theorem}[Generated replication trimer]
\label{p5:thm:generation-selection}
Let \(T_g\) be the complementary trimer generated by
Theorem~\ref{p5:thm:generated-color-family-incidence}.  Its cyclic orientation
defines a free transitive
\(Z_3\)-action.  The resulting replication carrier is
\[
G:=\mathbb C^{T_g}\cong\mathbb C[Z_3]\cong\mathbb C^3,
\]
unique up to cyclic relabelling and unitary equivalence.  Its interpretation
as three replicated copies and the trivial gauge action follow from
complementarity and the GP2 provenance theorem.
\end{theorem}

\begin{proof}
A cyclic trimer is a free transitive \(Z_3\)-set, so its complex permutation
carrier is the regular representation.  The incidence theorem identifies it
as the untouched multiplicity support, while
Lemma~\ref{p5:lem:replication-not-gauge-support} derives identity extension
from operation provenance.
\end{proof}

\begin{corollary}[Finite Weyl closure of the family carrier]
\label{p5:cor:family-weyl-closure}
In the orbit basis of Theorem~\ref{p5:thm:generation-selection}, define the clock operator \(Z\) by
\[
Ze_j=\omega^j e_j,
\qquad
\omega=e^{2\pi i/3},
\]
while \(X\) is the cyclic shift.  Then
\[
X^3=Z^3=1,
\qquad
ZX=\omega XZ,
\]
and the nine operators
\[
X^aZ^b,
\qquad a,b\in\ZZ_3,
\]
form a trace-orthogonal basis of \(\End(G)\cong M_3(\CC)\).  Thus the generated
replication carrier has a closed finite rank-three operator calculus.
\end{corollary}

\begin{proof}
This is the rank-three Weyl algebra of Paper~I-A's trimer Weyl-algebra proposition, in the finite-quantum language of Paper~IV's rank-three Weyl-basis theorem.  It records the operator closure of the family carrier after Theorem~\ref{p5:thm:generation-selection}; it does not choose Yukawa matrices, mixing angles, or physical phases.  In particular, the existence of the closed calculus
\(\End(G)\cong M_3(\CC)\) does not generate a relation such as
\([Y_f,X]=0\), does not force \(Y_f\propto I\), and does not identify the
structural trimer basis with a physical mass eigenbasis.
\end{proof}

% BEGIN FRAMEWORK v3.029: GP3 MATTER OWNER BRIDGE
\begin{theorem}[Exact Clifford--color--family algebra and faithful amplification classes]
\label{p5:thm:v329-matter-factor-envelope}
Let $W=W_L\oplus W_{R+}\oplus W_{R-}\cong\CC^4$,
$Q=\ell\oplus C\cong\CC\oplus\CC^3$, and $G\cong\CC^3$.
On the specified independently factor-local generated grammar, include
(i) all four normalized quaternionic Clifford-solder directions on $W$,
(ii) the $\ell/C$ central incidence projections and the color
$\mathfrak{su}(3)$ matrices on $C$, and (iii) the family Weyl operators
$X,Z$ on $G$, each extended by the identity on the other two factors.
The generated unital complex star-algebra on the canonical carrier
$F=W\otimes Q\otimes G$ is
\[
 \mathcal B_{5A}^{\mathrm{alg}}
 =M_4(\CC)\otimes(\CC\oplus M_3(\CC))\otimes M_3(\CC)
 \cong M_{12}(\CC)\oplus M_{36}(\CC).
\]
In particular $\dim_\CC\mathcal B_{5A}^{\mathrm{alg}}=1440$ and
$\dim_\CC F=48$. Every finite faithful star-representation of this
abstract algebra is unitarily equivalent to
\[
 \mathcal H_{k_\ell,k_C}
 =(\CC^{12}\otimes M_\ell)\oplus
   (\CC^{36}\otimes M_C),
 \qquad k_\ell=\dim M_\ell\ge1,\quad k_C=\dim M_C\ge1.
\]
This theorem classifies a stated operator grammar; it does not establish
physical occurrence of every generator as a test, detector, process, or
scalar action port, nor a GP3 whole-candidate comparison.
\end{theorem}
\begin{proof}
Writing $\Gamma(H)=\bigl(\begin{smallmatrix}
 0&\Phi(H)\\ \Phi(H)^\dagger&0
\end{smallmatrix}\bigr)$, the four real quaternionic coordinate
directions obey the complex $\mathrm{Cl}_4$ relations. Their sixteen
ordered monomials span $\operatorname{End}(W)\cong M_4(\CC)$.
The central $\ell/C$ projections and the irreducible generated color
$\mathfrak{su}(3)$ produce $\CC\oplus M_3(\CC)$ on $Q$.
The nine Weyl monomials $X^aZ^b$ span $M_3(\CC)$ on $G$.
The stipulated independent factor-local identity extensions give the
full tensor product, whose two summands are the displayed blocks.
The faithful representation classification is the standard one for
a finite direct sum of full matrix algebras.
\end{proof}
% END FRAMEWORK v3.029


% BEGIN FRAMEWORK v3.029: GP3 MATTER OWNER BRIDGE
\begin{proposition}[Finite process descent in the represented matter algebra]
\label{p5:prop:v329-matter-process-descent}
Let $\mathcal B$ be a specified represented copy of
$M_{12}(\CC)\oplus M_{36}(\CC)$ on
$H=(\CC^{12}\otimes M_\ell)\oplus(\CC^{36}\otimes M_C)$.
Write $E_\mathcal B$ for the trace-preserving conditional expectation
obtained by partial traces over the two multiplicity spaces.
Let $\Phi^*:\operatorname{End}(H)\to\operatorname{End}(H)$ be a
\emph{specified} unital completely positive Heisenberg map and
$F_1,\ldots,F_{1440}$ any real Hermitian basis of $\mathcal B$.
For any $\tau>0$,
\[
 \sum_{j=1}^{1440}
 {\cal H}_\tau(\Phi^*(F_j),E_\mathcal B\Phi^*(F_j))=0
 \quad\Longleftrightarrow\quad \Phi^*(\mathcal B)\subseteq\mathcal B.
\]
This is a finite mathematical descent criterion; applying it to a
physical process requires a separately verified source-incidence and
action-owner certificate identifying the actual process generator.
\end{proposition}
\begin{proof}
The expectation is a trace-preserving Hilbert--Schmidt orthogonal
projection onto the represented algebra. A completely positive map
preserves Hermitian adjoints. The finite positive-time heat--Bregman
converse of Theorem~\ref{p2a:thm:v267-heat-bregman-converse}
shows that each nonnegative summand vanishes exactly when its
arguments agree. Vanishing on a real Hermitian basis is equivalent
to invariance of the complex linear span, namely $\mathcal B$.
\end{proof}
% END FRAMEWORK v3.029




\begin{remark}[Exact status of the three-generation statement]
\label{p5:rem:finite-carrier-three-generations}
Paper~I-A proves both rank three and, through the axial/opposite-trimer frame,
which trimer is consumed by the \(1+3\) component.  The remaining trimer is
therefore the generated multiplicity carrier.  A fourth copy requires
enlarging or replacing \(G\); it is not available in the minimal generated
incidence.  Thus ``three generations'' is exhaustive only relative to the
native generated carrier and its frozen incidence tests.  It is not a
theorem that every mathematically larger or externally adjoined matter model
is impossible.
\end{remark}



\begin{theorem}[Generated composition of color and replication supports]
\label{p5:thm:canonical-color-family}
The oriented Hurwitz link supplies the incidence and complementary
rank-three supports
\[
C=\mathbb C^{T_c},
\qquad
G=\mathbb C^{T_g}.
\]
Under GP2 relation-free composition and the GP2/GP3 multiplicity sandwich,
their simultaneous realization
is
\[
C\otimes G,
\]
as a special case of Theorem~\ref{p5:thm:component-independence-composition}.
\end{theorem}

\begin{proof}
Their independently generated commuting test families satisfy Paper~IV's
relation-free composition theorem.  GP2 no-omission and Paper~II-B,
Proposition~\ref{p2:prop:gp3-multiplicity-bound}, give exactly one unmarked
copy of every generated component, hence the displayed carrier.
\end{proof}

\begin{remark}[Triality octets are support geometry, not extra gauge sectors]
\label{p5:rem:triality-octets-support-not-gauge}
Paper~I-A's triality octets delimit the inherited finite support.  They neither create additional color groups nor identify octets with families.  The physical strong octet is the traceless summand in \(C\otimes C^\ast=\mathbf1\oplus\mathbf8\) on the color support fixed by Theorem~\ref{p5:thm:color-selection}.
\end{remark}



\begin{remark}[Trimer incidence and finite quantum realization]
\label{p5:rem:paper5-support-selections}
The maps \(T_c\mapsto Q\) and \(T_g\mapsto G\) are generated by axial
orthogonality and complement.  Paper~IV realizes their simultaneous carrier
as a tensor product from the generated independent test families.
\end{remark}

\subsection{Framework-native ingredients relevant to the Standard Model}
\label{p5:sec:part1-native}

\begin{table}[!htbp]
\centering
\caption[Upstream imports and Paper-V-A closure]{Upstream imports and Paper--V-A closure.}
\label{p5:tab:native-ingredients}
\scriptsize
\setlength{\tabcolsep}{3pt}
\renewcommand{\arraystretch}{1.2}
\begin{tabularx}{\textwidth}{>{\raggedright\arraybackslash}p{0.28\textwidth}
                                >{\raggedright\arraybackslash}X
                                >{\raggedright\arraybackslash}p{0.26\textwidth}}
\toprule
\textbf{Item} & \textbf{Upstream theorem content} & \textbf{Paper--V-A status} \\
\midrule
Left--right carrier &
Diagonal round-fibre
\(\Spin(4)_{\rm fib}\cong SU(2)^a_L\times SU(2)^a_R\) &
Equivariant identification generated; local connection uses
\(\mathsf J_5^{\rm conn,der}\) \\

Hopf data &
For the provenance-fixed representative, primitive Chern class \(|c_1|=1\) and a
right Cartan half-spin spectrum &
Actual oriented token fixes the Cartan/oriented eigenline representative; no physical base-fibre phase is generated \\

Protected half-spin block &
\(W=(\frac12,0)\oplus(0,\frac12)\) and its
\(\Spin(4)\)-equivariant commutant &
Internal representation theorem; \(P_\pm\) and the character kernel are
derived; physical relative chirality is inherited from diagonal-fibre ancestry,
and physical structural hypercharge is generated by
Proposition~\ref{p5:prop:v187-pcad-yfix} \\

Quaternionic vector-space split &
\(\mathbb H_\CC=\CC1\oplus\operatorname{Im}(\mathbb H)_\CC
\cong\mathbf1\oplus\mathbf3\) &
Axial/opposite-trimer orthogonality derives
\(Q=\ell\oplus C\) \\

Two cyclic trimers &
On the actualized \(\mathsf{Supp}_{2T}\) role, two rank-three supports with identical trimer
grammar &
Orthogonality/complement derives component/multiplicity incidence;
\(\mathfrak{su}(3)\) closes on \(C\) \\

Quantum composition &
Tensor representation from distinct generated commuting test families, GP2
relation-free composition, GP2 no-omission, and the GP3 multiplicity bound &
Tensor form and blockwise multiplicity one derived \\

Physical fermion carrier &
Paper~IV supplies the finite obstruction and CAR/Fock calculus after
spin--statistics &
Lorentz-Weyl correlation uses the derived even ancestry relation; statistics uses
\(\mathsf R_5^{\rm SSQFT}\), and conjugate-sector occurrence uses
\(\mathsf R_5^{\rm PCT}\) \\

Scalar role &
Round coordinate representation and multiplicity-one protected intertwiner type &
One complex doublet, its conjugate, and \(Q,G\) blindness derived without a coordinate scale \\

Finite real/spectral datum &
No full finite spectral-triple \(J_F\) or \(D_F\) is required &
Functorial \(\bar F\) availability derived; \(F\oplus\bar F\) is not a finite
occurrence theorem; Majorana operator remains optional in
Sec.~\ref{p5:subsec:ncg-assumption-discharge} \\

\bottomrule
\end{tabularx}
\TableNotes{The ledger separates finite theorem outputs from generated hypercharge,
generated chirality ancestry, separately typed member occurrences, the derived
joint-connection certificate, and the analytic
SSQFT/PCT domains.  No observed charge, mass, or mixing value is used.}

\end{table}

\subsubsection{The protected \texorpdfstring{\(m_p=1\)}{mp = 1} momentum carrier}
\label{p5:sec:part1-native-carrier}

On Paper~II-B's generated and co-realized \(m_p=1\) momentum carrier,
Papers~I-A--II supply the protected abstract module
\begin{equation}
\label{p5:eq:part1-K-split}
W=W_L\oplus W_R,
\qquad
W_R=W_{R+}\oplus W_{R-}.
\end{equation}
Paper~III independently identifies the diagonal-fibre
\(\Spin(4)_{\rm fib}\) isometry carrier.  Equivariance identifies \(W\) with
its generated internal half-spin representation.  Physical Lorentz chirality
is carried by \(S_L,S_R\) and is correlated with \(W_L,W_R\) by the generated
even graded ancestry relation, as in
\eqref{p5:eq:physical-one-particle-factorization}.  Thus no internal
half-spin label is silently identified with spacetime chirality.

\subsubsection{Quaternionic left/right symmetry and the provenance-fixed Hopf \texorpdfstring{\(U(1)\)}{U(1)}}
\label{p5:sec:part1-native-hopf}

Paper~I-A \cite{Kugelmann2026a} identifies \(S^3\) with the group of unit
quaternions and exhibits commuting left and right \(SU(2)\) actions.
Paper~II's Hopf theorem fixes the primitive Chern class \(|c_1|=1\).  The upstream provenance-fixed representative of the unmarked oriented Hopf
class supplies the embedded Cartan, handedness, and positive half-spin eigenline
up to conjugate equivalence.  These are coordinates of the inherited
geometric token, not an additional Paper--V-A axis selector.

\begin{remark}[Hopf half-unit convention]
\label{p5:rem:forced-hopf-half-unit-import}
After the right-factor and handedness conventions are fixed, \(T_R^3\) is
normalized with eigenvalues \(\pm\frac12\).  The candidate character-kernel
\emph{line} below is invariant under common rescaling and becomes the physical
structural line by Proposition~\ref{p5:prop:v187-pcad-yfix}.  Writing its generator with coefficient one on
this fixed Cartan is the inherited-coordinate convention
\(Y=T_R^3+\frac12(B-L)\); an equivalent rescaling accompanied by the inverse
coupling rescaling describes the same local \(U(1)\) action.  The primitive
integer charge lattice later fixes the compact parametrization up to
inversion.  Thus the theorem output is the projective line
\(\mathfrak u(1)_Y=\RR(T_R^3+\tfrac12(B-L))\); coefficient one on
\(T_R^3\) is a generator/coupling convention, not a further physical selector.
\end{remark}

\subsubsection{The quaternionic split \texorpdfstring{$\mathbb H_{\mathbb C}=\mathbf 1\oplus \mathbf 3$}{HC=1+3}}
\label{p5:sec:part1-native-qsplit}

As a complexified quaternionic algebra,
\begin{equation}
\label{p5:eq:part1-HC-split}
\mathbb H_{\mathbb C}
=
\mathbb C\,1 \oplus \operatorname{Im}(\mathbb H)_{\mathbb C}
\cong \mathbf 1\oplus \mathbf 3.
\end{equation}
Theorem~\ref{p5:thm:generated-color-family-incidence} identifies this
decomposition with \(Q=\ell\oplus C\) by the axial/opposite-trimer
orthogonality relation.

\subsubsection{Protected commutant and absent finite real structure}
\label{p5:sec:part1-native-JGamma}

Papers~I-A--II supply the protected grading/half-spin decomposition and its
equivariant commutant.  Conjugate orientation functorially supplies the
representation \(\bar F\) from \(F\), but does not prove simultaneous
occurrence of \(F\oplus\bar F\).  No full finite spectral-triple real
structure \(J_F\) or finite Dirac operator \(D_F\) is claimed.  A Majorana
operator is optional interaction data in the finite-bimodule interface.


\begin{corollary}[Protected Fierz commutant without completion selection]
\label{p5:cor:protected-fierz-abelian-line}
On Paper~II-B's co-realized \(m_p=1\) finite cell, the protected representation
\[
W=\left(\tfrac12,0\right)\oplus\left(0,\tfrac12\right)
\]
has
\[
\boxed{\End_{\Spin(4)}(W)=\CC 1\oplus \CC\Gamma_5.}
\]
This commutant statement requires no member of the raw coordinate-completion
family.  If the explicit normalized Clifford--Fierz basis is invoked, the
current Paper~II contract requires only the nondegenerate protected stratum
\((\mathrm{ND}_{\rm sp}):\varrho_{\rm sp}>0\), on which every compatible
nonzero completion yields the same represented protected even
\(\Cl(0,3)\) algebra and hence the usual basis
\[
\End(W)
=
\langle
1,\Gamma_a,\Gamma_{ab},\Gamma_5\Gamma_a,\Gamma_5
\rangle.
\]
Consequently the non-scalar \(\Spin(4)\)-equivariant line is
\(\CC\Gamma_5\), up to normalization and sign.  This is an internal commutant
statement.  It is neither a finite real structure nor the physical
hypercharge line; the full diagonal candidate plane and its parent-derived
stabilizer are computed below.
\end{corollary}

\begin{proof}
The two half-spin summands are irreducible and inequivalent, so Schur's lemma
gives a two-dimensional equivariant commutant generated by their projectors.
Rewriting them as \((1\pm\Gamma_5)/2\) gives
\(\CC1\oplus\CC\Gamma_5\).  The optional normalized Clifford--Fierz
realization is completion-independent on \(\mathrm{ND}_{\rm sp}\); it does
not introduce a selector.  No gauge interpretation follows from this
calculation.
\end{proof}

\begin{remark}[Oriented-Hopf ancestry of the right-handed split]
\label{p5:rem:oriented-hopf-marked-ancestry}
Paper~II-E, Theorem~\ref{p2e:thm:oriented-hopf-marked-commutants},
computes the protected pointwise commutant after retaining the oriented
Hopf token already inherited here:
\[
 \End_{SU(2)_L\times U(1)_R}(W)
 =\CC P_L\oplus\CC P_{R+}\oplus\CC P_{R-}\cong\CC^3.
\]
This is the ancestry of the distinguished right-handed projectors used
below.  It coexists with the unmarked
\(\End_{\Spin(4)}(W)=\CC1\oplus\CC\Gamma_5\) of
Corollary~\ref{p5:cor:protected-fierz-abelian-line}, because the invariance
groups differ.  With the algebraic handedness convention
\(T_R^3=L_3^-\), the upstream geometric left-fibre charge is
\(Q_{H,L}=-2T_R^3\): the \(Q_{H,L}=+1\) eigenline is therefore
\(W_{R-}\), not \(W_{R+}\).  Reversing the inherited Hopf orientation
reverses this dictionary.

Paper~II-E, Theorem~\ref{p2e:thm:marked-hopf-weighted-trace-locus},
gives the protected grading-weighted trace-zero locus
\(\RR1\cup\RR Q_{H,L}\).  This finite protected calculation is not a
selection of hypercharge.  Tensoring with the color/lepton carrier and
retaining the actual scalar/conjugate comparison gives the six-block
centralizer, its plane, and the pre-Y anomaly calculation in
Lemmas~\ref{p5:lem:full-diagonal-centralizer}
and~\ref{p5:lem:intrinsic-abelian-centralizer-plane} and
Proposition~\ref{p5:prop:pre-y-trace-anomaly-plane}.  Their candidate plane
and the parent-derived hypercharge line remain unchanged.

This is a pointwise symmetry restriction.  The midpoint connection retains
full restricted \(SU(2)\) holonomy by Paper~I-C,
Theorem~\ref{p1b:thm:homogeneous}; choosing the Hopf axis does not reduce
that connection to \(U(1)\) or make its magnetic eigenlines parallel.
Likewise geometric Hopf grading alone is not superselection: the coordinate
charge-step rule is Paper~II-D,
Proposition~\ref{p2d:prop:coordinate-hopf-charge-steps}, while the accessible
algebra and dynamical conditions belong to Paper~IV,
Corollary~\ref{p4:cor:hopf-reeb-superselection}.
\end{remark}

\subsubsection{The minimal scalar/intertwiner slot}
\label{p5:sec:part1-native-scalar}
\label{p5:sec:part1-higgs}

The full event envelope contains many scalar representation modules.  The
following theorem identifies the one generated by the scale-free coordinate
representation and its protected multiplicity-one intertwiner type.

\begin{theorem}[Coordinate-type-derived one-field Dirac intertwiner]
\label{p5:thm:minimal-scalar-left-right-intertwiner}
On the canonical protected half-spin block, the scale-free
\((\tfrac12,\tfrac12)\) coordinate type has one retained intertwiner in each
right-to-left channel.  Its real self-adjoint form is one field and its
conjugate.  Thus
\begin{equation}
\label{p5:eq:part1-Higgs-slot}
H:W_{R-}\longrightarrow W_L,
\qquad
H\sim(\mathbf1,\mathbf2)_{h}\otimes1_G .
\end{equation}
Its conjugate presentation is
\begin{equation}
\label{p5:eq:part1-Higgs-conj}
\widetilde H=i\sigma_2\bar H
\sim(\mathbf1,\mathbf2)_{-h}\otimes1_G .
\end{equation}
The distinguished-line character kernel gives the candidate value \(h=1/2\);
it is the physical scalar hypercharge by Proposition~\ref{p5:prop:v187-pcad-yfix}.
This realization is unique up to nonzero field rescaling and unitary
equivalence within the generated coordinate module.  The theorem does not
claim that arbitrary scalar
singlets, weak triplets, colored \(\operatorname{Hom}(\ell,C)\) modules,
family-dependent modules, or independently typed additional doublets.
\end{theorem}

\begin{proof}
The source is a weak singlet and the target is the fundamental weak doublet,
so the Hom module is a weak doublet.  The coordinate exists before \(Q,G\)
composition as a representation type, so GP2 extends it by their identities.  Pseudoreality represents
the second column by \(i\sigma_2\bar H\), and self-adjointness leaves four
real coordinates, exactly one complex doublet.  The character-kernel weight
follows from Theorem~\ref{p5:thm:exact-oriented-vector-hypercharge}, while its
physical hypercharge interpretation follows from
Corollary~\ref{p5:cor:yfix-hypercharge}.  The other listed
modules exist algebraically but are not generated coordinate descendants.
No concrete reduced matrix element or coordinate normalization is used.
\end{proof}

\subsubsection{Native order-three structure and multiplicity-three geometry}
\label{p5:sec:part1-native-z3}

Paper~I-A \cite{Kugelmann2026a} contains both:
\begin{enumerate}[label=(\alph*), leftmargin=2em]
\item an order-three action on the quaternionic/projective line, and
\item the full multiplicity-three kinematic machinery
\[
M\cong \mathbb C^3,\qquad X^3=1,\qquad F_3,\qquad 3=1\oplus 2.
\]
\end{enumerate}
With \(T_g\) the complement of the incidence trimer,
Theorem~\ref{p5:thm:generation-selection} derives the regular trimer
representation as the family space,
\[
G=\mathbb C^{T_g}\cong \mathbb C[Z_3]\cong \mathbb C^3.
\]

\label{p5:part:anchor-free-charge}

\paragraph{Orientation.}
This part reconstructs the visible finite representation from
\(\mathsf G_5^{\rm fin}\) together with the explicitly displayed members of
the generated parent-relation hypercharge theorem.  It uses no observed charge value.  The candidate
plane and character kernel are finite outputs; physical structural hypercharge is generated by
Proposition~\ref{p5:prop:v187-pcad-yfix}, while chirality is inherited from
\eqref{p5:eq:chi5-derived}.

\subsection{Structural operators and their action}
\label{p5:sec:part1-action-table}

\subsubsection{Operator ledger}

The action ledger records the finite operators after the geometric bridge
theorems and charge calculation.  It distinguishes generated operators,
optional interactions, the separately typed member-occurrence boundary,
generated chirality ancestry, the derived joint-connection certificate, and analytic
theorem domains.
\begin{table}[H]
\centering
\caption[Structural Standard-Model operators]{Structural Standard-Model operators.}
\label{p5:tab:action-summary}
\renewcommand{\arraystretch}{1.2}
\begin{tabularx}{\textwidth}{>{\raggedright\arraybackslash}p{0.16\textwidth}
                                >{\raggedright\arraybackslash}p{0.22\textwidth}
                                >{\raggedright\arraybackslash}X
                                >{\raggedright\arraybackslash}p{0.22\textwidth}}
\toprule
\textbf{Object} & \textbf{Acts on} & \textbf{Action} & \textbf{Commutes with} \\
\midrule
\(SU(2)_L\), \(T_L^a\) &
\(W_L\) &
left weak action; trivial on \(\ell\oplus C\) and \(G\) &
\(P_\ell,P_C,X\) \\

\(T_R^3\) &
\(W_R\) &
derived right half-unit; trivial on \(\ell\oplus C\) and \(G\) &
\(P_\ell,P_C,X\) \\

\(SU(3)_c\), \(\lambda_A\) &
\(C\) &
color action on the core triplet; trivial on \(W\) and \(G\) &
\(T_L^a,T_R^3,X,P_\ell\) \\

\(X\) &
\(G\) &
cyclic family shift \(Xe_i=e_{i+1}\) &
all gauge/core operators \\

\(P_\ell,P_C\) &
\(\ell\oplus C\) &
project onto lepton singlet and color triplet core sectors &
\(T_L^a,T_R^3,X\) \\

\(B-L\) &
\(\ell\oplus C\) &
\(\frac13\) on \(C\), \(-1\) on \(\ell\); trivial on \(W,G\) &
\(T_L^a,T_R^3,X\) \\

\(Y\) &
\(W\otimes(\ell\oplus C)\) &
\(T_R^3+\frac12(B-L)\); trivial on \(G\) &
\(X\) \\

\(Q\) &
\(W\otimes(\ell\oplus C)\) &
\(T_L^3+Y\); trivial on \(G\) &
\(X\) \\

\(H\) &
\(W_R^{-}\to W_L\) &
generated coordinate weak-doublet intertwiner; family blind by GP2
provenance &
\(SU(3)_c,X\) \\

\(Y_u,Y_d,Y_e,Y_\nu\) &
\(G\) &
family endomorphisms; gauge singlets &
all gauge/core operators \\

\(D_M\) &
\(P_\ell P_{R+}\otimes G\) &
optional Majorana operator on the neutral right-lepton support &
\(\begin{gathered}SU(3)_c,\,SU(2)_L,\\ U(1)_Y\end{gathered}\) \\
\bottomrule
\end{tabularx}
\TableNotes{The ledger records the action generated on
\(\mathsf G_5^{\rm fin}\), with \(Y,Q\), the charged matter names, and the
neutral Majorana support typed relative to the generated hypercharge line, with visible
chirality derived by \eqref{p5:eq:chi5-derived}.  Local connection components use
\(\mathsf J_5^{\rm conn,der}\); a Majorana operator is optional downstream data.}

\end{table}

\subsection{visible carrier internal-response realization}
\label{p5:sec:v3102-visible-internal-response}

\begin{corollary}[Visible matter as an internal same-action response carrier]
\label{p5:cor:v3102-visible-internal-response}
The already derived visible one-particle carrier
\[
 \mathcal H_{\rm 1p}^{\rm matt}
 =(S_L\otimes F_L)\oplus(S_R\otimes F_R)
\]
may be consumed as $J_P$ in any independently licensed same-action physical
boundary split.  On the regular branch its quadratic matter response is then
\[
 \boxed{
 K_{\rm matt}
 =\operatorname{Schur}_{J_{\rm rest}}D^2S_\star
 \big|_{\mathcal H_{\rm 1p}^{\rm matt}}.}
\]
Accordingly the physical visible particle poles are poles of an internal
same-action boundary response on the already constituted particle carrier.
This does not identify an arbitrary structural boundary-energy packet or
complete boundary stress with visible matter energy; that stronger statement
requires a generated common-owner relation.
\end{corollary}
% END VISIBLE-INTERNAL-RESPONSE-REALIZATION v3.102

% BEGIN STRUCTURAL-BL2-GEOMETRIC-COEFFICIENT v3.103
\subsection{Consequence of the native graded matter vacuum ledger}
\label{p5:sec:v2216-matter-vacuum-ledger}

\begin{corollary}[Generated matter ranks do not cancel the vacuum coefficient]
\label{p5:cor:v2216-matter-vacuum-ledger}
On the unbroken, zero-fermion-background, parity-even smooth-parent
quadratic branch, the already generated $48$ internal Weyl coordinates,
four real Higgs coordinates and $12$ gauge generators, together with the
actual quotient Jacobian, give
\[
 \boxed{a_0^{\rm weighted}=-34,\qquad
 a_{2,\rm geom}^{\rm weighted}=0,\qquad
 a_{2,\rm Higgs\ mass}^{\rm weighted}=4\Lambda_g^2,}
\]
with the common $(4\pi)^{-2}\int\sqrt g$ factor suppressed.
The last term retains the actual Higgs mass squared $-2\Lambda_g^2$.
Thus the cancellation of the geometric Einstein coefficient is not a
vanishing volume coefficient or a solution of the physical vacuum
matching problem.
\end{corollary}
\begin{proof}
Import Paper~IV, Theorem~\ref{p4:thm:v2216-graded-matter-uv}, on the
same admitted owner. Its ranks are fixed here by component composition,
Theorem~\ref{p5:thm:component-independence-composition}, the coordinate
Higgs, Theorem~\ref{p5:thm:generated-coordinate-higgs}, and the least
self-adjoint heat carrier,
Theorem~\ref{p5:thm:least-self-adjoint-heat-carrier}. The real Higgs,
gauge, quotient and Weyl determinant weights are respectively
$1/2,1/2,-1,-1/4$ on ranks $4,48,12,192$, as in the existing
one-loop bridge, Theorem~\ref{p5:thm:part4-standard-one-loop-bridge}.
Their rank contributions sum to $2+24-12-48=-34$; their geometric
$R$ contributions sum to $0-2R-2R+4R=0$.
The conformal Higgs mass contribution is $4\Lambda_g^2$.
No physical mirror, conjugate or PCT family is added by the analytic
Clifford completion or by this import.
\end{proof}

\begin{remark}[Scope of the matter ledger]
\label{p5:rem:v2216-matter-vacuum-frontier}
The local UV coefficient calculation retains its common proper-time
domain. The tachyonic unbroken Higgs Gaussian needs an actual stable
member or an adequate high-scale regulator; it is not a positive
unregulated infrared determinant. Actual masses, thresholds and global
chiral phase, the field measure, the metric/scalaron gauge and ghost
blocks, and all interacting mixed terms retain their own contributions.
Paper~IV, Remark~\ref{p4:rem:v2216-graded-matter-scope}, states the
optional common-cutoff $41/2$ diagnostic and its uncontrolled mass
omission explicitly. It is not a physical vacuum prediction. The
bare/effective-action and factor-$96$ owner distinctions are unchanged;
neither a fitted subtraction nor a complete $\Gamma$ follows from this
ledger. Paper~IV,
Proposition~\ref{p4:prop:v2216-finite-heat-integrability}, additionally
requires the actual configuration domain and measure before applying its
exact finite effective-action flow.
\end{remark}
% END NATIVE-MATTER-VACUUM-LEDGER v2.216

% BEGIN FRAMEWORK v2.226: KINEMATIC CLOSURES
\subsection{Coordinate ancestry of the protected matter operators}
\label{p5:sec:v226-protected-matter}
\begin{corollary}[The protected chiral split needs no separate operator admission]
\label{p5:cor:v226-protected-matter}
On the intrinsic momentum fibre the protected chirality and its two
projectors are the degree-six coordinate words of
Theorem~\ref{p1:thm:v226-chirality}. Any already licensed matter embedding
or odd solder transports these equal operators into its declared associated
owner. The generated stabilizer and positive kinetic kernel therefore retain
their earlier operator ancestry without requiring a new scalar vacuum token.
\end{corollary}
\begin{proof}
Use the exact operator identities before the existing embedding, then
transport every action and source leg by that same map. This replaces an
operator presentation, not the physical matter embedding or solder itself.
\end{proof}
The native even algebra is $M_{11}\oplus M_4$; an odd solder can belong to
a larger graded field envelope. The native 137 coefficient tests cannot
stand in for an uncomputed odd leg, the complete matter representation,
charge normalization or a regulator-dependent anomaly calculation.
Optimal-radius localization also cannot by itself read protected chiral
populations (Theorem~\ref{p2c:thm:v226-chirality-radius}). These are domain
statements, not additional symmetry-breaking or vacuum requirements on the
already generated stabilizer.

% END FRAMEWORK v2.226: KINEMATIC CLOSURES

% BEGIN FRAMEWORK v2.227: DOWNWARD CLOSURES
\subsection{Canonical symmetry without an additional vacuum selector}
\label{p5a:sec:v227-canonical-symmetry}
For an actual finite Hamiltonian with the claimed unitary symmetry, the
normalized full-ground density and Gibbs density already have that
symmetry by Theorem~\ref{p1b:thm:v227-state-symmetry}. A finite invariant
state is therefore not an independent selector on these branches.
This applies to the supplied regulated matter/source theory; it does not
choose a pure orientation in a degenerate vacuum or prove a continuum
measure, Ward limit, or regulated chiral anomaly cancellation.

The symmetry commutant remains larger than the observable centre.
In particular the protected chirality is not central in $M_4$, since it
anticommutes with the Clifford coordinates. Only actual reducing central
block scalars have the neutral Liouvillian invariance of
Theorem~\ref{p1b:thm:v227-central-readouts}. Actual source blocks linking
the sectors read their relative energies and retain those coefficients.

% END FRAMEWORK v2.227: DOWNWARD CLOSURES


% BEGIN FRAMEWORK v2.229: ACTUAL CHARGED SOURCE SCOPE
\subsection{Derived structural identifications and full carrier}
\label{p5:sec:part1-selected}

\paragraph{No numerical charge anchoring.}
The construction below does not use observed Standard-Model hypercharges such as
\(Y(Q_L)=1/6\), \(Y(u_R)=2/3\), or observed electric charges as inputs.  The
abelian line is fixed from the generated geometry by:
\begin{enumerate}[label=(\roman*), leftmargin=2em]
\item the convention-fixed right half-unit \(T_R^3\),
\item the traceless quark--lepton splitter \(B-L\) on the \(\mathbf 3\oplus\mathbf 1\) block,
\item the distinguished one-dimensional module \(L_0=W_{R+}\otimes\ell\),
\item the generated coordinate doublet and its conjugate, and
\item the GP2-derived identity action on \(G\).
\end{enumerate}
The numerical charge list is therefore an output of the oriented carrier,
generated parent relation, distinguished-line character kernel, and the inherited
normalization convention.

\begin{table}[H]
\centering
\caption[Derived Standard-Model identifications]{Derived Standard--Model identifications.}
\label{p5:tab:selected-identifications}
\renewcommand{\arraystretch}{1.2}
\begin{tabularx}{\textwidth}{>{\raggedright\arraybackslash}p{0.18\textwidth}
                                >{\raggedright\arraybackslash}X
                                >{\raggedright\arraybackslash}p{0.28\textwidth}}
\toprule
\textbf{Structural item} & \textbf{Framework object} & \textbf{Status and proof location} \\
\midrule
Weak group \(SU(2)_L\) &
generated diagonal-fibre action on \(W_L\), coherently paired with Lorentz chirality &
Subsec.~\ref{p5:sec:part1-native-hopf};
Thm.~\ref{p5:thm:gp2-forced-visible-sector-admissibility};
Thm.~\ref{p5:thm:part1-minimal-sector-algebra} \\

Right half-unit \(T_R^3\) &
Hopf-fixed \(U(1)\subset SU(2)_R\) half-unit &
Subsec.~\ref{p5:sec:part1-native-hopf};
Rem.~\ref{p5:rem:forced-hopf-half-unit-import};
Eq.~\eqref{p5:eq:part1-TR3-full} \\

Color triplet \(C\) &
generated \(\mathbf 3\) block in \(\mathbb H_{\mathbb C}\), with exact
trimer Lie closure; local realization uses \(\mathsf J_5^{\rm conn,der}\) &
Thm.~\ref{p5:thm:color-selection};
Subsec.~\ref{p5:sec:part1-color-lepton};
Thm.~\ref{p5:thm:part1-minimal-sector-algebra} \\

Lepton singlet \(\ell\) &
complementary \(\mathbf 1\) block in \(\mathbb H_{\mathbb C}\) &
Subsec.~\ref{p5:sec:part1-native-qsplit};
Thm.~\ref{p5:thm:color-selection};
Subsec.~\ref{p5:sec:part1-color-lepton} \\

Family space \(G\) &
regular representation \(\mathbb C[Z_3]\cong \mathbb C^3\) &
Thm.~\ref{p5:thm:generation-selection};
Subsec.~\ref{p5:sec:part1-family-space};
Prop.~\ref{p5:prop:core-family-separation} \\

Family generator \(X\) &
regular \(Z_3\) shift on \(G\) &
Subsec.~\ref{p5:sec:part1-native-z3};
Thm.~\ref{p5:thm:generation-selection};
Eq.~\eqref{p5:eq:part1-X} \\

\(B-L\) operator &
\(
\frac13 P_C-P_\ell
\)
&
Subsec.~\ref{p5:sec:part1-BminusL};
Lem.~\ref{p5:lem:part1-BminusL-unique};
Lem.~\ref{p5:lem:intrinsic-abelian-centralizer-plane} \\

Hypercharge \(Y\) &
\(
T_R^3+\frac12(B-L)
\)
&
Generated by Proposition~\ref{p5:prop:v187-pcad-yfix};
Subsec.~\ref{p5:sec:part1-hypercharge};
Lem.~\ref{p5:lem:geometric-majorana-line};
Lem.~\ref{p5:lem:hypercharge-normalization-unique};
Thm.~\ref{p5:thm:part1-anchor-free-hypercharge} \\

Electric charge \(Q\) &
\(
T_L^3+Y
\)
&
Subsec.~\ref{p5:sec:part1-electric-charge};
Eq.~\eqref{p5:eq:part1-Q};
Thm.~\ref{p5:thm:part1-em-stabilizer} \\

Component algebra and physical gauge action &
\(\mathcal A_F^{\rm comp}
=(\mathbb H_w\oplus\mathbb C_R)\widehat\otimes_{\mathbb R}
(\mathbb C_\ell\oplus M_3(\mathbb C)_C)\); its derivation algebra is finite
constitution and its local realization uses \(\mathsf J_5^{\rm conn,der}\) &
Subsec.~\ref{p5:sec:part1-finite-algebra};
Eq.~\eqref{p5:eq:part1-finite-algebra};
Thm.~\ref{p5:thm:part1-minimal-sector-algebra} \\

Higgs doublet &
one-field left--right intertwiner; reconstructed \(Y=\frac12\) after the
hypercharge theorem &
Subsec.~\ref{p5:sec:part1-native-scalar};
Subsec.~\ref{p5:sec:part1-finite-algebra};
Eq.~\eqref{p5:eq:part1-Higgs-slot};
Cor.~\ref{p5:cor:scalar-hypercharge-after-reconstruction};
Prop.~\ref{p5:cor:finite-core-ewsb-topology} \\

Matter sectors &
\(Q_L,L_L,u_R,d_R,e_R,\nu_R\) on the full carrier &
Uses the generated structural hypercharge line;
Subsec.~\ref{p5:sec:part1-full-carrier};
Cor.~\ref{p5:cor:part1-matter-names-after-Y};
Subsec.~\ref{p5:sec:part1-representation-content};
Prop.~\ref{p5:prop:three-generation-content} \\
\bottomrule
\end{tabularx}
\TableNotes{Each row is a finite theorem output on the domain stated in the
text.  The charge values are not used as inputs.}

\end{table}

\subsubsection{Derivation of color and lepton sectors}
\label{p5:sec:part1-color-lepton}

By Theorem~\ref{p5:thm:color-selection}, the triplet block in
\(
\mathbb H_{\mathbb C}\cong \mathbf 1\oplus \mathbf 3
\)
is identified with the color triplet representation:
\begin{equation}
\label{p5:eq:part1-color-split}
\ell\oplus C
\cong
\mathbf 1\oplus \mathbf 3.
\end{equation}
We denote the corresponding orthogonal projectors by
\begin{equation}
P_\ell^2=P_\ell,\qquad P_C^2=P_C,\qquad P_\ell P_C=0,\qquad P_\ell+P_C=1.
\end{equation}
This is exactly the type of internal sector decomposition formalized abstractly in Paper~IV \cite{Kugelmann2026e}, where orthogonal projectors define legitimate quantum subspaces and block-supported operator actions.

\subsubsection{The full structural carrier}
\label{p5:sec:part1-full-carrier}

The one-copy structural core matter module is
\begin{equation}
\label{p5:eq:part1-finite-core-carrier}
F_{\rm core}
=
(W_L\oplus W_{R+}\oplus W_{R-})\otimes(\ell\oplus C).
\end{equation}
Under GP2 relation-free composition and the GP2/GP3 multiplicity sandwich,
the internal carrier is
\begin{equation}
\label{p5:eq:part1-full-carrier}
F=F_{\rm core}\otimes G
=
(W_L\oplus W_{R+}\oplus W_{R-})\otimes(\ell\oplus C)\otimes G.
\end{equation}
By the derived even ancestry relation, the physical Lorentz-chiral matter carrier is the
separate direct sum
\[
\mathcal H_{\rm 1p}^{\rm matt}
=(S_L\otimes F_L)\oplus(S_R\otimes F_R)
\]
defined in Theorem~\ref{p5:thm:component-independence-composition}.

\begin{proposition}[Core--family factor separation]
\label{p5:prop:core-family-separation}
By the generation-blindness theorem, all core gauge generators act
trivially on \(G\), while the family generator \(X\) acts trivially on
\(
(W_L\oplus W_R)\otimes(\ell\oplus C).
\)
Hence
\begin{equation}
\label{p5:eq:part1-core-family-comm}
[\rho_{\rm gauge}(a),X]=0
\qquad
\text{for every gauge generator }a.
\end{equation}
\end{proposition}

\begin{proof}
Lemma~\ref{p5:lem:replication-not-gauge-support} derives the factorized form
\[
\rho(a)=\rho_{\rm core}(a)\otimes 1_G,
\]
while the family generator acts as
\[
X=1_W\otimes 1_{\ell\oplus C}\otimes X_G.
\]
The two displayed tensor actions commute identically.  This proposition does not assert that an extra family connection occurs.
\end{proof}

\subsection{Derived structural Standard--Model outputs}
\label{p5:sec:part1-outputs}

\subsubsection{Full representation content}
\label{p5:sec:part1-representation-content}

\begin{proposition}[Matter content on the generated three-generation carrier]
\label{p5:prop:three-generation-content}
Assume
\(\mathsf G_5^{\rm fin}\).  The full carrier
\(F=F_{\rm core}\otimes G\) has matter sectors
\begin{align}
Q_L &\sim (\mathbf 3,\mathbf 2)_{1/6}\otimes G, &
L_L &\sim (\mathbf 1,\mathbf 2)_{-1/2}\otimes G, \\
u_R &\sim (\mathbf 3,\mathbf 1)_{2/3}\otimes G, &
d_R &\sim (\mathbf 3,\mathbf 1)_{-1/3}\otimes G, \\
\nu_R &\sim (\mathbf 1,\mathbf 1)_{0}\otimes G, &
e_R &\sim (\mathbf 1,\mathbf 1)_{-1}\otimes G.
\end{align}
Thus the generated rank-three replication carrier supports three
gauge-equivalent copies of one complete Standard--Model generation with a
right-handed neutrino.
\end{proposition}

\begin{proof}
For one core copy, define
\begin{align}
Q_L&:=P_LP_C F_{\rm core}, & L_L&:=P_LP_\ell F_{\rm core}, \\
u_R&:=P_{R+}P_C F_{\rm core}, & d_R&:=P_{R-}P_C F_{\rm core}, \\
\nu_R&:=P_{R+}P_\ell F_{\rm core}, & e_R&:=P_{R-}P_\ell F_{\rm core}.
\end{align}
The hypercharge values follow from \eqref{p5:eq:part1-Y}, and the family extension is obtained by tensoring each sector with \(G\).
\end{proof}

\subsubsection{Force sectors}
\label{p5:sec:part1-forces}

\paragraph{Strong sector.}
\label{p5:sec:part1-strong}
The reconstructed \(SU(3)_c\) action is the traceless unitary action on the
incidence trimer \(C\).  Its pair space decomposes as
\begin{equation}
\label{p5:eq:part1-color-pair-space}
C\otimes C^\ast\cong\mathbf3\otimes\bar{\mathbf3}=\mathbf1\oplus\mathbf8,
\end{equation}
so the active gauge directions form the adjoint octet.  In particular,
\begin{equation}
\label{p5:eq:part1-color-89}
\dim(C\otimes C^\ast)=9,
\qquad
\dim\mathfrak{su}(3)_c=8.
\end{equation}

\paragraph{Electroweak sectors.}
\label{p5:sec:part1-weak}
\label{p5:sec:part1-hypercharge-force}
\label{p5:sec:part1-em}
The generated weak action is on \(W_L\).  The parent-relation PCAD theorem
reconstructs the structural hypercharge direction
\(Y=T_R^3+\tfrac12(B-L)\).  If a nonzero neutral scalar direction is
supplied downstream, its candidate stabilizer is
\(Q=T_L^3+Y\).  Their reciprocal-position gauge fields and continuum
dynamics use \(\mathsf J_5^{\rm conn,der}+\mathsf R_5^{\rm SSQFT}\).


\section{Matter roles, conjugation and rooted family geometry}
\label{p5a:sec:organized-family}
\subsection{Native source-operation closure criterion}
\label{p5:sec:part1-scope}

\subsubsection{Generated finite package and residual status}
The finite supports below are inherited on the theorem domains stated in
Papers~I-A--IV and closed by the bridge proofs in this paper.  To prevent a
remaining occurrence from being hidden inside the word ``derived'', we state
the entire Paper--V-A boundary first.

\begin{definition}[Typed finite source profile]
\label{p5:def:primitive-finite-source-operator}
Let \(\mathfrak F_{\le4}\) be fixed.  A \emph{typed finite source profile} is
a generated carrier together with the role determined by its incidence,
grading, stabilizer, or functorial action.  GP2 preserves the prior geometry
under relation-free composition and GP3 removes kernels or unmarked
multiplicity.  GP2 or GP3 in isolation asserts neither a nonzero local
connection nor a Lorentzian QFT realization.  The former is nevertheless
derived when the explicit nonflat Paper-III connection is fed through the
licensed structure-group, associated-bundle, and operator constructors below.
\end{definition}

\begin{definition}[Paper--V-A finite package and residual envelope]
\label{p5:def:source-operation-envelope}
The notation \(\mathsf G_5^{\rm fin}\) records the following generated
theorem outputs; it is not a conjunction of assumptions.
\begin{enumerate}[label=(\roman*), leftmargin=2.2em]
\item \textbf{Oriented incidence.}  Paper~I-A's axial/opposite-trimer frame
      theorem generates
      \(Q=\ell\oplus C\cong\CC1\oplus\operatorname{Im}(\HH)_\CC\) and leaves
      the complementary trimer \(G\cong\CC^3\).  Orientation reversal gives
      the conjugate equivalent representative, not an alternative assignment.
\item \textbf{Left--right and tensor carrier.}  The protected half-spin block
      gives \(W=W_L\oplus W_R\), with
      \(W_R=W_{R+}\oplus W_{R-}\) fixed by the oriented Hopf Cartan.
      Paper~IV's relation-free tensor-product theorem gives
      \(F=W\otimes Q\otimes G\).  On each generated component block, GP2
      no-omission requires multiplicity at least one, while Paper~II-B's GP3
      unmarked-multiplicity theorem requires multiplicity at most one.
      Hence every displayed simple block occurs exactly once; no ambient full
      endomorphism algebra is promoted to a generated test algebra.
\item \textbf{Finite gauge candidates.}  The coordinate action on \(W_L\)
      and trimer inner-derivation closure on \(C\) generate
      \(\mathfrak{su}(2)_L\oplus\mathfrak{su}(3)_c\).  The generation-blind
      Abelian centralizer is the plane
      \(\operatorname{span}_{\mathbb R}\{T_R^3,B-L\}\), and the character of
      \(L_0=W_{R+}\otimes\ell\) has the exact kernel
      \(\mathbb R(T_R^3+\tfrac12(B-L))\).  The generated parent quadratic has
      a nonzero projected component in \((L_0^*)^{\otimes2}\) with weight
      \((-1,+2)\); Proposition~\ref{p5:prop:v187-pcad-yfix} below applies PCAD
      to that parent relation and identifies this kernel with the physical
      structural Abelian subgroup.
\item \textbf{Scalar type and finite symmetries.}  The scale-free coordinate type and
      its multiplicity-one protected intertwiner space form one complex weak
      doublet \(H\) and its pseudoreal conjugate \(\widetilde H\), acting as
      the identity on \(Q\otimes G\).  The separate external and internal
      gradings give two equal-rank invariant projectors
      \(P_\pm=\tfrac12(1\pm\gamma_x\otimes\Gamma_W)\).  Functorial complex
      conjugation makes \(\bar F\) available from \(F\), but neither a choice
      between \(P_+\) and \(P_-\) nor simultaneous occurrence of
      \(F\oplus\bar F\) follows from that fact.
\item \textbf{Typed kinetic theorem on the generated joint family.}  For the
      derived connection on the generated joint bundle, Paper~III-F proves a common ordinary
      carrier trace for the real, derivative-free, parity-even, positive,
      quadratic, uninserted, single-trace kinetic functionals.  A theta term,
      a nonscalar central insertion, or a functional outside this domain is
      not covered by that theorem.
\end{enumerate}

Let \(\mathbb T_5\) be the connected two-torus integrating the Abelian
candidate plane.

\begin{proposition}[Structural hypercharge from the generated parent quadratic]
\label{p5:prop:v187-pcad-yfix}
The joint Hopf--axial constructor already generates a nowhere-zero equivariant
quadratic
\[
 m_{\mathcal O}\in
 \Gamma_{\rm eq}(\mathcal O,N_{\mathcal O}^{*\otimes2}),
 \qquad \|m_{\mathcal O}\|^2=\frac12.
\]
The commuting Hopf/color projection of the complete geometric comparison has
a nonzero component in \((L_0^*)^{\otimes2}\).  Under the commuting Abelian
action its weight in the \((T_R^3,B-L)\) plane is \((-1,+2)\).  By Parent
Correspondence and Action Descent this component is the descendant of the
already generated parent relation, not a newly selected late pointing.  Its
connected stabilizer satisfies
\[
 -a+2b=0,
\]
and hence
\[
 \boxed{Y=T_R^3+\frac12(B-L).}
\]
Thus the structural G13/hypercharge obligation is derived on
\(\mathsf G_5^{\rm fin}\).  The theorem does not generate a Majorana mass, a
Higgs vacuum, a spacetime background field, a numerical Yukawa member, or the
G0 occurrence of any particular dynamical configuration.
\end{proposition}
\begin{proof}
The parent quadratic and the commuting projection are generated upstream in
the common rooted comparison geometry.  The projected component is nonzero
and has character \((-1,+2)\), so a torus generator
\(aT_R^3+b(B-L)\) fixes it exactly when \(-a+2b=0\).  PCAD transports this
stabilizer relation from the generated parent to its declared structural
descendant.  Normalizing the inherited \(T_R^3\) coefficient to one gives the
displayed generator.  No member-selection operation enters the constructor
word.
\end{proof}

The legacy symbol \(\mathrm{YFIX}_5\) is therefore retired as a live
structural premise.  For backward reference we record
\begin{equation}
\label{p5:eq:yfix5}
 \mathsf R_5^{\rm rep}:=\varnothing.
\end{equation}
If a downstream theorem needs a particular nonzero vector or field in
\(L_0\), a Higgs vacuum, a source member, or a spacetime background, that
occurrence must be named explicitly in the theorem's actual target sort; it
does not define hypercharge.  The relative external/internal grading is also
not an independent bridge.  Paper~III-B's generated diagonal-fibre graph has
an orientation-preserving spin lift, and Paper~III-E composes the same
protected constituents with the local Lorentz/null-spinor role.  The
resulting ancestry map \(J_\Delta:S_x\to W_{\mathrm{spin}}\), on the
complex internal spin carrier $W_{\mathrm{spin}}=W_L\oplus W_R$
(and not the seven-real-dimensional response $W_7$), satisfies
\begin{equation}
\label{p5:eq:chi5-derived}
 J_\Delta\gamma_x=\Gamma_WJ_\Delta.
\end{equation}
Hence the old symbol \((\mathrm{CHI}_5)\) is discharged by upstream ancestry
and is retained only as a legacy label in older editions.

The local gauge datum is the theorem certificate
\begin{equation}
\label{p5:eq:joint5-occ}
\begin{aligned}
({\rm JOINT}_5^{\rm der}):\quad
&\text{the canonical nonflat weak anchor, full candidate-gauge connection}\\
&\text{family, and generated odd slots occur in one rooted Dirac family}.
\end{aligned}
\end{equation}
Define
\[
 \mathsf J_5^{\rm conn,der}:=({\rm JOINT}_5^{\rm der}),
 \qquad
 \mathsf R_5^{\rm conn}:=\mathsf J_5^{\rm conn,der}
 \quad\text{(legacy notation only)}.
\]
Theorem~\ref{p5:thm:canonical-full-joint-connection} proves this certificate
from the generated prefix, so it is excluded from the residual envelope.  In the
represented carrier its joint connection is
\[
\nabla^F
=\nabla^W\otimes1_Q\otimes1_G
 +1_W\otimes\nabla^Q\otimes1_G .
\]
Paper~III-F, Theorem~\ref{p3f:thm:joint-same-datum-relation-closure}, then
derives the common factor connection, curvature, and profile functorially.
Thus the former
factorwise \(({\rm CONN}_{\rm occ})+({\rm JOINT}_5)\) split supplies no
residual data: the nonzero anchor and JOINT relations are theorem
consequences.  The uninserted common-trace
theorem applies only on the kinetic domain stated in item~(v); that domain is
a theorem qualifier, not a selector.

\(\mathsf R_5^{\rm SSQFT}\) is instead an analytic theorem domain.  The
local Lorentzian cone, time orientation, and positive polar orientation are
imported from Papers~III-A/III-E and are not re-assumed here.  What remains in
this domain are the invariant-vacuum, translation/evolution representation,
spectrum-condition relative to that already generated future cone, locality,
operator-domain, and renormalization hypotheses required for spin--statistics and the
stated four-dimensional QFT calculation.  It is not a member selector.  On this domain the half-spin visible carrier is fermionic by spin--statistics,
thereby discharging Paper~IV's previously unresolved statistics family on the
Paper--V-A branch.  The local background-gauge/heat-kernel realization is
licensed there, while the determinant powers are fixed upstream by Paper~IV
and the universal one-loop weights are derived below from the resulting
graded resolution-boundary response.  The claim-specific refinement
\begin{equation}
\label{p5:eq:r5-pct}
\mathsf R_5^{\rm PCT}\subseteq\mathsf R_5^{\rm SSQFT}
\end{equation}
means a local covariant positive-energy field realization satisfying the
locality, common invariant domain, covariance, vacuum, spectrum, and
analyticity hypotheses of the PCT theorem used in that claim
\cite{Haag1996}.  On this narrower domain the antiunitary PCT operation
actualizes the conjugate physical sector.  This is an analytic completion
theorem, not a finite derivation of \(F\oplus\bar F\).  We abbreviate
\begin{equation}
\label{p5:eq:omega5-envelope}
\mathsf R_5
=
\mathsf R_5^{\rm SSQFT},
\qquad
\mathsf R_5^{\rm PCT}\subseteq\mathsf R_5^{\rm SSQFT}.
\end{equation}
Only the clauses actually named in a theorem belong to its premise.  In
particular, the finite carrier and structural hypercharge use no representation
residual; hypercharge is generated by Proposition~\ref{p5:prop:v187-pcad-yfix},
while chirality uses the derived relation \eqref{p5:eq:chi5-derived}.  Local
gauge curvature and coupling matching use the derived
\(\mathsf J_5^{\rm conn,der}\); Fock/CAR, anomaly-as-a-local-QFT statement,
and one-loop running use \(\mathsf R_5^{\rm SSQFT}\); and physical
conjugate-sector completion uses \(\mathsf R_5^{\rm PCT}\).  These are named
closure targets for later construction, not terminal postulates.  Any
particular nonzero field, source, vacuum or background occurrence must be
stated in its own target sort.  The graded bridge is
already closed by \eqref{p5:eq:chi5-derived}, and the connection bridge is
closed by Theorem~\ref{p5:thm:canonical-full-joint-connection}; neither its
anchor nor its factor relations may be separately re-assumed.  Closing
\(\mathsf R_5^{\rm SSQFT}\)
requires construction of the stated Lorentzian positive-energy local-QFT
domain and verification of the spin--statistics hypotheses on it; once those
hold, the Fermi branch is theorem-level rather than selected.

For compatibility with the earlier source ledger, the post-composition
problems remain denoted
\begin{equation}
\label{p5:eq:completion-obligation-packet}
\mathcal B_5=\{\mathcal B_Y,\mathcal B_H,\mathcal B_\chi,\mathcal B_{\rm tr}\},
\end{equation}
but their statuses are now typed rather than conflated:
\(\mathcal B_H\) is closed by the finite scalar theorem,
\(\mathcal B_Y\) is closed by Proposition~\ref{p5:prop:v187-pcad-yfix},
\(\mathcal B_\chi\) is closed by \eqref{p5:eq:chi5-derived}, and
\(\mathcal B_{\rm tr}\) is a theorem consequence of
\(({\rm JOINT}_5^{\rm der})\) on the displayed kinetic domain.  A Majorana
operator, if later supplied, is an optional interaction; neither its existence
nor the neutrality of its carrier is used to choose hypercharge.
\end{definition}

\begin{remark}[Statistics chronology]
\label{p5:rem:statistics-derived-later}
Paper~IV retains the full Young-sector atlas and therefore leaves physical
statistics unresolved generically.  The domain \(\mathsf R_5^{\rm SSQFT}\)
is not a statistics selector; it supplies the Lorentzian local,
positive-energy spin--statistics theorem domain on which the visible
half-spin matter carrier is forced into the Fermi/exterior-power branch.
Thus the Paper--IV statistics interface is \emph{DERIVED-LATER} on this
Paper--V-A branch, while Paper~IV remains historically unresolved outside it.
\end{remark}

\begin{definition}[Derivation-status criterion]
\label{p5:post:structural-criterion}
\label{p5:def:structural-criterion}
A statement is \emph{derived} when its hypotheses are contained in the
displayed prefix and the bridge theorems defining
\(\mathsf G_5^{\rm fin}\), and the conclusion follows without a named
residual unless that residual appears explicitly in the theorem.  It is
\emph{orientation-relative unique} when every
realization on the provenance-fixed oriented branch is equivalent and reversal gives
the conjugate representative.  It is an \emph{interface datum} when changing
it leaves the same finite carrier representation realizable.  No conclusion
is called anchor-free unless every finite role used in it is generated and
structural hypercharge, generated chirality ancestry, separately typed member
occurrences, the derived joint-connection certificate, and analytic theorem
domains are reported separately.
\end{definition}

\begin{remark}[Standard--Model naming and the assignment boundary]
\label{p5:rem:visible-sector-target-convention}
The words ``color'', ``family'', ``gauge'', and ``Higgs'' name generated
roles by their representation-theoretic action: the quaternionic trimer
supports the \(A_2\) closure, the complementary trimer is the untouched
multiplicity support, the derived joint connection localizes the generated
derivations, and the protected coordinate representation is the scalar
intertwiner type.  A ``Majorana'' mass instead names optional downstream
interaction data.  Hypercharge and the charge table are consequences of the
generated parent relation and character-kernel stabilizer, and do not
assume that interaction.
\end{remark}

\begin{remark}[Imported source-support motivation for the visible carrier]
\label{p5:rem:visible-source-support-residue}
The macroscopic homogeneous-source problem motivates a finite internal source
analysis, but no closed-FLRW geometry is selected or proved in Paper~V-A.
Paper~III owns the gravitational action and metric response, while Paper~VI-A
later realizes their closed-FLRW branch.  The charged, chiral, replicated, and
scalar-coupled finite carrier is derived from \(\mathfrak F_{\le4}\) by the
bridge theorems below.  Majorana masses, vacuum data, and numerical couplings
remain optional downstream interactions or calibrations.
\end{remark}

\begin{proposition}[Conditional homogeneous-source interface]
\label{p5:thm:local-flrw-source-support-residue}
Given the homogeneous and isotropic scalar-clock branch later realized in Paper~VI-A,
\[
 ds^2=-dt^2+a(t)^2\gamma^{(+1)}_{ij}dx^idx^j,
\]
the unperturbed macroscopic source supported by the Paper--V-A carrier has the perfect-fluid form
\[
 T^\mu{}_{\nu}=\operatorname{diag}(-\rho c^2,p,p,p).
\]
This is an imported conditional interface, not a Paper--V-A selection of the metric, curvature sign, gravitational action, or cosmological branch.  It supplies no particle catalogue and no internal charge assignment.
\end{proposition}

\begin{proof}
On the stated imported homogeneous and isotropic branch, any unperturbed
rank-two source is invariant under the spatial isotropy group and therefore
has one temporal eigenvalue and one triply degenerate spatial eigenvalue,
giving the displayed perfect-fluid form.  The internal finite representation
is reconstructed independently from the oriented incidence geometry below.
\end{proof}

\begin{proposition}[Macroscopic source form does not replace the finite incidence proof]
\label{p5:prop:visible-source-residue-forces-sector-problem}
Let the upstream carrier and Paper--IV finite operator bridge be fixed.  The
macroscopic source problem by itself does not determine the oriented
axial/opposite-trimer frame, the quark--lepton \(1\oplus3\) incidence, or a
local connection.  The first two follow from Paper~I-A's finite orthogonality
theorem; the last follows independently from the constructor word defining
\(\mathsf J_5^{\rm conn,der}\).
\end{proposition}

\begin{proof}
Homogeneity and isotropy constrain only the spacetime tensor form of the
unperturbed source.  They are invariant under changes of finite internal
representation having the same macroscopic stress tensor.  Hence they cannot
replace the independent oriented finite-incidence proof or generate a
connection.  That separate result is
Theorem~\ref{p5:thm:canonical-full-joint-connection}.
\end{proof}

\subsection{Mixed flavor readouts from the existing weak-current maps}
\label{p5:sec:v199-weak-current-readouts}
Apply the kinetic and scalar-field normalization of the admitted action
branch before forming the left generation Grams
\(A=\widehat Y_u\widehat Y_u^\dagger\) and
\(B=\widehat Y_d\widehat Y_d^\dagger\) on the common \(G\).
Right Grams have the same separate spectra but cannot replace this pair in
mixed traces without the corresponding transport.  On the left weak doublet,
with the color spectator temporarily omitted, put
\[
 \mathsf A=P_u\otimes A,\quad \mathsf B=P_d\otimes B,\quad
 T_+=|u\rangle\langle d|\otimes I_G,\quad T_-=T_+^\dagger.
\]
The last two maps are the internal coefficients of the existing weak current.
A change of weak frame transports the whole tuple together.

\begin{theorem}[Same-carrier mixed-loop constructor]
\label{p5:thm:v199-weak-current-moments}
On \(\mathcal H_Q=(\mathbb C u\oplus\mathbb C d)\otimes G\),
\begin{align}
 \operatorname{Tr}_{\mathcal H_Q}
 (\mathsf A^pT_+\mathsf B^qT_-)&=\operatorname{tr}_G(A^pB^q),
 \qquad 0\le p,q\le2,\label{p5:eq:v199-current-moment}\\
 \frac1{6i}\operatorname{Tr}_{\mathcal H_Q}
 [\mathsf A,T_+\mathsf BT_-]^3
 &=\frac1{6i}\operatorname{tr}_G[A,B]^3.
 \label{p5:eq:v199-current-cp}
\end{align}
Zero powers mean the identity on \(\mathcal H_Q\).  Restoring the color
spectator multiplies both traces by three.
\end{theorem}
\begin{proof}
For positive powers, \(T_+(P_d\otimes B^q)T_-=P_u\otimes B^q\).
For zero powers the two currents themselves supply the block restriction.
The commutator is supported on the up block and restricts there to \([A,B]\).
Taking traces proves both identities.
\end{proof}
The constructor places $A,B$ on the same generation carrier $G$.
Once this pair is formed, its generic rank-three invariant calculus is
Paper~1B, Theorem~\ref{p1b:thm:rank-three}; the current identities above
supply the framework realization of those mixed traces.


The construction uses retained restrictions, current coefficients, adjoints,
composition and trace.  It supplies operator readouts of an admitted Yukawa
member; it does not insert mixed terms into an uninserted scalar trace or
assert a measured correlator equals its internal coefficient.  Paper~V-B
uses this packet to reconstruct and test the full generic Dirac orbit in
Theorem~\ref{p6:thm:v199-flavor-reconstruction}.  Its determinant-source
probe in Proposition~\ref{p6:prop:v199-mixed-resolvent} still requires the
actual Clifford insertion, regulator and effective-action weight before it
can be identified with a physical fermion determinant coefficient.
This is a readout constructor on its admitted action/inserted-source
domain. It does not reopen the rooted charged source occurrence or synthesis
closed by Corollary~\ref{p5:cor:v397-charged-preterminal-assembly-closure}.
An actual determinant/channel evaluation still retains its insertion,
regulator, physical state and every limiting domain that it consumes.


% BEGIN v2.36 LOCAL CORE PARENT-INDIFFERENCE
\subsection{rooted neutral response/transport sibling ownership}
\label{p5:sec:v3100-neutral-response-transport-sibling}

The protected neutral $R^+$ lepton branch and the rooted $D_4$ triality source
already carry two distinct descendants in the same source ancestry: the
positive central determinant response
\[
 R_{D_4}^{\rm fam}=xI_G,\qquad x^3=192,
\]
and the rooted family-transport cell
\[
 \mathscr T_{C_3}=\{I_G,X_G,X_G^{-1}\},\qquad X_G^3=I_G,
 \qquad X_G^\dagger=X_G^{-1}.
\]
The response coefficient and transport are different typed outputs; neither is
reconstructed from the other.

\begin{theorem}[Neutral response/transport sibling operation jet]
\label{p5:thm:v3100-neutral-response-transport-sibling-jet}
On the occurring rooted neutral source word, retain the already declared
one-pass operation labels
\(\mathsf{Inv},\mathsf{Dir},\mathsf{Gram}\) before terminal scalarization.
Preterminal relation preservation then retains the operation-matched sibling
pair
\[
 \boxed{
 \mathfrak J_{56}^{\rm sib}
 =\bigl(
 (R^+,X_G^+),
 (R,X_G),
 (R^\dagger R,X_G^\dagger X_G)
 \bigr)
 =\bigl(
 (x^{-1}I_G,X_G^{-1}),
 (xI_G,X_G),
 (x^2I_G,I_G)
 \bigr).}
\]
No numerical response coefficient selects a transport element: the pairing is
by the frozen operation role of one common rooted source word.
\end{theorem}

\begin{proof}
The central response and rooted $C_3$ cell are generated descendants of the
same occurring triality-source ancestry, while the inverse/direct/closed-Gram
roles are the already frozen constituents of the terminal neutral word.
Paper~II-A preterminal relation preservation forbids replacing a complete
operator-valued sibling by its later scalar descendant before all relations
consumed by the same source have been formed.  Hence the three operation roles
act on the retained response/transport sibling packet before terminal
readout.  Positivity of $R=xI_G$ gives $R^+=x^{-1}I_G$ and
$R^\dagger R=x^2I_G$; unitarity and order three of $X_G$ give
$X_G^+=X_G^{-1}$ and $X_G^\dagger X_G=I_G$.  No extra relation, coefficient or
family-unitary member is inserted.
\end{proof}

\begin{corollary}[Canonical neutral regular-module incidence]
\label{p5:cor:v3100-neutral-regular-module-incidence}
Let $g_0$ be the rooted family seed and
$g_j=X_G^jg_0$, $j\in\mathbb Z_3$.  On the neutral role space
\(\mathscr L_{56}=\operatorname{span}\{e_{\rm inv},e_{\rm dir},e_{\rm Gram}\}\),
define
\[
 \boxed{
 W_\nu e_{\rm inv}=X_G^{-1}g_0=g_2,\qquad
 W_\nu e_{\rm dir}=X_Gg_0=g_1,\qquad
 W_\nu e_{\rm Gram}=g_0.}
\]
Then $W_\nu$ is unitary.  Rooted cycle/reversal changes transform the role and
transport presentations simultaneously, so no physical $U(3)$ selector
survives.  This is current/incidence data only; it does not define or
rediagonalize any neutral mass.
\end{corollary}
% END NEUTRAL-RESPONSE-TRANSPORT-SIBLING v3.100

% BEGIN QUARK-CURRENT-INCIDENCE-CLOSURE v3.101
\subsection{quark weak-current incidence and role-space invariants}
\label{p5:sec:v3101-quark-current-incidence-closure}

\begin{theorem}[Generated quark weak-current incidence]
\label{p5:thm:v3101-generated-quark-current-incidence}
Paper~V-B Release v5.94 supplies the generated normalized terminal-role
incidences on the full-rank current-normalization domain stated there,
\[
 W_u:\mathscr L_u^{\rm role}\xrightarrow{\cong}G_L,
 \qquad
 W_d:\mathscr L_d^{\rm role}\xrightarrow{\cong}G_L.
\]
The left-handed quark weak-doublet current already has family coefficient
\(\mathscr J_W=I_G\) (with color a spectator).  Therefore its pullback to the
already-massed terminal role spaces is
\[
 \boxed{C_{ud}:=W_u^\dagger\mathscr J_WW_d=W_u^\dagger W_d.}
\]
This is a generated unitary current/incidence map; no Yukawa diagonalization,
mass eigenframe or additional family-unitary selector is required.  Hence
\[
 \boxed{
 \texttt{QUARK\_WEAK\_CURRENT\_INCIDENCE}=\texttt{CLOSED}.}
\]
This status closes ownership and the dependent current construction.
Numerical entries and physical flavor comparisons consume the actual
source evaluations, rank determinants and invariant contractions of
Paper~V-B, Proposition~\ref{p6:prop:v2211-quark-current-output-certificate};
the legacy CKM benchmark is not substituted for those outputs.
\end{theorem}

\begin{corollary}[Former CKM structural target is exhausted]
\label{p5:cor:v3101-former-ckm-target-exhausted}
The former relative mass-eigenframe object carries no independent native
obligation.  The charged quark masses are terminal geometric readouts and the
flavor-changing content is the current map $C_{ud}$.  Thus
\[
 \boxed{
 \texttt{CKM\_RELATIVE\_FRAME}=\texttt{RETIRED\_OBSOLETE\_ONTOLOGY}.}
\]
When a conventional coordinate section is useful one may define only after
this closure
\[
 V_{\rm CKM}^{\rm eff}:=[C_{ud}]_{\rm std}.
\]
\end{corollary}

\begin{definition}[Native terminal-role weight operators]
\label{p5:def:v3101-native-role-weight-operators}
On the already generated charged terminal role spaces define
\[
 Q_u:=\sum_{r\in\mathcal I_u}(m_{u,r}^{\rm fin})^2\Pi_{u,r}^{\rm role},
 \qquad
 Q_d:=\sum_{r\in\mathcal I_d}(m_{d,r}^{\rm fin})^2\Pi_{d,r}^{\rm role}.
\]
These are bookkeeping operators on the terminal role spaces; their entries are
the already generated finite mass coefficients and are not eigenvalues of a
later family operator.
\end{definition}

\begin{theorem}[Current-based mixed quark invariants]
\label{p5:thm:v3101-current-based-quark-invariants}
For $p,q\ge0$ the native mixed current moments are
\[
 \boxed{
 \mathcal M_{pq}^{(q)}
 :=\operatorname{tr}_{\mathscr L_u}
 \left(Q_u^pC_{ud}Q_d^qC_{ud}^\dagger\right).}
\]
The CP-odd current invariant is
\[
 \boxed{
 \chi_q^{\rm cur}
 :=\frac1{6i}\operatorname{tr}_{\mathscr L_u}
 \left[Q_u,C_{ud}Q_dC_{ud}^\dagger\right]^3.}
\]
These are invariant under simultaneous role rephasings/presentation changes.
If one subsequently chooses the conventional SM coordinate dictionary, they
reduce to the usual mixed-mass and Jarlskog-type expressions.  That reduction
is a dictionary theorem and does not alter native mass ownership.
\end{theorem}


% BEGIN FRAMEWORK v2.240: OBSERVABLE-MINIMAL PHYSICAL QUARK CURRENT
\begin{theorem}[Four mixed quark moments and one orientation test]
\label{p5:thm:v240-five-quark-current-invariants}
On the actual full-rank native current domain of
Theorem~\ref{p5:thm:v3101-generated-quark-current-incidence}, let
$Q_s=\operatorname{diag}(x_{s,1},x_{s,2},x_{s,3})$ in the already
ordered terminal-role coordinates, with pairwise distinct $x_{s,i}$.
Set $P_{ij}=|(C_{ud})_{ij}|^2$ and
\[
 (V_s)_{pi}=x_{s,i}^p\quad (p=0,1,2),\qquad
 M_{pq}=\operatorname{tr}(Q_u^pC_{ud}Q_d^qC_{ud}^\dagger).
\]
Then the exact current probabilities are
\begin{equation}\label{p5:eq:v240-quark-vandermonde-current}
 \boxed{P=V_u^{-1}MV_d^{-T}.}
\end{equation}
Unitarity fixes $M_{00}=3$, the two single-sector pairs
$M_{p0}=\operatorname{tr}Q_u^p$ and
$M_{0q}=\operatorname{tr}Q_d^q$ for $p,q=1,2$.
Thus only the four actual mixed contractions
$M_{11},M_{12},M_{21},M_{22}$ are needed to determine $P$.
For a generic nonzero-CP unitary current, the sign of the already
defined $\chi_q^{\rm cur}$ supplies the remaining
complex-conjugation orientation of its fixed ordered rephasing class.
These five invariants are \emph{tests to execute on the actual}
$E_u,E_d$; they cannot be inferred from a candidate parent
$H_+,H_-$ whose incidence relation is not established.
\end{theorem}
\begin{proof}
Expanding the mixed trace in the fixed terminal role bases gives
$M_{pq}=\sum_{ij}x_{u,i}^p P_{ij}x_{d,j}^q$, or
$M=V_uPV_d^T$. Distinct terminal weights make both Vandermonde
matrices invertible. Unitarity gives the row and column sums of $P$.
For a generic $3\times3$ unistochastic matrix with nonzero area,
its unitarity triangles fix the Jarlskog magnitude and leave the
two conjugate orientations, separated by the sign of $\chi_q^{\rm cur}$.
\end{proof}
\begin{corollary}[Exact generating function and residue test for the native current]
\label{p5:cor:v240-native-current-stieltjes}
For $z,w>0$, $R_s(z)=(zI+Q_s)^{-1}$ and
$\mathcal G_{\rm cur}(z,w)=\operatorname{tr}(R_u(z)C_{ud}
R_d(w)C_{ud}^\dagger)$, one has
\[
 \boxed{\mathcal G_{\rm cur}(z,w)=
       \sum_{ij}\frac{P_{ij}}{(z+x_{u,i})(w+x_{d,j})}.}
\]
Every $P_{ij}$ is the corresponding double residue of its
meromorphic continuation. This is an exact readout of the
\emph{actual} current when $C_{ud}$ has been compiled; it does not
identify the current with a raw candidate eigenframe.
\end{corollary}
\begin{proof}
Insert the diagonal resolvents and the already generated current.
Distinct role weights separate their poles, whose double residues
are exactly the squared current moduli.
\end{proof}
% END FRAMEWORK v2.240: OBSERVABLE-MINIMAL PHYSICAL QUARK CURRENT

\begin{remark}[Supersession of v1.99 Yukawa-Gram flavor language]
\label{p5:rem:v3101-supersede-yukawa-gram-flavor}
Section~\ref{p5:sec:v199-weak-current-readouts} and the older
$Y_sY_s^\dagger$ relative-eigenframe discussion remain valid algebra inside a
conventional Standard--Model/Yukawa coordinate realization.  They no longer
constitute the native flavor ontology and may not be used to generate terminal
masses, $C_{ud}$, $C_{e\nu}$, or a family selector.  The active native formulas
are the terminal-role/current expressions above.
\end{remark}
% END QUARK-CURRENT-INCIDENCE-CLOSURE v3.101

% BEGIN VISIBLE-INTERNAL-RESPONSE-REALIZATION v3.102
\subsection{Rooted source descendants after charged assembly closure}
\label{p5:sec:v393-rooted-assembly-target}
\begin{corollary}[Rooted source actuality discharges the descendant antecedents]
\label{p5:cor:v393-rooted-source-descendants}
On the occurring rooted Hurwitz--$D_4$/triality branch and the local
clocked uninserted action domain of
Theorem~\ref{p5:thm:v392-rooted-srcval-incidence}, one has
\[
 \operatorname{Occ}(\mathbb T_*)
 \Longrightarrow
 \operatorname{CoReal}_{\rm src}\{I_G,C_{1,G},C_{2,G}\}
 \quad\Longleftrightarrow\quad
 \operatorname{CoReal}_{\rm src}\{I_G,X_G,X_G^{-1}\}.
\]
The rank-$2,4,6$ primitive-owner descendants are consumed in that same
rooted source owner. No additional primitive source-occurrence token is needed.
\end{corollary}
\begin{proof}
Theorem~\ref{p5:thm:v392-rooted-srcval-incidence} supplies the specific
nonzero member $D_{F,*}^{\rm geom}$, with square $I$.
It therefore supplies the nonzero-source antecedent of
Theorems~\ref{p5:thm:v381-c3-source-cell-corealization}
and~\ref{p5:thm:v382-upstream-generated-source-cell} on their common rooted
domain. Their cell bases are related by the invertible change of generators
between $X_G,X_G^{-1}$ and $C_{1,G},C_{2,G}$.
Theorem~\ref{p5:thm:v375-d4-triality-source-consumption}, with the same
antecedent supplied, consumes the already derived primitive descendants.
This composes existing constructors; it does not choose a new family member.
\end{proof}

Paper I-C's root-normalized joint-quadratic restriction
(\texttt{p1b:thm:v135-\allowbreak hopf-\allowbreak sector-\allowbreak phase}) gives the structural sector typing
\[
 R^+\longmapsto R_G,\qquad R^-\longmapsto R_G^\dagger.
\]
Here the $R^\pm$ slots use the retained $T_R^3$ convention. The internal
Hopf interchange keeps the rooted family cycle fixed; reversing the whole
rooted tuple also reverses that cycle. Parent-relation descent preserves
these types without supplying a new physical action term.


The rooted charged zero-order packet is now supplied by
Theorem~\ref{p5:thm:v397-charged-preterminal-srcval-realization};
there is no remaining charged SVD assembly or independent $J_s$ target.
The superseded candidate matching ledger is no longer an active condition.
% END ROOTED-FERMION-DESCENDANTS v3.93


% BEGIN CO-CONSUMPTION-FERMION-FRONTIER v3.94
\subsection{neutral-role co-consumption and charged scope}
\label{p5:sec:v394-coconsumption-fermion-frontier}

\begin{corollary}[Neutral terminal-role handoff after co-consumption]
\label{p5:cor:v394-neutral-role-handoff}
On the rooted $D_4$ branch, the central determinant response
\[
 R_{D_4}^{\rm fam}=192^{1/3}I_G
\]
remains a generated descendant of the consumed triality/source parent.  The
physical use of its complete inverse/direct/closed-Gram role packet is no
longer proved by first inserting rank-one neutral family projectors.  Paper~V-B
v5.87 instead constructs the complete three-role response jet in its own typed
role space and closes its simultaneous use by Paper~II-A generated
co-consumption.  A subsequent role-to-family embedding is a distinct
downstream object.
\end{corollary}


The neutral role-to-family map retains its distinct output sort; the native
energy/incidence packet is supplied by Paper~V-B,
Theorem~\ref{p6:thm:v593-neutral-energy-incidence-packet}.
The former charged typed-parent residual is no longer an occurrence gate; the current charged
instance is closed by Corollary~\ref{p5:cor:v397-charged-preterminal-assembly-closure}.
% END CO-CONSUMPTION-FERMION-FRONTIER v3.94



% BEGIN CHARGED-LEPTON-DIRECTION-CLOSURE v3.107
\subsection{charged-lepton singlet endpoint and direction closure}
\label{p5:sec:v3107-charged-lepton-direction-closure}
\noindent\textbf{Current scope.} The in-plane open status in this historical release is superseded by main-text Release v3.108.  Its mass/incidence separation and the tau/muon direction result remain controlling inputs.
Paper~V-B v5.96 evaluates the retained tau primal--dual owner before terminal
scalarization.  Its canonical common-parent family return is the rooted
$C_3$-fixed line, so
\[
 \boxed{P_\tau=Q_0.}
\]
The tau mass coefficient remains the already derived $E_{36}$ terminal
coefficient: mass ownership and left-family endpoint are distinct sibling
readouts.  Since the electron is an $E_{40}$ role and the normalized
three-role current evaluation is full rank, the muon has the opposite
$E_{36}$ direction.  Thus
\[
 \boxed{
 \iota_\tau\in Q_0\End(G)Q_\perp,\qquad
 \iota_\mu\in Q_\perp\End(G)Q_0,}
\]
while the electron and muon normalized left-family lines span $Q_\perp G$.

The tau/muon direction result is retained. The in-plane incidence is
now evaluated by Theorem~\ref{p5:thm:v3108-lepton-hopf-incidence-parent}
and Corollary~\ref{p5:cor:v3108-lepton-incidence-status}.
The superseded open-plane/no-consumer statement is no longer an active gate.
% END CHARGED-LEPTON-DIRECTION-CLOSURE v3.107



% BEGIN CHARGED-LEPTON-HOPF-PCAD-INCIDENCE-CLOSURE v3.108
\subsection{charged-lepton Hopf--hierarchy incidence descent}
\label{p5:sec:v3108-lepton-hopf-pcad-incidence}
The remaining charged-lepton datum after Release v3.107 is a relation-valued
left-current incidence, not a missing mass coefficient.  This distinction
allows the already generated negative-Hopf family relation to be consumed in
its native structural sort without promoting it to a new Yukawa/source value.

\begin{theorem}[Negative-Hopf family relation is the charged-lepton incidence parent]
\label{p5:thm:v3108-lepton-hopf-incidence-parent}
Let the joint rooted Hopf--Hurwitz relation of Paper~I-C be
\[
 \mathscr R=P_+\otimes R_G+P_-\otimes R_G^\dagger,
 \qquad R_G=\operatorname{diag}(1,-i,i).
\]
Paper~V-A's physical charged-lepton right support is the already generated
$R^-$ source slot, hence its inherited family relation is
\[
 \boxed{R_-:=R_G^\dagger.}
\]
For the charged-current incidence sibling, apply the already generated
hierarchy/family block projections.  On the v3.107 placement these are
\[
 \pi_\mu(A)=Q_\perp A Q_0,
 \qquad
 \pi_\tau(A)=Q_0 A Q_\perp .
\]
Typed Parent Correspondence and Action Descent therefore retains
\[
 \boxed{A_\mu^{\rm rel}=Q_\perp R_-Q_0,
 \qquad A_\tau^{\rm rel}=Q_0R_-Q_\perp.}
\]
Both blocks are nonzero and rank one.  These are structural incidence
siblings of the charged roles.  They do not replace the direct charged mass
compiler and do not assert that $R_-$ is itself a physical mass/source
coefficient.
\end{theorem}
\begin{proof}
The joint quadratic generates $\mathscr R$ before the right-Hopf support slots
and family-current descendants are separated, and Paper~V-A identifies those
physical right supports by the same projectors $P_\pm$.  Hence restriction to
the charged $R^-$ slot gives the frozen structural relation $R_-$.  The
charged hierarchy/family bridge and Release v3.107 already type the muon and
tau incidence siblings as $Q_\perp\End(G)Q_0$ and
$Q_0\End(G)Q_\perp$, respectively.  The displayed maps are therefore licensed
linear projections on the domains used by the incidence word.  Paper~II-A's
PCAD theorem retains the image of the generated parent relation under exactly
such licensed projections.  No change of target to \(\mathsf{SrcVal}\) occurs.
Since $R_-c$ is not proportional to $c$, both off-diagonal compressions are
nonzero; each has one-dimensional source or target support and hence rank one.
\end{proof}

\begin{corollary}[No independent charged-lepton plane selector remains]
\label{p5:cor:v3108-lepton-incidence-status}
After canonical left-current normalization, the muon range is the normalized
range of $Q_\perp R_-Q_0$, the tau range remains $Q_0G$, and the electron line
is the unique orthogonal complement of the muon line in $Q_\perp G$.
Consequently
\[
 \boxed{\begin{gathered}
 \texttt{LEPTON\_ROLE\_TO\_FAMILY\_INCIDENCE}\\
 =\texttt{CLOSED\_LEFT\_CURRENT\_BY\_HOPF\_HIERARCHY\_PCAD}.
 \end{gathered}}
\]
This relation-level result does not evaluate physical pole/rest or detector
claims. A genuinely different generic fermion target must specify its own
source sort; the rooted charged source contraction itself is already closed
by Corollary~\ref{p5:cor:v397-charged-preterminal-assembly-closure}.
\end{corollary}
% END CHARGED-LEPTON-HOPF-PCAD-INCIDENCE-CLOSURE v3.108



% BEGIN INLINE SOURCE: Paper5A-ChargedSourceApplication-Main-v1_0.tex
% Framework v2.168: integrated finite charged-source derivations.
\subsection{Charged source occurrence and the remaining pole application}
\label{p5:sec:v2168-charged-source-application}
On the independently certified hierarchy-arrow domain, the occurring charged source is the synthesis
$R_{\rm act}=\sum_r\varepsilon_r\Phi_{\rm HG}(A_r^{\rm hier})$
of the retained terminal coefficients and hierarchy/incidence siblings.
Theorems~\ref{p5:thm:v397-charged-preterminal-srcval-realization} and
\ref{p5:cor:v397-charged-preterminal-assembly-closure} establish its
charged occurrence and assembly on that certified domain. A05 consumes this source; it does not
reactivate A01 or introduce an independent source-incidence selector.

\begin{corollary}[Two sufficient finite application routes]
\label{p5:cor:v2168-charged-source-application-routes}
On a certified same-action source/kinetic/constraint domain, the actual
source may be tested by either of the following claims, at its stated scope.
\begin{enumerate}
\item For the complete tagged free response, the exact condition is
$R_{\rm act}R_{\rm act}^\dagger=UM^2U^\dagger$ on the uninserted chart,
or the corresponding normalized return on an inserted chart. V-B's
unit-support transfer graph is an equivalent sufficient certificate when
proved for the actual exporter.
\item For an interacting pole-only claim, use the full action-generated
$H$, physical source and stationary quotient of IV. The actual complete
Schur kernel at the fixed terminal energy must have a source-visible zero;
the full charge-sector rank test is equivalent and also applies where an
eliminated complement is singular.
\end{enumerate}
The second route does not require the stronger free-response identity.
\end{corollary}
\begin{proof}
The first statement is III-F,
Theorem~\ref{p3f:thm:v2168-left-return-test}, together with V-B,
Theorem~\ref{p6:thm:v2168-unit-support-equivalence}. The second is IV's
spectral source decomposition and
Theorem~\ref{p4:thm:v2168-faithful-stationary-poles}, with its actual
state-support visibility test where faithfulness is not available.
Both consume the occurring source; neither defines a terminal coefficient
through a later spectral value.
\end{proof}

For the free route, the first unevaluated exporter equality is the actual
unit-support graph or equivalent left return, rather than equality to one
particular geometric right frame. For the dynamic route, the physical
rest/source/constraint chart is certified first; then evaluate the full
$S(\varepsilon_r)$ of V-C,
Theorem~\ref{p5c:thm:v2168-complete-weak-memory}. A current ray or gauge
normalization does not prove this kernel vanishes. The complete action-generated root remains A05's finite task. A root
proved on the exact actual support includes physical visibility; a
faithful-carrier root alone retains that support check. The named particle
interpretation and consumed rest/refinement/resonance/charged infrared
limits remain on their exact domains.

% END INLINE SOURCE: Paper5A-ChargedSourceApplication-Main-v1_0.tex



% v2.169: generated action/support constructors and consuming scope.

% BEGIN INLINE SOURCE: Paper5A-ChargedSupportApplication-Main-v1_1.tex
% Framework v2.169: Batches 7--13, integrated at the owning source.
\subsection{Current charged hierarchy-arrow origin and source-value scope}
\label{p5:sec:v228-charged-origin-correction}
\begin{definition}[Charged hierarchy-arrow origin certificate]
\label{p5:def:v228-hierarchy-origin}
Write $\mathsf{Cert}_{\rm hier}$ for an independently certified licensed
result-independent constructor on the occurring rooted parent for every
$A_{s,r}^{\rm hier}$ consumed by the charged packet, retaining its directed
incidence, native metric and coefficients before terminalization.
The certificate may be given by a directed expression or by an equivalent
proved whole-family constructor. It is an input-origin obligation,
not another source-member G0 token.
\end{definition}
\begin{proposition}[Current scope of charged assembly and source occurrence]
\label{p5:prop:v228-charged-origin-scope}
The hierarchy sibling export in
Proposition~\ref{p6:prop:v589-mass-incidence-sibling-export} is conditional
on an independently generated hierarchy arrow. The inverse channel
presentation and linear sum in
Theorem~\ref{p6:thm:v590-charged-channel-representatives} construct the
assembly on that generated-packet domain. The justified source handoff is
\[
 \operatorname{Occ}(\mathbb T_*)\land\mathsf{Cert}_{\rm hier}
 \ \Longrightarrow\
 \operatorname{CoReal}_{\mathsf{SrcVal}}(D_{F,\rm ch}^{\rm geom}).
\]
The displayed assembly chain does not itself discharge
$\mathsf{Cert}_{\rm hier}$. Earlier root-only charged closure wording is
read with this current antecedent. The explicit primitive member of
Definition~\ref{p5:def:v392-rooted-geometric-finite-odd} retains its
independent geometric source-existence closure.
\end{proposition}
\begin{proof}
Apply Proposition~\ref{p2a:prop:v228-assembly-origin} to the conditional
export, its presentation isomorphism and the linear assembly. Once its
input origins are certified, the existing source consumer requires no
duplicate occurrence assumption. A later assembly formula cannot supply
an earlier input's missing origin certificate by citing itself.
\end{proof}
The compiler explicitly supplies $R_G,X,K_7,K_9$. Its common commuting
subalgebra is $\mathbb C[X]$; the full input algebra also admits additional
noncommuting Hopf operations. Therefore a particular source word must be
identified before using Theorem~\ref{p1b:thm:v228-character-defect} or
the constant-readout fibre of
Proposition~\ref{p1b:prop:v228-opposite-odd-fibre}.
These furnish concrete sufficient closure routes with no fitted phase.
The structural negative-Hopf current relations retain their stated domains.

% END FRAMEWORK v2.228: CHARGED ORIGIN AND CHARACTER TESTS

% BEGIN INLINE SOURCE: Paper5A-HistoricalChargedAssembly-Appendix-v1_0.tex
% Framework v2.187: superseded charged assembly/incidence diagnostics.
\subsection{Actual charged source incidence after the rooted hierarchy classification}
\label{p5:sec:v229-charged-origin-scope}
\begin{corollary}[Geometric hierarchy arrows are not an automatic action word]
\label{p5:cor:v229-charged-action-origin-scope}
The correct rooted negative-Hopf compressions and the rank-one electron
left-support criterion are given by
Theorem~\ref{p6:thm:v229-rooted-negative-hopf-hierarchy-blocks}.
Theorem~\ref{p6:thm:v229-opposite-odd-positive-converse}
classifies all positive commuting three-role orthonormal returns.
These close geometric arrow existence and a conditional exhaustive
mass-only return, but do not themselves prove the occurrence of the
physical noncentral charged source on those arrows.  The existing
conditional handoff
\[
 \operatorname{Occ}(\mathbb T_*)\land\mathsf{Cert}_{\rm hier}
 \Longrightarrow
 \operatorname{CoReal}_{\mathsf{SrcVal}}(D_{F,{\rm ch}}^{\rm geom})
\]
is unchanged.  To discharge $\mathsf{Cert}_{\rm hier}$ evaluate the
actual same-action preterminal mixed chiral derivative, with right word,
comparison ports and supports intact, as required by
Definition~\ref{p2a:def:v229-action-incidence-certificate}.
No physical quark-origin or interacting stationary pole follows from
these purely family-level certificates.
\end{corollary}
% END FRAMEWORK v2.229: ACTUAL CHARGED SOURCE SCOPE

% BEGIN FRAMEWORK v2.242: SIGNED GEOMETRY AND CHARGED ORIGIN
\begin{remark}[Explicit phase lift and unchanged charged-action certificate]
\label{p5:rem:v242-signed-provider-origin}
The signed rooted phase port now has an explicit unique real-vector lift
with the rooted spinor frames fixed
(Corollary~\ref{p1:cor:v242-joint-phase-vector-lift}); the full seed
also prevents distinct channel phases from being absorbed into either
spinor leg alone. The complete transverse Gram and the explicit family
quadratic-band embedding are further geometric descendants. None changes
the antecedent \(\mathsf{Cert}_{\rm hier}\): physical source use,
the positive-path word and its normalization must follow from the actual
same-action mixed chiral derivative. In particular the new Gram value
\(5\) is not an identification with a rank-five primitive owner.
\end{remark}
% END FRAMEWORK v2.242: SIGNED GEOMETRY AND CHARGED ORIGIN


\subsubsection{Generated replication space}
\label{p5:sec:part1-family-space}

By Theorem~\ref{p5:thm:generation-selection}, the family space is
\begin{equation}
\label{p5:eq:part1-G}
G=\mathbb C^{T_g}\cong \mathbb C[Z_3]\cong \mathbb C^3.
\end{equation}
In the language of Paper~IV \cite{Kugelmann2026e}, this is the internal finite quantum factor carrying the family degrees of freedom, on which family operators act while gauge operators remain generation-blind.
Choose the standard basis
\[
\{e_0,e_1,e_2\}
\]
and define the cyclic generator
\begin{equation}
\label{p5:eq:part1-X}
Xe_0=e_1,\qquad Xe_1=e_2,\qquad Xe_2=e_0,\qquad X^3=1.
\end{equation}
Then the canonical family projectors are
\begin{equation}
\label{p5:eq:part1-family-projectors}
P^{(1)}=\frac{1+X+X^2}{3},
\qquad
P^{(2)}=1-P^{(1)},
\end{equation}
so that \(G\) carries the canonical real/Hermitian family split
\[
G\cong 1\oplus 2.
\]


\section{Parent stabilizer, hypercharge and faithful global symmetry}
\label{p5a:sec:organized-charge}
\subsection{rooted structural \texorpdfstring{$B-L=2$}{B-L=2} neutral coefficient}
\label{p5:sec:v3103-structural-bl2-geometric-coefficient}

\begin{theorem}[Generated rooted charge-two family coefficient]
\label{p5:thm:v3103-generated-bl2-family-coefficient}
Let $m_{\mathcal O}$ be the already generated joint Hopf--axial quadratic of
Proposition~\ref{p5:prop:v187-pcad-yfix}.  Paper~I-A fixes its rooted positive
Hopf coefficient to
\[
 \mu_0=\frac{-1+i}{2},
 \qquad |\mu_0|^2=\frac12,
\]
while Paper~I-C's root-normalized restriction of the same joint quadratic to
the ordered rooted trimer gives
\[
 R_G=\operatorname{diag}(1,-i,i).
\]
Therefore the common rooted structural descendant on the unique neutral-right
support is the concrete symmetric family coefficient
\[
 \boxed{
 M_{BL2}^{\rm geom}:=\mu_0R_G
 =\frac12\operatorname{diag}(-1+i,\,1+i,\,-1-i).}
\]
It obeys
\[
 (M_{BL2}^{\rm geom})^{\mathsf T}=M_{BL2}^{\rm geom},
 \qquad
 M_{BL2}^{\rm geom}(M_{BL2}^{\rm geom})^\dagger=\frac12I_G,
\]
\[
 \boxed{\|M_{BL2}^{\rm geom}\|_{\rm HS}^2=\frac32},
 \qquad
 \boxed{\det M_{BL2}^{\rm geom}=\frac{1+i}{4}\ne0.}
\]
The inherited parent character is $(-1,+2)$ in the
$(T_R^3,B-L)$ plane, hence this structural coefficient has $Y=0$ and source
weight $\Delta(B-L)=+2$.  Complete rooted reversal sends it to its complex
conjugate.
\end{theorem}

\begin{proof}
The phase restriction is diagonal in the ordered rooted-trimer presentation,
so it is symmetric under the transpose used by the Majorana pairing slot.
The displayed norm and determinant follow directly from $|\mu_0|^2=1/2$ and
$\det R_G=1$.  The charge statement is the already generated parent character,
and $Y=-1+\tfrac12(2)=0$.
\end{proof}

\begin{corollary}[Structural coefficient closed; Majorana action occurrence remains open]
\label{p5:cor:v3103-bl2-structural-vs-action}
The optional Majorana support of
Theorem~\ref{p5:thm:part1-yukawa-majorana-unique} no longer requires an
arbitrary family matrix merely to exhibit a native structural candidate:
$M_{BL2}^{\rm geom}$ is generated and full rank.  However the current local
fermion action still contains no constant odd/Majorana offset, so
\[
 \boxed{
 \texttt{BL2\_GEOMETRIC\_COEFFICIENT}=\texttt{CLOSED\_STRUCTURAL},}
\]
while
\[
 \boxed{
 \texttt{BL2\_BOUNDARY\_ACTION\_OCCURRENCE}=\texttt{OPEN\_ACTION\_OWNER}.}
\]
Neither PCAD nor co-consumption may be run backwards to promote the structural
coefficient to a physical source value without a licensed action constructor
whose required output contains the neutral-right pairing block.
\end{corollary}

\begin{remark}[CP scope]
\label{p5:rem:v3103-bl2-cp-scope}
The raw coordinate phase of $\mu_0$ is not by itself an intrinsic CP source:
Paper~I-A proves that the normalized nonzero quadratic defines a flat Hermitian
connection and is locally representable by a fixed positive real coefficient.
A physical CP-odd history therefore requires an occurring relative or
history-dependent mark in the same action owner, not merely the displayed
constant complex coordinate.
\end{remark}
% END STRUCTURAL-BL2-GEOMETRIC-COEFFICIENT v3.103

% BEGIN BL2-BOUNDARY-PAIR-COREALIZATION v3.105
\subsection{co-realized \texorpdfstring{$B-L=2$}{B-L=2} neutral boundary pair}
\label{p5:sec:v3105-bl2-boundary-pair-corealization}

\begin{definition}[Rooted neutral pair member with its Weyl/CAR type]
\label{p5:def:v3105-rooted-bl2-pair-member}
On the constituted neutral-right one-particle factor
$\mathcal K_{N,R}=S_R\otimes G\cong\mathbb C^2\otimes\mathbb C^3$,
let $\varepsilon_R\in\Lambda^2 S_R^*$ be the determinant-one right-Weyl
alternating form, normalized in an oriented spin frame by
$\varepsilon_R=\left(\begin{smallmatrix}0&1\\-1&0\end{smallmatrix}\right)$.
For a symmetric family coefficient $M\in\operatorname{Sym}^2G^*$, define
\[
 \iota_{\rm pair}(M)(s\otimes g,t\otimes h)
 :=\varepsilon_R(s,t)M(g,h),\qquad
 \Delta_{BL2}^{\rm geom}:=\iota_{\rm pair}(M_{BL2}^{\rm geom}).
\]
In a family-first interleaved basis,
$\Delta_{i\alpha,j\beta}=M_{ij}(\varepsilon_R)_{\alpha\beta}$.
It is a nonzero coefficient in $\Lambda^2\mathcal K_{N,R}^*$, of source
weight $\Delta(B-L)=+2$ and $Y=0$; reversal supplies its adjoint $-2$
coordinate. This defines the algebraic support insertion, not by itself an
action-owner occurrence or normalization.
\end{definition}
\begin{proposition}[Faithful neutral-Weyl pairing intertwiner]
\label{p5:prop:v382-majorana-car-spinor-intertwiner}
The map $\iota_{\rm pair}:\operatorname{Sym}^2G^*\to
\Lambda^2(S_R\otimes G)^*$ is injective, has skew image, and is
covariant under determinant-one right-spin transformations. On the rooted
full-rank coefficient its rank is six, and in the normalized spin frame
$\Delta\Delta^\dagger=\frac12 I_6$.
\end{proposition}
\begin{proof}
Interchanging the arguments changes the sign of $\varepsilon_R$ while
$M$ is symmetric. The entries $\Delta_{i\uparrow,j\downarrow}=M_{ij}$
recover $M$ and prove injectivity. The spin identity
$U^T\varepsilon_R U=\varepsilon_R$ for $U\in SL(S_R)$ proves covariance.
The displayed rooted diagonal matrix has
$MM^\dagger=\frac12I_3$, while
$\varepsilon_R\varepsilon_R^\dagger=I_2$; the full tensor coefficient
therefore has rank six and the stated singular-value identity.
Its CAR pair operator is nonzero because the coefficient is skew.
\end{proof}


\begin{theorem}[K16 neutral pair-cell co-realization]
\label{p5:thm:v3105-k16-bl2-pair-corealization}
On the already constituted neutral-right matter support, combine:
\begin{enumerate}[label=(\roman*),leftmargin=2em]
\item the occurring rooted Hopf--axial/family parent that generates
$M_{BL2}^{\rm geom}$;
\item the predeclared finite fermionic pairing coordinate of Paper~IV;
\item the closed K16 direct particle/internal-boundary owner; and
\item the owner-preserving support insertion
$\iota_{\rm pair}$ of Definition~\ref{p5:def:v3105-rooted-bl2-pair-member}.
\end{enumerate}
The mathematical tensor insertion is fixed by Proposition~\ref{p5:prop:v382-majorana-car-spinor-intertwiner}. The following same-owner physical conclusion additionally requires the action-owned K16 boundary functional to consume this particular skew coefficient, with its licensed insertion and coefficient; that consumption is not implied merely by the CAR tensor type.
Then the complete rooted pair cell
\[
 \mathfrak C_{BL2}
 =\bigl(\Delta_{BL2}^{\rm geom},
        (\Delta_{BL2}^{\rm geom})^\dagger\bigr)
\]
is a generated nonredundant input of the occurring K16 neutral quadratic
boundary owner \emph{when the stated same-owner consumption certificate holds}.  Paper~II-A,
Theorem~\ref{p2a:thm:v3109-generated-member-owner-corealization}, and
Paper~IV, Corollary~\ref{p4:cor:v4101-generated-pair-member-handoff}, therefore
give the same-owner boundary pairing
\[
 \boxed{
 H_{\partial,BL2}
 =\frac12\bigl(
  \mathfrak a^\dagger\Delta_{BL2}^{\rm geom}\mathfrak a^\dagger
 +\mathfrak a(\Delta_{BL2}^{\rm geom})^\dagger\mathfrak a
 \bigr)}
\]
as a co-realized required constituent.  Consequently
\[
 \boxed{
 \texttt{BL2\_BOUNDARY\_ACTION\_OCCURRENCE}
 =\texttt{CONDITIONAL\_OWNER},}
\]
\[
 \boxed{
 \texttt{BL2\_ACTION\_NORMALIZATION}
 =\texttt{CONDITIONAL\_NORMALIZATION}.}
\]
There is no second occurrence token and no independent multiplicative
pairing efficiency.
\end{theorem}

\begin{proof}
Theorem~\ref{p5:thm:v3103-generated-bl2-family-coefficient} supplies a
specific nonzero rooted family coefficient rather than an unresolved family
member.  The neutral-right pairing support and the finite quadratic pairing
coordinate are declared before that value is evaluated.  The K16 owner is independently constituted, but its consumption of the specific
skew tensor is an additional typed same-owner premise; this proof derives
only the asserted handoff conditional on that occurrence certificate.  Generated co-consumption retains the $+2$ and rooted-reversal
$-2$ siblings jointly.  Replacing the generated pair by zero deletes a
required nonredundant constituent; replacing it by
$\lambda\Delta_{BL2}^{\rm geom}$ changes the generated normalization unless a
separate deformation is generated.  Neither operation is the same owner
presentation.
\end{proof}

\begin{remark}[Pair ownership and the consumed normal-current word]
\label{p5:rem:v2167-pair-current-scope}
Faithful pairing support and its generated family coefficient certify the
nonzero algebraic pair; physical K16 consumption and its normalization
retain the same-owner condition of Theorem~\ref{p5:thm:v3105-k16-bl2-pair-corealization}. The actual normal-current handoff on the same field identity
is a different consumed map: it must retain the mixed return identified in
Paper~V-C, Corollary~\ref{p5c:cor:v2167-cpbl2-handoff-frontier}.
The pair certificate alone does not assert current-algebra intertwining.
This scope adds no second pair occurrence, state preparation or efficiency.
\end{remark}


\begin{corollary}[Dirac-only no-go becomes a comparison subbranch]
\label{p5:cor:v3105-bl2-dirac-no-go-scope}
The historical Dirac-only branch of
Corollary~\ref{p5:cor:v3103-bl2-structural-vs-action} remains a valid
charge-preserving comparison branch, but its open-occurrence status is superseded only on a K16 branch with
verified same-owner consumption as specified in
Theorem~\ref{p5:thm:v3105-k16-bl2-pair-corealization}.  On that certified branch the boundary
owner contains a nonzero $|\Delta(B-L)|=2$ pair coordinate, so the
charge-preserving-parent hypothesis used by the Dirac-only no-go does not hold
on this branch.
\end{corollary}


% BEGIN PROBLEM-4-K16-ACTION-HYPERCHARGE v3.188
\begin{theorem}[Action-consumed Hopf--axial quadratic selects the finite conserved hypercharge]
\label{p5:thm:v270-k16-action-hypercharge-ward}
On the already occurring K16 neutral-right pair owner of
Theorem~\ref{p5:thm:v3105-k16-bl2-pair-corealization}, let
\(\Delta=\Delta_{BL2}^{\rm geom}\ne0\),
\(A_\Delta=\mathfrak a^\dagger\Delta\mathfrak a^\dagger\), and
\(H_{\partial,BL2}=(A_\Delta+A_\Delta^\dagger)/2\).
Use the already generated candidate Abelian plane
\(\mathfrak t_{\rm cand}=\RR T_R^3\oplus\RR(B-L)\), acting
on the same neutral-right CAR carrier by second quantization.
For \(X=aT_R^3+b(B-L)\), the one-particle support
\(W_{R+}\otimes\ell\otimes G\) has charge \(a/2-b\), hence
\begin{equation}\label{p5:eq:v270-k16-torus-commutator}
\boxed{[d\Gamma(X),A_\Delta]=(a-2b)A_\Delta,\qquad
[d\Gamma(X),H_{\partial,BL2}]=\frac{a-2b}{2}
(A_\Delta-A_\Delta^\dagger).}
\end{equation}
Consequently the connected candidate-Abelian stabilizer of the
\emph{actual fixed coefficient in this occurring finite action} is exactly
\[
\boxed{\mathfrak{stab}_{\mathfrak t_{\rm cand}}H_{\partial,BL2}
=\RR\bigl(T_R^3+\tfrac12(B-L)\bigr)=\RR Y.}
\]
On the finite Ward domain where the normal part \(H_0\) commutes with
\(d\Gamma(X)\), the same action gives the exact pair-charge readout
\[
\boxed{\frac{d}{dt}\langle d\Gamma(X)\rangle_t
=\frac{a-2b}{\hbar}\operatorname{Im}\langle A_\Delta\rangle_t.}
\]
Thus the owned pair cannot source \(Y\), but can source \(B-L\)
with pair charge \(-2\). If \(H_0\) is number-preserving but does not
commute with every candidate \(X\), its commutator cannot cancel the
fermion-number-\(\pm2\) pair term; the pair still forces
\(a-2b=0\). On a \(Y\)-covariant normal companion domain, the complete
owner's connected candidate-Abelian stabilizer is therefore \(\RR Y\).
This is a new \emph{observable-level} converse and same-owner consistency
certificate for the structural PCAD selection of
Proposition~\ref{p5:prop:v187-pcad-yfix}. It does not claim a bulk
constant Majorana mass, a selected state or history, or a unique
full-spacetime reduction of every unreduced candidate-torus connection.
\end{theorem}
\begin{proof}
On \(W_{R+}\) one has \(T_R^3=+1/2\), and on
\(\ell\) one has \(B-L=-1\), while the candidate torus acts trivially
on \(G\). The CAR identity
\([d\Gamma(X),\mathfrak a_i^\dagger]=(a/2-b)\mathfrak a_i^\dagger\)
therefore gives the two identities
\eqref{p5:eq:v270-k16-torus-commutator} term by term.
The faithful predeclared pairing insertion and nonzero generated
\(M_{BL2}^{\rm geom}\) guarantee \(A_\Delta|0\rangle\ne0\),
whereas \(A_\Delta^\dagger|0\rangle=0\). Applying the second commutator
to the Fock vacuum proves that its vanishing forces \(a-2b=0\);
the converse is immediate. Paper~IV,
Theorem~\ref{p4:thm:v4102-pair-charge-source}, with
\(q_{\rm pair}=a-2b\), proves the Ward identity. Finally, a
number-preserving normal term cannot cancel a number-changing
\(\pm2\) operator. No second coefficient, pointing, phase selection,
or action occurrence is introduced.
\end{proof}

\begin{proposition}[Boundary stabilizer does not eliminate a bulk candidate direction]
\label{p5:prop:v270-boundary-bulk-torus-firewall}
The fixed K16 pair reduces the \emph{simultaneous finite action}
connected candidate-Abelian stabilizer to \(U(1)_Y\), but the same
statement does not imply that all local candidate-torus gauge
fluctuations on an unrelated spacetime chart take values in
\(\RR Y\). If the pair is supported on a proper boundary locus
\(S\subsetneq U\), any smooth gauge transformation
\(g(x)=\exp(f(x)X_\perp)\) with
\(f\in C_c^\infty(U\setminus S)\) and
\(X_\perp\notin\RR Y\) preserves the fixed boundary pair while acting
nontrivially on the unreduced bulk candidate connection. The separately
constituted local \(Y\)-reduced connection of
Theorem~\ref{p5:thm:canonical-full-joint-connection} and its currents
remain exact; a claim of \emph{necessity} for the full bulk family
requires a licensed spacetime comparison/connection port, rather than
this finite boundary theorem alone.
\end{proposition}
\begin{proof}
The displayed transformations equal the identity on \(S\), so preserve
the pair; the unreduced bulk Dirac and gauge action is covariant under
the full candidate torus and hence admits these nontrivial interior
transformations. This constructs the required counterexample. A
separately restricted \(Y\)-bundle is a valid derived descendant,
but restricting the whole unreduced affine family is not a consequence
of the K16 pairing commutator.
\end{proof}
% END PROBLEM-4-K16-ACTION-HYPERCHARGE v3.188

\begin{remark}[Boundary pairing, not terminal Majorana mass]
\label{p5:rem:v3105-boundary-not-terminal-majorana}
The theorem does not promote a constant bulk $D_M$ into the local fermion mass
operator and does not alter the already generated terminal neutral masses.
It closes the mathematical K16 \emph{internal-boundary pairing} tensor
on the constituted neutral-right support; action occurrence is conditional
on the stated same-owner consumption certificate.  Type-I/seesaw/Takagi mass dictionaries
remain conditional external comparison structures unless separately invoked.
\end{remark}

\begin{corollary}[New baryogenesis frontier]
\label{p5:cor:v3105-bl2-new-frontier}
After pair occurrence is certified, the next native baryogenesis obligation is
phase/CP ownership:
\[
 \boxed{
 \texttt{BL2\_PHASE\_LOCK}=\texttt{OPEN\_OWNER\_BINDING}.}
\]
The constant complex coordinate of $M_{BL2}^{\rm geom}$ is not by itself a
CP-odd trajectory.  The density-matrix/washout history and final $Y_{B-L}$
remain downstream.
\end{corollary}
% END BL2-BOUNDARY-PAIR-COREALIZATION v3.105

\subsubsection{Derivation of \texorpdfstring{$B-L$}{B-L}}
\label{p5:sec:part1-BminusL}

The internal quaternionic split is already fixed as
\[
\ell\oplus C\cong \mathbf 1\oplus\mathbf 3.
\]
Any generation-blind abelian operator constant on the triplet block and on the singlet block has the form
\(
aP_C+bP_\ell
\).
The unique traceless splitter on the native \(3\oplus1\) block is obtained from
\[
3a+b=0.
\]
Trace zero fixes only the eigenvalue ratio.  We choose the representative
whose triplet eigenvalue is the inverse triplet multiplicity,
\(a=1/3\), \(b=-1\).  Hence
\begin{equation}
\label{p5:eq:part1-BminusL}
B-L=1_W\otimes\left(\frac13\,P_C-P_\ell\right)\otimes 1_G.
\end{equation}

\begin{lemma}[Unique quark--lepton splitter direction]
\label{p5:lem:part1-BminusL-unique}
The traceless generation-blind block-scalar subspace on \(C\oplus\ell\) is
the one-dimensional line
\(\RR(\frac13P_C-P_\ell)\).  The displayed vector fixes a normalization and
the convention that the triplet block is positive; trace zero alone fixes
neither overall scale nor sign.
\end{lemma}

\begin{proof}
The space of such operators is two-dimensional before imposing trace zero:
\(aP_C+bP_\ell\).  Since \(\dim C=3\) and \(\dim\ell=1\), tracelessness is exactly
\(3a+b=0\).  The solution space is therefore one-dimensional.  The
normalization \(a=1/3\) makes the summed triplet contribution one unit, and
the sign convention selects the displayed representative.
\end{proof}

\begin{theorem}[Conditional Pati--Salam parent completion]
\label{p5:thm:pati-salam-parent-carrier}
Let the internal quark--lepton carrier be
\[
F_{\rm int}=\ell\oplus C,
\qquad
\dim_{\mathbb C}\ell=1,
\qquad
\dim_{\mathbb C}C=3.
\]
If one separately requires the full skew-Hermitian off-diagonal completion
\[
\mathfrak m=
\left\{
\begin{pmatrix}
0&-B^\dagger\\ B&0
\end{pmatrix}:B\in\operatorname{Hom}(\ell,C)
\right\},
\]
together with the split stabilizer, then their Lie closure is
\(\mathfrak{su}(F_{\rm int})\) and the minimal connected parent is
\[
SU(F_{\rm int})\cong SU(4)_c.
\]
The stabilizer of the splitting is
\[
S(U(C)\times U(\ell))\cong S(U(3)\times U(1)),
\]
with Lie algebra
\[
\mathfrak{su}(3)_c\oplus\mathfrak u(1)_{B-L}.
\]
The abelian generator is the traceless quark--lepton splitter
\[
B-L=\frac13P_C-P_\ell.
\]
Combining this optional completion with the diagonal-fibre left--right carrier
gives the local Pati--Salam parent \cite{PatiSalam1974}.  On the standard
chiral carrier its faithful connected global form is
\[
\frac{SU(2)_L\times SU(2)_R\times SU(4)_c}
{\langle(-1_2,-1_2,-1_4)\rangle}.
\]
The off-diagonal completion and its gauging are additional parent-role data,
not consequences of GP2 or of the split stabilizer.
\end{theorem}

\begin{proof}
The block-diagonal stabilizer and \(\mathfrak m\) span every traceless
skew-Hermitian block matrix on \(\ell\oplus C\); their commutators close to
\(\mathfrak{su}(4)\).  Minimality follows because any connected
special-unitary parent containing all these directions has the same Lie
algebra.  On
\((\mathbf2,\mathbf1,\mathbf4)\oplus
(\mathbf1,\mathbf2,\mathbf4)\), the displayed diagonal central element acts
trivially, and the fundamental representations show that no larger central
kernel remains.
\end{proof}


\begin{corollary}[Diagonal-fibre source of the left--right carrier]
\label{p5:cor:diagonal-fibre-left-right-origin}
\label{p5:cor:born-antidiagonal-left-right-origin}
Paper~III's ratio fibration has diagonal vector fibre \(S^3_{\rm diag}\)
with connected spin-isometry carrier
\[
\Spin(4)_{\rm fib}\cong SU(2)^a_L\times SU(2)^a_R.
\]
The Born-odd subspace \(\mathfrak k_-\) is instead the tangent module of the
relative quotient and is not a gauge Lie algebra.  Under the generated
equivariant identification, the diagonal-fibre carrier supplies the finite
left--right factor.  Its kinetic realization and \(g_L=g_R\)
require Paper~III's explicit \(({\rm GRED})\) package.
\end{corollary}

\begin{proof}
This is exactly Paper~III's carrier-level gauge--gravity split together with
the symmetric-pair bracket relation
\([\mathfrak k_-,\mathfrak k_-]\subseteq\mathfrak k_+\).
\end{proof}

\subsubsection{Split-phase baryon number and perturbative proton stability}
\label{p5:sec:split-phase-proton-stability}

The Pati--Salam result above is a conditional parent-completion statement.
In the declared split phase the active gauge algebra is the stabilizer of the
\(C\oplus\ell\) decomposition, not the full off-diagonal \(SU(4)_c\) algebra.
This distinction is exactly what fixes the proton-stability status of the
minimal finite core.

Define the quark and lepton counting operators on the matter carrier by
\begin{equation}
\label{p5:eq:split-phase-BL-counting-operators}
B:=1_W\otimes \frac13 P_C\otimes 1_G,
\qquad
L:=1_W\otimes P_\ell\otimes 1_G,
\end{equation}
and extend them additively to multiparticle states.  The algebraic conjugate
representation carries the opposite charges; its physical-sector occurrence
is licensed on \(\mathsf R_5^{\rm PCT}\).  Then \(B-L\) is precisely the
traceless splitter \eqref{p5:eq:part1-BminusL}.

\begin{theorem}[Perturbative baryon stability for the declared split-phase vertices]
\label{p5:thm:split-phase-proton-stability}
Assume
\(\mathsf G_5^{\rm fin}
+\mathsf R_5^{\rm SSQFT}+\mathsf R_5^{\rm PCT}\).  For the declared
split-phase renormalizable field
content, every perturbative
interaction generated by
\[
\mathfrak g_{\rm split}
:=\mathfrak{su}(2)_L\oplus\mathfrak{su}(3)_c\oplus\mathfrak u(1)_Y
\]
together with the finite Dirac/Yukawa blocks preserves baryon number \(B\).
An optional Majorana block may violate \(L\) by two units, but it acts only on the
neutral lepton line and still has \(\Delta B=0\).  Consequently no perturbative
local operator generated by the minimal split-phase finite core can mediate proton decay.
\end{theorem}

\begin{proof}
The \(SU(3)_c\) factor acts inside \(C\), the \(SU(2)_L\) factor acts on the
Hopf/chiral factor \(W\), and the residual abelian generator is
\(Y=T_R^3+\frac12(B-L)\).  Hence every gauge generator in
\(\mathfrak g_{\rm split}\) commutes with \(P_C\), and therefore with \(B\).
The charged-fermion Yukawa blocks pair left and right modes inside the same
quark or lepton sector, so they also commute with \(B\).  The neutrino Dirac
block pairs lepton modes with lepton modes and commutes with \(B\).  The
Majorana block acts on the derived neutral right-handed lepton line; it violates
lepton number but contains no quark projector and therefore also commutes with
\(B\).  Since products and commutators of \(B\)-preserving vertices are again
\(B\)-preserving, every perturbative operator generated by the minimal
split-phase finite core has \(\Delta B=0\).  Proton decay has \(\Delta B=-1\),
so it is absent in this sector.
\end{proof}

\begin{corollary}[Status of off-diagonal parent exchange]
\label{p5:cor:offdiagonal-parent-exchange-nonminimal}
The only \(SU(4)_c\) gauge directions capable of changing quark number into
lepton number are the off-diagonal complement
\[
\mathfrak m_{4/(3+1)}
:=
\left\{
\begin{pmatrix}
0&-Z^\dagger\\ Z&0
\end{pmatrix}
:Z\in\operatorname{Hom}(\ell,C)
\right\}
\subset
\mathfrak{su}(4)_c .
\]
These directions are absent from the declared split stabilizer
\(S(U(3)\times U(1))\).  Therefore any proton-decay operator obtained by
integrating out a dynamical vector in \(\mathfrak m_{4/(3+1)}\) belongs to a
separate parent-completion role, not to the declared Standard--Model finite
core.  GP2--GP3 do not prove that such a role is impossible.
\end{corollary}

\begin{proof}
The block-diagonal stabilizer of \(C\oplus\ell\) is exactly
\(\mathfrak{su}(3)_c\oplus\mathfrak u(1)_{B-L}\).  The remaining generators of
\(\mathfrak{su}(4)_c\) are precisely maps from \(C\) to \(\ell\) and from
\(\ell\) to \(C\).  Those maps do not commute with \(P_C\), hence they do not
preserve \(B\).  Since the split phase retains only the stabilizer algebra, the
off-diagonal maps are absent unless a separate massive parent-completion field
is introduced.
\end{proof}

\begin{remark}[Higher-dimensional baryon-violating operators]
\label{p5:rem:higher-dimensional-proton-decay-nonminimal}
The theorem classifies the renormalizable split-phase vertices explicitly
declared here.  It does not assert that every gauge-invariant
higher-dimensional Standard--Model operator is impossible.  Operators such
as \(QQQL/\Lambda^2\), whether supplied directly or generated by a
parent-completion exchange field, lie outside that declared vertex class.
Nonperturbative electroweak effects are outside the perturbative statement as
well.
\end{remark}

\subsubsection{Geometric derivation of hypercharge}
\label{p5:sec:part1-hypercharge}

\begin{theorem}[Haar averaging and its exact scope]
\label{p5:thm:primitive-local-ray-singlet-criterion}
Let a compact group \(K\) act unitarily on a finite carrier by \(\rho\).
Then
\begin{equation}
\label{p5:eq:local-ray-invariant-projector}
\Pi_{\rm inv}= \int_K \rho(k)\,dk
\end{equation}
is the orthogonal projector onto the invariant subspace.  On every nontrivial
irreducible subrepresentation \(V_\lambda\),
\begin{equation}
\label{p5:eq:local-ray-projection-kills-nonsinglet}
\Pi_{\rm inv}|_{V_\lambda}=0 .
\end{equation}
This theorem does not imply that an elementary gauge-covariant field must be
an invariant ray, and it does not select a matter line, a real structure, or a
gauge representation.
\end{theorem}

\begin{proof}
Haar invariance makes \(\Pi_{\rm inv}\) self-adjoint and idempotent, and its
range is exactly the invariant subspace.  Schur orthogonality makes it zero
on a nontrivial irreducible representation.  The final sentence is a scope
statement: gauge-covariant elementary fields transform in representations and
need not lie in the invariant state subspace.
\end{proof}

\begin{corollary}[Generated distinguished right-lepton line]
\label{p5:cor:supplied-self-conjugate-line-neutrality}
\label{p5:lem:geometric-majorana-line}
The right-lepton block is
\begin{equation}
\label{p5:eq:local-fixed-right-lepton-block}
P_RP_\ell F_{\rm core}=P_{R+}P_\ell F_{\rm core}\oplus P_{R-}P_\ell F_{\rm core}.
\end{equation}
The oriented Hopf Cartan and axial quaternionic component therefore generate
the one-dimensional tensor module
\begin{equation}
\label{p5:eq:r0-local-fixed-line}
L_0:=P_{R+}P_\ell F_{\rm core}=W_{R+}\otimes\ell,
\end{equation}
with replicated carrier
\begin{equation}
\label{p5:eq:r0-family-support}
P_{R+}P_\ell\otimes G.
\end{equation}
The connected character kernel of this line consists of the generators that
act trivially on \(L_0\).  No choice of a vector, phase, or norm is needed.  No
Majorana bilinear or finite real structure is used.
\end{corollary}

\begin{proof}
This is the distinguished module used in
Theorem~\ref{p5:thm:exact-oriented-vector-hypercharge}.  The infinitesimal
kernel of a one-dimensional representation is the kernel of its character;
it is independent of a representative vector and requires no global Hopf
section.
\end{proof}

\begin{lemma}[Full diagonal centralizer before the scalar constraint]
\label{p5:lem:full-diagonal-centralizer}
\label{p5:lem:one-abelian-slot-derived}
On \(\mathsf G_5^{\rm fin}\), put
\[
P_{\sigma q}:=P_\sigma\otimes P_q\otimes1_G,
\qquad
\sigma\in\{L,R+,R-\},\quad q\in\{\ell,C\}.
\]
The generation-blind Hermitian operators that commute with
\(\mathfrak{su}(2)_L\oplus\mathfrak{su}(3)_c\) and preserve all six displayed
projectors are exactly
\[
\mathfrak d
=
\bigoplus_{\sigma,q}\RR P_{\sigma q}
\cong\RR^6.
\]
Consequently \(\dim(\mathfrak d/\RR1)=5\).  In particular,
\(P_L=P_{L\ell}+P_{LC}\) is an allowed direction and is not, even modulo the
identity, in \(\operatorname{span}\{T_R^3,B-L\}\).
\end{lemma}

\begin{proof}
Each \(P_{\sigma q}F\) is irreducible under the retained nonabelian action
once the six projectors are required to be preserved.  Schur's lemma makes a
commuting Hermitian operator scalar on every block.  Conversely every such
block-scalar operator commutes with the action and projectors.  Removing the
one common identity reduces the dimension from six to five.  Direct
comparison of block eigenvalues proves the final counterexample.
\end{proof}

\begin{lemma}[Abelian plane under coordinate-conjugate covariance]
\label{p5:lem:intrinsic-abelian-centralizer-plane}
On \(\mathsf G_5^{\rm fin}\), let
\(\mathfrak d_{\rm 1H}\subseteq\mathfrak d\) consist of those \(D\) for which
the modules
\[
\mathcal S_-=\operatorname{Hom}(W_{R-},W_L)\otimes1_Q\otimes1_G,
\qquad
\mathcal S_+=\operatorname{Hom}(W_{R+},W_L)\otimes1_Q\otimes1_G
\]
are one field and its conjugate and therefore have opposite \(D\)-weights.
Then
\[
\mathfrak d_{\rm 1H}
=\RR1\oplus\RR T_R^3\oplus\RR(B-L).
\]
Consequently
\begin{equation}
\label{p5:eq:intrinsic-abelian-centralizer-plane}
\mathfrak d_{\rm 1H}/\RR1
=\operatorname{span}\{[T_R^3],[B-L]\}.
\end{equation}
Equivalently, its traceless representatives lie in
\(\operatorname{span}\{T_R^3,B-L\}\).
\end{lemma}

\begin{proof}
Write the six block eigenvalues of \(D\) as
\[
x_{Lq},\quad x_{+q},\quad x_{-q},
\qquad q\in\{\ell,C\}.
\]
If the weight of \(\mathcal S_-\) is \(h\), one-field conjugacy gives
\[
x_{Lq}-x_{-q}=h,\qquad
x_{Lq}-x_{+q}=-h
\quad(q=\ell,C).
\]
Thus \(x_{+q}=x_{Lq}+h\) and \(x_{-q}=x_{Lq}-h\), leaving the three
parameters \(x_{L\ell},x_{LC},h\).  Put
\[
a=2h,\qquad
b=\frac34(x_{LC}-x_{L\ell}),\qquad
c=x_{L\ell}+b.
\]
Since \(B-L=-1\) on \(\ell\), \(1/3\) on \(C\), and \(T_R^3\) has right
eigenvalues \(\pm1/2\), direct substitution gives
\[
D=c1+aT_R^3+b(B-L).
\]
Conversely every such operator gives opposite scalar weights.  Both
\(T_R^3\) and \(B-L\) are traceless on \(F_{\rm core}\), while \(1\) is not,
so quotienting by \(\RR1\) is canonically represented by the traceless plane.
\end{proof}

% BEGIN M8 v1.78 SM01
% Results: B03-MAJORANA-Y, B04-COMPLETE-OWNER, B08-HOPF-PROJECTOR, B08-COLOR-PROJECTOR, B08-JOINT-PROJECTION, B08-PROJECTED-SEED, B08-RETENTION-ORDER, B08-PHASE-SENSITIVE-OWNER, B11-CHARGED-TARGET
\subsubsection{The joint comparison and its physical target}

The geometric comparison now exists upstream:
Paper I-A constructs \(m_{\mathcal O}\) on the full joint orbit,
and Paper II-C exports its complete geometric constructor word.
The former structural hypercharge residual is closed by Proposition~\ref{p5:prop:v187-pcad-yfix}; the remaining question is instead the physical
transport and retention.  A shared host does not identify the
quaternionic coordinate action with the axial/color action on
\(Q=\ell\oplus C\).

\begin{proposition}[Commuting color and Hopf projections]
\label{p5:prop:commuting-hopf-color-comparison-projections}
\label{p5:prop:joint-hopf-color-projection}
Let \(Z=\operatorname{diag}(1,-1)\) on \(W_R\), and let
\(\mathcal L_Z(S)=Z^TS+SZ\) on symmetric spinor forms.  Its
weight projectors are
\[
 \Pi_+=\mathcal L_Z(\mathcal L_Z+2)/8,\qquad
 \Pi_0=1-\mathcal L_Z^2/4,\qquad
 \Pi_-=\mathcal L_Z(\mathcal L_Z-2)/8.
\]
On \(\operatorname{Sym}^2Q^*\otimes\operatorname{Sym}^2W_R^*\),
color averaging is restriction of both \(Q\) arguments to \(\ell\).
It commutes with \(\Pi_+\).  Their joint image is the complex line
\[
 (L_0^*)^{\otimes2},\qquad L_0=W_{R+}\otimes\ell.
\]
In the compatible axial coordinates, with the parent unit
transported as in Paper I-A, the actual seed has projection
\[
 \frac{-1+i}{2}\,z_0^2s_+^2\ne0.
\]
This is a geometric component calculation.  If this component
has the physical target action, its \((T_R^3,B-L)\) weights are
\((-1,+2)\).
\end{proposition}
\begin{proof}
The entries \(S_{++},S_{+-},S_{--}\) have
\(\mathcal L_Z\)-eigenvalues \(2,0,-2\), so polynomial
interpolation gives the three projectors.
The color decomposition of \(\operatorname{Sym}^2Q^*\)
is \(\ell^{*2}\oplus(\ell^*\otimes C^*)\oplus
\operatorname{Sym}^2C^*\); only the first summand is an
\(SU(3)\) singlet.  The two projections act on separate
tensor factors and commute.  Evaluating the quadratic matrix
of Paper I-A at \(h=(1+i+j+k)/2\) gives the displayed coefficient.
Restriction to the axial arguments retains the \(+2\) dual
axial weight, and the positive spin dual square has weight
\(-1\).  Neither projection changes these target characters.
\end{proof}

\begin{proposition}[Complete retention and the quadratic Y criterion]
\label{p5:prop:complete-retained-quadratic-y-criterion}
\label{p5:thm:quadratic-incidence-yfix-criterion}
Work in the one-field/conjugate covariance domain of
Lemma~\ref{p5:lem:intrinsic-abelian-centralizer-plane}.
If a licensed physical comparison \(0\ne m\in(L_0^*)^2\)
is retained coherently, its infinitesimal character equation is
\[
 X=aT_R^3+b(B-L),\qquad a-2b=0.
\]
It has the same connected kernel as the distinguished-line
character, namely \(\RR Y\), \(Y=T_R^3+(B-L)/2\).
The full physical conclusion also requires that every other
retained companion preserve the proposed subgroup.
A nonzero \(M\in(L_0^*)^2\otimes\operatorname{Sym}^2G^*\)
has the same gauge equation if the family action is trivial.
\end{proposition}
\begin{proof}
The character of \(L_0\) is \(\chi_0(X)=a/2-b\).
The coefficient \(m\) transforms by \(e^{-2i\chi_0(X)}\);
differentiating its fixed-value equation gives \(a-2b=0\).
A nonzero family tensor with gauge-trivial family action has
the same character.  The automorphism group of a retained
diagram is the intersection of the stabilizers of all its
data, which proves the companion qualification.
\end{proof}

In particular, retaining a projected composite as an output and
retaining the full seed before reading its component are different
diagrams.  Fixed charged factors can impose two independent phase
conditions even when their mixed product imposes only one.
For a line vector \(n\), invariance of
\(\operatorname{Re}m(n,n)\) for all \(n\) requires
\(e^{2i\theta}=1\); invariance of \(|m(n,n)|^2\) permits every
phase.  Averaging over the full parent torus annihilates \(m\).
Thus norms, a moving covariant coefficient, or a component of an
overmarked seed do not by themselves establish a physical member occurrence; structural hypercharge is supplied separately by the parent-relation PCAD theorem.
The explicit failed target substitutions are collected in
Appendix~\ref{p5:app:joint-incidence-physical-port-tests}.
% END M8 v1.78 SM01


% BEGIN M8 v1.78 SM02
% Results: B03-TORUS, B07-CONJUGATE-TYPING, B09-Y-KERNEL, B11-TRACE-COMPLEMENT, B11-ORBIT-FORM, B11-CANDIDATE-GLOBAL-KERNEL, B11-PRIMITIVE-CHARACTER, B11-INTEGER-LATTICE, B11-QUADRATIC-STABILIZER, B11-POWER-COMPONENTS, B11-CLOCK-DISTINCTION, B12-PHYSICAL-DETERMINANT-CHARACTER, B12-RELATIVE-STABILIZER, B13-CENTRE-PARITY
\subsubsection{Electric charge}
\label{p5:sec:part1-electric-charge}

With the weak third component \(T_L^3\) acting on \(W_L\),
\[
T_L^3=T_L^3\otimes 1_{\ell\oplus C}\otimes 1_G,
\]
the structurally generated hypercharge branch gives the conventional charge
operator
\begin{equation}
\label{p5:eq:part1-Q}
\boxed{Q=T_L^3+Y=T_L^3+T_R^3+\frac12(B-L).}
\end{equation}
The load-bearing electroweak stabilizer is fixed directly by the occurring
chiral source rather than by postulating an independently occurring scalar
vacuum orbit.

\begin{theorem}[Electromagnetic stabilizer of an occurring chiral solder]
\label{p5:thm:v377-occurring-solder-em-stabilizer}
Assume the generated Standard--Model charge representation and an actually
occurring nonzero odd zero-order source member in the Paper--V-A source role.
By Theorem~\ref{p5:thm:v375-d4-triality-source-consumption}, its primitive
chiral factor contains the generated unit solder \(c(\hat h)\).  Use
Theorem~\ref{p5:thm:v377-unit-solder-normalized-scalar} to write its normalized
weak columns as \(\widetilde{\widehat H}\) and \(\widehat H\).

Put \(U_H:=\Phi(\widehat H)\in SU(2)\) and define the source-aligned
weak Cartan and electromagnetic generator by
\[
 T_{L,H}^3:=U_HT_L^3U_H^\dagger,
 \qquad Q_H:=T_{L,H}^3+Y.
\]
For a general electroweak generator
\[
 X=x_1T_L^1+x_2T_L^2+x_3T_L^3+bY,
\]
invariance of any nonzero allowed left--right source slot is equivalent to
\[
 \sum_{a=1}^3x_aT_L^a=bT_{L,H}^3.
\]
Hence
\[
 \boxed{\mathfrak{stab}_{\rm EW}=\mathbb R Q_H
 =\operatorname{Ad}_{U_H}\mathfrak u(1)_Q.}
\]
In the source-aligned weak frame \(U_H^\dagger\widehat H=(0,1)^T\),
this is exactly \(x_1=x_2=0\), \(x_3=b\), and \(Q_H\) has the
diagonal representative \(Q=T_L^3+Y\).  The frame change transports
the complete source/current tuple; it is not a new source pointing.
Since color acts trivially on the solder and as the identity on the quark
mass multiplicity,
\[
 \boxed{\mathfrak{stab}_{\rm SM}
 =\mathfrak{su}(3)_c\oplus\mathbb R Q_H.}
\]
No scalar potential, scalar magnitude, stationary point, or vacuum selector
is consumed.
\end{theorem}

\begin{proof}
The unit-solder theorem gives
\(U_H(0,1)^T=\widehat H\) and
\(U_H(1,0)^T=\widetilde{\widehat H}\).
Conjugate the complete source and weak generators by \(U_H^\dagger\).
Hypercharge commutes with this weak change of frame.  In that aligned frame
the off-diagonal weak generators \(T_L^{1,2}\) rotate a nonzero typed weak
source column into the orthogonal weak component, while hypercharge acts
diagonally on the corresponding right singlet; hence invariance forces
\(x_1=x_2=0\).  The four allowed chiral charge slots then give
\[
\begin{aligned}
 u:&\quad \frac{x_3}{2}+\frac b6=\frac{2b}{3},\\
 d:&\quad-\frac{x_3}{2}+\frac b6=-\frac b3,\\
 \nu:&\quad\frac{x_3}{2}-\frac b2=0,\\
 e:&\quad-\frac{x_3}{2}-\frac b2=-b.
\end{aligned}
\]
Each nonzero slot independently yields \(x_3=b\) in this aligned frame.
Conjugating back proves that the stabilizer is \(\mathbb R Q_H\).
Equivalently, writing \(T_L^a=\sigma_a/2\) gives
\(T_{L,H}^3=\frac12\sum_an_a\sigma_a\), where
\(n_a=-\widehat H^\dagger\sigma_a\widehat H\); thus
\(x_a=bn_a\), rather than \(x_1=x_2=0\) in an arbitrary weak frame.
Color commutes with the solder and acts only on the
already generated quark multiplicity, giving the displayed Standard--Model
stabilizer.
\end{proof}

\begin{corollary}[Conventional scalar-coordinate presentation of the same stabilizer]
\label{p5:prop:part1-em-stabilizer}
\label{p5:thm:part1-em-stabilizer}% compatibility alias with environment prefix
If one chooses the conventional lower-component representative of the already
generated normalized solder line and transports every current/source with
that choice, its scalar-coordinate stabilizer is
\(Q=T_L^3+Y\).  This is a coordinate presentation of
Theorem~\ref{p5:thm:v377-occurring-solder-em-stabilizer}, not an independent
nonzero-vacuum premise.
\end{corollary}

\begin{proposition}[Invariant polynomial for the generated scalar]
\label{p5:cor:finite-core-ewsb-topology}
Assume
\(\mathsf G_5^{\rm fin}\) and impose strict
large-field stability on the four-dimensional degree-\(\le4\) scalar
potential.  The full polynomial invariant ring for one complex weak doublet
is generated by \(H^\dagger H\).  Its real degree-\(\le4\) potential
therefore has, up to an additive constant, the form
\begin{equation}
\label{p5:eq:finite-core-scalar-potential}
V(H)=-\mu_H^2 H^\dagger H+\lambda_H(H^\dagger H)^2,
\qquad \lambda_H>0.
\end{equation}
This proposition does not select \(\mu_H^2>0\), a nonzero scalar orbit, or a
physically unbroken subgroup.
\end{proposition}

\begin{proof}
The first fundamental invariant of one \(SU(2)\) doublet with nonzero
hypercharge is its Hermitian norm \(H^\dagger H\); every polynomial invariant
is a polynomial in that norm.  Truncation at canonical degree four leaves its
linear and quadratic powers.  Strict large-field stability makes the
quartic coefficient positive.
\end{proof}

\begin{remark}[Papers--V-B--V-C native electroweak handoff]
\label{p5:rem:finite-core-ewsb-handoff}
Paper~V-A fixes the generated scalar representation, the normalized
Clifford/solder direction, and the exact source stabilizer
\(Q=T_L^3+Y\).  Paper~V-B places the conventional scalar comparison coordinate
on the native \(E_{32}\) hierarchy owner and treats it as a
readout/normalization coordinate, not as a separately generated physical
vacuum.  Paper~V-C may study scalar stationary/effective-action branches for
claims which explicitly require equilibrium, scalar Hessians, or cosmological
history; those branches are not prerequisites for the native electroweak
stabilizer or gauge-mass kernel.  In particular no theorem-level identity
\(v=8\rho_VE_{32}\) is used as the cause of electroweak reduction.
\end{remark}


\section{Chiral descent, anomalies and global spin conditions}
\label{p5a:sec:organized-chirality}
\subsection{Chirality ancestry and the member-occurrence boundary}
\label{p5:sec:g4-update}

The upstream diagonal-fibre ancestry removes the former relative-chirality
selector.  An odd external/internal bridge would require reversing orientation
on exactly one owner and is not another realization of the generated graph.
Consequently every theorem in this paper which used $({\rm CHI}_5)$ only to
select $P_+$ may now consume the derived relation
\eqref{p5:eq:chi5-derived} instead.

Structural hypercharge is already closed by parent-relation PCAD.  Separately,
the nontrivial torus character on $L_0$ admits no nonzero vector natural under
the unbroken parent torus:
if $v$ were such a vector, $t\cdot v=v$ for all $t$ would imply
$(\chi_0(t)-1)v=0$, hence $v=0$ for some $t$ with $\chi_0(t)\ne1$.
Thus a genuine nonzero background/field/member occurrence remains a separate
G0/member statement in its actual target sort, while the connected structural
hypercharge stabilizer is already generated from the parent relation by PCAD.


\subsection{The generated representation fixes the anomalous export coefficients}
The three-family matter content of V-A,
\texttt{p5:prop:three-generation-content}, gives the left-chiral charge
table below. Right-handed fields are represented by their conjugates for
this coefficient calculation; this is no physical trace doubling.
\[
\begin{array}{c|r|r|r|r|r}
 \hbox{slot}&\hbox{dimension}&Y&B&L&\hbox{weak copies}\\\hline
 Q_L&6&1/6&1/3&0&3\\
 L_L&2&-1/2&0&1&1\\
 u_R^c&3&-2/3&-1/3&0&0\\
 d_R^c&3&1/3&-1/3&0&0\\
 e_R^c&1&1&0&-1&0\\
 \nu_R^c&1&0&0&-1&0
\end{array}
\]
With $T(\mathbf2)=1/2$, define
$\mathcal A_{X,2}=\sum X T(\mathbf2)$, retaining the color and family
multiplicities, and $\mathcal A_{X,Y^2}=\sum \dim(\hbox{slot})XY^2$.
Direct substitution gives the exact finite coefficients
\begin{equation}
 \begin{array}{c|rr}
 X&\mathcal A_{X,2}&\mathcal A_{X,Y^2}\\\hline
 B&3/2&-3/2\\
 L&3/2&-3/2\\
 B-L&0&0
 \end{array}.
 \label{p5a:v2171:eq:exportcoefficients}
\end{equation}
The color and mixed gravitational coefficients for these vector charges
vanish with the displayed right-neutral slot retained.
For example per family
$\mathcal A_{B,Y^2}=1/18-4/9-1/9=-1/2$ and
$\mathcal A_{L,Y^2}=1/2-1=-1/2$.

Thus, on a licensed regulated quantum anomaly domain, baryon and lepton
export have the same gauge source, while this contribution to $B-L$
cancels. Both the weak and hypercharge gauge terms must be retained when
they occur. The finite coefficient does not select a topological history,
fix an anomalous current expectation, or supply a sphaleron rate.
VII-C, \texttt{p12:prop:v248-finite-quantum-ward-antecedent} and
\texttt{p12:cor:v236-ward-dependent}, require the actual regulated identity,
including its Jacobian, contact and regulator terms, before transport.

For a unit weak topological transition with the anomaly/spectral-flow
bridge actually established, each of the $3+1$ weak doublets per family
has one index unit. In the orientation chosen to produce particles,
\begin{equation}
 (\Delta B,\Delta L,\Delta(B-L))=(3,3,0).
 \label{p5a:v2171:eq:unitexport}
\end{equation}
This is an evaluated selection rule. It does not assert that one such
transition has occurred. The familiar schematic twelve-fermion selection
word has nine quark and three lepton slots, with each family carrying
hypercharge $3(1/6)-1/2=0$. A licensed insertion representing it would
have charges $(3,3,0)$; its coefficient and time history cannot be chosen.


\begin{corollary}[Finite baryon retention and the export requirement]
\label{p5a:v2171:cor:baryon-export-requirement}
On the complete finite baryon-equivariant split-plus-BL2 descendant of
Theorem~\ref{p3f:v2171:thm:baryonstationary}, the inherited neutral root has
$q_B(t)=0$. A nonzero baryon readout therefore consumes a generated
physical baryon-export law outside that theorem's conservation hypotheses,
such as the actual regulated anomalous current on its licensed domain.
\end{corollary}
\begin{proof}
The lepton pair has zero baryon charge. Apply the equivariant stationary
descent and the complete charge balance of Theorem~\ref{p4:v2171:thm:fullcharge}.
\end{proof}
The finite coefficient table fixes relative representation weights; it
does not evaluate an anomaly history or a sphaleron rate.
Appendix~\ref{p5a:v2171:app:gaugehistoryscope} records why the Cartan
midpoint and the optional Euclidean kink do not supply that physical history.

% END INLINE SOURCE: Paper5A-VectorChargeExport-Main-v1_0.tex


% BEGIN NEUTRAL-CORE-FAMILY-TYPING v2.208
\subsection{Literature-verified global anomaly and phase-transport domains}
\label{p5:sec:v2210-lit-global-domain}
The finite representation, chirality ancestry, hypercharge line, faithful
global group and charged source closure above are retained.  This section
narrows a quantum export obligation; it does not constitute a new matter
carrier, fermion source, full weak-sector CP symmetry or regulator.

\begin{theorem}[Ordinary spin--\(Z_6\) global gauge-anomaly obstruction]
\label{p5:thm:v2210-lit-anomaly-obstruction}
Let the claimed four-dimensional quantum export have precisely the faithful
connected group
\[
 G_Q=\bigl(SU(3)\times SU(2)\times U(1)\bigr)/\mathbb Z_6
\]
of Proposition~\ref{p5:prop:faithful-global-sm-group}, with the same quotient
embedding and representations.  Suppose its background category consists
of orientable spin manifolds with principal \(G_Q\)-bundles, and its
regulated chiral anomaly is the ordinary Dai--Freed anomaly functional.
Assume the full occupied chiral content has vanishing local anomaly
polynomial; for the constituted three-generation carrier the relevant
representation coefficients are already those of
Theorem~\ref{p5:thm:part1-anomaly-cancellation}.
Then the bordism obstruction to this ordinary global gauge anomaly vanishes:
\[
 \boxed{\Omega_5^{\rm Spin}(BG_Q)=0,
 \qquad \mathcal A_Q(X)=1\quad
 \text{for every closed spin five-background }X.}
\]
This conclusion does not construct a native local fermion measure or
identify its regulator with the stated anomaly functional.
\end{theorem}
\begin{proof}
Davighi--Gripaios--Lohitsiri compute the displayed bordism group for this
exact quotient in Section~4.5 and Table~1
\cite{LitSMGlobalAnomaly2020}.  On the declared anomaly domain, vanishing
of the local anomaly polynomial makes the exponentiated eta anomaly a
bordism-invariant character
\(\mathcal A_Q:\Omega_5^{\rm Spin}(BG_Q)\to U(1)\).  A character of the
zero group has only the value one.  The generated group and coefficient
theorems supply the finite identifications used here; the stated background
category and regulated anomaly identification are the exact remaining
export-domain checks.  No \(SU(5)\) completion or additional scalar
expectation is used.
\end{proof}

\begin{proposition}[Path-independent determinant-line transport criterion]
\label{p5:prop:v2210-lit-measure-transport}
On one connected component of the actual occupied gauge-background space
modulo gauge transformations, suppose the regulated determinant line
\(\mathcal L\) has a unitary connection with zero curvature and trivial
holonomy on every closed loop.  Then an assigned unit vector in one fibre
extends to a unique parallel unit section.  Changing that one reference
vector multiplies the section by a constant phase; all ratios of the section
under background transport are unchanged.
\end{proposition}
\begin{proof}
Transport the reference vector along a path to the requested background.
Two such paths differ by a closed loop and give the same answer by trivial
holonomy.  Parallel transport preserves its norm and supplies local smooth
parallel sections; these agree on overlaps.  Uniqueness follows by transport
from the reference fibre.  A second reference differs by one constant unit
scalar, which cancels from every transport ratio.
\end{proof}

\begin{corollary}[Scope of the global-anomaly narrowing]
\label{p5:cor:v2210-lit-anomaly-remainder}
On the exact domain of
Theorem~\ref{p5:thm:v2210-lit-anomaly-obstruction}, no separate ordinary
spin--\(G_Q\) global-anomaly obstruction remains.  The physical measure
claim still consumes the identification of its determinant-line transport,
locality, regulator contacts and quantum Ward identities.  Its normalization
and weights on disconnected physical sectors are not selected by
Proposition~\ref{p5:prop:v2210-lit-measure-transport}.  A boundary, an
additional gauge/discrete symmetry, a pin or spin--gauge background, or a
different anomaly functional is outside this theorem and retains its own
obstruction calculation.
\end{corollary}
The electroweak reconstruction theorem of Kadoh--Kikukawa
\cite{LitSMChiralMeasure2008} offers a separate conditional method: its
local currents must satisfy smoothness, locality, covariance and local/global
integrability, on its admissible \(SU(2)\times U(1)\) lattice sectors with
vanishing \(U(1)\) magnetic flux.  The native finite Dirac operator and
faithful quotient do not automatically supply that regulator or its
configuration space.  This citation therefore supplies a reconstruction
route, not a discharged native-measure premise.

\begin{proposition}[Complete determinant-phase transport test]
\label{p5:prop:v2210-lit-strong-phase-jet}
Let \(D_u(t),D_d(t)\) be differentiable full-rank quark-source matrices
along an actual declared transport.  In a continuous local branch of the
determinant phase,
\[
 \boxed{\frac{d}{dt}\arg\det(D_uD_d)
 =\operatorname{Im}\operatorname{Tr}
 \bigl(D_u^{-1}\dot D_u+D_d^{-1}\dot D_d\bigr).}
\]
In particular, if the complete transported rooted source retains conjugate
orientation-core determinants and positive polar determinants, its product
determinant stays positive and this phase derivative vanishes.  This
criterion tests the complete source derivative, not an after-the-fact
anomalous chiral rephasing.
\end{proposition}
\begin{proof}
Jacobi's determinant formula gives
\(d\log\det D/dt=\operatorname{Tr}(D^{-1}\dot D)\).
Add the two identities and take the imaginary part.  Under the additional
rooted relation, the determinant product is a product of positive polar
determinants times \(|\det A_+(t)|^2\), so its phase is constant zero.
All contacts in the actual \(\dot D\) must be retained before this test.
\end{proof}

\begin{remark}[Finite strong CP is retained; the quantum transport is tested]
\label{p5:rem:v2210-lit-finite-cp-retained}
Paper~V-C already fixes the finite native sector character to one and the
rooted quark determinant orientation to be positive.  The quantum extension
requires transport of that phase together with its actual regulated
Jacobian/measure and effective-action terms.  Neither a nonzero weak-current
quartet nor absence of a bare inserted theta port by itself evaluates every
such quantum term.  The phase-jet proposition is an explicit sufficient
source test where its conjugate-parent hypotheses hold; it does not assert
those hypotheses for all weak-current quantum corrections.  No new free
strong-phase selector, axion or Majorana interaction is inserted.
\end{remark}

\paragraph{Current local obligation ledger.}
The finite \(\mathsf G_5^{\rm fin}\), named roles and charged source
assembly remain closed.  For the ordinary spin--\(Z_6\) quantum anomaly
claim, the new first check is the exact regulated anomaly/background descent;
once checked, the global obstruction is discharged by
\ref{p5:thm:v2210-lit-anomaly-obstruction}, while the actual measure/Ward
certificate remains.  Charged physical roots retain
\ref{p5:sec:v2169-charged-support-application} and the canonical source test
of Paper~V-B, \ref{p6:sec:v2210-lit-physical-source-test}.  The harmonic
neutral branch separately retains the result-independent K56-to-harmonic
mass/incidence constructor of
\ref{p6:sec:v2208-neutral-rank-word}; charged closure and the historical
\(D_4\) response do not supply it.
% END LITERATURE-VERIFIED STANDARD-MODEL DOMAINS v2.210


% BEGIN FRAMEWORK v2.224: particle geometry and source comparison
\subsection{Global quadratic line and spin descent}
\label{p5:app:global-quadratic-line-spin-descent}

This appendix assumes an actual global line realization
over \(M_8\), with its base maps and target actions.
It does not infer that realization from the different
parameter orbit of Paper I-A.

\begin{proposition}[Global quadratic line and its two signs]
\label{p5:prop:global-m8-quadratic-sign-class}
Let \(N\to M_8\) be a Hermitian line with an everywhere
nonzero \(m\in\Gamma(N^{*2})\).  Then \(N\) is trivial
as a complex line, but its normalized real line can have
\[
 w_1=\alpha=(a_t,a_E)\in H^1(M_8,\ZZ_2)=\ZZ_2^2.
\]
It has a nonzero global real frame exactly when
\(\alpha=0\).  A unitary \(N\)-connection preserving \(m\)
requires constant \(\|m\|\) and has holonomies
\((-1)^{a_t},(-1)^{a_E}\) around the labeled scalar cycles.
\end{proposition}
\begin{proof}
Apply Paper~1B, Theorem~\ref{p1b:thm:global-quadratic-line}, to
the global Hermitian line and nonzero quadratic section in the statement.
It gives trivial complex Chern class, the two scalar-cycle sign bits,
and the preserving flat connection precisely at constant norm.
All four sign classes are represented here by
$m_{ab}=\sigma e^{i(a\phi_t+b\phi_E)}(e^*)^2$.
Complex topological triviality therefore retains the stated possible
nontrivial flat holonomies.
\end{proof}

\begin{corollary}[Exchange-compatible global quadratic-line signs]
\label{p5:cor:exchange-compatible-quadratic-line-signs}
On the global constant-norm quadratic-line branch above, suppose exchange
of the labeled scalar cycles has a parallel involutive lift to this same
line, preserving its Hermitian and quadratic structures. Then
\[
 a_t=a_E,\qquad (a_t,a_E)\in\{(0,0),(1,1)\}.
\]
Both classes admit such lifts, and the residual real fibre lift has two
signs. This corollary does not select the actual global class.
\end{corollary}
\begin{proof}
Parallel exchange conjugates the cycle holonomies:
\(J U_tJ^{-1}=U_E\). In rank one they are scalars, so
\((-1)^{a_t}=(-1)^{a_E}\).
Conversely an invariant character gives an exchange-equivariant flat
line by Paper~I-C, Theorem~\ref{p1b:thm:equivariant-flat}.
On the real quadratic line the involutive fibre constant is \(+1\)
or \(-1\). The two invariant characters supply the asserted realizations.
\end{proof}
Unequal characters can be exchanged after taking the direct sum of two
lines; that rank-two construction is not a same-line lift.


\begin{proposition}[Relative determinant and Poincare parity]
\label{p5:prop:relative-determinant-poincare-parity}
Let \(\mathcal M=N^{*2}\) be the relative determinant,
and let \(\mathcal P\) be the scalar Poincare line with
primitive \(c_1(\mathcal P)=\eta\).  Then
\[
 c_1(\mathcal M)=-2c_1(N).
\]
Thus \(\mathcal M\) cannot equal an odd power of
\(\mathcal P\).  If
\(\mathcal M=\mathcal P^k\otimes A\), then
\[
 c_1(A)=-2c_1(N)-k\eta;
\]
under the everywhere-nonzero comparison it becomes
\(c_1(A)=-k\eta\).  If \(c_1(N)=d\eta\), a nonzero
section of \(N^{*2}\otimes\mathcal P^k\) requires
\(k=2d\).
\end{proposition}
\begin{proof}
The tensor-line and parity calculation is Paper~1B,
Theorem~\ref{p1b:thm:global-quadratic-line}, with the same primitive
choice $c_1(\mathcal P)=\eta$. Additivity gives
$c_1(A)=-2c_1(N)-k\eta$ and the nonzero-section condition $k=2d$.
For the actual connection on $A_{12}\otimes\ell^{12}$, the endpoint
test remains $F_A+12F_\ell=0$ together with trivial tensor holonomy.
These conditions are necessary by parallelism and sufficient by
path-independent parallel transport from one nonzero value.
\end{proof}

The positive determinant-density line has a positive real
action, since \(\det_{\RR}U=|\det_{\CC}U|^2\).
It has no complex phase to supply this relative
determinant.  Also \(H^2(SO(3),\ZZ)=\ZZ_2\) for the
joint parameter orbit, unlike \(H^2(M_8,\ZZ)\).
A hosted parameter line is not automatically its
pullback along a specified physical base map.

\begin{proposition}[Conditional spin and matter sign incidence]
\label{p5:prop:quadratic-sign-spin-matter-incidence}
On a constituted common base, use the full spin/matter
tensor role for which the central spin sign \(s_0\)
and \(p\) both act by minus one on their factors.
Then
\[
 (\operatorname{Spin}_x\times K_2)/\langle(s_0,p)\rangle
       \simeq\operatorname{Spin}_x\times K_1.
\]
For a \(K_2\)-bundle with sign class \(\alpha\),
\[
 S_s\otimes F_P\simeq S_{s+\alpha}\otimes F_{\mathrm{SM}}.
\]
For the mathematical position bundle
\(\pi_x^*TS_x^3\to M_8\), the four spin twists all
restrict to the unique spin structure on \(S_x^3\).
\end{proposition}
\begin{proof}
Using $K_2=K_1\times\{1,p\}$, the homomorphism
$(s,k,p^\varepsilon)\mapsto(s\,s_0^\varepsilon,k)$ has exactly the
displayed diagonal kernel. Write the actual matter transition as
$\varepsilon_{ij}g_{ij}$, $\varepsilon_{ij}\in\{\pm1\}$. Moving the
same sign into the spin transition gives the stated matter-bundle
isomorphism, compatible with Clifford multiplication and the transported
connections. This sign-line operation is the spin-twist conclusion of
Paper~1B, Theorem~\ref{p1b:thm:global-quadratic-line};
Theorem~\ref{p1b:thm:spin} gives the two scalar-cycle twist labels.
For the pulled-back position bundle the twist torsor is likewise
$H^1(M_8;\mathbb Z_2)$; restriction kills both labels because
$H^1(S_x^3;\mathbb Z_2)=0$.
\end{proof}

The isomorphism concerns this common full tensor role.
A separately fixed spin owner, boundary condition or
individual factor can retain the sign.  No physical
\(\operatorname{Spin}(8)\) premise is used.
A nowhere-zero smooth comparison family on a fixed
labeled base preserves \(\alpha\), since its winding
class is homotopy-invariant; transport covering \(f\)
instead carries it by \(f^*\).
Local Noether identities do not choose \(\alpha\).
Global spectra can distinguish it: on a scalar circle
the trivial and sign holonomies give respectively
integer and half-integer modes, with different
zero-mode content.  Thus the sign-absorption isomorphism
does not leave a separately fixed spectrum unchanged.
% END M8 v1.78 TOP01


\subsubsection{Anomaly cancellation}
\label{p5:sec:part1-anomalies}

\begin{theorem}[Generationwise anomaly-coefficient cancellation]
\label{p5:thm:part1-anomaly-cancellation}
On \(\mathsf G_5^{\rm fin}\), all finite gauge and mixed
gauge--gravitational anomaly coefficients vanish for each generation, and
the \(SU(2)\) doublet count is even.  Using the derived
\(\mathsf J_5^{\rm conn,der}\) and assuming
\(\mathsf R_5^{\rm SSQFT}\), these identities imply
absence of the corresponding perturbative local anomalies and of the global
\(SU(2)\) Witten anomaly.  Hence all three replicated generations are anomaly
free on that local-QFT domain.
\end{theorem}

\begin{proof}
Using left-chiral fields
\[
\begin{aligned}
Q_L&:(\mathbf3,\mathbf2)_{1/6}, &
L_L&:(\mathbf1,\mathbf2)_{-1/2}, &
\nu_R^c&:(\mathbf1,\mathbf1)_0,\\
u_R^c&:(\bar{\mathbf3},\mathbf1)_{-2/3}, &
d_R^c&:(\bar{\mathbf3},\mathbf1)_{1/3}, &
e_R^c&:(\mathbf1,\mathbf1)_1,
\end{aligned}
\]
the nontrivial sums are
\begin{align}
SU(3)^2U(1)_Y:&\quad
2\left(\frac16\right)\frac12-\frac23\frac12+\frac13\frac12=0,\\
SU(2)^2U(1)_Y:&\quad
3\left(\frac16\right)\frac12-\frac12\frac12=0,\\
\mathrm{grav}^2U(1)_Y:&\quad
6\left(\frac16\right)+2\left(-\frac12\right)+3\left(-\frac23\right)+3\left(\frac13\right)+1=0,\\
U(1)_Y^3:&\quad
6\left(\frac16\right)^3+2\left(-\frac12\right)^3+3\left(-\frac23\right)^3+3\left(\frac13\right)^3+1=0.
\end{align}
The \(SU(3)^3\) coefficient cancels between two triplet and two conjugate-triplet
contributions, while the perturbative \(SU(2)^3\) coefficient vanishes because
the doublet is pseudoreal.  Each generation contains \(3+1=4\) left weak
doublets, so the Witten parity condition is even
\cite{Witten1982SU2Anomaly}.  Tensoring with the family carrier repeats these
finite coefficient identities three times; their local-QFT interpretation is
on the residual domain stated in the theorem.
\end{proof}

% BEGIN FRAMEWORK v3.029: GP3 MATTER OWNER BRIDGE
\begin{proposition}[Common-copy branch from the complete solder and anomalies]
\label{p5:prop:v329-solder-anomaly-commoncopy}
Consider a hypothetical physical extension of the generated matter
owner which represents the complete factor algebra of
Theorem~\ref{p5:thm:v329-matter-factor-envelope}, including
the full unitary Clifford-solder relations on the represented
left/right sectors. Suppose the Standard Model charge table is
retained and the source's local-QFT anomaly-cancellation criterion
applies on $\mathsf R_5^{\rm SSQFT}$.
Then the additional matter-sector multiplicities necessarily satisfy
\[
 m_Q=m_u=m_d=k_C,\qquad
 m_L=m_e=m_\nu=k_\ell,\qquad
 k_C=k_\ell=k\in\ZZ_{\ge1}.
\]
This proposition does not select $k=1$ and is not a proof that every
physically occurring GP3 candidate has the stipulated full-solder
and local-QFT properties.
\end{proposition}
\begin{proof}
The full Clifford matrix algebra acts irreducibly on each $W$ factor
and commutes with each spectator multiplicity space; its faithful
representation classification makes the right and left multiplicities
equal separately in the color and lepton blocks. Equivalently, the
extended unit solder $S_H=(\widetilde H\otimes J_+,H\otimes J_-)$,
with both $S_H^\dagger S_H=I$ and $S_HS_H^\dagger=I$, makes
$J_\pm$ unitary. The $SU(2)^2U(1)_Y$ anomaly coefficient with the
source charge table is $(k_C-k_\ell)/4$ per unamplified family.
Its vanishing gives $k_C=k_\ell$. All remaining listed anomaly
coefficients vanish on the resulting common-copy branch.
\end{proof}
% END FRAMEWORK v3.029



\section{Finite Dirac geometry and charged source-value constitution}
\label{p5a:sec:organized-source}
\subsection{Generated Yukawa--Dirac owner family and optional charged-member strata}
\label{p5:sec:v186-yukawa-dirac-owner-family}

The earlier G15 firewall distinguished the complete family operator grammar
from the missing physical weighted Yukawa/source owner.  The distinction
remains correct for the operator grammar \emph{alone}; the complete owner
family is now generated by solving the already typed support, grading,
coordinate-linearity, gauge-covariance, and adjoint constraints as one
Paper-II-A exact-solution family.

\begin{definition}[Constrained Yukawa--Dirac family]
\label{p5:def:v186-constrained-yukawa-dirac-family}
Put
\[
 \mathfrak Y
 :=\operatorname{End}(G)_u\oplus\operatorname{End}(G)_d
   \oplus\operatorname{End}(G)_e\oplus\operatorname{End}(G)_\nu
 \simeq M_3(\mathbb C)^4.
\]
Let $\mathscr D_Y$ be the complete family of finite operators $\Phi(H)$ on the
matter carrier that are self-adjoint after dagger completion, odd for the
left/right grading, coordinate-linear in the generated pair $(H,\widetilde
H)$, preserve the color/lepton projectors, connect only the generated
right-singlet to left-doublet support slots, and are covariant for the full
one-field/conjugate candidate Abelian plane of
Lemma~\ref{p5:lem:intrinsic-abelian-centralizer-plane}.  No numerical hypercharge member or
numerical Yukawa member is included in this definition.
\end{definition}

\begin{theorem}[Generated weighted Yukawa--Dirac owner family]
\label{p5:thm:v186-generated-yukawa-dirac-owner-family}
On the generated finite matter carrier, the map
\[
 \Theta_Y:\mathfrak Y\longrightarrow\mathscr D_Y
\]
defined by
\begin{align}
 \Theta_Y(Y_u,Y_d,Y_e,Y_\nu)(H)
 ={}&H P_{R-}\otimes(P_C\otimes Y_d+P_\ell\otimes Y_e)\nonumber\\
 &+\widetilde H P_{R+}\otimes(P_C\otimes Y_u+P_\ell\otimes Y_\nu)
 +\mathrm{h.c.}
 \label{p5:eq:v186-yukawa-owner-map}
\end{align}
is a bijection.  Its inverse consists of the four generated block-projection
readouts onto the orthogonal $(R\pm,C/\ell)$ support slots followed by the
canonical identification of the residual family factor with
$\operatorname{End}(G)$.  Hence $\mathscr D_Y$ is the complete generated
weighted dagger-quiver/Dirac owner family; no unique numerical member is
required for structural closure.

The same finite-carrier trace owns its scalar action rows.  In particular the
common action readouts are exactly the already generated invariants
\[
 a_Y=3\operatorname{tr}Y_u^\dagger Y_u
     +3\operatorname{tr}Y_d^\dagger Y_d
     +\operatorname{tr}Y_e^\dagger Y_e
     +\operatorname{tr}Y_\nu^\dagger Y_\nu,
\]
\[
 b_Y=3\operatorname{tr}(Y_u^\dagger Y_u)^2
     +3\operatorname{tr}(Y_d^\dagger Y_d)^2
     +\operatorname{tr}(Y_e^\dagger Y_e)^2
     +\operatorname{tr}(Y_\nu^\dagger Y_\nu)^2.
\]
Thus the positive polar parts of the edge weights and the color multiplicity
weights are family data of the same owner, not additional selectors.
\end{theorem}
\begin{proof}
Odd grading forces every coordinate-linear Dirac arrow to connect one of
the two generated right-singlet supports to the left weak doublet.  The
color/lepton projectors forbid crossed $C\leftrightarrow\ell$ arrows, and the
generated coordinate pair assigns $H$ to the $R-$ arrow and $\widetilde H$
to the $R+$ arrow.  Hence there are exactly the four displayed support slots.
Before applying parent-relation hypercharge descent, write a candidate Abelian generator as
$X=aT_R^3+b(B-L)$.  The $B-L$ weight cancels across every same-$Q$
left--right arrow, while the two coordinate modules have opposite
$T_R^3$ weights $\pm a/2$; consequently all four slots are covariant for the
entire candidate plane, not merely its Y ray.

The constructor word used here is the frozen composition of these generated
projector restrictions, tensor extension by the already generated family
endomorphism carrier, coordinate insertion by $(H,\widetilde H)$, finite
direct-sum assembly, and adjoint closure.  Its target sort is
$\mathscr D_Y$ of
Definition~\ref{p5:def:v186-constrained-yukawa-dirac-family}, declared before
solving.  The four source blocks are mutually separated by the generated
projectors $P_C,P_\ell,P_{R+},P_{R-}$.  Projecting
\eqref{p5:eq:v186-yukawa-owner-map} onto one such block and reading the
remaining $G$ factor recovers exactly one of $Y_u,Y_d,Y_e,Y_\nu$.  Therefore
$\Theta_Y$ is injective, while completeness of the four support slots makes it
surjective.  The resulting word is a predeclared finite exact constraint
constructor, so
Corollary~\ref{p2a:cor:v186-generated-constrained-owner-family} retains the
whole family as generated owner content.  The formulas for $a_Y,b_Y$ are
Corollary~\ref{p5:cor:v175-yukawa-scalar-action-family}; hence their weights
come from the same finite trace.
\end{proof}

\begin{corollary}[G15 owner-family closure and residual member/source-word firewall]
\label{p5:cor:v186-g15-owner-family-split}
The structural G15 obligation to construct the physical weighted
Yukawa--Dirac/source-row \emph{family} is discharged by
Theorem~\ref{p5:thm:v186-generated-yukawa-dirac-owner-family}.  The local
scalar jet compiler and, on the named clocked uninserted branch, the local
fermion jet placement are also derived.  At the historical v1.86 stage the action-owned occurrence of a sector-specific
normalized source member was still recorded as residual G15 debt.  In the
historical v3.88 graph that occurrence was the G0 fermion-source/member
boundary. The current rooted primitive and charged instances are supplied by
Theorems~\ref{p5:thm:v392-rooted-srcval-incidence}
and~\ref{p5:thm:v397-charged-preterminal-srcval-realization}; their mixed
readouts are dependent consequences rather than additional selectors.
The charged comparison's $25$-label exclusion consumption is supplied by
Theorem~\ref{p5:thm:v398-charged-rooted-comparison-owner-consumption}.
An actual preparation, physical pole/rest identification, vacuum or quantum
measure/CP claim retains only its own consumed constructor/state/domain;
none is a universal antecedent of this finite charged source.
\end{corollary}

\begin{proposition}[Optional charged-member strata after Yukawa-owner closure]
\label{p5:prop:v186-abelian-pointing-strata}
The generated Yukawa--Dirac owner family preserves the full candidate plane
\(\operatorname{span}\{T_R^3,B-L\}\) and therefore does not itself perform
the parent-relation hypercharge descent.  Structural hypercharge is already
generated independently by Proposition~\ref{p5:prop:v187-pcad-yfix}.

If a later physical claim separately supplies a charged member
\(m_{\rm occ}\in(L_0^*)^{\otimes2}\), the occurrence strata
\[
 \mathfrak C_0:=\{m_{\rm occ}=0\},\qquad
 \mathfrak C_{\ne0}:=\{m_{\rm occ}\ne0\}
\]
classify that member statement only.  On \(\mathfrak C_{\ne0}\) a member of
weight \((-1,+2)\) has the same connected stabilizer
\(-a+2b=0\), but this is a consistency check on the already generated
\(U(1)_Y\), not its structural origin.  A nonzero Majorana quadratic is one
possible downstream member and remains optional unless separately generated.
\end{proposition}
\begin{proof}
Candidate-plane covariance follows from
Theorem~\ref{p5:thm:v186-generated-yukawa-dirac-owner-family}.  The stabilizer
of any separately supplied nonzero member of weight \((-1,+2)\) is the kernel
\(-a+2b=0\).  Proposition~\ref{p5:prop:v187-pcad-yfix} derives the same line
without requiring such a member occurrence.
\end{proof}
% END OWNER-FAMILY v1.86

\subsection{Predeclared zero-order chiral source target}
\label{p5:sec:v2187-predeclared-chiral-source-target}
The general four-slot source grammar is declared before evaluating a
particular member. On $\mathsf{Cert}_{\rm hier}$ its rooted charged
$u,d,e$ instance is supplied by
Theorem~\ref{p5:thm:v397-charged-preterminal-srcval-realization}.
Availability of the fourth slot does not require an additional bulk neutral
source for the charged packet; native N56 energy/current roles retain their
separate typed constructor. Historical alphabet/offset candidates are not active source prerequisites.

\begin{definition}[Same-action direct chiral-source compiler target]
\label{p5:def:v369-direct-mass-compiler-target}
A direct chiral-source constructor is a result-independent generated
same-parent map
\[
 \mathcal C_{\rm mass}:
 (R^\pm,C/\ell;K_7,K_9,R_G,X,\mathscr E_H,c(\hat h))
 \longrightarrow
 (D_u,D_d,D_e,D_\nu),
\]
whose outputs occupy the four already generated support slots of
\eqref{p5:eq:v186-yukawa-owner-map} as zero-order chiral maps.  The compiler
must fix the finite return word, its relative rooted transport against the
Clifford hierarchy, and any action-owned functional calculus before scalar
trace moments erase this information.  It may not be selected by comparison
with observed masses or mixings.
\end{definition}

% BEGIN SAFEGUARD-HISTORICAL-G15-PROVENANCE v3.88
% The following release sequence is provenance and is superseded as a live dependency graph by Release v3.85--v3.88.


\subsection{rooted geometric finite-Dirac member and source-value incidence}
\label{p5:sec:v392-rooted-geometric-srcval}

The generic source-port/value firewall of Release v3.89 remains valid: a
generated nonzero port does not by itself actualize a nonzero member of a
later physical role.  The rooted triality branch is stronger.  The occurring
Hurwitz--$D_4$ parent generates one specific nonzero finite odd operator before
any source-family member is selected, and the already generated
Yukawa--Dirac block readouts identify the owner coordinates of that operator.
Paper~III-F then consumes that independently generated member in the
pre-existing local fermion-action constructor.  No implication
$\mathsf{SrcPort}\Rightarrow\mathsf{SrcVal}$ is used.

\begin{definition}[Rooted geometric finite odd operator]
\label{p5:def:v392-rooted-geometric-finite-odd}
On the occurring rooted Hurwitz--$D_4$ domain of
Theorem~\ref{p5:thm:v375-d4-triality-source-consumption}, take the canonical
axial direction $h_*=-e_1$.  Paper~I-C identifies the protected chiral and
family restrictions of the same rooted triality contraction as
\[
 c(h_*):W_R\to W_L,
 \qquad I_G:G\to G.
\]
Relation-free tensor extension on the already generated matter carrier gives
\[
 \boxed{T_*:=c(h_*)\otimes I_Q\otimes I_G.}
\]
Let
\[
 \boxed{D_{F,*}^{\rm geom}:=
 \begin{pmatrix}0&T_*\\T_*^\dagger&0\end{pmatrix}.}
\]
\end{definition}

\begin{proposition}[Rooted primitive member is nonzero and full rank]
\label{p5:prop:v392-rooted-member-full-rank}
The rooted Clifford factor is unitary on its protected chiral slot, hence
\[
 T_*^\dagger T_*=I,
 \qquad \boxed{(D_{F,*}^{\rm geom})^2=I},
 \qquad \boxed{\ker D_{F,*}^{\rm geom}=0}.
\]
Thus the rooted geometry generates a specific nonzero finite odd operator
before any member of the ambient Yukawa--Dirac family is chosen.
\end{proposition}

\begin{theorem}[Owner-coordinate readout of the rooted geometric member]
\label{p5:thm:v392-rooted-member-owner-readout}
Let $\widehat H_*$ be the unique normalized coordinate representative of the
same rooted solder supplied by
Theorem~\ref{p5:thm:v377-unit-solder-normalized-scalar}, so that
\[
 \Phi(\widehat H_*)=c(h_*).
\]
Apply the four generated block-projection readouts inverse to
Theorem~\ref{p5:thm:v186-generated-yukawa-dirac-owner-family} to
$D_{F,*}^{\rm geom}$.  Because $T_*$ is the identity on $Q\otimes G$, every
$(R\pm,C/\ell)$ block readout leaves the same family factor $I_G$.  Therefore
the unique owner coordinate is
\[
 \boxed{Y_*=(I_G,I_G,I_G,I_G)}
\]
and
\[
 \boxed{\Theta_Y(Y_*)(\widehat H_*)=D_{F,*}^{\rm geom}.}
\]
Hence $Y_*$ is a readout of one rooted geometric operator, not a point selected
from $M_3(\mathbb C)^4$ by GP3, simplicity, a VEV, or phenomenological input.
\end{theorem}

\begin{theorem}[Rooted fermion source-incidence closure]
\label{p5:thm:v392-rooted-srcval-incidence}
On the local clocked uninserted classical branch of Paper~III-F, the rooted
operator $D_{F,*}^{\rm geom}$ satisfies the Paper~II-A source-incidence
certificate.  Its fields are:
\begin{enumerate}
\item occurring parent: the occurring rooted Hurwitz--$D_4$/triality frame;
\item licensed geometric constructor: axial triality restriction,
      relation-free identity extension and dagger completion;
\item finite owner target: the generated Yukawa--Dirac family $\mathscr D_Y$,
      with owner coordinate fixed by the generated block readouts;
\item nonzero proof: $(D_{F,*}^{\rm geom})^2=I$;
\item physical action owner: the Paper~III-F local uninserted fermion-jet
      constructor;
\item required same-action output: the independently generated member
      $D_{F,*}^{\rm geom}$ itself.
\end{enumerate}
Consequently
\[
 \boxed{\operatorname{Occ}(\mathbb T_*)
 \Longrightarrow
 \operatorname{CoReal}_{\mathsf{SrcVal}}(D_{F,*}^{\rm geom}).}
\]
No second G0 source-member occurrence token is required for primitive source
existence on this rooted branch.
\end{theorem}

\begin{proof}
The first four certificate fields follow from the rooted triality ancestry,
Proposition~\ref{p5:prop:v392-rooted-member-full-rank}, and
Theorem~\ref{p5:thm:v392-rooted-member-owner-readout}.  The last two fields
are supplied by Paper~III-F's independently-generated-member handoff to the
local fermion jet: a finite Dirac/Yukawa member which is itself a required
output of a licensed constructor on an occurring root is consumed by the
already constituted fermion-action slot without a duplicate occurrence token.
Canonical normalization on that branch gives
\[
 \widehat D_{F,*}=D_{F,*}^{\rm geom}\neq0.
\]
By Definition~\ref{p5:def:v389-srcport-srcval}, this is a physical finite-action
$\mathsf{SrcVal}$.  The generic port/value firewall is untouched because the
proof never promotes a port merely from nonvanishing; it independently
generates the nonzero member and then places it in the action.
\end{proof}

\begin{corollary}[Primitive source-existence boundary retired]
\label{p5:cor:v392-primitive-source-boundary-retired}
The live status is
\[
 \boxed{\texttt{FERMION\_SRCVAL\_INCIDENCE = CLOSED\_GEOMETRIC}.}
\]
The historical token
\[
 \texttt{G0\_FERMION\_SOURCE\_MEMBER}
\]
is retired as a prerequisite for primitive source existence on the rooted
branch.  The charged noncentral source instance is supplied by
Theorem~\ref{p5:thm:v397-charged-preterminal-srcval-realization}.
A genuinely different fermion target must name its own source sort and
consumed observable; physical pole/rest identification remains separately typed.
\end{corollary}

\begin{remark}[No VEV or scalar-background premise]
\label{p5:rem:v392-no-vev-source-origin}
The normalized $\widehat H_*$ is only the coordinate presentation of the
already generated unit Clifford solder, $\Phi(\widehat H_*)=c(h_*)$.  No scalar
stationary point, nonzero spacetime background field, VEV magnitude, fitted
Yukawa coefficient or $v/\sqrt2$ mass law is used.
\end{remark}

\begin{remark}[No overgeneration]
\label{p5:rem:v392-no-canonical-operator-overgeneration}
This theorem does not assert that every canonical odd endomorphism occurs in
the action.  The rooted member is distinguished by the complete origin word:
occurring triality parent, same-parent chiral/family restrictions,
relation-free tensor extension, exact owner readout, nonzero/full-rank proof,
and an already licensed fermion-action consumer.  Diagnostic antiprism,
Schur/Feshbach and comparison operators remain behind their existing
consumption firewalls unless they independently satisfy an equally complete
origin certificate.
\end{remark}
% END ROOTED-GEOMETRIC-SRCVAL-CLOSURE v3.92

% BEGIN FERMI-PRETERMINAL-PARENT-CONTRACTION v3.95
\subsection{charged preterminal parent-contraction target}
\label{p5:sec:v395-preterminal-parent-contraction}
The following factorization and generated-member criteria preserve the
charged owner information before noninjective terminalization.
Their actual rooted instance is supplied by
Theorem~\ref{p5:thm:v397-charged-preterminal-srcval-realization}; the
superseded release-stage residuals are not active preconditions.

\begin{definition}[Charged preterminal parent contraction]
\label{p5:def:v395-preterminal-parent-contraction}
For a charged sector $s\in\{u,d,e\}$, let $\widehat{\mathscr C}_s$ denote the rooted tuple of
owner-resolved charged channels retained with their native metrics,
orientation/Hodge/triality/shell marks and terminal amplitudes.  A
\emph{charged preterminal parent contraction} is a licensed occurring
constructor whose required finite output is
\[
 \mathfrak A_s(\widehat{\mathscr C}_s)
 \in\operatorname{Hom}(F_{R,s},F_{L,s})
\]
before any noninjective scalar terminal readout of the individual channels.
\end{definition}

\begin{proposition}[Sector factorization removes an independent family embedding]
\label{p5:prop:v395-sector-factorization-preterminal}
If Definition~\ref{p5:def:v395-preterminal-parent-contraction} occurs, then the
already generated sector decomposition forces
\[
 \boxed{
 \mathfrak A_s
 =\Sigma_{\sigma_s}\otimes P_{Q_s}\otimes R_s,
 \qquad R_s:G_R\to G_L,}
\]
with $(\sigma_s,Q_s)=(+,C),(-,C),(-,\ell)$ for
$s=u,d,e$, respectively.  The neutral $\nu$ role-space/PMNS handoff retains
its existing separately typed downstream boundary.  Hence $R_s$ is the unique remaining family
coefficient of the actual action block.  Any support isometry later denoted
$J_s$ is a coordinate/polar readout of this already generated incidence and is
not an additional CORE constructor.
\end{proposition}

\begin{proof}
The generated Yukawa--Dirac owner family of
Theorem~\ref{p5:thm:v186-generated-yukawa-dirac-owner-family} is factor-local
on the protected chiral and color/lepton blocks, leaving $G$ as the family
factor.  Once an actual member in that block is generated, projection onto the
nonzero spin/internal factor identifies the unique coefficient on $G$.
No map from an external three-space is required.
\end{proof}



\begin{remark}[No Higgs/VEV reopening]
\label{p5:rem:v395-no-higgs-reopening}
This retyping concerns only the native finite fermion action.  It introduces
no Higgs VEV, no fitted Yukawa matrix and no scalar-background selector.
\end{remark}
% END FERMI-PRETERMINAL-PARENT-CONTRACTION v3.95

% BEGIN CHARGED-GENERATED-MEMBER-CLOSURE-CRITERION v3.96
\subsection*{generated-member route for charged preterminal occurrence}
\label{p5:sec:v396-generated-member-route}
Earlier occurrence gaps were closed by generating the particular physical
member from an occurring root and only afterward reading its coordinates in a
larger owner family.  The charged preterminal construction has the same required
shape, fulfilled by Theorem~\ref{p5:thm:v397-charged-preterminal-srcval-realization}.
Selecting a member of the Yukawa--Dirac owner first is not the licensed
closure route.

\begin{definition}[Charged geometric preterminal member]
\label{p5:def:v396-charged-geometric-preterminal-member}
For sector $s\in\{u,d,e\}$ let $\widehat{\mathscr C}_s$ be the generated
owner-resolved charged packet.  A \emph{charged geometric preterminal member}
is a result-independent licensed contraction
\[
 \boxed{
 \mathfrak K_s:\widehat{\mathscr C}_s
 \longrightarrow R_s^{\rm geom}\in\operatorname{End}(G)}
\]
formed before terminal readout, with all generated owner metrics, rooted
orientation/Hodge/triality/shell marks and required cross relations retained.
Its finite odd completion is
\[
 \boxed{
 D_{F,s}^{\rm pre,geom}
 =\Sigma_{\sigma_s}\otimes P_{Q_s}\otimes R_s^{\rm geom}+\mathrm{h.c.}}
\]
with the already generated sector factors.
\end{definition}

\begin{theorem}[Generated-member occurrence route]
\label{p5:thm:v396-generated-member-occurrence-route}
Assume Definition~\ref{p5:def:v396-charged-geometric-preterminal-member} is
satisfied by a licensed geometric contraction of the occurring rooted parent.
Then $D_{F,s}^{\rm pre,geom}$ is an independently generated finite fermion
member in the existing action owner.  Its owner-family coordinate is read only
after the member has been generated.  The existing Paper-III-F
independently-generated-member handoff therefore consumes it without an
additional G0/source-member occurrence token.  Paper II-A preterminal relation
preservation and generated co-consumption then retain all mixed blocks,
metrics and current coefficients before any noninjective terminal readout.
\end{theorem}

\begin{proof}
The generated matter carrier and factor-local odd-source algebra already
provide the unique sector tensor factors $\Sigma_{\sigma_s}\otimes P_{Q_s}$.
A generated $R_s^{\rm geom}$ therefore determines the displayed finite odd
member without selecting a coordinate from the generic owner family.  This is
the same logical order as the rooted primitive source closure: generate the
specific finite member first, read the generic owner coordinate afterward,
and then apply the existing finite-action consumer.  The II-A chronology
forbids terminalizing the owner channels before their complete contraction.
\end{proof}



\begin{remark}[No owner-family member selection]
\label{p5:rem:v396-no-owner-selection}
It is not sufficient to choose a mathematically available
$R\in\operatorname{End}(G)$ and evaluate the generic Yukawa--Dirac owner on
it.  The closure route must independently generate $R_s^{\rm geom}$ from the
owner-resolved charged parent, as in the primitive rooted-source precedent.
\end{remark}
% END CHARGED-GENERATED-MEMBER-CLOSURE-CRITERION v3.96

% BEGIN CHARGED-PRETERMINAL-SRCVAL-CLOSURE v3.97
\subsection*{charged preterminal source-value closure}
\label{p5:sec:v397-charged-preterminal-srcval-closure}
The direct mass compiler target of
Definition~\ref{p5:def:v369-direct-mass-compiler-target} is already the four
zero-order chiral source slots. On the independently certified
hierarchy-arrow domain, Paper~V-B v5.90 assembles, for each charged
sector, the mass/incidence sibling packet and its preterminal family
coefficient. The assembly retains the origin obligation
$\mathsf{Cert}_{\rm hier}$; it does not generate an absent hierarchy arrow.

\begin{definition}[Charged preterminal family coefficient]
\label{p5:def:v397-charged-preterminal-family-coefficient}
For $s\in\{u,d,e\}$ let
\[
 \boxed{
 R_s^{\rm ch}:=R_s^{\rm pre}
 =\sum_{r\in\mathcal I_s}m^{\rm fin}_{s,r}\iota_{s,r}
 \in\operatorname{Hom}(G_R,G_L)\cong\operatorname{End}(G),}
\]
with $(m^{\rm fin}_{s,r},\iota_{s,r})$ the complete generated sibling packet
of Paper~V-B.  The incidence operator carries the inherited root, orientation,
Hodge, triality and shell information; $m^{\rm fin}_{s,r}$ is the already
fixed charged finite-mass coefficient.
\end{definition}

\begin{theorem}[Charged preterminal source-value realization]
\label{p5:thm:v397-charged-preterminal-srcval-realization}
On the occurring rooted charged branch, assume that the hierarchy row
arrows have the independent origin certificate of
Definition~\ref{p5:def:v228-hierarchy-origin}. On that generated-packet
domain, the frozen direct-mass constructor has required output
\[
 \boxed{
 D_{F,s}^{\rm ch}
 =\Sigma_{\sigma_s}\otimes P_{Q_s}\otimes R_s^{\rm ch}
 +\mathrm{h.c.}
 \in\mathsf{SrcVal}_{\rm odd,s},}
\]
with $(\sigma_s,Q_s)=(+,C),(-,C),(-,\ell)$ for $s=u,d,e$.
Equivalently the charged part of the complete finite source is
\[
 \boxed{D_{F,\rm ch}^{\rm geom}=\sum_{s\in\{u,d,e\}}D_{F,s}^{\rm ch}.}
\]
This is a generated image of the owner-resolved charged packet, not a selected
point of $M_3(\mathbb C)^4$.  Therefore
\[
 \boxed{
 \operatorname{Occ}(\mathbb T_*)\land\mathsf{Cert}_{\rm hier}
 \Longrightarrow
 \operatorname{CoReal}_{\mathsf{SrcVal}}
 (D_{F,\rm ch}^{\rm geom})}
\]
on the stated rooted/action domain, with no additional charged source-member
occurrence token.
\end{theorem}

\begin{proof}
On the independently certified hierarchy-arrow domain, Paper~V-B v5.90
assembles $R_s^{\rm ch}$ by a frozen linear synthesis of the complete
generated mass/incidence sibling packet.  Paper~I-A's signed triality
map and Paper~I-C's hierarchy bridge remove the old external support-to-$G$
identification: every incidence term already lies in the generated family
operator sort.  Theorem~\ref{p5:thm:v389-factorized-family-source} then fixes
the only factor-local odd operator in the sector to be
$\Sigma_{\sigma_s}\otimes P_{Q_s}\otimes R_s^{\rm ch}$, while
Definition~\ref{p5:def:v369-direct-mass-compiler-target} predeclares that
operator as the zero-order chiral mass-source target.  Paper~II-A sibling
co-consumption retains each $(m^{\rm fin}_{s,r},\iota_{s,r})$ before the
synthesis and its generated-image theorem promotes no new sort: the target
$\mathsf{SrcVal}_{\rm odd,s}$ is already declared by the direct-mass
constructor.  The existing root-generated finite-member handoff then supplies
physical occurrence without a duplicate G0 token.
\end{proof}

\begin{corollary}[Charged fermion preterminal assembly closure]
\label{p5:cor:v397-charged-preterminal-assembly-closure}
On the rooted charged branch,
\[
 \boxed{\begin{gathered}
 \texttt{FERMION\_PRETERMINAL\_PARENT\_CONTRACTION}\\
 =\texttt{CLOSED\_ON\_CERTIFIED\_PACKET}.
 \end{gathered}}
\]
The historical first-principle gaps remain retired:
\begin{itemize}
\item \texttt{FERMION\_SAME\_ACTION\_SVD\_ASSEMBLY},
\item \texttt{FERMION\_TYPED\_PARENT\_CONTRACTION},
\item an independent $J_s$.
\end{itemize}  No charged family-member selector or extra `SrcVal' postulate is
introduced on that domain. The input hierarchy-arrow origin certificate
remains required; this assembly formula does not discharge it.
\end{corollary}

\begin{corollary}[Mass/incidence semantics after closure]
\label{p5:cor:v397-mass-incidence-semantics}
The source consumes the complete pair
\[
 \boxed{(m^{\rm fin}_{s,r},\iota_{s,r})}
\]
for every charged role.  The first entry is the direct Clifford/split finite
mass coefficient and the second its family incidence.  No equation of the
form
\[
 m^{\rm fin}_{s,r}=\sqrt{\lambda_r(R_sR_s^\dagger)}
\]
is part of this construction.  Actual source-visible pole/rest identification and any consumed RG,
reporting or limiting comparison remain downstream and separately typed.
Native weak-current incidence is a different generated readout, including
the charged-lepton closure of Corollary~\ref{p5:cor:v3108-lepton-incidence-status};
it is not a reconstruction of the terminal mass coefficients.
\end{corollary}

\begin{remark}[Supersession of the v3.93--v3.96 charged boundary]
\label{p5:rem:v397-supersede-charged-boundary}
Remark~\ref{p5:rem:v393-assembly-boundary},
Corollary~\ref{p5:cor:v395-preterminal-frontier} and
Corollary~\ref{p5:cor:v396-charged-residual-retyped} are retained as historical
chronology.  Those statements are superseded by Corollary~\ref{p5:cor:v397-charged-preterminal-assembly-closure}.
Their charged open-boundary statements, including the intervening typed-parent
label, are superseded by
Corollary~\ref{p5:cor:v397-charged-preterminal-assembly-closure}.
This closure covers the rooted $u,d,e$ zero-order source packet. It does not
evaluate the complete physical source-cyclic pole polynomial, identify an
interacting root with a named particle/rest observable, or construct a
separately activated bulk high-neutral source.
\end{remark}
% END CHARGED-PRETERMINAL-SRCVAL-CLOSURE v3.97

% BEGIN FRAMEWORK v2.245: FINITE CHARGED ACTION-INCIDENCE TEST
\begin{proposition}[Four finite invariants certify one complete three-arrow chiral member]
\label{p5:prop:v245-charged-four-invariant-source-test}
Work on a common canonically normalized three-family Hilbert--Schmidt
owner in which the geometric polar Hopf arrows $B_r$ of
Theorem~\ref{p6:thm:v240-canonical-polar-hopf-triple} are all placed,
with $\Tr(B_r^\dagger B_s)=\delta_{rs}$. Let
$M_s=Z_{L,s}^{-1/2}D_{{\rm ch},s}Z_{R,s}^{-1/2}$ be the
\emph{independently evaluated actual} charged mixed action jet, and let
$m_{s,r}^{\rm fin}$ be the already generated coefficients. Then
\[
 \boxed{M_s=\sum_{r=1}^3m_{s,r}^{\rm fin}B_r}
\]
if and only if the following three oriented contractions and one
nonnegative norm test hold:
\[
 \boxed{\Tr(B_r^\dagger M_s)=m_{s,r}^{\rm fin}\ (r=1,2,3),
 \qquad \Tr(M_s^\dagger M_s)
 =\sum_{r=1}^3|m_{s,r}^{\rm fin}|^2.}
\]
This is a complete \emph{finite equality criterion}, not an evaluation
of $M_s$ and not an unconditional proof of $\mathsf{Cert}_{\rm hier}$.
\end{proposition}
\begin{proof}
The arrow Gram orthonormality follows from
Eq.~\eqref{p6:eq:v240-hopf-full-left-return} and the rank-one returns.
Orthogonal Hilbert--Schmidt projection onto their span has coefficients
$\Tr(B_r^\dagger M_s)$. With the three prescribed values,
Pythagoras gives
$\|M_s-\sum_rm_{s,r}^{\rm fin}B_r\|_{\rm HS}^2
=\Tr(M_s^\dagger M_s)-\sum_r|m_{s,r}^{\rm fin}|^2$.
It vanishes precisely for equality of the operators.
\end{proof}

% BEGIN CHARGED PHYSICAL-SOURCE EXCLUSION TESTS (v3.216)
\begin{corollary}[Charged E36 no-go for a C3-only actual source]
\label{p5:cor:charged-c3-source-no-go}
On the canonical charged-lepton Hilbert--Schmidt owner and its
rooted $E_{32}\oplus E_{36}\oplus E_{40}$ orthogonal operator-block
decomposition, let
$M_{\rm ch}=m_e^{\rm fin}B_e+m_\mu^{\rm fin}B_\mu
 +m_\tau^{\rm fin}B_\tau$
use the fixed normalized polar arrows. If $M$ is built only
from source coefficients in the rooted $C_3$-preserving algebra,
including regular split-preserving reductions as in
the rooted commutant-stability theorem of Paper~III-A, then
\[
 \|M-M_{\rm ch}\|_{\rm HS}^2\ge
 |m_\mu^{\rm fin}|^2+|m_\tau^{\rm fin}|^2.
\]
This is a no-go for the stated source-constructor class, not
for the full signed triality geometry or an independently generated
noncentral physical action insertion.
\end{corollary}
\begin{proof}
The two normalized $E_{36}$ polar arrows occupy the opposing
off-diagonal $Q_0/Q_\perp$ blocks, each of squared
Hilbert--Schmidt norm one. The source $M$ has zero projection onto
both blocks. Orthogonality and the Pythagorean theorem give the bound.
\end{proof}


\begin{proposition}[Fixed-frame signed-triality channel obstruction for the charged source]
\label{p5:prop:charged-signed-return-distance}
In the normalized signed-Fock family basis and on the same
canonical charged-lepton Hilbert--Schmidt owner, let
\(\mathcal K_{I_4}(U)=\left.\mathcal W_0^\dagger(U\otimes I_4)
\mathcal W_+\right|_G\), \(U\in M_3(\mathbb C)\), be the
mathematical full-target channel return of Paper~I-B.
Its image is the hyperplane \(M_{32}=M_{23}\) in one-based indices.
For the existing oriented charged-lepton polar-arrow packet
\(M_{\rm ch}=\sum_{r=e,\mu,\tau}m_r^{\rm fin}B_r\), the fixed-frame
formula gives
\[
 (M_{\rm ch})_{32}-(M_{\rm ch})_{23}
 =\frac{i}{\sqrt2}(m_\tau^{\rm fin}-m_\mu^{\rm fin}).
\]
Consequently the exact orthogonal Hilbert--Schmidt distance is
\[
 \boxed{\min_{U\in M_3(\mathbb C)}
 \|\mathcal K_{I_4}(U)-M_{\rm ch}\|_{\rm HS}^2
 =\frac14|m_\tau^{\rm fin}-m_\mu^{\rm fin}|^2.}
\]
This is a mathematical exclusion for the stated \emph{channel-only}
contraction and fixed common family frame, not an obstruction to a
separately licensed target-spinor-sensitive physical operator or to
quark roles with different source maps. The rank-nine seeded
contraction is mathematically available, but on unitary mixers its
singular values remain equal and no physical source incidence follows.
\end{proposition}
\begin{proof}
The full-target signed-Fock return has rank eight with image
\(\{M:M_{32}=M_{23}\}\), by direct multiplication of the source-owned
transverse coefficients. Substitution of the existing normalized
charged-lepton polar arrows gives the displayed oriented entry
contrast. The perpendicular projection of any matrix onto the
one-dimensional orthogonal complement of the hyperplane is
\(\frac12(M_{32}-M_{23})(E_{32}-E_{23})\), whose squared norm is
\(\frac12|M_{32}-M_{23}|^2\). Since the image fills the entire
hyperplane, the lower bound is attained for some \(U\).
The rooted seed result is multiplication by fixed unitary diagonal
factors and hence preserves singular values for unitary inputs.
\end{proof}

% END CHARGED PHYSICAL-SOURCE EXCLUSION TESTS

% END FRAMEWORK v2.245: FINITE CHARGED ACTION-INCIDENCE TEST

% BEGIN HURWITZ-GRAMMAR ROOT-SENSITIVE CHARGED SOURCE TESTS v3.211
\subsection{Root-sensitive hierarchy-arrow certificates for the charged chiral jet}
\label{p5:sec:v3211-root-sensitive-charged-jet}
The finite hierarchy geometry and the generated charged-current relation have
separate target sorts.  The current relation of
Theorem~\ref{p5:thm:v3108-lepton-hopf-incidence-parent} is a structural
incidence parent, whereas the source-value assertion of
Theorem~\ref{p5:thm:v397-charged-preterminal-srcval-realization} continues to
require the independent input-origin certificate
\(\mathsf{Cert}_{\rm hier}\) of
Definition~\ref{p5:def:v228-hierarchy-origin}.  The following exact
statements supply diagnostics on an independently evaluated fermion jet;
they do not identify a candidate family matrix as an actual source member.

\begin{proposition}[Commutator extraction of the two rooted negative-Hopf crossings]
\label{p5:prop:v3211-hopf-crossing-commutator}
On the existing rooted charged-lepton family chart put
\[
 c=\frac1{\sqrt3}(1,1,1)^{\mathsf T},\quad
 Q_0=cc^\dagger,\quad Q_\perp=I_G-Q_0,\quad
 W=R_-=\operatorname{diag}(1,i,-i).
\]
The commutators \(J=[Q_0,W]\) and \(K=[Q_0,J]\) determine both directed
\(E_{36}\) structural crossing blocks without any additional geometric
relation or terminal mass coefficient:
\[
 \boxed{Q_0WQ_\perp=\tfrac12(K+J),\qquad
 Q_\perp WQ_0=\tfrac12(K-J).}
\]
Both blocks have rank one and squared Hilbert--Schmidt norm \(8/9\).
This is a relation-level construction on the charged-current chart,
not an equality to the independently required physical chiral source.
\end{proposition}
\begin{proof}
Writing \(B=Q_0WQ_\perp\) and \(A=Q_\perp WQ_0\), the projection
identities give \(J=B-A\) and \(K=B+A\), proving the extraction formulas.
Since \(W\) is unitary and \(\langle c,Wc\rangle=1/3\),
\[
 \|A\|_{\rm HS}^2=1-|\langle c,Wc\rangle|^2=\tfrac89,
 \qquad \|B\|_{\rm HS}^2=\tfrac89.
\]
Each crossing has one-dimensional source or target and is nonzero.
\end{proof}

\begin{corollary}[Primitive family-blind mixed jets have no \(E_{36}\) crossing]
\label{p5:cor:v3211-primitive-mixed-jet-centrality}
Fix the primitive rooted member \(Y_*=(I_G,I_G,I_G,I_G)\) of
Theorem~\ref{p5:thm:v392-rooted-member-owner-readout}.  On a
family-blind bosonic/kinetic chart that inserts no other noncommuting
family relation in its zero-order chiral coefficient, the fermion mixed
variation of the printed action remains proportional to \(I_G\) on each
primitive family block.  It therefore obeys
\[
 \boxed{Q_0\mathscr D_{LR,*}Q_\perp
       =Q_\perp\mathscr D_{LR,*}Q_0=0.}
\]
Regular classical elimination of stationary bosons does not create a new
quadratic fermion coefficient on this fixed-member chart.  Thus the
primitive identity member alone cannot realize the nonzero crossings of
Proposition~\ref{p5:prop:v3211-hopf-crossing-commutator} as chiral mass
arrows by these operations.  The corollary excludes neither the already
proved charged-current incidence nor a different, source-owned noncentral
fermion contraction.
\end{corollary}
\begin{proof}
On the stipulated chart the family coefficient of the primitive bilinear
is scalar, hence its commutator with \(Q_0\) vanishes and both off-diagonal
corner compressions vanish.  The precise mixed-variation/regular-elimination
identity is the same-action fermion-jet proposition of Paper~III-F;
it adds no independent bilinear member when the eliminated bosonic variables
vary only at quadratic fermion order.  The nonzero commutator crossings
above consequently require a separately generated family-sensitive input
in this source sort.
\end{proof}

\begin{theorem}[A single nonnegative hierarchy-corner defect for the actual jet]
\label{p5:thm:v3211-hierarchy-corner-defect}
Let \(R:G_R\to G_L\) be an \emph{independently evaluated}, canonically
normalized three-family chiral source coefficient on a fixed common rooted
chart, and identify its left/right root projectors only as permitted by that
chart.  Put
\[
 D=Q_0RQ_0,\quad A=Q_\perp RQ_0,\quad
 B=Q_0RQ_\perp,\quad C=Q_\perp RQ_\perp,
 \qquad N=\Tr(R^\dagger R)>0,
\]
\[
 \sigma_C=\frac12\left\{
  [\Tr(C^\dagger C)]^2-\Tr[(C^\dagger C)^2]\right\},
\]
where in \(\sigma_C\) the nonzero block \(C\) is regarded on the
rank-two augmentation.  The dimensionless defect
\[
 \boxed{\Delta_{\rm hier}(R)=
   \frac{\|D\|_{\rm HS}^2}{N}
  +\frac{\|C^\dagger A\|_{\rm HS}^2+
          \|CB^\dagger\|_{\rm HS}^2+\sigma_C}{N^2}\ge0}
\]
vanishes if and only if
\[
 D=0,\qquad \rank C\le1,\qquad
 C^\dagger A=0,\qquad CB^\dagger=0.
\]
On the separately required nonzero-block domain \(A,B,C\ne0\), these
conditions are equivalent to the normalized polar factors of \(A,B,C\)
having mutually orthogonal initial and final one-dimensional supports.
Consequently their sum is a unitary three-arrow transport and the
unweighted left return splits into three orthogonal rank-one blocks.

Vanishing of \(\Delta_{\rm hier}\) tests this \emph{coarse corner
structure}; it does \emph{not} prove equality to the particular
Hopf-polar arrows, determine their source-owned phases or terminal
coefficients, or discharge \(\mathsf{Cert}_{\rm hier}\).
\end{theorem}
\begin{proof}
Every displayed term is nonnegative.  Since the two (possibly zero) singular
values of the augmented block \(C\) have squares \(u,v\ge0\),
\(\sigma_C=uv\), which vanishes precisely when \(\rank C\le1\).
For \(A,B,C\ne0\), the two Gram conditions state that the final line
of \(A\) is orthogonal to the final line of \(C\), and that the initial
line of \(B\) is orthogonal to the initial line of \(C\), respectively.
The initial line of \(A\) is the singlet, while the final line of \(B\)
is the singlet; the other two lines on each side exhaust the augmentation.
The three polar factors therefore have orthogonal initial and final
rank-one projections summing to identity.  Their sum is unitary.  The
converse follows by the same orthogonality and the vanishing singlet
corner.  No assertion about the actual origin of \(R\) enters this proof.
\end{proof}

% BEGIN FRAMEWORK INTEGRATION v3.024: EXACT ROOTED HIERARCHY DISTANCE CERTIFICATES
\begin{corollary}[Primitive identity hierarchy defect and comparison-scale invariance]
\label{p5:cor:v3024-primitive-hierarchy-defect}
On the fixed common rooted three-family chart of
Theorem~\ref{p5:thm:v3211-hierarchy-corner-defect}, let
\(R=\lambda I_G\) with \(\lambda\in\mathbb C\setminus\{0\}\).
Then
\[
 \boxed{\Delta_{\rm hier}(\lambda I_G)=\frac49,\qquad
 \Delta_{\rm hier}(\lambda R')=\Delta_{\rm hier}(R')
 \ \text{for every nonzero }R'\text{ and }\lambda\ne0.}
\]
Thus a common nonzero scalar comparison does not turn the primitive
family-identity source into a zero-defect three-arrow source and cannot
by itself certify which role-specific polar comparison operation is
consumed.  This statement concerns the fixed normalized source chart;
it neither excludes a separately generated noncentral charged member
nor evaluates the native chiral action derivative.
\end{corollary}
\begin{proof}
For \(R=\lambda I_G\), the four source corners are
\(D=\lambda Q_0\), \(A=B=0\), \(C=\lambda Q_\perp\).
Consequently \(N=3|\lambda|^2\) and
\(\sigma_C=|\lambda|^4\), so the two contributions to the original
defect are \(1/3\) and \(1/9\).  Under the replacement
\(R'\mapsto\lambda R'\), \(D,A,B,C\) are multiplied by
\(\lambda\), \(N\) by \(|\lambda|^2\), and the quartic
numerator by \(|\lambda|^4\).  Both ratios are unchanged.
\end{proof}

\begin{corollary}[Sharp fixed-chart separation from the primitive identity source]
\label{p5:cor:v3024-rooted-identity-separation}
Let \(R\ne0\) obey \(\Delta_{\rm hier}(R)=0\) on the same
fixed three-family chart, and let \(\lambda\ne0\).  Then
\[
 \boxed{\|R-\lambda I_G\|_{\rm HS}\ge\sqrt2\,|\lambda|,
 \qquad
 \frac{\|R-\lambda I_G\|_{\rm HS}}
 {\|\lambda I_G\|_{\rm HS}}\ge\sqrt{\frac23}.}
\]
The bounds are attained in the degenerate zero-crossing class and are
sharp infima when all three hierarchy-arrow blocks are required to be
nonzero.  They are distances in a fixed Hilbert--Schmidt chart, not
physical energy barriers or source-generation prohibitions.
\end{corollary}
\begin{proof}
Zero defect requires \(D=0\) and \(\operatorname{rank}C\le1\)
by Theorem~\ref{p5:thm:v3211-hierarchy-corner-defect}.
The four rooted corners are mutually Hilbert--Schmidt orthogonal,
so \(\|D-\lambda Q_0\|_{\rm HS}^2=|\lambda|^2\).
The augmentation compression \(C\) has a unit kernel vector
\(v\in\operatorname{ran}Q_\perp\), whence
\(\|(C-\lambda Q_\perp)v\|^2=|\lambda|^2\) and
\(\|C-\lambda Q_\perp\|_{\rm HS}^2\ge|\lambda|^2\).
All other block contributions are nonnegative.
For sharpness choose a root-adapted orthonormal basis
\(e_0,e_1,e_2\), with \(e_0\) the rooted singlet, and take
\(R_0=\lambda|e_2\rangle\langle e_2|\).
It has zero defect and achieves equality.  For
\(R_\varepsilon=R_0+\varepsilon|e_1\rangle\langle e_0|
 +\varepsilon|e_0\rangle\langle e_1|\), the three blocks
\(A,B,C\) are all nonzero when \(\varepsilon\ne0\), and their
orientations make all defect terms vanish, while
\(\|R_\varepsilon-\lambda I_G\|_{\rm HS}^2
 =2|\lambda|^2+2|\varepsilon|^2\).  Let
\(\varepsilon\to0\).
\end{proof}
% END FRAMEWORK INTEGRATION v3.024


\begin{proposition}[Unrooted spectral moments do not determine hierarchy incidence]
\label{p5:prop:v3211-spectral-moment-orientation-nogo}
Let \(g=(1,-1,0)^{\mathsf T}/\sqrt2\) and
\(h=(1,1,-2)^{\mathsf T}/\sqrt6\), orthogonal to the rooted singlet
\(c\).  Set
\[
 R_0=cc^\dagger+2gg^\dagger+3hh^\dagger,
 \qquad R_1=VR_0V^\dagger,
\]
where \(V\) rotates the \((c,g)\)-plane by
\(\cos\theta=3/5,\ \sin\theta=4/5\) and fixes \(h\).  These
positive operators satisfy for every \(k\ge1\)
\[
 \Tr[(R_0^\dagger R_0)^k]
 =\Tr[(R_1^\dagger R_1)^k]=1+4^k+9^k,
\]
but have different rooted incidence:
\[
 \boxed{Q_0R_0Q_\perp=0,\qquad
   \|Q_0R_1Q_\perp\|_{\rm HS}^2=\frac{144}{625}.}
\]
Thus scalar spectral traces alone do not certify \(\mathsf{Cert}_{\rm hier}\)
or the mass-source consumption of the structural Hopf arrows.
\end{proposition}
\begin{proof}
Unitary conjugation preserves every listed singular-value moment.  In the
\((c,g,h)\)-basis the rotation produces a mixed \((c,g)\) coefficient of
absolute value \(|(2-1)(3/5)(4/5)|=12/25\); squaring yields \(144/625\).
The unrotated operator is block diagonal relative to \(Q_0\), so its
rooted crossing vanishes.  The counterexample concerns unrooted scalar
readouts, not an independently supplied noncommuting same-action parent.
\end{proof}

\begin{remark}[Source-origin status of the finite hierarchy tests]
\label{p5:rem:v3211-hurwitz-source-status}
The charged-lepton structural current-incidence closure of
Corollary~\ref{p5:cor:v3108-lepton-incidence-status} remains intact.
The charged mass-source handoff of
Theorem~\ref{p5:thm:v397-charged-preterminal-srcval-realization} continues
to have its declared \(\mathsf{Cert}_{\rm hier}\) antecedent; neither the
commutator representation nor \(\Delta_{\rm hier}\) evaluates the
independent actual preterminal source member.  The nine charged comparison
port assignments, K56 neutral mass-source selection, and interacting
pole/rest identifications remain separately typed obligations.
\end{remark}

% BEGIN v3.058 SHARP COMMUTING-HOPF SOURCE OBSTRUCTION
% Proposed mathematical insertion for Paper V-A after
% p5:rem:v3211-hurwitz-source-status. Baseline v3.217.
% STATUS: sharp restricted-operator no-go; no action-source promotion.
\begin{proposition}[Sharp Hopf-commutant obstruction to three distinct free left responses]
\label{p5:prop:v3217-hopf-commutant-left-return-nogo}
On the fixed common rooted charged-lepton family chart let
\[
 W=\operatorname{diag}(1,i,-i),\quad
 c=\frac1{\sqrt3}(1,1,1)^{\mathsf T},\quad
 P_0=cc^\dagger,\quad Q=I-P_0,\quad s^2=\frac89,
\]
\[
 P_\mu=s^{-2}QWP_0W^\dagger Q,\qquad
 P_e=Q-P_\mu,\qquad P_\tau=P_0.
\]
For any real positive squared finite coefficients
$E=(m_e^{\rm fin})^2$, $M_\mu=(m_\mu^{\rm fin})^2$,
$T=(m_\tau^{\rm fin})^2$, define the \emph{target} left return
$K_\star=EP_e+M_\mu P_\mu+TP_0$.
Then the exact least Hilbert--Schmidt error of any canonically normalized
chiral source coefficient $N\in M_3(\mathbb C)$ that commutes with
$W$ in this common source/target chart is
\[
 \boxed{\inf_{[N,W]=0}\|NN^\dagger-K_\star\|_{\rm HS}^2
  =\frac{11}{24}(E-M_\mu)^2
   +\frac16(E+M_\mu-2T)^2.}
\]
The infimum is attained within the commutant. The right side
vanishes only if $E=M_\mu=T$. Hence for distinct finite
coefficients the exact full three-role \emph{free left} response
cannot arise from any source word contained only in
$\mathbb C[W]$, even allowing arbitrary polynomial coefficients.
This is a fixed-common-chart no-go for a restricted action-source
class, not a proof of any actual charged source or an exclusion of
noncommuting rooted action words.
\end{proposition}
\begin{proof}
The eigenvalues $1,i,-i$ of $W$ are distinct. Thus its commutant
on this common source/target chart consists exactly of diagonal
matrices $N$, and every $NN^\dagger$ is diagonal and nonnegative.
The nearest diagonal matrix to $K_\star$ is its diagonal part;
its entries are strictly positive because $K_\star\succ0$,
so that projection is an attainable $NN^\dagger$ with diagonal
$N$. Since each diagonal $N$ is a quadratic polynomial in the
simple-spectrum $W$, this also proves sharpness within $\mathbb C[W]$.
Using the generated polar projectors, the upper-triangular entries
of $K_\star$ in the $W$-eigenframe are
\[
 \begin{split}
 (K_\star)_{12}&=\frac{4T-3E-M_\mu+3i(E-M_\mu)}{12},\\
 (K_\star)_{13}&=\frac{4T-3E-M_\mu-3i(E-M_\mu)}{12},\\
 (K_\star)_{23}&=\frac{4(T-M_\mu)+3i(E-M_\mu)}{12}.
 \end{split}
\]
Consequently the squared error of the diagonal projection is
\[
 4|(K_\star)_{12}|^2+2|(K_\star)_{23}|^2
 =\frac{11}{24}(E-M_\mu)^2
  +\frac16(E+M_\mu-2T)^2.
\]
Both squares must vanish for equality, which forces all three
squared coefficients to be equal.
\end{proof}
\begin{remark}[Observable scope of the noncommutative source requirement]
\label{p5:rem:v3217-hopf-commutant-source-scope}
The obstruction concerns the \emph{positive return} of a source that
commutes with $W$ in one common rooted left/right family frame. The canonical uninserted kinetic factor of Paper III-F preserves this
commuting property, but an independently generated noncentral internal
kinetic metric could change it under normalization. A
source-valued operator produced by a differently typed chiral
intertwiner need not commute with this $W$. The construction of
$P_e,P_\mu,P_\tau$ is upstream geometry, not a license to place them
in $\delta^2S/(\delta\bar\psi_L\delta\psi_R)$.
The restricted conclusion neither discharges
$\mathsf{Cert}_{\rm hier}$ nor supplies the actual mixed Hessian.
It isolates an experimentally relevant necessary kind of source
noncommutativity for the complete three-role free left response.
\end{remark}

% END v3.058 SHARP COMMUTING-HOPF SOURCE OBSTRUCTION
% END HURWITZ-GRAMMAR ROOT-SENSITIVE CHARGED SOURCE TESTS v3.211


% BEGIN CHARGED-COMPARISON-OWNER-CONSUMPTION v3.98
\subsection*{charged same-action comparison-owner consumption}
\label{p5:sec:v398-charged-comparison-owner-consumption}
The charged finite source of Release v3.97 remains unchanged.  The additional
result here concerns the finite comparison factor already used by the charged
spectrum readout, not the origin of the raw Clifford/split coefficient.

\begin{theorem}[Charged rooted comparison-owner consumption]
\label{p5:thm:v398-charged-rooted-comparison-owner-consumption}
On the rooted charged branch and on the stationary local quartic domain of
Corollary~\ref{p5:cor:v180-exact-icosian-higgs-profile-action}, the same action
already consumes the complete $m_{\rm cmp}=2$ scalar comparison band.
Paper~I-A, Corollary~\ref{p1:cor:v569-unique-quartic-comparison-support},
therefore fixes its finite exceptional comparison support to
$I\cong2I$, while the charged root supplies the occurring Hurwitz core
$H\cong2T\subset I$.

Let $X=I/H$.  Then $|X|=5$.  Taking the already declared Paper--V-B
source--target cell-label carrier as the predeclared relation-comparison target,
Paper~II-A, Corollary~\ref{p2a:cor:v3104-rooted-relation-corealization},
co-realizes
\[
 \mathcal L_{25}=\ell^2(X_L\times X_R)
\]
with the rooted identity label $(x_0,x_0)$ of the occurring $H$-member.  The
unique relative-comparison reduction is
\[
 Q_{\rm cmp}=I-|x_0,x_0\rangle\langle x_0,x_0|,
\]
and for the invariant label state
\[
 \boxed{\eta_H
 =\langle\Omega_{25},Q_{\rm cmp}\Omega_{25}\rangle
 =1-\frac1{25}=\frac{24}{25}.}
\]
Thus the charged same-action comparison consumes the already constituted
$25$-label exclusion owner with no additional occurrence or cell selector.
\end{theorem}

\begin{remark}[Comparison support, not a new microscopic symmetry]
\label{p5:rem:v398-comparison-support-firewall}
The theorem identifies the finite support used by this same-action comparison.
It does not assert that every microscopic field owner is globally $2I$-
symmetric, and it does not replace the local Hurwitz $2T$ source owner by
$2I$.  The charged raw mass coefficient, family incidence and preterminal
source of Release v3.97 are unchanged.
\end{remark}

\begin{corollary}[Charged exclusion-owner occurrence is dependent]
\label{p5:cor:v398-charged-exclusion-owner-dependent}
On the stated rooted charged comparison branch, the former conditional clause
``if the Higgs/Yukawa word consumes $\mathfrak E_H$'' is discharged.  The
remaining charged finite-comparison question is the owner-provenance grading
of the already occurring exclusion factor, supplied by Paper~V-B v5.91.
\end{corollary}
% END CHARGED-COMPARISON-OWNER-CONSUMPTION v3.98


% BEGIN FRAMEWORK v2.221: CHARGED COMPARISON SOURCE OWNERSHIP SCOPE
\begin{proposition}[Scope of the charged rooted comparison-owner theorem]
\label{p5:prop:v221-charged-comparison-source-scope}
Theorem~\ref{p5:thm:v398-charged-rooted-comparison-owner-consumption}
identifies a rooted finite comparison owner and its scalar expectation
$\eta_H=24/25$.  Its conclusion alone does not identify which
normalized chiral frame is consumed by the mass-producing left or
right source leg.  Paper~I-C,
Corollary~\ref{p1b:cor:v221-hurwitz-principal-angle}, supplies
normalized scalar comparison frames $J_0,J_1$ with overlap
$\sqrt{\eta_H}I$ on the 14-mode quadratic band, and
Paper~III-F, Proposition~\ref{p3f:prop:v221-relative-chiral-comparison-frame},
calculates their possible chiral source responses.  In particular,
relative versus common rooted chiral consumption gives inequivalent
canonically normalized mass operators.

The displayed charged odd source of Release v3.97,
\[
 D_{F,s}^{\rm ch}=\Sigma_{\sigma_s}\otimes P_{Q_s}
 \otimes\sum_r m_{s,r}^{\rm fin}\iota_{s,r}+\mathrm{h.c.},
\]
starts with terminal coefficients that already contain their
comparison grades.  Hence this downstream equation may transport the
coefficients, but cannot independently derive their primitive
rooted-port insertions.  The latter require a mixed action derivative
or an explicitly generated preterminal contraction identifying the
chiral frames and exterior owner sorts for each physical role.
\end{proposition}
\begin{remark}[Physical charged-owner certificate still outstanding]
\label{p5:rem:v221-charged-source-port-certificate}
This is a scope refinement, not a retraction of the existing
$25$-label geometry, the charged occurrence theorem, or the nine
finite numerical coefficients.  Full mass ownership requires the
actual preterminal source to select (i) a comparison-port operation,
(ii) its chiral frame placement, and (iii) the relevant exterior or
reciprocal owner operation \emph{without reading the comparison
exponent back out of the final mass table}.  The finite rest/pole
and current-residue tests remain separately typed.
\end{remark}
% END FRAMEWORK v2.221: CHARGED COMPARISON SOURCE OWNERSHIP SCOPE



% BEGIN DIRECT-MASS-SM-ONTOLOGY v3.99
\subsection{direct-mass Standard--Model ontology synchronization}
\label{p5:sec:v399-direct-mass-sm-ontology}

\begin{theorem}[Generated chiral slots are not native mass generators]
\label{p5:thm:v399-yukawa-slot-mass-firewall}
The four generated chiral source slots customarily denoted
$Y_u,Y_d,Y_e,Y_\nu$ remain valid coordinates of the gauge-compatible odd
source/interactions.  On the active direct-mass branch, however, the native
finite masses are generated upstream by their terminal geometric constructors
and shell/hierarchy readouts.  No theorem of Paper~V-A identifies those masses
with singular values of $vY_f/\sqrt2$, nor is a Higgs-VEV magnitude an upstream
mass selector.
\end{theorem}

\begin{corollary}[Weak-current rather than mass-eigenframe flavor target]
\label{p5:cor:v399-current-incidence-flavor-target}
For claims involving quark or lepton flavor change, the native target is the
family coefficient of the generated weak current between the already-massed
terminal role spaces.  A CKM or PMNS matrix may be introduced downstream as a
coordinate representative of that normalized current map when a conventional
comparison is useful; it is not a primitive relative mass-eigenframe required
by the finite carrier.
\end{corollary}

\begin{remark}[Historical eigenframe language]
Earlier references in this paper to Yukawa singular spectra or relative
left-eigenframes are retained as operator diagnostics and Standard--Model
comparison language.  They are superseded as statements of native mass
ontology by Theorem~\ref{p5:thm:v399-yukawa-slot-mass-firewall}.
\end{remark}
% END DIRECT-MASS-SM-ONTOLOGY v3.99

% BEGIN NEUTRAL-RESPONSE-TRANSPORT-SIBLING v3.100
\subsection{Reuse of the occurring source in the physical-support calculation}
\label{p5:sec:v2169-charged-support-application}

The physical source is still the already synthesized $R_{\rm act}$;
Theorem~\ref{p5:thm:v397-charged-preterminal-srcval-realization} and its
assembly corollary are not reopened. IV now generates a terminating tagged
polynomial on the actual stationary support. Its terminal root already
includes visibility there, by
Theorem~\ref{p4:v2169:b10:thm:support}. Thus the dynamic route consumes
$H,q_r,P$ and the fixed $\varepsilon_r$, with no independent self-energy,
density weights or extra overlap premise after that same-support root.
The free-return certificate remains an alternative stronger response claim.
If the state/physical chart is already certified, the first unproved finite
calculation is the actual complete $p_{r,P}$ and its source-visible tagged
root set.  The equality $p_{r,P}(\varepsilon_r)=0$ is then tested only for
the stronger terminal-coordinate--equals--pole claim; it is not imposed when
constructing the root.  State stationarity where unproved, named particle
interpretation and the limits actually consumed retain their own ordered
domains.

% END INLINE SOURCE: Paper5A-ChargedSupportApplication-Main-v1_1.tex



% v2.171: conversation proofs, with actual history obligations retained.

% BEGIN INLINE SOURCE: Paper5A-VectorChargeExport-Main-v1_0.tex
% Framework v2.171: integrated conversation proof on its stated domain.
\subsection{Scalar core terminalization and the neutral source-value scope}
\label{p5:sec:v2208-neutral-core-family}
The family factor $G$ in the constituted matter carrier and its coefficient
space $\End(G)$ are the ones already used in
\eqref{p5:eq:part1-Dirac-Yukawa}. The geometric bridge
$\Phi_{\rm IR}(A)=J_{\rm IR}^\dagger A J_{\rm IR}$ is
\ref{p1b:thm:v2208-ir-operator-bridge}.

There are two differently typed uses of $P_0$. In that geometric bridge it
is an operator element transported out of one geometric corner. In the
coefficient carrier below it acts on the Clifford factor, leaving the
separate family coefficient factor available.

\begin{theorem}[Family coefficients survive scalar core terminalization]
\label{p5:thm:v2208-terminal-tensor-extension}
The spectator extension of the core map is
\[
 \mathcal T_{56}=P_0\otimes\operatorname{id}_{\operatorname{End}(G)}:
 \mathcal A_{\mathbb C}\otimes\operatorname{End}(G)
 \longrightarrow\mathbb C\mathbf1\otimes\operatorname{End}(G).
\]
Its rank is nine, and for every family coefficient $B$,
\[
 \boxed{\mathcal T_{56}(\mathbf1\otimes B)=\mathbf1\otimes B.}
\]
Consequently scalar terminalization of the core alone imposes no rank-one
or equal-mass condition on a family coefficient.
\end{theorem}
\begin{proof}
Expand the coefficient as $\sum_\alpha a_\alpha\otimes B_\alpha$ in a
Clifford basis beginning with $a_0=\mathbf1$. The projector retains
exactly $\mathbf1\otimes B_0$. The nine family matrix units are linearly
independent in that range and are each fixed. Its rank is
$\rank P_0\,\dim\operatorname{End}(G)=1\cdot9$. In particular every
$3\times3$ coefficient can occur in the range of this map. Whether a
specific coefficient is generated and consumed is a separate constructor
statement, not an additional constraint supplied by this tensor identity.
\end{proof}
\noindent\textbf{Type distinction.}
$\Phi_{\rm IR}(P_0)=Q_c$ is one transported geometric element.
It is neither $\mathcal T_{56}$ nor a derived physical neutrino mass
coefficient. Applying that geometric identity as though it projected the
matter's independent family factor would be a type error. This theorem
establishes the tensor consequence of the core action, without claiming
that scalar projection is the entire mass compiler.


\begin{theorem}[Inherited binding on a certified new neutral export]
\label{p5:thm:v2208-neutral-boundary-export}
Suppose a result-independent licensed neutral constructor has as required
inputs and outputs the full harmonic/local-port packet of
\ref{p1b:sec:v2208-neutral-boundary-geometry}, its operation-matched response
and incidence siblings, and their predeclared odd-source insertion.
Then its required source value and every nonzero constituent used by the
same action inherit that occurrence. No second $B$-binding token or
independent normalization efficiency is introduced.
\end{theorem}
\begin{proof}
The geometry theorem fixes the evaluation, polar normalization and
projectors before source evaluation. Apply
\ref{p2a:thm:v3109-generated-member-owner-corealization} and
\ref{p2a:thm:v3101-generated-coconsumption} to the licensed required
output and insertion, and \ref{p3a:cor:v183-action-coconsumption} to
the constituents used by that action. This proves inherited occurrence;
it does not construct the hypothesized neutral origin arrow.
\end{proof}
The factor-local odd neutral slot then has its existing form
\[
 D_{F,\nu}=\Sigma_+\otimes P_\ell\otimes R_\nu+\mathrm{h.c.}
\]
by \ref{p5:thm:v389-factorized-family-source}. Deriving $R_\nu$ from a
licensed new terminal mass/incidence packet remains the actual target.
The charged realization theorem does not supply this neutral packet.

\begin{proposition}[Boundary-only consumption has no kernel response]
\label{p5:prop:v2208-boundary-only-kernel}
For the map $B$ of \ref{p1b:thm:v2208-boundary-coisometry}, if an action
on $\mathcal H_{12}$ factors as $S[f]=s[Bf]$, then every derivative
with a direction in $\ker B$ vanishes. In particular a quadratic return
$B^\dagger K_\partial B$ annihilates that kernel.
\end{proposition}
\begin{proof}
Since $BQ=0$, $S[f+tQv]=S[f]$ for all $t$; differentiate this identity.
The quadratic identity follows by multiplying by $Q$ on either side.
\end{proof}
This scoped action control does not exclude a licensed collective full-owner
trace, whose domain and consumed response differ from a boundary-only action.

The K16 internal-boundary pairing in
\ref{p5:thm:v3105-k16-bl2-pair-corealization} retains its own action
scope. As stated in \ref{p5:rem:v3105-boundary-not-terminal-majorana},
that pairing is not a terminal mass dictionary. Its production history or
normal-return evaluation is not an antecedent of the K56 geometric and
finite-coefficient constructions above. Downstream interacting readouts
retain their own actual-action domains.
% END NEUTRAL-CORE-FAMILY-TYPING v2.208

% BEGIN LITERATURE-VERIFIED STANDARD-MODEL DOMAINS v2.210
\subsection{Joint-incidence physical port tests}
\label{p5:app:joint-incidence-physical-port-tests}

These tests keep their stated input actions and marks.  Failure
of one substitution is not a no-go theorem for all \(M_8\)
geometry.

\begin{proposition}[Coordinate, color and axial target tests]
\label{p5:prop:joint-incidence-linear-port-tests}
The physical coordinate role is \(W_L\otimes W_{R+}^*\)
for the positive column, not \(W_{R+}^*\).  It is neutral
under \(B\), whereas \(\ell\) has \(B\)-weight \(-1\).
Consequently a fixed equivariant map
\(\ell\to V_{\mathrm{coord}}\) is zero; a nonzero element
of that Hom space would itself have \(B\)-weight \(+1\).
Moreover
\[
 \operatorname{Hom}_{SU(2)_L}(Q,V_{\mathrm{coord}})=0,\qquad
 \dim_{\CC}\operatorname{Hom}_{SU(3)}(Q,\operatorname{End}W_R)=4.
\]
Every map in the second space factors through \(\ell\).
\end{proposition}
\begin{proof}
Weights in a Hom space are target minus source weights.
The central element of \(SU(2)_L\) is \(+1\) on \(Q\)
and \(-1\) on the coordinate doublet, proving the first
vanishing.  Color is trivial on \(\operatorname{End}W_R\)
and nontrivial on \(C\), so a color intertwiner vanishes
on \(C\); the arbitrary image of a generator of \(\ell\)
has dimension \(\dim\operatorname{End}W_R=4\).
No positive symmetric power of a single defining weak
doublet has an \(SU(2)\)-invariant covector.  A two-column
weak determinant instead cancels the opposite right
weights and leaves the axial-square condition.
\end{proof}

\begin{proposition}[Stabilizers of stronger fixed seeds]
\label{p5:prop:full-seed-stabilizer-tests}
In the thirteen-dimensional candidate Lie algebra
\(\mathfrak{su}(2)_L\oplus\mathfrak{su}(3)_c
\oplus\RR T\oplus\RR B\), fixing the full quaternionic
Clifford map gives connected stabilizer dimension one.
Fixing both axial columns gives dimension nine; fixing
one axial column gives dimension ten.  These statements
use the displayed physical actions on the proposed fixed
seed, without additional source marks.
\end{proposition}
\begin{proof}
Normalize \(c(h)=I\), put
\(r(a)=\operatorname{diag}(e^{ia/2},e^{-ia/2})\) and
\(g_Q=e^{-ib}1_\ell\oplus e^{ib/3}U\).
Full invariance is
\[
 u\,c(q)=c(g_Qq)r(a)\quad\hbox{for all }q.
\]
At \(q=h\), \(u=e^{-ib}r(a)\); its determinant forces
\(b=0\) in the connected group.  On the complement,
\(U=\operatorname{Ad}_{r(a)}\), so only \(a\) is free.
For both axial columns the same equation fixes \(b=0\)
and \(u=r(a)\) but leaves color free, giving \(8+1\).
For one column it becomes
\(ue_+=e^{i(a/2-b)}e_+\), fixing a weak Cartan in terms
of two free abelian parameters and leaving color free,
giving \(8+2\).  The norm of a column is instead neutral.
\end{proof}

The full unital quaternion multiplication also fails to
intertwine the defining physical \(SU(3)\): on the imaginary
three-space it retains the metric and cross product, whose
unitary complex-linear stabilizer is \(SO(3)\).
Fixing the actual parent unit as a physical vector in
\(\ell\oplus C\), whose nonzero color component is fixed,
reduces \(SU(3)\) to \(SU(2)\).  If multiplication and the
axial mark are both fixed, the remaining rotations are
\(SO(2)\).  These are stronger retention signatures than
the projected quadratic.  For a real-quaternionic quadratic,
\(E^TSE=\bar S\); a nonzero positive diagonal block has
a nonzero negative companion.  The equations from the
two opposite Hopf weights impose rank-two phase
constraints if both are fixed on the same axial line.
That companion is not the true conjugate
\(W_{R-}\otimes\bar\ell\).

\begin{proposition}[Neutral operator closure and first order]
\label{p5:prop:neutral-joint-word-comparison-boundary}
Every member of the displayed gauge/coordinate joint family
commutes with \(B_F\).  Its per-family finite operator
envelope is contained in
\(\operatorname{End}(\CC^4)\oplus\operatorname{End}(\CC^{12})\),
of dimension \(160\), and cannot create a charged comparison
cross block by neutral closure alone.
For the toy algebra \(\CC\oplus\CC\), represented as
\(\rho(a)=\operatorname{diag}(a_+,a_-)\) on the
two-line comparison, the conjugate-swap opposite action
satisfies
\[
 [[D_m,\rho(a)],\rho^o(b)]
     =-\Delta(a)\Delta(b)D_m,\qquad
 \Delta(a)=a_+-a_-.
\]
For \(m\ne0\), a first-order compatible complex subalgebra
must therefore act equally on the two scalar summands.
\end{proposition}
\begin{proof}
In the joint family the background makes \(B_F\) parallel,
the gauge representation centralizes it, and the coordinate
odd field is blind to \(Q\).  Hence each term in
\(D_{5,B}+c_g(a)+\gamma_g\otimes\Phi\) commutes with \(B_F\).
Its eigenspaces have dimensions four and twelve.
A reducing projection commuting with the represented
algebra and \(D\) continues to commute under sums,
products, adjoints, commutators, resolvents and functional
calculus on their domains.  In a block-diagonal conjugate
completion the conjugate swap exchanges reducing blocks;
an expression \(A+JAJ^{-1}\) still leaves the absent cross
block zero.  Direct multiplication of the two-by-two
matrices gives the double commutator.  Taking \(b=a\)
shows that a complex subalgebra admitting \(D_m\ne0\)
must have \(\Delta(a)=0\) for all its elements.
\end{proof}

The possible \(B\)-weights of a complex bilinear on \(F\)
are \(-2,-2/3,+2/3\), none zero.  Thus a neutral family
alone supplies no such pre-Y bilinear.  This test says
nothing about a different licensed geometric operation
with the required charged output.

For the higher-invariant alternative, \(U_{12,L}\) gives
\(L^{12}\simeq A_{12,L}\).  A nonzero endpoint
\(\kappa:A_{12,L}\otimes\ell^{12}\to\CC\) is exactly
equivalent, by composition, to a nonzero
\(m_{12}\in(L_0^*)^{12}\).
The pair has phase equations
\(\gamma=12\alpha,\ \gamma+12\beta=0\), hence
\(\alpha+\beta=0\).  Independently fixing an image
frame imposes \(\gamma=0\) as an additional equation.
Squaring the degree-six and degree-eight invariants
and pairing with fixed axial dual powers gives
degree-twelve and degree-sixteen geometric candidates;
their physical weights would be \((-6,12)\) and
\((-8,16)\), but the endpoint/action certificate is
not generated by this weight calculation.

\subsubsection{Optional action tests}
\label{p5:app:optional-comparison-action-tests}
If the mathematical comparison direct sum is actually
included in a physical joint operator with its proved trace,
and the local scalar-action hypotheses apply, then
\(\operatorname{Tr}D^2=8\) is angularly constant and
a positive quartic coefficient multiplies
\(\operatorname{Tr}D^4=2Q(n)\), favoring the six axes.
The required operator inclusion and coefficient are
antecedents here.  All these traces and heat readouts are
invariant under \(m\mapsto e^{-2i\theta}m\), and hence
cannot select its coherent phase.

If a phase kinetic action is separately constituted with
\(\delta A_X=d\lambda,\ \delta\phi=10\lambda\), integration
by parts in its action variation gives the sector identity
\(\operatorname{div}E_{A_X}-10E_\phi=0\), with contributions
from every other charged field added in a complete owner.
This conditional identity does not imply that such an
action occurs, that \(m\) is nonzero, or that it is parallel.
% END M8 v1.78 SM03



% BEGIN M8 v1.78 TOP01
% Results: B05-GLOBAL, B08-POINCARE-COMPATIBILITY, B11-POSITION-SPIN-QUOTIENT, B12-POSITIVE-DENSITY-TYPE, B12-POINCARE-PARITY, B12-CLOCK-CANCELLATION, B12-QUADRATIC-SIGN-CLASS, B12-PARALLEL-HOLONOMY, B12-SPIN-SIGN-INCIDENCE, B12-POSITION-PULLBACK, B12-DYNAMICAL-PRESERVATION
\subsection{Owner-saturation firewall for the remaining member-occurrence and Yukawa-source ports}
\label{p5:sec:v184-owner-completion-correction}
The GP-derived joint-owner saturation theorem of Paper~II-A applies only after
constructor provenance and the physical target sort have already been
certified.  The v1.83 synchronization overstated this theorem in two places.
The corrected statements below retain the new mathematical reductions without
promoting an unlicensed cross-sort map or source word.

\begin{proposition}[Late fixed-descendant transport does not replace parent-relation PCAD]
\label{p5:prop:v184-yfix-port-firewall}
If one starts only from already separated descendants and asks for a new
physical transport into \(L_0\), common rooting, point evaluation on the
hosted orbit, and abstract fibre identification do not by themselves certify
such a late cross-sort map.  This fixed-input no-go is retained.

It does not apply to Proposition~\ref{p5:prop:v187-pcad-yfix}: that theorem
changes the primitive object by retaining the already generated parent
quadratic and descending its nonzero projected relation by PCAD before the
descendants are separated.  Structural hypercharge is therefore closed,
whereas a particular nonzero field/background/member occurrence remains a
separate G0/member statement.
\end{proposition}

\begin{proposition}[Operator grammar alone does not license the physical Yukawa source word]
\label{p5:prop:v184-yukawa-family-firewall}
The generated family incidence returns and normalized two-return words span
\[
 \operatorname{End}(G)=M_3(\mathbb C),
\]
so they provide a complete finite operator grammar.  That grammar by itself
does not produce a weighted dagger-quiver/source owner.  Release v1.86 closes
that \emph{family-level} target differently: the already generated support,
grading, coordinate-linearity, covariance and adjoint constraints define the
predeclared exact-solution family of
Theorem~\ref{p5:thm:v186-generated-yukawa-dirac-owner-family}, and
Paper~II-A retains that complete family as generated owner content.  This does
not select a numerical Yukawa member, an occurrence, a pole member, or the
separate 25-label source-word consumption.
\end{proposition}

\begin{remark}[Corrected framework status]
\label{p5:rem:v184-g13-g15-status}
The descendant-first classification is retained, and the structural G13/hypercharge
obligation is closed by PCAD in Proposition~\ref{p5:prop:v187-pcad-yfix}:
the relevant nonzero relation descends from the generated parent quadratic.  A mass, vacuum,
background-field member or other dynamical occurrence remains separately G0.
The structural G15 weighted Yukawa--Dirac owner family is now generated by
Theorem~\ref{p5:thm:v186-generated-yukawa-dirac-owner-family}; the local
same-action scalar jet and the clocked uninserted fermion jet are also
compiled.  The remaining G15 residual is sector-specific normalized
source/member occurrence together with action-owned mixed invariant readouts;
25-label consumption, RG/pole matching, vacuum selection and regulated
quantum CP/measure data remain downstream and are not selected by the owner
family.
These updates do not alter the generated character kernel or any conditional
downstream theorem.
\end{remark}
% END OWNER-COMPLETION v1.84 AUDIT CORRECTION



\subsection{charged primitive-owner compiler and contracted G15 residual}
\label{p5:sec:v372-charged-primitive-owner-g15}

The Paper-I owner results now remove the charged node-support ambiguity from
G15.  They do not by themselves prove occurrence of a physical normalized
Yukawa/source member.

\begin{theorem}[Charged primitive-owner support]
\label{p5:thm:v372-charged-primitive-owner-support}
On the charged finite grammar, after all already generated refinements, the
primitive non-scalar owner types are
\[
 \boxed{u^\perp,\qquad V,\qquad\Lambda^2V}
\]
of ranks \(2,4,6\), respectively.  More explicitly, the rank-one branch is
scalar; the rank-three \(q=5\) branch is annihilated by the
\(C_6\to C_3\) owner map; ranks four and six are the triality-vector and
coordinate-pair bivector owners; and ranks eight, twelve and sixteen are
nonprimitive/composite parent blocks on the charged branch.  Hence every
owner-natural charged node descendant factors through
\[
 \boxed{P_2\oplus P_4\oplus P_6}.
\]
\end{theorem}
\begin{proof}
Import Theorem~\texttt{p1b:thm:v125-\allowbreak q5-\allowbreak rank2-\allowbreak owner},
Theorems~\texttt{p1:thm:v558-\allowbreak triality-\allowbreak tetrad-\allowbreak vector-\allowbreak owner} and
\texttt{p1:thm:v558-\allowbreak coordinate-\allowbreak pair-\allowbreak bivector-\allowbreak owner}, and the generated
Hurwitz refinements summarized in Paper~V-B.  The statement is about support
and owner type only.
\end{proof}

\begin{corollary}[Charged G15 compiler layer is derived; same-action source occurrence remains]
\label{p5:cor:v372-g15-contracted-residual}
The charged finite compiler no longer requires independent row-support choices
for \(u,\mu,s\), nor independent radicand choices for the charged split words.
The remaining G15 gate is
\[
 \boxed{\texttt{SAME\_ACTION\_PRIMITIVE\_OWNER\_SOURCE\_OCCURRENCE}}.
\]
Concretely, one must prove that the actual finite fermion source/Hessian uses
the generated primitive-owner incidence maps as its same-action couplings and
has the finite odd chiral-source target sort before Feshbach evaluation.  Only
then may one form an origin certificate of the form
\[
 \operatorname{Orig}_{\Sigma}
 \bigl(w_{\rm mass};D_\star,P_L\operatorname{End}(F)P_R\bigr).
\]
The relative family member, CP state and pole/RG handoffs remain downstream.
\end{corollary}

\begin{remark}[Mass-PCAD target after support contraction]
\label{p5:rem:v372-mass-pcad-target}
The sharpened target word is
\[
 D_\star\longrightarrow\mathbb T_{\rm Hur/Cl}
 \longrightarrow\mathbb D_\star^{\rm fin}
 \xrightarrow{\operatorname{Fesh}_{Q_{\rm int}}}
 \Xi_{\rm mass}
 \longrightarrow P_L\operatorname{End}(F)P_R.
\]
Paper~I-C proves the odd Schur target shape once the paired same-owner parent
is present.  This release does not assert that physical same-action parent;
therefore G15 remains open at the narrower gate above.
\end{remark}
% END CHARGED-PRIMITIVE-OWNER-G15 v3.72

% BEGIN D4-TERMINAL-NEUTRAL-OCCURRENCE v3.73
\subsection{Hurwitz--\texorpdfstring{\(D_4\)}{D4} terminal-neutral occurrence target}
\label{p5:sec:v373-d4-terminal-neutral-occurrence}

Paper~I-A now supplies, on the same rooted Hurwitz--\(D_4\) geometry, the
closed oriented-frame determinant return
\[
 J_{D_4}^{\dagger}J_{D_4}=192
\]
and the unique positive central rank-three lift
\[
 R_{D_4}=xI_G,\qquad x=192^{1/3},
\]
whenever the determinant-line scalar is lifted centrally to a rank-three
carrier \(G\).  This is mathematical capacity until the physical terminal
neutral source is proved to consume it.

\begin{definition}[Terminal-neutral \(D_4\) determinant occurrence gate]
\label{p5:def:v373-n56-d4-occurrence-gate}
Write
\[
 \mathsf{N56}_{D_4}^{\rm occ}
\]
for the same-action statement that the actual \(K=56\) neutral fermion source
uses the Paper-I-A oriented-frame determinant return as its terminal finite
response and uses its minimal positive central lift on the already generated
rank-three family carrier \(G\).  Concretely the typed word is
\[
 D_4\text{ Hurwitz cell}
 \longrightarrow\mathcal F_{D_4}
 \longrightarrow\mathscr D_{D_4}
 \xrightarrow{\ J_{D_4}^{\dagger}J_{D_4}\ }
 \det G
 \xleftarrow{\ \Lambda^3\ }
 R_{\nu,56}\in Z(\End G)^+ .
\]
The gate includes physical occurrence of the terminal response but does not
include an absolute dimensionful normalization or a choice of neutral
mass-eigenbasis.
\end{definition}

\begin{proposition}[Consequence of the terminal-neutral occurrence gate]
\label{p5:prop:v373-n56-d4-central-response}
Conditional on \(\mathsf{N56}_{D_4}^{\rm occ}\),
\[
 \boxed{R_{\nu,56}=xI_G,\qquad x^3=192,\qquad x=192^{1/3}.}
\]
Thus no independent numerical terminal-neutral shape coefficient is inserted
at this stage.
\end{proposition}

\begin{proof}
Apply Corollary~\texttt{p1:cor:v559-positive-central-rank3-lift} to the
physical family carrier supplied by the occurrence gate.
\end{proof}

\begin{remark}[What remains physical]
\label{p5:rem:v373-neutral-residual}
This release does not prove \(\mathsf{N56}_{D_4}^{\rm occ}\).  It also does
not fix the dimensionful neutral normalization, the mass-eigenprojectors, or
the PMNS frame.  The result narrows the old terminal-neutral coefficient
problem from several unrelated numerical weights to one typed same-action
occurrence statement plus those downstream normalization/frame obligations.
The charged G15 gate of Corollary~\ref{p5:cor:v372-g15-contracted-residual}
remains unchanged.
\end{remark}
% END D4-TERMINAL-NEUTRAL-OCCURRENCE v3.73

% BEGIN D4-NEUTRAL-SOURCE-CONTRACTION v3.74
\subsection{\texorpdfstring{\(D_4\)}{D4} triality-source consumption and the G15 occurrence contraction}
\label{p5:sec:v375-d4-triality-source-consumption}

Papers~I-A v5.61 and I-B v1.26 now identify a common finite parent behind the
oriented Hurwitz-frame determinant response and the protected Clifford source
channel.  This permits the framework's existing action-ownership rules to be
applied to the former G15 same-action occurrence gate.

\begin{definition}[Rooted \(D_4\) triality mass-source constructor]
\label{p5:def:v375-rooted-d4-triality-mass-source}
Let $\Theta_{\rm Spin(8)}$ be the signed, metric-normalized coefficient
tensor of Paper~I-A, Theorem~\ref{p1:thm:v214-normalized-spin8-triality}.
The unsigned tagged incidence vector $\Theta_{D_4}^{\rm ord}$ records its
support; its counting norm is not used as a coupling normalization.
Use the spectator coefficient extraction $\operatorname{Ext}_B$ of
Paper~I-C, Corollary~\ref{p1b:cor:v126-triality-restriction-is-clifford}.
On an occurring rooted Hurwitz frame, contract the signed tensor's
vector slot with the already constituted Higgs/vector direction \(\hat h\),
restrict to the protected chiral block, and apply the already generated
sector/family projectors.  Denote the resulting typed odd zero-order source
word by
\[
 \mathcal C_{\rm TRI}(D_\star)
 \in P_L\End(F)P_R\oplus P_R\End(F)P_L.
\]
With this specified extraction, Paper~I-C computes its primitive chiral
factor to be exactly $c(\hat h)$. A unit-singlet contraction is a different
map and carries its explicitly computed factor $\sqrt2$; it is not substituted
for the coefficient extraction in this constructor.
\end{definition}

\begin{theorem}[D4 triality-source consumption]
\label{p5:thm:v375-d4-triality-source-consumption}
Assume the already generated finite fermion-source role and an actually
occurring nonzero odd zero-order mass/source member in that role.  Then the
claim-complete source/Hessian owner cannot omit
\(\mathcal C_{\rm TRI}(D_\star)\).  In every GP2-exact realization of that same
source role,
\[
 \boxed{\mathcal C_{\rm TRI}(D_\star)\text{ is retained as a generated,
 nonredundant source port}.}
\]
Consequently its generated rank-\(2,4,6\) primitive-owner descendants are
co-realized in the same action/source owner.  No second physical occurrence
token is required for those descendants.
\end{theorem}

\begin{proof}
Paper~I-C v1.26 proves that the protected restriction of the generated
triality parent is the already normalized Clifford source channel
\(c(\hat h)\), and that the rank-\(2,4,6\) owner types are descendants of the
same rooted parent.  Hence the triality contraction is not merely a
mathematically available Hom-space: it has licensed constructor provenance in
the declared odd source sort.

For the mass/source claim, the consumed zero-order source and Hessian belong to
the claim-complete action owner of Paper~III-A.  Omitting the nonzero protected
triality/Clifford contraction changes that consumed source/Hessian.  Therefore
it is not in the collectively invisible omitted sector.  The generated-action
constructor closure theorem forbids deleting a generated source port from the
same claim-complete owner.  Equivalently, Paper~II-A's no-idle licensed
constructor corollary and generated joint-owner saturation theorem retain each
required nonredundant output in its declared sort.  The rank-\(2,4,6\)
descendants then follow by finite generated-image inheritance.
\end{proof}

\begin{corollary}[Retirement of the independent G15 same-action occurrence gate]
\label{p5:cor:v375-g15-source-occurrence-retired}
The former gate
\[
 \texttt{SAME\_ACTION\_PRIMITIVE\_OWNER\_SOURCE\_OCCURRENCE}
\]
is not an independent theoretical assumption once the triality ancestry is
included.  Conditional only on G0 actuality of a nonzero fermion mass/source
member, same-action consumption of the generated primitive-owner compiler is a
GP2/action-closure descendant.

This does \emph{not} let GP2 manufacture a nonzero source member: if the zero
member is compatible with all frozen relations, nonzero actuality remains G0.
Nor does this corollary by itself determine a dimensionful mass normalization,
a noncentral family eigenframe/CKM or PMNS matrix, CP state/measure, or pole/RG
matching.
\end{corollary}

\begin{corollary}[No independent electroweak-vacuum occurrence token for the solder]
\label{p5:cor:v377-no-independent-ew-vacuum-token}
Conditional only on the existing G0 actuality of a nonzero odd mass/source
member, Theorem~\ref{p5:thm:v375-d4-triality-source-consumption} and
Theorem~\ref{p5:thm:v377-unit-solder-normalized-scalar} co-realize the
normalized electroweak solder direction in the same source owner.  No second
\(\mathsf H_{\rm vac}\) occurrence token is required to constitute its
stabilizer.  This does not select a scalar stationary history, a pole mass, or
a physical scalar vertex.
\end{corollary}

\begin{corollary}[Neutral \(K=56\) consumption is dependent, not independent]
\label{p5:cor:v375-n56-consumption-dependent}
On an actually occurring nonzero neutral source coordinate, the relation
\(\mathsf{N56}_{D_4}^{\rm consume}\) of v3.74 is a descendant of the same
triality-source theorem: the global closed contraction of the consumed parent
is the already generated central family response
\[
 R_{D_4}^{\rm fam}=192^{1/3}I_G.
\]
Hence \(\mathsf{N56}_{D_4}^{\rm consume}\) is removed as an independent
neutral assumption.  The live neutral obligations are the dimensionful
normalization owner and the noncentral relative-eigenframe/PMNS theorem.
\end{corollary}

\begin{remark}[What G15 still names]
\label{p5:rem:v375-g15-retyped-frontier}
The source-consumption obstruction of G15 is discharged.  The label G15 is
retained only as a convenient umbrella for any still-unproved \emph{noncentral
rooted family-member evaluation} and its downstream CP/pole/RG readouts.  It
must no longer be described as a missing same-action occurrence of the
rank-\(2,4,6\) primitive owner maps.  A future explicit evaluation of the
rooted triality contraction may contract that remaining family-member layer
further.
\end{remark}
% END D4-TRIALITY-SOURCE-CONSUMPTION v3.75


% BEGIN D4-ROOTED-HALFSTEP-FRONTIER v3.76
\subsection{rooted source orientation covariance and determinant phase}
\label{p5:sec:v378-rooted-source-orientation-covariance}

The v3.75--v3.76 source theorem leaves one primitive noncentral rooted line,
\(\RR C_{2,G}\), equivalently the two oriented half-step transports
\(X_G\) and \(X_G^\dagger\).  Their common rooted parent also fixes how a
complete orientation reversal acts on every orientation-odd descendant.

\begin{theorem}[Rooted source orientation covariance]
\label{p5:thm:v378-rooted-source-orientation-covariance}
Let
\[
 X_G=\begin{pmatrix}0&0&1\\1&0&0\\0&1&0\end{pmatrix},
 \qquad
 R_G=\operatorname{diag}(1,-i,i),
 \qquad
 C_{2,G}=-i(X_G-X_G^\dagger).
\]
The complete rooted orientation reversal acts simultaneously as
\[
 R_G\longmapsto R_G^\dagger,
 \qquad
 X_G\longmapsto X_G^\dagger,
 \qquad
 C_{2,G}\longmapsto-C_{2,G}.
\]
Consequently the two claim-complete rooted orientation cores are a conjugate
pair.  The finite source grammar contains no additional independent relative
\(U(1)\) determinant-phase mark.
\end{theorem}
\begin{proof}
The first two transformations are the rooted Hopf/half-step reversals already
proved upstream.  The third follows algebraically from the definition of
\(C_{2,G}\).  Any extra relative phase between the two reversed source cores
would require a new noncentral orientation datum not generated by the rooted
triality parent, contrary to the v3.76 exhaustion and GP2/no-insertion.
\end{proof}

\begin{corollary}[Quark determinant phase is a rooted descendant]
\label{p5:cor:v378-quark-determinant-phase-descendant}
On a nonzero full-rank quark-source branch, write the two rooted source blocks
as
\[
 D_u=P_u^{(L)}A_+P_u^{(R)},
 \qquad
 D_d=P_d^{(L)}A_-P_d^{(R)},
\]
where all \(P\)'s are positive polar/radial factors and the rooted orientation
cores obey \(A_-=\overline{A_+}\).  Then
\[
 \boxed{\det(D_uD_d)>0,
 \qquad
 \frac{\det(D_uD_d)}{|\det(D_uD_d)|}=1.}
\]
This fixes only the determinant orientation.  The numerical radial spectrum,
relative eigenframe, CKM/PMNS member and pole/RG realization remain downstream.
\end{corollary}
\begin{proof}
Positive factors have positive real determinants, while
\(\det A_- = \overline{\det A_+}\).  Hence
\[
 \det(D_uD_d)=\Bigl(\prod\det P\Bigr)|\det A_+|^2>0.
\]
\end{proof}

\begin{remark}[Weak CP is not the determinant phase]
\label{p5:rem:v378-weak-cp-separate}
The corollary does not require the two positive Hermitian family operators to
commute.  Weak-family CP resides in relative eigenframe invariants such as
\(\operatorname{tr}[H_u,H_d]^3\), not in an independently free overall quark
determinant phase.  That relative-frame/member problem remains under G15.
\end{remark}
% END ROOTED-SOURCE-ORIENTATION-COVARIANCE v3.78


% BEGIN HSD4-BORN-Q5-RETYPE v3.80
\subsection{rooted \texorpdfstring{$C_3$}{C3} source-cell co-realization and G15 closure}
\label{p5:sec:v381-c3-source-cell}

The v3.75 source-consumption theorem, v3.76 identity/even-transport theorem,
and v3.80 Born-odd q=5 theorem are descendants of one rooted triality/source
owner.  The former residual treated one oriented half-step as though the
physical source had to choose it.  The correct finite object is instead the
complete rooted transport cell carried by that source owner.

\begin{definition}[Rooted \(C_3\) source transport cell]
\label{p5:def:v381-c3-source-cell}
On an actually occurring nonzero odd fermion-source branch of
Theorem~\ref{p5:thm:v375-d4-triality-source-consumption}, freeze the family
transport presentation
\[
 \boxed{\mathcal C_{C_3}^{\rm src}:=(I_G,C_{1,G},C_{2,G})},
\]
where
\[
 C_{1,G}=X_G+X_G^\dagger,
 \qquad
 C_{2,G}=-i(X_G-X_G^\dagger).
\]
The three ports are not imported from the ambient algebra
\(\End(G)\): \(I_G\) and \(C_{1,G}\) are the axial/even descendants of
Proposition~\ref{p5:prop:v376-even-rooted-transport-generated}, while
\(C_{2,G}\) is the same-root Born-odd q=5 descendant of
Theorem~\ref{p5:thm:v380-hsd4-geometry-closed}.  Their common occurrence
owner is the retained rooted triality/source port of v3.75.
\end{definition}

\begin{lemma}[Nonredundancy of the three source-cell ports]
\label{p5:lem:v381-c3-source-cell-nonredundant}
With the Hilbert--Schmidt pairing on the rooted trimer,
\[
 \operatorname{tr}(I_GC_{1,G})=
 \operatorname{tr}(I_GC_{2,G})=
 \operatorname{tr}(C_{1,G}C_{2,G})=0,
\]
and
\[
 \operatorname{tr}(I_G^2)=3,
 \qquad
 \operatorname{tr}(C_{1,G}^2)=
 \operatorname{tr}(C_{2,G}^2)=6.
\]
Hence the identity, orientation-even and orientation-odd ports are linearly
independent.  They are also distinguished by their frozen source signatures:
zero family displacement, reversal-even transport, and Born/reversal-odd
q=5 transport respectively.  Deleting any one changes the generated rooted
transport presentation.
\end{lemma}
\begin{proof}
For \(X_G^3=I_G\) one has
\(\operatorname{tr}X_G=\operatorname{tr}X_G^2=0\) and
\(X_G^\dagger=X_G^{-1}=X_G^2\).  Direct expansion gives the displayed
orthogonality and norms.  The source-signature distinction follows from
v3.76 and v3.80 and is therefore fixed before evaluation.
\end{proof}

\begin{theorem}[GP2 co-realization of the rooted source cell]
\label{p5:thm:v381-c3-source-cell-corealization}
Conditional only on the already separate G0 actuality of a nonzero odd
fermion-source member, the frozen source presentation
\(\mathcal C_{C_3}^{\rm src}\) is a finite generated-constituent
presentation in the sense of Paper~II-A.  Therefore GP2 exactness, typed host
faithfulness and no-omission give
\[
 \boxed{\operatorname{CoReal}_{\rm src}
 \{I_G,C_{1,G},C_{2,G}\}.}
\]
No additional family selector, occurrence token, sign-line coefficient or
new geometric principle is required.
\end{theorem}
\begin{proof}
Theorem~\ref{p5:thm:v375-d4-triality-source-consumption} retains the rooted
triality contraction as a generated nonredundant port of the occurring
claim-complete source owner.  Proposition~\ref{p5:prop:v376-even-rooted-transport-generated}
and Theorem~\ref{p5:thm:v380-hsd4-geometry-closed} generate the three declared
family-transport descendants from that same rooted port.  The preceding lemma
proves nonredundancy, and no port is introduced only by ambient availability.
Paper~II-A's GP2 co-realization theorem for finite generated-constituent
presentations therefore retains all three simultaneously in their native
source sort.
\end{proof}

\begin{corollary}[Equivalent cyclic transport presentation]
\label{p5:cor:v381-c3-transport-presentation}
The exact basis change
\[
 X_G=\frac12(C_{1,G}+iC_{2,G}),
 \qquad
 X_G^{-1}=X_G^\dagger=\frac12(C_{1,G}-iC_{2,G})
\]
is invertible.  Hence the same realized source cell has the equivalent rooted
cyclic presentation
\[
 \boxed{\mathscr T_{C_3}^{\rm src}=\{I_G,X_G,X_G^{-1}\}.}
\]
This is one co-realized finite source cell, not three independently prepared
physical tokens.
\end{corollary}

\begin{corollary}[Rooted conjugation orbit]
\label{p5:cor:v381-c3-conjugation-orbit}
For any rooted trimer reflection \(S_c e_k=e_{c-k}\), indices modulo three,
\[
 S_cX_GS_c^{-1}=X_G^{-1},
 \qquad
 S_cI_GS_c^{-1}=I_G.
\]
Thus the realized transport cell has the unique inversion-orbit split
\[
 \boxed{\{I_G\}\sqcup\{X_G,X_G^{-1}\}.}
\]
There is no invariant distinction between which member of the two-cycle is
called the matter orientation; complete rooted reversal exchanges them.
\end{corollary}

\begin{corollary}[Finite geometric conjugation and physical retyping]
\label{p5:cor:v381-geometric-conjugation}
Let \(J_F:F\to\bar F\) be the canonical anti-linear core conjugation and set
\[
 \Theta_c=(1\otimes S_c)J_F.
\]
Then
\[
 \boxed{(F,X_G)\stackrel{\Theta_c}{\longleftrightarrow}
 (\bar F,X_G^{-1})},
 \qquad
 \Theta_c I_G\Theta_c^{-1}=I_G.
\]
Core conjugation, not family inversion alone, reverses gauge weights: a
\(U(1)\) weight \(y\) goes to \(-y\), and a color \(\mathbf3\) goes to
\(\bar{\mathbf3}\).  After structural realization one may therefore name the
nontrivial inversion pair matter/antimatter orientations and the fixed
identity member the neutral/fixed \emph{transport} orientation.  Full electric
neutrality or Majorana self-conjugacy additionally requires the core
representation to satisfy \(F_0\simeq\bar F_0\); for an Abelian weight this
forces \(y_0=0\).
\end{corollary}

\begin{corollary}[Conjugate mass spectra agree]
\label{p5:cor:v381-conjugate-mass-spectrum}
If \(T:F_R\to F_L\) is a finite mass map and
\(T^c=\Theta_LT\Theta_R^{-1}\), then
\[
 (T^c)^\dagger T^c
 =\Theta_R(T^\dagger T)\Theta_R^{-1}.
\]
Hence particle and conjugate-particle singular spectra coincide, in particular
\[
 \boxed{m_{\bar p}=m_p}.
\]
No second independent antimatter Yukawa or mass parameter is introduced.
This statement does not determine the numerical matter spectrum.
\end{corollary}

\begin{corollary}[Closure of the finite G15 transport selector]
\label{p5:cor:v381-g15-transport-closed}
The former selector-type gate
\[
 \mathsf{D4\_Q5\_COMPLEMENT\_TRANSPORT\_OCCURRENCE}:\quad U_Q=X_G
\]
is retired as a live finite closure target.  It was a choice of one oriented
representative of a source cell whose complete rooted transport content is
co-realized.  Record instead
\[
 \boxed{\mathsf{G15}_{C_3}^{\rm src}
 =\mathsf{CLOSED\_DEPENDENT\_ON\_G0\_NONZERO\_FERMION\_SOURCE\_ACTUALITY}.}
\]
When an oriented Feshbach block is written with \(U_Q=X_G\), that equality is
only a matter-orientation representative convention; complete rooted reversal
uses \(X_G^{-1}\).  Relative CKM/PMNS eigenframes, G0 actuality itself,
Majorana occurrence and pole/RG matching remain under their separately typed
downstream owners and are not reopened by this closure.
\end{corollary}

\begin{remark}[Retyping of the old identity countermodel]
\label{p5:rem:v381-identity-fixed-branch}
The old observation that Born covariance also permits an identity family
factor is no longer a nonclosure certificate.  It is the required fixed member
\(I_G^{-1}=I_G\) of the complete source cell.  Symmetry still does not select
one oriented member; the point is that no such one-member selection is a
physical premise after source-cell co-realization.
\end{remark}
% END G15-C3-SOURCE-CELL v3.81




% BEGIN G15-UPSTREAM-GENERATED-CELL-HARDENING v3.82
\subsection{upstream-generated \texorpdfstring{$q=5$}{q=5} transport-cell hardening}
\label{p5:sec:v382-upstream-generated-q5-cell}

Paper~I-C v1.30 packages the rooted Hurwitz six-cycle as a regular $C_3$
local system and defines, before physical evaluation, the single generated
transport object
\[
 \mathfrak C_{q5}^{\rm geom}
 =(I_G,C_{1,G},C_{2,G}),
\]
with equivalent cyclic presentation
\[
 \{I_G,X_G,X_G^{-1}\}.
\]
This strengthens the logical interface used in v3.81: the source cell is not
formed by retrospectively collecting three independently available
endomorphisms.

\begin{theorem}[Upstream-generated composite status of the G15 source cell]
\label{p5:thm:v382-upstream-generated-source-cell}
On an actually occurring nonzero odd fermion-source branch of
Theorem~\ref{p5:thm:v375-d4-triality-source-consumption}, the v3.81 source cell
is the source-owned realization of the pre-existing Paper-I-B constructor
$\mathfrak C_{q5}^{\rm geom}$.  Its three projections are respectively the
already generated axial identity, reversal-even transport and Born-odd
$q=5$ transport:
\[
 \boxed{
 \jmath_{\rm src}(\mathfrak C_{q5}^{\rm geom})
 =\mathcal C_{C_3}^{\rm src}
 =(I_G,C_{1,G},C_{2,G}).}
\]
Consequently the finite generated-constituent hypothesis used in
Theorem~\ref{p5:thm:v381-c3-source-cell-corealization} is satisfied by an
upstream-generated composite, not by post-evaluation tuple formation.
\end{theorem}

\begin{proof}
Paper~I-C v1.30 defines $\mathfrak C_{q5}^{\rm geom}$ from the rooted
Hurwitz-link six-cycle, its two edge orbits and the aligned axial
identification.  Therefore the identity, even and odd ports are projections
of one constructor output before any physical source member is evaluated.
Theorem~\ref{p5:thm:v375-d4-triality-source-consumption} supplies the occurring
rooted source owner.  Proposition~\ref{p5:prop:v376-even-rooted-transport-generated}
identifies the source images of the identity/even projections, and
Theorem~\ref{p5:thm:v380-hsd4-geometry-closed} identifies the source image of
the odd projection.  All three therefore descend from one already generated
mathematical transport object in one retained source owner.  Applying the
Paper-II-A generated-constituent theorem to that frozen composite is thus
legitimate without introducing a new sum, product or tuple after evaluation.
\end{proof}

\begin{corollary}[G15 closure is unchanged and proof-hardened]
\label{p5:cor:v382-g15-closure-hardened}
The status
\[
 \boxed{
 \mathsf{G15}_{C_3}^{\rm src}
 =\mathsf{CLOSED\_DEPENDENT\_ON\_G0\_NONZERO\_FERMION\_SOURCE\_ACTUALITY}}
\]
from Corollary~\ref{p5:cor:v381-g15-transport-closed} is retained.  The v3.82
result strengthens only the generated-composite premise of the proof.  It does
not derive G0 source actuality, a numerical Yukawa member, relative CKM/PMNS
frames, Majorana occurrence, or pole/RG matching.
\end{corollary}

\begin{remark}[No use of the primitive-owner retraction in G15]
\label{p5:rem:v382-no-owner-retraction-in-g15}
Paper~I-C v1.30 also proves a canonical scalar-owner retraction for finite
unitary associated modules.  That map is not used in the G15 closure and is
not needed to define the primitive gravitational curvature, which is already
fixed independently by the momentum-curvature theorem
\(U=(3/\pi)\mathcal K_p\), \(\mathcal K_p=6/R_{\rm mom}^2\).  The retraction remains an independent mathematical
capacity rather than an action-trace replacement.
\end{remark}
% END G15-UPSTREAM-GENERATED-CELL-HARDENING v3.82


% END SAFEGUARD-HISTORICAL-G15-PROVENANCE v3.88

% BEGIN FERMION-SOURCE-BOUNDARY v3.85
\subsection{fermion-source structural closure and the G0 actuality boundary}
\label{p5:sec:v385-fermion-source-boundary}

The v3.75 triality theorem and the v3.81--v3.82 rooted $C_3$ source-cell
theorems remain valid.  The present historical release isolates the generic
port-only firewall.  The normalized triality/Clifford source port is generated.
Paper~II-C shows that the raw vanishing protected-coordinate completion face
is not in the physical fixed-radius coordinate essential image.  Generation of
an owner family and a nonzero port still does not, by itself, actualize a
nonzero physical source value.  Release v3.92 later supplies an independent
geometry-first nonzero-member constructor on the rooted branch; the generic
firewall stated here remains valid for unresolved families.

\begin{theorem}[Fermion-source port-only zero-member firewall]
\label{p5:thm:v385-fermion-source-zero-branch}
Generation of the owner family $\mathfrak Y$ together with a nonzero normalized
$\mathsf{SrcPort}$ does not, by GP1--GP3 alone, imply a nonzero
$\mathsf{SrcVal}$.  If no independent licensed constructor generates a
specific nonzero member as a required output, the zero member remains an
admissible realization of the unresolved physical source role.
\end{theorem}

\begin{proof}
Theorem~\ref{p5:thm:v186-generated-yukawa-dirac-owner-family} generates
the complete owner family $\mathfrak Y\simeq M_3(\mathbb C)^4$, including
its zero member, without selecting a numerical point.  Paper~II-A's
zero-member test, Corollary~\ref{p2a:cor:v399-nonzero-port-zero-coordinate},
and the typed cross-relation firewall therefore prohibit promotion of a
nonzero normalized source port to source-value actuality when no independent
required-output constructor fixes a member.  Release v3.92 supplies exactly
the allowed exception on the rooted triality branch: the nonzero finite member
is generated independently from the occurring root and is then consumed by
the already constituted fermion-action constructor.  Thus the present theorem
remains the generic firewall but no longer furnishes a countermodel to that
completed rooted word.
\end{proof}

\begin{corollary}[G15 structural retirement]
\label{p5:cor:v385-g15-structural-retirement}
The finite G15 source problem is structurally closed:
\[
 \boxed{\texttt{G15\_STRUCTURAL = CLOSED}.}
\]
Conditional on an actually occurring nonzero odd source member, the existing
triality-source theorem and rooted generated-constituent theorems consume the
complete $C_3$ source/transport cell and its generated rank-$2,4,6$ structural
descendants without an additional selector or occurrence gate.  The nonzero
physical source/member itself is retyped as
\[
 \boxed{\texttt{G0\_FERMION\_SOURCE\_MEMBER = OCCURRENCE\_BOUNDARY}.}
\]
It is not a residual G15 constructor target.
\end{corollary}

\begin{remark}[Supersession of the old ``what G15 still names'' wording]
\label{p5:rem:v385-g15-wording-superseded}
Remark~\ref{p5:rem:v375-g15-retyped-frontier} is retained as historical
provenance for the derivation path, but its remaining-family-member wording is
superseded by Corollary~\ref{p5:cor:v385-g15-structural-retirement}.  The
presentation-free physical family/source member is an actual G0 specialization
of the generated source-owner role, not a missing structural selector.
\end{remark}

\begin{remark}[No promotion of diagnostic mass words]
\label{p5:rem:v385-no-diagnostic-promotion}
This retyping does not promote any exploratory Schur/Feshbach family word into
the actual physical mass member, does not insert the $24/25$ exclusion scalar
into fermion masses, and does not identify terminal structural words with pole
masses.  Such identifications require their own predeclared ancestry/cross-sort
constructors.  Majorana occurrence and pole/RG certificates remain separately
typed.
\end{remark}
% END FERMION-SOURCE-BOUNDARY v3.85


% BEGIN TRANSITIVE-FERMION-SOURCE-HANDOFF v3.87
\subsection{transitive fermion-source handoff}
\label{p5:sec:v387-transitive-fermion-source-handoff}

The structural result of Theorem~\ref{p5:thm:v186-generated-yukawa-dirac-owner-family}
and the actuality boundary of
Corollary~\ref{p5:cor:v385-g15-structural-retirement} together determine the
correct downstream dependency.  The generated owner family and the actual
member must not be counted as two independent physical choices.

\begin{theorem}[Actual source member fixes a unique generated owner member]
\label{p5:thm:v387-actual-source-unique-owner-member}
Let an actually occurring nonzero finite fermion-source member be given on the
already generated matter carrier.  Under the source-slot typing of
Theorem~\ref{p5:thm:v186-generated-yukawa-dirac-owner-family}, the block
projection inverse of the bijection
\(\Theta_Y:\mathfrak Y\to\mathscr D_Y\) fixes a unique member of the generated
Yukawa--Dirac/source-owner family, up to the presentation equivalences already
quotiented in the construction.  No second family selector, mass constructor,
or scalar-vacuum/VEV datum is required for this finite source-member
identification.
\end{theorem}
\begin{proof}
Theorem~\ref{p5:thm:v186-generated-yukawa-dirac-owner-family} proves that
\(\Theta_Y\) is bijective because the four generated support projectors recover
the four source blocks separately.  Corollary~\ref{p5:cor:v385-g15-structural-retirement}
retypes the remaining nonzero member as the G0 occurrence boundary rather than
as a missing structural selector.  Supplying that occurrence therefore selects
one point of an already generated family; applying the proved inverse gives the
unique corresponding owner member.  A scalar stationary branch is a distinct
claim-scoped problem and does not enter this inverse map.
\end{proof}

\begin{corollary}[Downstream mass, mixing and QCD source ownership]
\label{p5:cor:v387-downstream-source-ownership}
Conditional on the actual G0 source member, finite spectra, relative left
frames and weak-CP invariants are deterministic readouts of that member.  A
renormalized running quark mass/source operator at scale \(\mu\) is therefore a
transported readout of the same microscopic source owner once the applicable
RG/threshold certificate is supplied.  A physical pole additionally consumes
the separate continuation, sheet/domain, isolation and residue certificates;
those pole conditions are not prerequisites for the existence of the running
mass/source operator itself.
\end{corollary}

\begin{remark}[No VEV re-entry]
\label{p5:rem:v387-no-vev-reentry}
The native source-member chain above is the occurring finite solder/source
chain.  Paper~V-C's scalar stationary/effective-action branch remains
claim-scoped.  Consequently a Higgs-vacuum expectation value is not reinserted
as an upstream mass selector by this downstream handoff.
\end{remark}
% END TRANSITIVE-FERMION-SOURCE-HANDOFF v3.87

% BEGIN SAFEGUARD-SOURCE-MASS-CHAIN v3.88
\subsection{source-port/value typing and algebraic gap candidate}
\label{p5:sec:v389-source-port-value-factor}

\begin{definition}[Source port and source value]
\label{p5:def:v389-srcport-srcval}
A $\mathsf{SrcPort}$ is a generated normalized odd direction/intertwiner
available to an already occurring source.  A $\mathsf{SrcVal}$ is an actual
odd zero-order operator member of the physical finite action.  The implication
$\mathsf{SrcPort}\Rightarrow\mathsf{SrcVal}$ is forbidden unless the
Paper~II-A source-incidence certificate is discharged.
\end{definition}

\begin{theorem}[Factor-local form of every coordinate-linear family source]
\label{p5:thm:v389-factorized-family-source}
On a fixed component sector, the generated multiplicity-one chiral slot and
relation-free family factor imply that every generated coordinate-linear
family-dependent odd source has the form
\[
 \boxed{T_{\sigma q}=\Sigma_\sigma\otimes P_q\otimes R_{\sigma q}},
 \qquad R_{\sigma q}\in\End(G).
\]
The family factor modulates the chiral edge multiplicatively; it is not an
independent additive even source.
\end{theorem}
\begin{proof}
Apply Paper~IV's relation-free tensor multiplication corollary to the odd
multiplicity-one chiral intertwiner and the even family endomorphism.  The
component projector is fixed by the already generated sector support.
\end{proof}

\begin{proposition}[Primitive rooted owner member and algebraic full rank]
\label{p5:prop:v389-primitive-rooted-owner-member}
Let $h_*=-e_1$ be the rooted axial direction and let $H_*$ be its normalized
coordinate representative.  The explicit family member
\[
 Y_*=(I_G,I_G,I_G,I_G)\in\mathfrak Y
\]
has the rooted source word
\[
 T_*:=c(h_*)\otimes I_Q\otimes I_G,
\]
with dagger completion in the full Dirac owner.  On its directed chiral slot,
\[
 \boxed{T_*^\dagger T_*=I},
\]
so for
$D_*^{\rm seed}=\left(\begin{smallmatrix}0&T_*\\T_*^\dagger&0\end{smallmatrix}\right)$,
\[
 \boxed{(D_*^{\rm seed})^2=I},\qquad \boxed{\ker D_*^{\rm seed}=0}.
\]
This is an algebraic full-rank member of the generated owner family; the
proposition does not assert that this member occurs physically.
\end{proposition}

\begin{corollary}[Family-factor gap criterion]
\label{p5:cor:v389-family-factor-gap-criterion}
For $T_s=c(h_*)\otimes P_s\otimes R_s$,
\[
 T_s^\dagger T_s=I\otimes P_s\otimes R_s^\dagger R_s,
\]
so on the declared sector
\[
 \boxed{\ker T_s=0\iff\ker R_s=0.}
\]
Consequently primitive chiral gap geometry and noncentral family-rank
selection are separate problems.
\end{corollary}

\begin{remark}[Generic source-port firewall and rooted exception]
\label{p5:rem:v389-open-srcval-incidence}
The generic implication
$\mathsf{SrcPort}\Rightarrow\mathsf{SrcVal}$ remains forbidden.  Release
v3.92 does not use it: the rooted triality geometry independently generates
one nonzero finite Dirac member, the generated owner block readouts identify
its unique coordinate, and Paper~III-F places that member in the physical
local fermion action.  Thus the rooted primitive source-incidence target is
closed while the generic firewall is retained.
\end{remark}
% END FERMION-SOURCE-PORT-VALUE-FACTOR v3.89


% BEGIN FERMION-RESIDUAL-BOUNDARY-SPLIT v3.90
\subsection{fermion residual-boundary split and same-parent target}
\label{p5:sec:v390-fermion-residual-boundary-split}

\begin{definition}[Resolved primitive incidence and remaining family obligation]
\label{p5:def:v390-two-fermion-obligations}
With Release v3.92 the two formerly separated obligations have status
\[
 \begin{gathered}
 \boxed{\mathsf{FERMION\_SRCVAL\_INCIDENCE}=\mathsf{CLOSED\_GEOMETRIC}},\\[4pt]
 \boxed{\mathsf{FERMION\_FAMILY\_COMPILER}=\mathsf{OPEN\_PHYSICAL}}.
 \end{gathered}
\]
The live problem begins with the same-parent noncentral responses
\[
 I_G\longmapsto R_u,R_d,R_e,R_\nu
\]
in Theorem~\ref{p5:thm:v389-factorized-family-source}, and specifically with
same-action co-ownership of structural magnitudes and physical left frames.
This obligation is distinct from RG transport and pole certification.
\end{definition}

\begin{proposition}[No downstream repair of source origin or family assembly]
\label{p5:prop:v390-no-downstream-source-repair}
Canonical normalization, regular hierarchy dressing, basis/scheme change and
finite regular RG transport cannot replace a geometric source-origin
certificate and cannot select a noncentral $R_s$ from the generated owner
family.  On the rooted branch Release v3.92 supplies the primitive nonzero
member upstream by independent geometry-first generation.  G16 therefore
remains downstream of the still-open same-action family assembly.
\end{proposition}

\begin{remark}[Sharper same-parent family target]
\label{p5:rem:v390-same-parent-family-target}
Paper~III-F supplies the conditional paired-Feshbach form
\[
 R_s=-\Delta_s^{-1}A_{L,s}U_sA_{R,s}^\dagger
\]
whenever the corresponding blocks actually occur in the finite fermion
action.  This formula narrows the family compiler to a same-parent action
selection problem; it is not itself an occurrence theorem.  In particular
this release does not promote any coefficient-free diagnostic word to the
physical family member.
\end{remark}
% END FERMION-RESIDUAL-BOUNDARY-SPLIT v3.90

% BEGIN ANTIPRISM-FRAME-PORT v3.91
\subsection{antiprism frame port and source-value firewall}
\label{p5:sec:v391-antiprism-frame-port}

\begin{proposition}[Imported full-rank noncentral family-frame port]
\label{p5:prop:v391-antiprism-frame-port}
Paper~I-C v1.33 constructs a canonical full-rank odd antiprism primitive
$\mathcal D_A$ from the rooted phase, native six-cycle adjacency and the
$q=5$ grading, with chiral singular values
\[
 \{\sqrt5-1,2,\sqrt5+1\}.
\]
It therefore supplies a generated noncentral mathematical family-frame port
compatible with the factorized source form of
Theorem~\ref{p5:thm:v389-factorized-family-source}.  This proposition imports
that port as available geometry only; it does not identify it with any
physical $R_s$.
\end{proposition}

\begin{remark}[No cross-sort premise minimization; rooted incidence already closed]
\label{p5:rem:v391-no-cross-sort-premise-minimization}
Paper~II-A no-idle closure still preserves outputs only in their predeclared
target sorts; premise minimization alone never promotes a
$\mathsf{SrcPort}$ to $\mathsf{SrcVal}$.  Release v3.92 closes the rooted
primitive incidence by a different route: independent geometry-first member
generation followed by the existing Paper~III-F action-placement constructor.
The live frontier is therefore same-action noncentral family assembly, not
primitive source existence.
\end{remark}

\begin{remark}[Narrowed family obligation]
\label{p5:rem:v391-narrowed-family-obligation}
With the noncentral frame primitive now available mathematically, the family
compiler should no longer be described as the unconstrained selection of four
arbitrary matrices.  The physical debt is same-action co-ownership/assembly
of sector magnitude data and an occurring noncentral frame into the positive
returns $H_s=R_sR_s^\dagger$.
\end{remark}
% END ANTIPRISM-FRAME-PORT v3.91


































% BEGIN CHARGED-ROLE-INCIDENCE-SCOPE-CORRECTION v3.106
\subsubsection{Yukawa and Majorana block structure}
\label{p5:sec:part1-yukawa}

On \(G\cong\mathbb C^3\), the general gauge-compatible Dirac block is
\begin{equation}
\label{p5:eq:part1-Dirac-Yukawa}
D_Y=
H P_{R-}\otimes\bigl(P_C\otimes Y_d+P_\ell\otimes Y_e\bigr)
+\widetilde H P_{R+}\otimes\bigl(P_C\otimes Y_u+P_\ell\otimes Y_\nu\bigr)
+\mathrm{h.c.},
\end{equation}
where \(Y_u,Y_d,Y_e,Y_\nu\in\End(G)\).

\begin{theorem}[Generated Yukawa slots and optional Majorana support]
\label{p5:thm:part1-yukawa-majorana-unique}
On \(\mathsf G_5^{\rm fin}\), for the generated
coordinate doublet, the only
gauge-compatible coordinate-linear Dirac slots are
\[
\bar Q_LHd_R,
\qquad
\bar Q_L\widetilde H u_R,
\qquad
\bar L_LHe_R,
\qquad
\bar L_L\widetilde H\nu_R.
\]
On \(\mathsf R_5^{\rm SSQFT}\) these are precisely the power-counting
renormalizable Yukawa interactions.
If a Majorana operator is supplied downstream, gauge invariance restricts its
support to
\begin{equation}
\label{p5:eq:part1-Majorana-slot}
D_M\ \text{supported on}\ P_\ell P_{R+}\otimes G.
\end{equation}
\end{theorem}

\begin{proof}
A coordinate-linear Dirac block must pair a left weak doublet with a right
singlet through \(H\) or \(\widetilde H\), while preserving the color/lepton
projectors.  The four displayed sectors are precisely the possibilities, and
their hypercharge sums are
\[
-\frac16+\frac12-\frac13=0,
\quad
-\frac16-\frac12+\frac23=0,
\quad
\frac12+\frac12-1=0,
\quad
\frac12-\frac12+0=0.
\]
All crossed color--lepton, left--left, and right--right Dirac blocks fail the
projector or generated coordinate-intertwiner test.  A Majorana block must be
a self-conjugate gauge singlet; the only neutral right support is the line of
Corollary~\ref{p5:lem:geometric-majorana-line}.  Its coefficients are
arbitrary family endomorphisms, symmetric in the Majorana pairing.  The
existence or nonzero value of this optional operator is not derived.
\end{proof}

\begin{corollary}[Co-realized Yukawa slots and downstream flavor family]
\label{p5:cor:yukawa-matrices-interface-data}
The four generated coordinate-linear Dirac support slots co-realize in the
finite matter carrier; Paper~V-A does not select one instead of the others.  If
a Majorana operator occurs, its support is uniquely restricted to the neutral
right-handed sector.  The numerical interaction data remain the competing
downstream family
\[
\mathfrak A_{\rm flavor}
=
\{(Y_u,Y_d,Y_e,Y_\nu,M_R):
Y_f\in\End(G),\ M_R=M_R^{\mathsf T}\},
\]
with the Majorana coordinate omitted when no such operator occurs.  Paper~V-A
does not determine the eigenvalues, phases, diagonalizers, or the occurrence
of a nonzero Majorana matrix.
\end{corollary}

\begin{corollary}[Family-valued local scalar action invariants]
\label{p5:cor:v175-yukawa-scalar-action-family}
Use the local odd-scalar action branch of
Paper~III-F, Theorem~\ref{p3f:thm:v175-native-local-scalar-action}, and the
Dirac-neutrino branch of \eqref{p5:eq:part1-Dirac-Yukawa} with no additional
constant odd/Majorana offset.  Define the weak-basis invariants
\[
\begin{aligned}
 a_Y:={}&\operatorname{tr}(Y_\nu^\dagger Y_\nu)
       +\operatorname{tr}(Y_e^\dagger Y_e)
       +3\operatorname{tr}(Y_u^\dagger Y_u)
       +3\operatorname{tr}(Y_d^\dagger Y_d),\\
 b_Y:={}&\operatorname{tr}\!\left((Y_\nu^\dagger Y_\nu)^2\right)
       +\operatorname{tr}\!\left((Y_e^\dagger Y_e)^2\right)
       +3\operatorname{tr}\!\left((Y_u^\dagger Y_u)^2\right)
       +3\operatorname{tr}\!\left((Y_d^\dagger Y_d)^2\right).
\end{aligned}
\]
Then the complete generated Yukawa block satisfies
\[
 \operatorname{Tr}_F\Phi_Y(H)^2=2a_YH^\dagger H,
 \qquad
 \operatorname{Tr}_F(D_\mu\Phi_YD^\mu\Phi_Y)
 =2a_Y(D_\mu H)^\dagger D^\mu H,
\]
and
\[
 \operatorname{Tr}_F\Phi_Y(H)^4=2b_Y(H^\dagger H)^2.
\]
These identities are retained over the entire flavor family and select no
Yukawa member.  On the subfamily \(a_Y>0\), the canonical field
\[
 H_c=\sqrt{\frac{a_Y}{8\pi^2}}\,H
\]
has the generated local Lorentzian scalar action
\[
\begin{aligned}
 S_{L,U_0}^{H}=\int_{U_0}d\mu_L\bigg[&
 -(D_\mu H_c)^\dagger D^\mu H_c
 -\frac{R_L}{6}H_c^\dagger H_c
 +2\Lambda_g^2H_c^\dagger H_c\\
 &-\lambda_Y(H_c^\dagger H_c)^2\bigg],
 \qquad
 \lambda_Y=8\pi^2\frac{b_Y}{a_Y^2}>0.
\end{aligned}
\]
The zero-Yukawa fibre \(a_Y=0\) has no kinetic normalization by this word.
The displayed scalar-action family does not own native fermion-mass
selection.  The generated support/grading/covariance/adjoint construction fixes
the admissible Yukawa--Dirac/source owner family, while the actual nonzero
sector member is supplied only by the typed G0 fermion-source occurrence.
A scalar stationary point or nonzero vacuum, when separately claimed, is a
claim-scoped scalar-dynamics statement and cannot feed back as a second
fermion-mass selector.
\end{corollary}

% BEGIN HM v1.79 YH01
\begin{theorem}[Canonical Higgs direction and common-action mass moments]
\label{p5:thm:v179-canonical-higgs-direction}
On the \(a_Y>0\), Dirac-neutrino and no-constant-offset branch of
Corollary~\ref{p5:cor:v175-yukawa-scalar-action-family}, put
\[
 Y_{s,c}=\sqrt{\frac{8\pi^2}{a_Y}}\,Y_s,\qquad
 Y_{\rm tot,c}=(I_3\otimes Y_{u,c})\oplus(I_3\otimes Y_{d,c})
                    \oplus Y_{e,c}\oplus Y_{\nu,c}.
\]
At an existing nonzero neutral background, the local radial chart
\(H_c=(0,(v+h)/\sqrt2)\) gives the actual odd-field derivative
\begin{equation}
\label{p5:eq:v179-higgs-odd-direction}
 X_H=\partial_h\Phi_Y
 =\frac1{\sqrt2}
 \begin{pmatrix}0&Y_{\rm tot,c}\\Y_{\rm tot,c}^\dagger&0\end{pmatrix}.
\end{equation}
The chiral block \(Y_{\rm tot,c}\) has dimension \(24\); the trace
\(\operatorname{Tr}_F\) is the physical finite-carrier trace already used
in the scalar action.  Then
\begin{equation}
\label{p5:eq:v179-higgs-moments}
 \operatorname{Tr}_F X_H^2=8\pi^2,\qquad
 \lambda_Y=\frac{\operatorname{Tr}_F X_H^4}{4\pi^2}
           =8\pi^2\frac{b_Y}{a_Y^2}.
\end{equation}
Thus the kinetic coefficient \(2C\operatorname{Tr}_F X_H^2\) is one.
At a nonzero stationary point of this local quartic block,
\begin{equation}
\label{p5:eq:v179-local-higgs-mass}
 \mu_R^2=\lambda_Yv^2,\qquad
 m_{h,\rm loc}^2=2\lambda_Yv^2
 =16\pi^2v^2\frac{b_Y}{a_Y^2}=4\Lambda_g^2-R/3 .
\end{equation}
This stationary-block identity neither proves a vacuum of the complete
effective action nor identifies \(\Lambda_g\) with a hierarchy anchor.
\end{theorem}
\begin{proof}
The canonical field change rescales all four Yukawa blocks by the same
factor, and
\(\operatorname{tr}(Y_{\rm tot,c}^\dagger Y_{\rm tot,c})=8\pi^2\).
Squaring the odd block gives two diagonal return blocks with the factor
\(1/2\), proving its second moment.  Its fourth moment is
\(\tfrac12(8\pi^2/a_Y)^2 b_Y\), proving the quartic relation.
The radial potential is
\(-\mu_R^2(v+h)^2/2+\lambda_Y(v+h)^4/4\).
Its first and second derivatives give
\eqref{p5:eq:v179-local-higgs-mass}; this is the compression of the
same local Hessian in
Theorem~\ref{p3f:thm:v179-normalized-scalar-hessian}.
\end{proof}

% Proposed placement: V-A, immediately after the canonical radial-Higgs
% action/mass discussion; the full-heat corollary may instead be placed III-F.
% This transports I-A providers into an existing action, not a new field.

\begin{proposition}[Sextic-profile evaluation of the existing radial action]
\label{p5:prop:v215-sextic-radial-evaluation}
Work on the canonical radial branch of
Theorem~\ref{p5:thm:v179-canonical-higgs-direction} and the local-action
domain of Theorem~\ref{p3f:thm:v175-native-local-scalar-action}. Use a
fixed round spatial slice \(S^3(a)\), a homogeneous existing radial
background \(v\), and a parallel radial internal line, so the spatial
covariant derivative reduces to the ordinary derivative. To evaluate a
finite position-band field at the following profile requires its actual
band \(m_x\geq6\), and the actual position/reciprocal comparison map if
the profile is transported from the rooted momentum cell. A degree-two
comparison label alone does not supply either condition.

Let \(\psi_6\) be the volume-normalized harmonic of
Proposition~\ref{p1:prop:v215-sextic-moments}. The already defined radial
potential
\[
 V_{\rm rad}(v+h)=-\frac{\mu_R^2}{2}(v+h)^2
                         +\frac{\lambda_Y}{4}(v+h)^4,
 \qquad\lambda_Y>0,
\]
evaluated at the existing field configuration \(h=q\psi_6\), gives
the exact fixed-slice restricted energy
\[
 E(q)-E(0)=
 \frac12\left(\frac{48}{a^2}-\mu_R^2+3\lambda_Yv^2\right)q^2
 +\frac{\lambda_Y}{4}\frac{427}{143V_a}q^4,
 \qquad V_a=2\pi^2a^3.
\]
In particular the cubic self-coefficient vanishes. At a nonzero
stationary point of this same local quartic block,
\[
 \mu_R^2=\lambda_Yv^2,\qquad
 \omega_6^2=\frac{48}{a^2}+2\lambda_Yv^2>0.
\]
The last identity has the local stationary-block domain already stated
in Theorem~\ref{p5:thm:v179-canonical-higgs-direction}; it does not select
a stationary point of the complete effective action.
\end{proposition}
\begin{proof}
Expand the native radial potential and integrate over the slice. The
linear term vanishes since \(\int\psi_6=0\); the cubic term vanishes
since \(\int\psi_6^3=0\). The unit quadratic norm, gradient norm
\(48/a^2\), and quartic norm \(427/(143V_a)\) are precisely the
I-A provider. Substitution of the existing stationary identity proves
the positive quadratic coefficient.
\end{proof}

\begin{remark}[Configuration, degeneracy, and nonlinear scope]
This calculation adds no interaction to the action and no field to its
carrier. It evaluates one configuration only when that configuration
belongs to the actual field domain; it makes no occupation claim. The
round linearized kinetic operator has the same eigenvalue on all
\(49\) degree-six harmonics, so the full-cell invariant does not select
a unique unstable or populated mode. At \(v=0\), a degree-six
instability would require \(\mu_R^2>48/a^2\), when the constant mode
is already unstable. A single-profile nonlinear ansatz is generally
not closed: \(3\lambda_Yv h^2\) sources the constant mode at order
\(q^2\), and higher powers have other harmonic components. Nor does
one selected profile supply homogeneous, isotropic stress. The
second-moment isotropy calculation is the full-multiplet identity of
Proposition~\ref{p1:prop:v215-degree-six-kernels}; an actual state and
its fourth moments remain needed for nonlinear stress. No independent
global \(F_4\) field symmetry is asserted.
\end{remark}

\begin{corollary}[Odd sextic jets of an already generated heat family]
\label{p3f:cor:v215-sextic-heat-parity}
Suppose the supplied round geometric, gauge, and homogeneous scalar
background, its actual finite window, and any trace localizer are
preserved by the orientation-preserving isometry \(O\) of
Proposition~\ref{p1:prop:v215-sextic-moments}. Require its bundle lift
also to fix the already constituted radial matrix \(X_H\). On the
existing family, when the same position-band and comparison-map
conditions admit the profile,
\[
 D(q)=D_A+\gamma_g\otimes(\Phi_0+qX_H\psi_6),
\]
every defined spectral or heat readout satisfies
\[
 \operatorname{Tr}f(D(q)^2/\Lambda^2)
          =\operatorname{Tr}f(D(-q)^2/\Lambda^2).
\]
The corresponding localized equality holds with a preserved localizer.
Thus every existing odd amplitude derivative at zero vanishes for the
exact readout, including the linear forcing and cubic self-vertex.
This statement does not impose positivity on its even derivatives or
remove mixed vertices with other profiles.
\end{corollary}
\begin{proof}
The spin lift exists for the orientation-preserving round isometry.
Under the stated background and internal-matrix hypotheses it fixes
\(D_A+\gamma_g\otimes\Phi_0\) and reverses the profile term.
Apply Theorem~\ref{p1:thm:v215-odd-profile-spectral-parity} on the
actual preserved window or trace-class domain.
\end{proof}

\begin{remark}[The already exact comparison has no sextic alias term]
The Hurwitz errors in
Proposition~\ref{p1:prop:v215-eight-alias} evaluate the specified
\(2T\) quadrature. On the existing degree-two quartic comparison,
the actual \(2I\) eleven-design integrates degrees six and eight
exactly, so both displayed alias errors vanish. They do not produce a
residual sextic correction in that comparison. Conversely, this
eleven-design is not sufficient by degree alone for the degree-twelve
norm of an actual \(H_6\) field; such a field calculation uses the
exact I-A Haar moments or a separately certified quadrature.

On the distinct native unit-preserving scalar-process branch,
\(L_0\mathbf1=0\) excludes a bare modification
\(L_\epsilon=L_0+\epsilon M_{H_6}\), unless \(\epsilon=0\): indeed
\(L_\epsilon\mathbf1=\epsilon H_6\). The existing joint-field
evaluation above does not make this scalar-generator modification.
\end{remark}




\begin{corollary}[Internal Yukawa participation invariant and sharp support bound]
\label{p5:cor:v181-yukawa-participation-invariant}
On the canonical Higgs branch of
Theorem~\ref{p5:thm:v179-canonical-higgs-direction}, define
\[
 \boxed{N_{\rm eff}:=\frac{a_Y^2}{b_Y}.}
\]
Then
\[
 \boxed{\lambda_Y=\frac{8\pi^2}{N_{\rm eff}},\qquad 1\le N_{\rm eff}\le24.}
\]
The upper endpoint is the total color-weighted chiral singular-value support
of the canonical finite Higgs/Yukawa block.  Spatial harmonic multiplicities
do not enter this invariant.
\end{corollary}
\begin{proof}
Expand the weights $(3,3,1,1)$ as repeated nonnegative squared singular
values $x_j$ on the $24$-entry finite chiral support.  Then
$a_Y=\sum_jx_j$ and $b_Y=\sum_jx_j^2$.  Positivity gives
$b_Y\le a_Y^2$, while Cauchy--Schwarz gives $a_Y^2\le24b_Y$.
The quartic identity is Eq.~\eqref{p5:eq:v179-higgs-moments}.
\end{proof}

\begin{corollary}[Normalized fourth moment of the homogeneous Yukawa ray]
\label{p5:cor:v2211-yukawa-ray-moment}
On the no-constant-offset Dirac branch of
Corollary~\ref{p5:cor:v175-yukawa-scalar-action-family}, let
\(\Phi_Y(H)\ne0\) and set
\(T:=\Phi_Y(H)/\sqrt{\operatorname{Tr}_F\Phi_Y(H)^2}\).
Then
\[
 \boxed{\operatorname{Tr}_FT^2=1,\qquad
 c_4:=\operatorname{Tr}_FT^4
 =\frac{b_Y}{2a_Y^2}=\frac1{2N_{\rm eff}},\qquad
 \frac1{48}\le c_4\le\frac12.}
\]
The primitive identity family \(Y_u=Y_d=Y_e=Y_\nu=I_G\) has
\(a_Y=b_Y=24\), \(N_{\rm eff}=24\) and \(c_4=1/48\).
These are moments of this homogeneous constituent; a constant offset,
additional source contact or complete effective-action term is not included.
\end{corollary}
\begin{proof}
The same-carrier identities give
\(\operatorname{Tr}_F\Phi_Y^2=2a_YH^\dagger H\) and
\(\operatorname{Tr}_F\Phi_Y^4=2b_Y(H^\dagger H)^2\).
Their quotient is \(b_Y/(2a_Y^2)\); apply the preceding sharp support
bound.  For four identity family blocks, color weights give
\(3\cdot3+3\cdot3+3+3=24\) in both moments.
\end{proof}

% Framework v2.217 helper integration.
\begin{corollary}[Color repetition and participation on a quark-only ray]
\label{p5:cor:v2217-color-participation}
For a nonzero canonical homogeneous Yukawa member with $Y_e=Y_\nu=0$,
write its nonzero squared quark singular values as $q_i\ge0$ before color
repetition. Then
\[
 N_{\rm eff}=3\frac{(\sum_iq_i)^2}{\sum_iq_i^2}\ge3,
\]
with equality precisely when one $q_i$ is nonzero.
\end{corollary}
\begin{proof}
The same color-weighted trace gives $a_Y=3\sum q_i$ and
$b_Y=3\sum q_i^2$. Since $\sum_{i\ne j}q_iq_j\ge0$, substitution into
$a_Y^2/b_Y$ proves both assertions.
\end{proof}
The quantitative quartic remainder is
Paper~III-F, Proposition~\ref{p3f:prop:v2217-quartic-omitted-term}.
For clarity, the separate zero-offset local comparison with
$\kappa_J=12U$, $L^2=12\kappa_J$, $\Lambda_g^2=96\pi^2U$ and
$\lambda_Y=8\pi^2/N_{\rm eff}$ has normalized radial potential
\[
 \frac{V_E}{12\kappa_J\Lambda_g^2}
   =(3/N_{\rm eff}-1)s^2-s,\qquad s=\sinh^2(\rho/L).
\]
Its nonzero finite radial minimum exists only for $N_{\rm eff}<3$,
at $s=N_{\rm eff}/[2(3-N_{\rm eff})]$, as follows by differentiating
this quadratic in $s\ge0$. Thus a quark-only ray cannot supply that
comparison minimum. The comparison deletes a constant term by hypothesis;
it is not the bare-vacuum polynomial retained in v2.216, a full vacuum
theorem, or evidence of a runaway in the exact completed action.


\begin{proposition}[Claim-scoped scalar/gauge shape-matching certificate]
\label{p5:prop:v2211-scalar-gauge-shape-certificate}
Suppose one actual common canonical local tree branch satisfies the
single-doublet gauge dictionary of Paper~V-C and the stationary scalar
Hessian identity of Theorem~\ref{p5:thm:v179-canonical-higgs-direction}.
Let \(v\ne0\) denote the radial coordinate of that existing branch,
not a fundamental mass-generation premise.  Then
\[
 \boxed{\frac{m_{h,\rm loc}^2}{m_{W,\rm loc}^2}
 =\frac{8\lambda_Y}{g_2^2}
 =\frac{64\pi^2}{g_2^2N_{\rm eff}}.}
\]
At the source-native common-heat coupling \(g_2^2=2\pi^2\), a further
identification of this ratio with the finite shell value \(17/7\) requires
\[
 \boxed{N_{\rm eff}=\frac{224}{17},\qquad
 \frac{b_Y}{a_Y^2}=\frac{17}{224}.}
\]
This is allowed by \(1\le N_{\rm eff}\le24\), but is not selected by
that inequality or primitive source occurrence.  The primitive identity
member instead gives ratio \(4/3\).  The actual source moments, common
scale and physical action transport must be evaluated before using this
certificate as a physical matching.  Offsets, nonminimal metric feedback
and quantum terms retain their complete-Hessian calculation.  This
certificate neither derives a stationary vacuum nor identifies finite
coefficients with interacting poles.
\end{proposition}
\begin{proof}
On the stated common branch,
\(m_{W,\rm loc}^2=g_2^2v^2/4\) and
\(m_{h,\rm loc}^2=2\lambda_Yv^2\).  Divide, cancel \(v^2\), and use
\(\lambda_Y=8\pi^2/N_{\rm eff}\).  Substitution of the separately
stated common-heat coupling and \(17/7\) gives the required moment.
The cancellation shows why no additional scalar magnitude is selected.
\end{proof}

\begin{definition}[External SM right/left Yukawa Gram coordinates]
\label{p5:def:v355-yukawa-left-right-grams}
For each Dirac sector $s\in\{u,d,e,\nu\}$ set
\[
\begin{aligned}
 H_s^R&:=Y_s^\dagger Y_s,\\
 H_s^L&:=Y_sY_s^\dagger.
\end{aligned}
\]
These Gram operators are retained as coordinates of the external
Standard--Model/\allowbreak Yukawa comparison dictionary.  They may be used for scalar
source moments on that dictionary, but they do not own the framework-native
terminal masses or flavor mixing.  Native charged flavor is carried by the
terminal-role incidence maps $W_s$ and their weak-current pullbacks.  The two
Gram operators have the same nonzero singular spectrum but remain differently
typed comparison objects.
\end{definition}

\begin{corollary}[Projective Yukawa participation factorization]
\label{p5:cor:v355-yukawa-projective-participation}
Let $(w_u,w_d,w_e,w_\nu)=(3,3,1,1)$ and define
\[
 a_s:=\operatorname{tr}H_s^R,\qquad
 b_s:=\operatorname{tr}(H_s^R)^2.
\]
For $a_s>0$ put $p_s:=a_s^2/b_s$, and put
\[
 \rho_s:=\frac{w_sa_s}{a_Y},\qquad \sum_s\rho_s=1.
\]
A sector with $a_s=0$ has $\rho_s=0$ and contributes zero below.  Then
\[
 \boxed{
 \frac1{N_{\rm eff}}
 =\sum_s\frac{\rho_s^2}{w_sp_s}
 =\frac{\rho_u^2}{3p_u}+\frac{\rho_d^2}{3p_d}
  +\frac{\rho_e^2}{p_e}+\frac{\rho_\nu^2}{p_\nu}.}
\]
Thus $N_{\rm eff}$ records only the distribution of singular-value weight
among and within sectors in the external Yukawa dictionary.  It is blind to
the generated weak-current incidences $C_{ud}$ and $C_{e\nu}$ that carry the
framework-native flavor data.

The sharp bound $1\le N_{\rm eff}\le24$ of
Corollary~\ref{p5:cor:v181-yukawa-participation-invariant} has the following
endpoint interpretation.  The upper endpoint occurs exactly when all
nonzero entries on the full $24$-entry color-expanded singular-square support
are equal and the full support is occupied.  The lower endpoint can occur
only when a single color-expanded entry is nonzero; in particular a single
nonzero quark singular value alone gives $N_{\rm eff}=3$ because of the
threefold color repetition.
\end{corollary}
\begin{proof}
Since $a_Y=\sum_sw_sa_s$ and $b_Y=\sum_sw_sb_s$,
\[
 \frac1{N_{\rm eff}}=\frac{b_Y}{a_Y^2}
 =\sum_s\frac{w_sa_s^2}{p_sa_Y^2}
 =\sum_s\frac{\rho_s^2}{w_sp_s}.
\]
The blindness statement follows because every quantity displayed is a
function only of the external-dictionary singular values and contains no
terminal-role current-incidence input.  The endpoint statements are the
equality cases of the positivity and Cauchy--Schwarz proof in
Corollary~\ref{p5:cor:v181-yukawa-participation-invariant}, with the quark
entries repeated three times by color.
\end{proof}

\begin{corollary}[Common positive Yukawa-scale cancellation]
\label{p5:cor:v355-common-yukawa-scale-cancellation}
For a common positive scalar $r>0$, replace every Dirac block by
$Y_s' = rY_s$.  Then
\[
 a_Y'=r^2a_Y,\qquad b_Y'=r^4b_Y,
\]
and therefore
\[
 \boxed{
 N_{\rm eff}'=N_{\rm eff},\qquad
 \lambda_Y'=\lambda_Y,\qquad
 Y_{s,c}'=Y_{s,c}.}
\]
Hence a factor such as a common $\sqrt{24/25}$ multiplying all four raw
Yukawa blocks cannot by itself survive as a canonically normalized mass or
quartic multiplier.  A physical effect requires a generated non-common
sector deformation, altered normalized source image, altered vacuum/Hessian,
or another asymmetric same-action role.
\end{corollary}
\begin{proof}
Substitution into the definitions of $a_Y,b_Y,N_{\rm eff},\lambda_Y$ and
$Y_{s,c}=\sqrt{8\pi^2/a_Y}\,Y_s$ gives the identities immediately.
\end{proof}

\begin{corollary}[External SM-coordinate kinetic-normalized fermion comparison]
\label{p5:cor:v179-canonical-fermion-comparison}
On the local clocked, uninserted classical branch of
Paper~III-F, Theorem~\ref{p3f:thm:v177-local-clocked-fermion-jet}, the
same-action fermionic jet has
\[
 (Z_{L,s},D_s,Z_{R,s})=c_F(I_L,D_{Y,s},I_R),
 \qquad c_F>0,
\]
so its canonical chiral operator is
\[
 \boxed{\widehat D_s=Z_{L,s}^{-1/2}D_sZ_{R,s}^{-1/2}=D_{Y,s}.}
\]
If one additionally adopts the conventional Standard--Model comparison
coordinate
\[
 M_s^{\rm SM,cmp}:=vY_{s,c}/\sqrt2,
\]
then at one common comparison slice, with
\((w_u,w_d,w_e,w_\nu)=(3,3,1,1)\),
\[
 \sum_s w_s\operatorname{tr}(M_s^\dagger M_s)=4\pi^2v^2,\qquad
 m_{h,\rm loc}^2
 =4\frac{\sum_s w_s\operatorname{tr}((M_s^\dagger M_s)^2)}
          {\sum_s w_s\operatorname{tr}(M_s^\dagger M_s)} .
\]
\end{corollary}
\begin{proof}
The local clocked uninserted compiler of
Theorem~\ref{p3f:thm:v177-local-clocked-fermion-jet} cancels the common
positive fermionic coefficient between the left kinetic, chiral and right
kinetic slots.  The common second- and fourth-moment formulas then follow by substitution
in \eqref{p5:eq:v179-higgs-moments} and
\eqref{p5:eq:v179-local-higgs-mass}.  This proves only an external SM
coordinate identity.  It does not identify $M_s^{\rm SM,cmp}$ with the
framework-native terminal mass coefficient or permit VEV/Yukawa back-generation.
\end{proof}

The \(Y_s\) remain the family of
Corollary~\ref{p5:cor:yukawa-matrices-interface-data}.  These derivatives
co-realize with that family and use no new G0 occurrence, but do not select
its eigenvalues.  The local clock-hermitizing kinetic pairing is supplied on
the named Paper~III-F branch; global determinant/Pfaffian, Majorana, PCT and
quantum-measure questions remain separate.  Masses reported at different
renormalization scales and physical pole values cannot be combined as the
single matching slice in the corollary.  Coupled scalar--metric and other
retained responses are handled by
Proposition~\ref{p3f:prop:v179-radial-einstein-response} and the complete
action reduction.
% END HM v1.79 YH01



\begin{remark}[Flavor presentation and occurrence firewall]
\label{p5:rem:flavor-presentation-occurrence-firewall}
The set \(\mathfrak A_{\rm flavor}\) is a raw coefficient atlas on the
generated support slots, not a count of physical flavor parameters.
Kinetic-term-preserving changes of family bases act on these coefficients by
the induced weak-basis action, and presentation-equivalent representatives
are to be identified only at the downstream physical-quotient stage.
Paper~V-B owns that quotient explicitly.  This flavor-level presentation
statement does not remove the anomalous strong-CP action/measure coordinate,
which has a distinct downstream owner.

Likewise, co-realization of an allowed support slot does not construct a
nonzero numerical coefficient on that support.  In particular, the optional
Majorana support separates a branch on which a nonzero Majorana interaction
actually occurs from a branch on which it does not; that branch distinction is
an occurrence question, not a consequence of the family Weyl calculus.
\end{remark}


% BEGIN DOWNSTREAM FLAVOR PRESENTATION IMPORT v1.1.
\begin{corollary}[Protected flavor mass and polynomial calculus transport]
\label{p5:cor:fullcarrier-flavor-mass-transport}
For any already licensed Hermitian three-flavor source/mass member
$M_f$ in the physical flavor-presentation atlas, and the mathematical
frame $Q$ and local lift $\mathscr K$ of
Theorem~\ref{p1b:thm:fullcarrier-protected-local-flavor},
\[
 \boxed{\mathscr K(M_f)^mQ=QM_f^m,\qquad
 Q^\dagger p(\mathscr K(M_f))Q=p(M_f)}
\]
for every $m\ge0$ and polynomial $p$.  Thus a spectral projector of
an \emph{already selected} $M_f$ is transported exactly when restricted
to $\operatorname{Ran}Q$.  This does not select $M_f$, its physical
eigenvalues or its occurrence, and does not imply the same spectral
projector on all $20$ vertex states.
\end{corollary}
\begin{proof}
Iterate $\mathscr K(M_f)Q=QM_f$ and use $Q^\dagger Q=I_3$;
polynomial spectral projectors of a finite Hermitian matrix follow by
interpolation on its distinct eigenvalues.
\end{proof}
% END DOWNSTREAM FLAVOR PRESENTATION IMPORT v1.1.


\section{Gauge, fermion and spectral-action exports}
\label{p5a:sec:organized-exports}
\subsection{Bridge criterion and universal one-loop form}
\label{p5:sec:part4-bridge-criterion}

\begin{definition}[Graded resolution-boundary bridge]
\label{p5:post:part4-bridge-criterion}
\label{p5:def:part4-bridge-criterion}
A one-loop coefficient is \emph{boundary-derived} on the present branch when
(i) the physical carrier and representation indices are fixed before the
calculation, (ii) the quadratic response blocks and gauge quotient descend
from the occurring action, (iii) their determinant powers are fixed by the
finite bosonic/Berezin integration owners, and (iv) the logarithmic
four-dimensional heat coefficient is evaluated on the SSQFT local-QFT domain.
The renormalization scale is the position of a resolution/spectral boundary;
it is not physical time and introduces no new particle or coupling member.
\end{definition}

\subsubsection{Graded heat-boundary one-loop coefficient theorem}
\begin{theorem}[Graded heat-boundary one-loop coefficient theorem]
\label{p5:thm:part4-standard-one-loop-bridge}
Let a fixed threshold interval of the reciprocal-position QFT interface
contain compact gauge factors \(G_i\), Weyl fermions \(f\) in
representations \(R_f\), and complex scalars \(s\) in representations
\(R_s\).  Use the Paper--III-F background-gauge response complex and the
Paper--IV determinant weights on the \(\mathsf R_5^{\rm SSQFT}\) local
four-dimensional heat-kernel domain.  Define
\begin{equation}
\label{p5:eq:part4-beta-sign-convention}
\mu\frac{dg_i}{d\mu}
=\frac{b_i}{16\pi^2}g_i^3+O(g_i^5),
\qquad
\frac{d\alpha_i^{-1}}{d\log\mu}
=-\frac{b_i}{2\pi}+O(\alpha_i).
\end{equation}
Then the logarithmic resolution-boundary response gives
\begin{equation}
\label{p5:eq:part4-standard-bi}
\boxed{
 b_i
 =
 -\frac{11}{3}C_2(G_i)
 +\frac{2}{3}\sum_{\text{Weyl }f}T_i(R_f)
 +\frac{1}{3}\sum_{\text{complex }s}T_i(R_s).
}
\end{equation}
For an Abelian factor the adjoint self-term vanishes because
$C_2(U(1))=0$.  No coefficient in
\eqref{p5:eq:part4-standard-bi} is an independently supplied response count.
\end{theorem}

\begin{proof}
\ProofPlan{Project the finite graded shell determinant onto the pure-gauge
$F^2$ heat coefficient and then convert that boundary coefficient into the
beta-function convention \eqref{p5:eq:part4-beta-sign-convention}.}

For a Laplace-type operator
$\Delta=-(\nabla^2+E)$, the pure internal-gauge part of the local
four-dimensional heat coefficient is
\begin{equation}
\label{p5:eq:v2140-a4-gauge-piece}
 a_4(\Delta)\big|_{F^2}
 =\frac1{(4\pi)^2}\int
 \operatorname{tr}\!\left(\frac12E^2+\frac1{12}\Omega_{\mu\nu}\Omega^{\mu\nu}\right)d\mu_g,
\end{equation}
with the convention already used in Paper~III-F
\cite{Vassilevich2003HeatKernel}.  Put
\[
 X_R:=\int\operatorname{Tr}_R(\mathcal F_{\mu\nu}^{\dagger}
 \mathcal F^{\mu\nu})d\mu_g
 =T(R)\int F^a_{\mu\nu}F^{a\mu\nu}d\mu_g,
\]
and write
$a_4|_{F^2}=(4\pi)^{-2}hX_R$.

For one real minimally coupled scalar, $E=0$ and
$\operatorname{tr}\Omega^2=-X_R$, hence
\[
 h_{s,\mathbb R}=-\frac1{12}.
\]
Paper~IV fixes the real-boson effective-action determinant weight to
$\omega=+1/2$; a complex scalar contains two real blocks and therefore has
$\omega_s=+1$.

For a Dirac carrier,
$E_F=\tfrac12\gamma^{\mu\nu}\mathcal F_{\mu\nu}$ and the Clifford traces used
already in Paper~III-F give, in the positive $X_R$ convention,
\[
 \operatorname{tr}_{\mathbb S}(E_F^2)=2X_R,
 \qquad
 \operatorname{tr}_{\mathbb S}(\Omega_F^2)=-4X_R.
\]
Thus
\[
 h_D=1-\frac13=\frac23.
\]
The paired Grassmann determinant has $\omega_D=-1/2$ after squaring the
first-order Dirac operator.  The parity-even $F^2$ response of one Weyl
carrier is half the Dirac response, equivalently $\omega_W=-1/4$ against the
same rank-four heat trace.

For the Yang--Mills one-form operator of Paper~III-F,
$E_\mu{}^\nu=2\operatorname{ad}(\mathcal F_\mu{}^\nu)$ in the pure-gauge
projection.  Therefore
\[
 \operatorname{tr}E^2=4X_{\rm ad},\qquad
 \operatorname{tr}\Omega^2=-4X_{\rm ad},\qquad
 h_1=\frac53.
\]
Its real bosonic determinant has $\omega_1=+1/2$.  The quotient Jacobian is a
paired Grassmann scalar in the adjoint representation, so
$h_0=-1/12$ and $\omega_0=-1$.  Hence the complete gauge-orbit response has
\begin{equation}
\label{p5:eq:v2140-gauge-11-12}
 \boxed{
 \omega_1h_1+\omega_0h_0
 =\frac12\frac53+(-1)\!\left(-\frac1{12}\right)
 =\frac{11}{12}.}
\end{equation}

Finally, the Paper--I-B heat shell gives
$\Tr\log\Delta=-\int ds\,s^{-1}\Tr e^{-s\Delta}$, so the $a_4$ term between
$\Lambda$ and $\mu$ is
$2\omega h(4\pi)^{-2}\log(\mu/\Lambda)X_R$.  Comparing with
$\frac1{4g^2(\mu)}\int F^2$ gives
\begin{equation}
\label{p5:eq:v2140-master-b-weight}
 \boxed{b=-4\omega h.}
\end{equation}
Therefore a complex scalar gives $+\frac13T(R)$, a Weyl fermion gives
$+\frac23T(R)$, and the one-form plus gauge-orbit quotient gives
$-\frac{11}{3}C_2(G)$.  Summing the independently fixed carrier blocks proves
\eqref{p5:eq:part4-standard-bi}.
\end{proof}

\begin{remark}[Analytic scope and non-circularity]
The local $a_4$ theorem, background-field regularity, and SSQFT
renormalization chart remain analytic-domain hypotheses; the result does not
claim that the bare finite Born--Hopf grammar alone proves four-dimensional
perturbative QFT.  What has been removed is the independent import of the
numbers $-11/3$, $2/3$, and $1/3$.  The carrier, hypercharge normalization,
rank-three family multiplicity, scalar role, action, quotient complex, and
determinant powers are all fixed before this coefficient evaluation.
\end{remark}

\subsection{Full unbroken Standard-Model coefficient vector}
\label{p5:sec:part4-full-sm-vector}

\subsubsection{Vector blocks}
The vector blocks are
\begin{equation}
\label{p5:eq:part4-vector-blocks}
\boxed{b^{\rm vector}=\left(0,-\frac{22}{3},-11\right),}
\end{equation}
since
\[
-\frac{11}{3}C_2(SU(2))=-\frac{22}{3},\qquad -\frac{11}{3}C_2(SU(3))=-11,\qquad C_2(U(1))=0.
\]

\subsubsection{Full coefficient vector}
\begin{proposition}[Evaluation of the Standard--Model one-loop vector]
\label{p5:prop:part4-sm-vector}
Under
\(\mathsf G_5^{\rm fin}
+\mathsf R_5^{\rm SSQFT}\), the full unbroken Standard--Model one-loop
coefficient vector, in the sign convention
\eqref{p5:eq:part4-beta-sign-convention}, is
\begin{equation}
\label{p5:eq:part4-sm-vector}
\boxed{b^{\rm SM}=\left(\frac{41}{10},-\frac{19}{6},-7\right).}
\end{equation}
\end{proposition}

\begin{proof}
Using Corollary~\ref{p5:cor:part4-three-generation-matter}, the Higgs block \eqref{p5:eq:part4-bH}, and the vector blocks \eqref{p5:eq:part4-vector-blocks},
\[
b^{\rm SM}=(4,4,4)+\left(\frac{1}{10},\frac16,0\right)+\left(0,-\frac{22}{3},-11\right).
\]
Therefore
\[
b_1=4+\frac{1}{10}=\frac{41}{10},\qquad b_2=4+\frac16-\frac{22}{3}=-\frac{19}{6},\qquad b_3=4-11=-7.
\]
\end{proof}


\paragraph{Boundary of the one-loop bridge.}
The Paper--V-A one-loop bridge stops with the unbroken Standard--Model coefficient vector.  Broken-phase threshold towers, carrier--host constants, electroweak endpoint packets, and numerical comparison laws are not part of the structural Standard--Model carrier theorem and are deferred to later comparison layers.

\subsection{Common-heat trace carrier and absolute gauge matching}
\label{p5:sec:common-heat-trace-carrier}

The finite Standard--Model carrier and the derived relative-even projector now
suffice to close the trace-carrier question left by the broad gauge kinetic
classification.  This section concerns only the source-native common-heat
branch of Paper~III-F; it does not alter the larger class of admissible
parity-even kinetic functionals.

\begin{theorem}[Least self-adjoint Clifford completion of the visible Weyl carrier]
\label{p5:thm:least-self-adjoint-heat-carrier}
Let
\[
 \mathbb S=S_L\oplus S_R,
 \qquad
 F
 =(W_L\oplus W_{R+}\oplus W_{R-})\otimes(\ell\oplus C)\otimes G,
\]
with the already derived relative grading
$\Gamma_{\rm rel}=\gamma_x\otimes\Gamma_W$ and
$P_\pm=\tfrac12(1\pm\Gamma_{\rm rel})$.
For the generated nondegenerate Clifford principal symbol $c(\xi)$,
\[
 \{c(\xi)\otimes1_F,\Gamma_{\rm rel}\}=0.
\]
Consequently the physical Weyl space $\operatorname{Ran}P_+$ is not closed
under a self-adjoint first-order Dirac datum.  The unique least complex
Clifford-stable carrier containing it is
\[
 \boxed{
 \mathcal H_{\rm heat}
 =\operatorname{Ran}P_+\oplus\operatorname{Ran}P_-
 =\mathbb S\otimes F,
 \qquad
 \dim_{\mathbb C}\mathcal H_{\rm heat}=4\cdot48=192.}
\]
The complementary $P_-$ space is the analytic Clifford partner required by
self-adjointness; it is not a second physical mirror-family occurrence.
\end{theorem}

\begin{proof}
The anticommutation identity follows from
$\{c(\xi),\gamma_x\}=0$ and the trivial action of the external Clifford
symbol on the internal grading.  Hence
$(c(\xi)\otimes1)P_+=P_-(c(\xi)\otimes1)$, so the symbol is not an
endomorphism of $\operatorname{Ran}P_+$.  Adding the complementary grading
sector closes the least faithful complex $\Cl_4$ module.  Paper~V-A's
no-hidden-multiplicity theorem excludes further unmarked internal copies.
Moreover, complex conjugation supplies $\bar F$ only as a
functorially available conjugate carrier; simultaneous finite
$F\oplus\bar F$ occurrence belongs to the separately typed
PCT/real-structure branch and is not inserted here.
\end{proof}


\begin{theorem}[Raw joint/scalar Einstein-coefficient transport]
\label{p5:thm:g24-raw-eh-transport}
On the source-native ordinary-trace common-heat branch, let
\[
 U:=\Lambda_{\rm EH}^2
 =\frac{\Lambda_g^2}{6(4\pi)^2}
 =\frac{\Lambda_g^2}{96\pi^2}
\]
be Paper~III-E's generated scalar-resolution Einstein curvature-scale owner,
and let
\[
 U_J^{\rm raw}:=8\kappa_J
\]
be Paper~III-F's raw represented joint-Dirac Einstein coefficient after the
already-derived positive pure-metric support and Born-diagonal routing.  Then
\[
 \boxed{
 U_J^{\rm raw}=\frac{\Lambda_g^2}{\pi^2}=96\,U.}
\]
Equivalently, the relative coefficient is the generated visible-Weyl heat
multiplicity
\[
 \boxed{
 \frac{U_J^{\rm raw}}{U}
 =\frac{192/12}{1/6}
 =96
 =2\cdot48
 =\dim_{\mathbb C}\operatorname{Ran}P_+.}
\]
No measured Newton datum, normalized finite trace, or additional carrier copy
enters this transport.
\end{theorem}

\begin{proof}
Paper~III-F gives
\[
 E_A|_{R_E}=-\frac14R_E1,
 \qquad
 a_2(D_A^2)|_{R_E}
 =-(4\pi)^{-2}\frac{\operatorname{tr}(1)}{12}
 \int R_E\,d\mu_E.
\]
Theorem~\ref{p5:thm:least-self-adjoint-heat-carrier} fixes
\(\operatorname{tr}(1)=\dim_{\mathbb C}\mathcal H_{\rm heat}=192\) for the
ordinary additive common-heat trace.  The same already-derived pure-metric
protected support/Born routing entering Paper~III-F's \(\kappa_J\) then gives
\[
 \kappa_J
 =\frac18(4\pi)^{-2}\frac{192}{12}\Lambda_g^2
 =\frac{\Lambda_g^2}{8\pi^2},
\]
so
\[
 U_J^{\rm raw}=8\kappa_J=\frac{\Lambda_g^2}{\pi^2}.
\]
Eliminating \(\Lambda_g^2\) with Paper~III-E's independently generated scalar
identity gives \(U_J^{\rm raw}=96U\).
\end{proof}

\begin{remark}[Owner and trace firewall for the factor 96]
\label{p5:rem:g24-raw-eh-owner-firewall}
The theorem evaluates the \emph{raw represented} joint coefficient.  It does
not replace the ordinary additive trace in the gauge sector by a normalized
trace.  Framework v2.88 independently fixes the primitive curvature at the
pre-amplification EH owner \(U=(3/\pi)\mathcal K_p\) with intrinsic momentum curvature \(\mathcal K_p=6/R_{\rm mom}^2\); therefore no physical
\(96\to1\) division or owner-retraction process is required to define that
curvature.  Any separately claimed represented-to-scalar dynamical process
would still require its own theorem.  The numerical coincidence of other
framework objects of cardinality 192 is not used in this proof.
\end{remark}

\begin{proposition}[Raw Newton coordinates and the coupled rescaling test]
\label{p5:prop:v2211-raw-newton-coupled-rescaling}
On one fixed metrological convention define the two coefficient coordinates
\[
 G_{J,\rm raw}:=\frac{c_{\rm op}^3}{2\pi\hbar U_J^{\rm raw}},
 \qquad
 G_{\rm scal}:=\frac{c_{\rm op}^3}{2\pi\hbar U}.
\]
Then \(G_{\rm scal}=96G_{J,\rm raw}\).  Neither coordinate is
identified with one measured physical Newton response by the raw heat
calculation alone.  Such an identification consumes the complete normalized
metric/source action map, including the actual source, gauge, scalar,
boundary and mixed metric jets.

In particular, multiplying the entire uninserted common action by a constant
\(s>0\), on the same fields and fixed cutoff, would give
\[
 \boxed{\kappa_J(s)=s\kappa_J,\qquad
 U_J^{\rm raw}(s)=sU_J^{\rm raw},\qquad
 g_i^2(s)=\frac{(g_i^{\rm heat})^2}{s}.}
\]
Thus \(U_J^{\rm raw}(s)=U\) on this merely rescaled family would force
\(s=1/96\) and simultaneously multiply every canonical gauge coupling
square by 96.  It is not a gravity-only normalization convention on the
fixed native action and is not selected here.
\end{proposition}
\begin{proof}
The coordinate ratio is the proved raw coefficient ratio.  The Einstein
coefficient is linear in the common action prefactor.  The gauge coefficient
before canonical matching is likewise multiplied by \(s\); the same
ordinary trace gives
\(s\,3k_ig_i^2(s)/(48\pi^2)=1/4\).
Divide by the corresponding native equation at \(s=1\).
The resulting simultaneous changes establish the coupled test.  They
supply no physical projection or compensating source/contact identity.
\end{proof}

\begin{corollary}[Source-native common-heat gauge normalization]
\label{p5:cor:common-heat-absolute-gauge-normalization}
On the common-heat branch of Paper~III-F, use the ordinary additive Hilbert
trace on $\mathcal H_{\rm heat}$.  The full visible carrier contains exactly
three family copies and Paper~V-A's one-generation trace indices are
\[
 k_3=2,\qquad k_2=2,\qquad k_Y=\frac{10}{3}.
\]
Let \(\mu_h\) denote the source-native common-heat matching datum and write
\(g_i^{\rm heat}:=g_i(\mu_h)\).  Canonical normalization of the Yang--Mills
kinetic terms gives
\[
 \boxed{k_i (g_i^{\rm heat})^2=4\pi^2.}
\]
Equivalently,
\[
 \boxed{
 (g_3^{\rm heat})^2=(g_2^{\rm heat})^2=2\pi^2,
 \qquad
 (g_Y^{\rm heat})^2=\frac65\pi^2,
 \qquad
 \sin^2\theta_W(\mu_h)=\frac38.}
\]
With $\alpha_s(\mu)=g_3^2(\mu)/(4\pi)$, the color coupling at the native
common-heat matching datum is
\[
 \boxed{\alpha_s^{\rm heat}:=\alpha_s(\mu_h)=\frac\pi2.}
\]
No identification of \(\mu_h\) with an observed renormalization scale, the
K16 hard-production scale, or a thermal-plasma scale is made here.  Any use of
physical \(g_i(\mu)\) away from \(\mu_h\) consumes the separately typed
G16 RG/threshold matching map.
\end{corollary}

\begin{proof}
Paper~III-F gives
$S_{\rm YM}^{\rm heat}=\frac1{48\pi^2}
\int\operatorname{Tr}_{F}(\mathcal F_{\mu\nu}^\dagger
\mathcal F^{\mu\nu})$.
The gauge generators act trivially on the generated family factor
$G\cong\mathbb C^3$, so
$\operatorname{Tr}_{F}(T_i^aT_i^b)=3k_i\delta^{ab}$.
Writing $\mathcal F_{\mu\nu}=-ig_iF^a_{i,\mu\nu}T_i^a$ and matching the
canonical coefficient $1/4$ gives
$3k_ig_i^2/(48\pi^2)=1/4$ and therefore the result.
\end{proof}

% BEGIN FRAMEWORK v3.036: POSITIVE RETAINED-TRACE GAUGE COMPARISON
\begin{proposition}[Sharp positive trace-two obstruction to unchanged weak heat normalization]
\label{p5:prop:v3036-positive-trace-two-weak-bound}
On the established $192$-dimensional heat carrier, let $T_2^a$ be
one native weak fundamental generator, zero on singlets, with
$\|(T_2^a)^2\|=1/4$ and
$\operatorname{Tr}_F[(T_2^a)^2]=3k_2=6$.
For $Q_2=I_{\mathbb S}\otimes(T_2^a)^2$ and every positive operator
$Z\succeq0$ satisfying $\operatorname{Tr}_{\mathcal H} Z=2$,
\begin{equation}\label{p5:eq:v3036-trace-two-gauge-bound}
 \boxed{0\le \frac{\operatorname{Tr}(ZQ_2)}{\operatorname{Tr}(Q_2)}
 \le\frac1{48}<1,\qquad \operatorname{Tr}(Q_2)=24.}
\end{equation}
For Paper~I-C's relative-graded primitive insertion
$Z_{\rm rp}=P_+/48$, with internal grading trace zero, the ratio is
exactly $1/96$, namely $\operatorname{Tr}(Z_{\rm rp}Q_2)=1/4$.
Thus an unchanged-generator, positive trace-two substitution in the
\emph{ordinary parity-even} gauge heat word cannot retain the native
weak kinetic coefficient.  This does not exclude a separately generated
sector-aware complete-action projection with contact, graded,
boundary or source terms.
\end{proposition}
\begin{proof}
Positivity gives $0\preceq Q_2\preceq I/4$. Hence
$0\le\operatorname{Tr}(ZQ_2)\le\operatorname{Tr}(Z)/4=1/2$,
while $\operatorname{Tr}(Q_2)=4\operatorname{Tr}_F(T_2^a)^2=24$.
For $Z_{\rm rp}$, expand $P_+=(I+\gamma_x\otimes\Gamma_W)/2$.
The graded term vanishes in this quadratic comparison because
$\operatorname{Tr}_{\mathbb S}\gamma_x=0$, so the marked trace is
$\operatorname{Tr}(Q_2)/96=1/4$.
The parity-odd graded $F\wedge F$ and mixed-boundary packets are
separate terms and cannot be substituted for the parity-even
$F_{\mu\nu}F^{\mu\nu}$ coefficient by a constant trace weight.
\end{proof}

\begin{remark}[Complete-action gauge restoration is a target, not a new term]
\label{p5:rem:v3036-gauge-restoration-obligation}
At fixed cutoff, generators and matching prescription, a uniform
$1/96$ parity-even heat replacement would give
$g_{i,\rm rp}^2=96(g_i^{\rm heat})^2$.
Recovering the native parity-even gauge word would require an actual
additional $(95/96)$ of its raw coefficient, with any graded and
boundary contacts transported. This equation specifies the comparison
target only; no such compensator is postulated or claimed to exist.
Physical source, scalar and gauge jets must be compared as one action.
\end{remark}
% END FRAMEWORK v3.036

% BEGIN FRAMEWORK v3.029: GP3 MATTER OWNER BRIDGE
\begin{proposition}[Absolute/relative common-copy gauge and heat routing]
\label{p5:prop:v329-commoncopy-heat-gauge-routing}
On the same-owner common-amplification branch $F_k=F\otimes\CC^k$
with the source's ordinary additive common-heat action, one unchanged
scalar doublet, unchanged gauge-vector sectors, and the same unbroken
threshold regime, the represented carrier and raw geometric coefficients are
\[
 \dim F_k=48k,\qquad
 \dim\mathcal H_{{\rm heat},k}=192k,\qquad
 U_{J,k}^{\rm raw}/U=96k.
\]
The one-generation gauge indices are multiplied by $k$,
\[
 (k_3,k_2,k_Y)=(2k,2k,10k/3),
\]
so at the source-native common-heat matching datum $\mu_h$,
\[
 g_3^2=g_2^2=2\pi^2/k,\qquad
 g_Y^2=6\pi^2/(5k),\qquad
 \sin^2\theta_W(\mu_h)=3/8.
\]
Using the source's GUT-normalized $b_1$ convention, the unbroken
one-loop beta coefficient vector for common charged-family
replication is
\[
 b(k)=(4k+1/10,\ 4k-43/6,\ 4k-11),
\]
with $b_1-b_2=109/15$, $b_1-b_3=111/10$, and $b_2-b_3=23/6$.
These are comparisons at fixed owner, scale/convention, and thresholds.
They do not imply experimentally calibrated absolute couplings or
physical $k=1$.
\end{proposition}
\begin{proof}
The same-owner ordinary trace on a replicated internal carrier is
$k$ times the trace on the base carrier, giving the represented
internal dimension, gauge weights and scalar Einstein coefficient.
The source-native gauge kinetic normalization is $k_i g_i^2=4\pi^2$.
The weak-angle ratio cancels the common factor $k$. At one loop, the
three replicated families contribute $4k$ to each of the normalized
coefficients, while the fixed one-doublet scalar and gauge-vector
terms retain their source constants. Taking pairwise differences
cancels $k$.
\end{proof}
% END FRAMEWORK v3.029


\begin{remark}[No PCT trace doubling and no normalized-trace freedom]
\label{p5:rem:no-pct-trace-doubling}
The closed spectral functional uses the ordinary finite trace.  A normalized
trace is not additive on licensed direct sums, and a weighted trace is an
inserted endomorphism.  Likewise the PCT theorem does not double the finite
bosonic trace carrier.  The common-heat result therefore has no continuous
trace-normalization parameter on the premise-minimal branch.
\end{remark}



% BEGIN SC v1.80 HIGGS-PROFILE-ACTION
\subsubsection{Finite position profiles of the canonical Higgs direction}
\label{p5:sec:v180-higgs-position-profile-action}

\begin{corollary}[Exact $2I$ evaluation of the canonical local Higgs action]
\label{p5:cor:v180-exact-icosian-higgs-profile-action}
Let $X_H$ be the canonical odd Higgs direction of
Theorem~\ref{p5:thm:v179-canonical-higgs-direction}.  On the stationary local
quartic domain of that theorem, set
\[
 \delta\Phi(t,x)=X_Hh(t,x),\qquad
 h(t,x)=V_a^{-1/2}\sum_{A=1}^{14}z_A(t)f_A(x),
 \quad f_A\in\mathcal P_x(m_{\rm cmp}),\quad m_{\rm cmp}=2.
\]
Then the inherited kinetic normalization of $X_H$ and the exact finite scalar
action theorem of Paper~III-F imply that its quadratic, cubic, and quartic
spatial tensors are reproduced exactly by uniform evaluation on the existing
$120$ vertices of $2I$.

The label $m_{\rm cmp}=2$ is position-side harmonic resolution used for this
comparison.  It does not alter the physical compact momentum result $m_p=1$,
does not truncate the growing reciprocal-position spectral window, and does
not supply a Yukawa
member, flavor-family assignment, pole mass, or stable vacuum.  Spatial
harmonic labels multiply the already generated Higgs direction and are not
substituted for the three family labels of the Yukawa matrices.
The exact spatial moments and their mode derivatives are Paper~1B,
Theorem~\ref{p1b:thm:design} and
Corollary~\ref{p1b:result:charge-linked-degree-two-spectra-and-polynomial-integrals}
\cite{Kugelmann2026Paper1B}. Paper~III-F supplies their use in the
normalized physical action on the Higgs direction fixed here.
\end{corollary}
% END SC v1.80 HIGGS-PROFILE-ACTION


% BEGIN OWNER-FAMILY v1.86
\subsubsection{Conditional action and continuation handoff}

\begin{proposition}[Retained comparison action handoff]
\label{p5:prop:retained-comparison-action-handoff}
Suppose the physical comparison port has been licensed with
weights \((-1,+2)\), the complete companion diagram is
retained, and the action owner is complete for the claim.
A fixed nonzero comparison imposes \(a-2b=0\) on symmetries
of that retained diagram.  Every such actual action symmetry
has its native Noether identities, which enter the generated
exact-continuation diagram on their stated domain.
For a declared local Hermitian line \(N=L_0\), in a unit frame
write \(m=\sigma e^{i\phi}\) and \(\nabla_N=d+iA_N\).  Then
\[
 \nabla m=e^{i\phi}
       [d\sigma+i\sigma(d\phi-2A_N)],\qquad
 A_N=A_R/2-A_B=5A_X.
\]
Parallel \(j_m\) requires \(d\phi=2A_N\), and parallel \(m\)
also requires \(d\sigma=0\).  In this case
\(d\phi=2A_N=10A_X\).
\end{proposition}
\begin{proof}
Differentiate the fixed-character equation of the preceding
proposition.  The complete-owner Noether theorem applies only
after all transformed sources and coefficients have been
included.  In the faithful two-line presentation, the
quadratic response varies by
\(i\lambda\langle\Psi,[D_m,Q_X]\Psi\rangle\), whose vanishing
for all \(\Psi\) is the same character equation.
On the dual square of a unitary line the connection is
\(d-2iA_N\), giving the local formula.  The Riesz construction
of \(j_m\) depends only on the normalized phase of \(m\),
which proves the parallelism assertions.  Finally the line
has charge zero under \(Y\) and five under \(X\).
\end{proof}

A coherently fixed structural incidence can restrict symmetries
without a new action term.  A dynamical comparison must instead
have its configuration, action and nonzero branch derived.
If the coefficient is transformed as a field, its own
variation participates in Noether; freezing it silently would
change the owner.  A norm-only term is phase-blind.
For a full fixed quadratic matrix the infinitesimal equation is
\(2\beta S+\alpha(ZS+SZ)=0\), so its companion blocks can impose
additional constraints beyond the one-line equation.

The parameter connection in Paper I-A and the nonflat Born
connections in III-F remain distinct until their typed
pullback/intertwiner is supplied.  Position transport carries
the existing target action, and neutral multipliers on the
two scalar circles leave the internal character equation
unchanged.  Noether alone implies neither nonzero \(m\),
parallelism, a phase kinetic term, a physical mass, a discrete
sign choice nor an actual history.  Optional action tests and
global descent are in the following appendices.
% END M8 v1.78 DYN02




\begin{theorem}[Closure ledger for the Paper--V-A finite constitution]
\label{p5:thm:reconstruction-package-sufficiency}
\label{p5:thm:gp2-forced-visible-sector-admissibility}
The following is the typed Paper--V-A closure ledger.
\begin{enumerate}[label=(\roman*), leftmargin=2em]
\item the diagonal-fibre spin-isometry supplies the internal
\(SU(2)_L\times SU(2)_R\) carrier, distinct from Lorentz chirality;
\item the incidence trimer has exact traceless Lie closure
\(\mathfrak{su}(3)\), while its local connection directions lie in
\(\mathsf J_5^{\rm conn,der}\);
\item the complementary trimer is \(G\cong\mathbb C[Z_3]\), and the
core gauge action is generation blind by provenance;
\item Paper~IV's relation-free theorem forces the tensor carrier; GP2
no-omission and Paper~II-B's GP3 multiplicity bound give multiplicity one;
\item coordinate/conjugate covariance and the distinguished line give the
candidate character-kernel ray
\(\mathbb R(T_R^3+\frac12(B-L))\), made structural-physical by
Proposition~\ref{p5:prop:v187-pcad-yfix};
\item the generated coordinate module is one weak doublet and its conjugate; and
\item the derived even ancestry fixes visible chirality; on
the generated hypercharge line, the charge table, faithful \(Z_6\) quotient,
anomaly cancellations, Yukawa slots, and representation trace weights follow.
The derived \(\mathsf J_5^{\rm conn,der}\) certificate gives the full local
connection family and the common kinetic trace on
Paper~III-F's real, derivative-free, parity-even, positive, quadratic,
uninserted single-trace domain; if
\(\mathsf R_5^{\rm SSQFT}\) occurs, Fermi/CAR and the graded heat-boundary one-loop
evaluation follow; physical conjugate-sector completion uses
\(\mathsf R_5^{\rm PCT}\).
\end{enumerate}
Thus the finite base, generated hypercharge, generated chirality, the derived
joint family, separately typed member occurrences, and analytic theorem domains
are not conflated.  Majorana
masses, Yukawa entries, the scalar vacuum, and numerical couplings are
downstream interaction or calibration data.
\end{theorem}

\begin{proof}
Apply, in order, the diagonal-fibre theorem, the trimer Lie-closure theorem,
the generated incidence theorem, Theorem~\ref{p5:thm:component-independence-composition},
Lemmas~\ref{p5:lem:intrinsic-abelian-centralizer-plane}
and~\ref{p5:lem:hypercharge-normalization-unique}, and the representation
calculations in Parts~II--IV, together with
Theorem~\ref{p5:thm:two-residual-sm-constitution}.
\end{proof}

\begin{lemma}[Exact scope of GP3 for additional finite actions]
\label{p5:lem:minimal-sector-enlargement-trichotomy}
Once a generated role has been fixed, GP3 may remove a kernel, a spectator
summand, or unmarked duplicate multiplicity that realizes the same typed
action.  It cannot exclude a faithful action belonging to a distinct
independently marked role---for example a family gauge action, another scalar
representation, an additional abelian factor, or an off-diagonal
parent-completion module.  Such an independently occurring marked role would
define an extension of the minimal generated constitution; GP3 neither
generates it nor proves its physical absence.
\end{lemma}

\begin{proof}
The first sentence is the kernel/multiplicity minimization content of GP3.
A faithful action in another role is not redundant relative to the original
role: it changes the typed observable or interaction structure.  Removing it
would therefore erase constituted data rather than quotient a redundancy.
GP3 contains no rule for deciding that the second role is physically absent.
\end{proof}

\begin{theorem}[Universal-property uniqueness of the component representation]
\label{p5:thm:part1-minimal-sector-algebra}
On the carrier \(F\) already fixed without multiplicity by
Theorem~\ref{p5:thm:component-independence-composition}, restrict the full
local factor representations to the two displayed component algebras.  Their
commuting joint representation is the tensor-product
representation of \(\mathcal A_F^{\rm comp}\) in
\eqref{p5:eq:part1-finite-algebra}--\eqref{p5:eq:part1-finite-rep}, with
identity action on \(G\).  This is a universal-property statement; it does not
derive either component algebra, identify the physical gauge subgroup, or
exclude other physical roles.
\end{theorem}

\begin{proof}
Theorem~\ref{p5:thm:component-independence-composition} fixes
\(F=W\otimes Q\otimes G\) from the actual generated semisimple component
algebra and the GP2/GP3 multiplicity sandwich.  Restrict its first two
component representations to \(\mathcal A_\chi\) and
\(\mathcal A_{q\ell}\).  Their images commute, so the
universal property induces the unique real-algebra representation
\(\rho_\chi\widehat\otimes_{\mathbb R}\rho_{q\ell}\), tensored with \(1_G\).
No carrier-resolution claim is made for these proper component subalgebras.
The stated limitations do not follow from this universal property and are
therefore recorded explicitly.
\end{proof}

\begin{remark}[Scope of finite-algebra uniqueness]
\label{p5:rem:finite-algebra-uniqueness-scope}
The theorem identifies the joint representation after its component roles
have been generated.  It gives no ``no-extra-sector'' theorem for separately
marked extensions.  In particular,
GP2 cannot classify an already available endomorphism module as an insertion
merely because that module is absent from the chosen Standard--Model field
content.
\end{remark}

\label{p5:part:gauge-anomaly-spectral-closure}

\paragraph{Orientation.}
This part collects exact consequences of the finite geometric package.
It records the representation content, force sectors, charges, anomaly
cancellation, scalar and Yukawa slots, trace weights, and the optional finite
spectral-datum interface, while displaying the generated hypercharge role and the
derived joint-connection certificate, and the analytic SSQFT/PCT domains wherever
their conclusions are used.

\subsubsection{Gauge normalization}
\label{p5:sec:part1-gauge-normalization}

\begin{proposition}[Exact representation trace weights]
\label{p5:prop:representation-trace-weights}
On \(\mathsf G_5^{\rm fin}\), tracing the squared
generators over one reconstructed generation gives
\begin{align}
k_3&=2T(\mathbf 3)+T(\mathbf 3)+T(\mathbf 3)=2,\\
k_2&=3T(\mathbf 2)+T(\mathbf 2)=2,\\
k_Y&=
6\left(\frac16\right)^2
+
3\left(\frac23\right)^2
+
3\left(\frac13\right)^2
+
2\left(\frac12\right)^2
+
1^2
+
0
=
\frac{10}{3}.
\end{align}
Hence
\begin{equation}
\label{p5:eq:part1-gauge-ratio}
k_3:k_2:k_Y=1:1:\frac53.
\end{equation}
\end{proposition}

\begin{proof}
The weak trace receives three colored copies of \(Q_L\) and one copy of
\(L_L\).  The color trace receives the two weak components of \(Q_L\) and
the two right singlets.  The displayed sum of squared hypercharges includes
the corresponding spectator multiplicities.  With
\(T(\mathbf2)=T(\mathbf3)=1/2\), the three values follow.
\end{proof}

\begin{corollary}[Typed coupling relation from the derived joint connection]
\label{p5:cor:conditional-common-kinetic-normalization}
Assume \(\mathsf G_5^{\rm fin}\), use the derived
\(\mathsf J_5^{\rm conn,der}\), and restrict the kinetic functional to the real,
derivative-free, parity-even (equivalently \(O(4)\)-even), positive,
quadratic, uninserted, single-trace domain of Paper~III-F.  All three
represented gauge factors are components of the one joint connection.
Paper~III-F,
Theorem~\ref{p3f:thm:joint-connection-uninserted-trace-closure}, therefore
gives one kinetic functional with a common prefactor and the same trace
convention.  Canonical
normalization then gives
\[
k_3g_3^2=k_2g_2^2=k_Yg_Y^2,
\]
and therefore
\[
g_3^2=g_2^2=\frac53\,g_Y^2,
\]
so
\begin{equation}
\label{p5:eq:part1-sin2}
\sin^2\theta_W=\frac{g_Y^2}{g_2^2+g_Y^2}=\frac38.
\end{equation}
A theta term, a nonscalar central weight, kinetic mixing, or any
factor-dependent inserted trace lies outside this conclusion.
\end{corollary}

\begin{proof}
The \(i\)-th curvature block enters the joint \(\Omega\) with the factor
\(g_iT_i\).  The uninserted trace action therefore gives its quadratic term
the coefficient \(C g_i^2 k_i\).  Mixed-factor traces vanish: the nonabelian
generators are traceless on their fundamental blocks, while \(Y\) is block
scalar for both nonabelian actions.  Canonical normalization with no
factor-dependent coefficient requires these three products to agree.
Substitution of \(k_3=k_2=2\) and \(k_Y=10/3\) gives the displayed coupling
relation and then \(\sin^2\theta_W=3/8\).
\end{proof}


\begin{corollary}[Conditional gauge-cone export]
\label{p5:cor:conditional-gauge-cone-export}
Assume \(\mathsf G_5^{\rm fin}\), use the derived
\(\mathsf J_5^{\rm conn,der}\), and impose the upstream \((\mathrm{CAUS})\)
domain.  On the native stationary-product, parity-even pure-gauge branch use
Paper~III-F, Theorem~\ref{p3f:thm:v174-generated-joint-gauge-phase-space-action};
on a broader branch retain the claim-specific same-variable Lorentzian
metric-realization premise.  Then every represented visible gauge component
inherits the same
metric principal characteristic cone,
\[
 \boxed{\operatorname{Char}(A_a)=\{k:g_{\rm op}^{\mu\nu}k_\mu k_\nu=0\}},
\]
modulo the gauge direction, and therefore
\[
 \boxed{c_{{\rm front},a}=c_{\rm op}.}
\]
This conclusion uses no absolute gauge calibration: the common action constant,
representation trace indices, and canonical field rescalings do not alter the
principal characteristic set.  On the native theorem domain both the
connection and local Yang--Mills action are generated outputs; only a broader
Lorentzian branch retains the same-variable realization premise.
\end{corollary}

\begin{proof}
The preceding corollary places all factor curvatures in the one uninserted
joint kinetic functional.  Paper~III-F's joint gauge principal-symbol
inheritance theorem then contracts every factor with the same generated
Lorentzian metric owner.  Positive internal trace factors multiply the
principal symbol but do not change its zero set.  The connection certificate
is supplied chronologically earlier by the explicit constructor word, not
inferred backward from this characteristic consequence.  The same-process
factor relations follow from Paper~III-F's same-datum closure theorem.
\end{proof}

\begin{corollary}[Conditional diagonal-fibre left--right normalization]
\label{p5:cor:born-spherical-left-right-normalization}
Under Paper~III's separate normalized gauge-reduction package
\(({\rm GRED})\), the round diagonal fibre gives
\[
g_L=g_R,
\qquad
\frac1{g_L^2}=\frac1{g_R^2}
=\frac{r_{\rm fib}^2}{16\pi G_B}
=\frac{f_2}{48}.
\]
This is a conditional weak left--right statement.  It neither supplies the
color or hypercharge gauge action nor proves
\(g_3=g_2=\sqrt{5/3}\,g_Y\).  That stronger relation follows on
\(\mathsf J_5^{\rm conn,der}\) from the preceding corollary.  A separately
calibrated diagonal-circle \(U(1)_{\rm fib}\) would be a different connection
occurrence and is not part of the projective visible connection.
\end{corollary}

\begin{proof}
Paper~III's diagonal-fibre bulk-reduction theorem gives the displayed
coefficient only under \(({\rm GRED})\), including its round-metric,
bulk-matching, cell-saturation, and generator-normalization clauses.  Its
theorem contains no color or hypercharge kinetic term, so by itself it gives
no cross-factor equality; that equality comes from the joint-connection trace
theorem.
\end{proof}

\begin{remark}[The former \(24\cdot7\) gauge boundary is superseded]
\label{p5:rem:born-gauge-trace-support}
The Hurwitz orbit size \(24\) and scalar-free local rank \(7\) remain valid
framework counts, but their Cartesian product is not the native bosonic heat
trace carrier.  The preceding theorem now supplies the actual least
self-adjoint carrier \(\mathbb S\otimes F\), the ordinary trace, and the
three-family multiplicity.  Consequently
\[
24\cdot7=168,
\qquad
\alpha_B^{-1}=42
\]
remains only a legacy numerical diagnostic and is not used as a boundary
condition or export.  The absolute source-native common-heat gauge
normalization is instead the computed Corollary~\ref{p5:cor:common-heat-absolute-gauge-normalization}; the broader kinetic-class branch may still carry its separately typed overall calibration.
\end{remark}

\begin{remark}[Status of the \(3/8\) weak-angle normalization]
\label{p5:rem:weak-angle-normalization-status}
The value \(\sin^2\theta_W=3/8\) is a consequence of the generated
representation, the generated hypercharge line and the theorem certificate
\(\mathsf J_5^{\rm conn,der}\), only on the real, derivative-free, parity-even,
positive, quadratic, uninserted single-trace kinetic domain.  The relative
trace weights are finite theorem outputs on the selected representation; the
common normalization follows from the typed trace theorem on the joint
connection family.  Theta terms, nonscalar central insertions, and
factor-dependent trace weights are outside the theorem.

Low-energy running and threshold comparison are downstream and are not used here as structural inputs.  In particular, this common-joint-action coupling
coordinate is not identified by name alone with a later broken-shell/tree-mass
reconstruction.  Any such same-role comparison requires an independently
generated running/matching map and a declared common scale and convention.
Thus a later tree-shell value such as \(\sin^2\theta_W=2/9\) is not a
same-stage contradiction of \eqref{p5:eq:part1-sin2}; an identity matching of
the two typed roles would itself require proof.
\end{remark}

\begin{theorem}[Geometric reconstruction of the minimal visible representation]
\label{p5:thm:conditional-sm-representation}
\label{p5:thm:part1-anchor-free-sm-carrier}
On \(\mathsf G_5^{\rm fin}\), the generated finite
carrier carries the
Standard--Model representation
\begin{equation}
Y=T_R^3+\frac12(B-L),
\qquad
Q=T_L^3+Y,
\end{equation}
and the projector sectors have exactly the hypercharge and electric-charge
spectrum in \eqref{p5:eq:projector-hypercharge-discharge}.  Its faithful
connected global group is the \(Z_6\) quotient, its anomalies cancel
generationwise at the coefficient level, and its one-doublet Yukawa slots are
the four displayed Dirac slots plus the unique optional Majorana support.  No
observed charge value is used to fit \(Y\).  Local gauge realization uses
\(\mathsf J_5^{\rm conn,der}\), and physical Fermi/QFT anomaly statements use
\(\mathsf R_5^{\rm SSQFT}\).  Physical conjugate-sector completion, when
claimed, uses \(\mathsf R_5^{\rm PCT}\).
\end{theorem}

\begin{proof}
Theorem~\ref{p5:thm:component-independence-composition} gives the tensor
carrier from relation-free composition and the GP2/GP3 blockwise multiplicity
sandwich.
Theorem~\ref{p5:thm:trimer-su3-lie-closure} computes the color-trimer Lie
closure on the incidence trimer.  Lemma~\ref{p5:lem:part1-BminusL-unique}
gives the normalized traceless \(1\oplus3\) splitter.  The full diagonal
centralizer is five-dimensional modulo the identity;
Lemma~\ref{p5:lem:intrinsic-abelian-centralizer-plane} reduces it to the
\(T_R^3,B-L\) plane using coordinate conjugacy.  The distinguished-line
character kernel then gives
\[
Y=T_R^3+\frac12(B-L)
\]
by Lemma~\ref{p5:lem:hypercharge-normalization-unique}.  The scalar stabilizer
gives \(Q=T_L^3+Y\).  Evaluating these operators yields the charge table;
Propositions~\ref{p5:prop:faithful-global-sm-group}
and~\ref{p5:prop:three-generation-content} and
Theorems~\ref{p5:thm:part1-anomaly-cancellation}
and~\ref{p5:thm:part1-yukawa-majorana-unique} give the remaining claims.
\end{proof}



\section{Consolidated geometric results and downstream limits}
\label{p5a:sec:organized-results}
\subsection{Conclusion}
\label{p5:sec:conclusion}
Paper~V-A now proves the finite Standard--Model constitution from the oriented
geometry.  Axial-frame incidence, relation-free composition, the coordinate
doublet type, nonabelian algebra, Abelian candidate plane and character
kernel, the two relative-parity projectors, conjugate-representation
availability, the even external/internal ancestry, and the blockwise
multiplicity-one carrier are derived.  Release v1.86 additionally generates the complete constrained weighted
Yukawa--Dirac owner family as a whole exact-solution family; this alone does not
select an arbitrary member of that generic family. On the occurring rooted
charged branch, however, Theorem~\ref{p5:thm:v397-charged-preterminal-srcval-realization}
supplies the particular zero-order source packet and
Theorem~\ref{p5:thm:v398-charged-rooted-comparison-owner-consumption}
supplies its comparison owner. A physical interacting pole/rest identification,
an additional field/background member or a vacuum is not supplied by this
finite source closure.  Physical structural hypercharge follows from
Proposition~\ref{p5:prop:v187-pcad-yfix}; visible chirality follows from
\eqref{p5:eq:chi5-derived}.  The local gauge datum is the derived joint family
\(({\rm JOINT}_5^{\rm der})\).  Paper~III-F then derives the former JOINT
factor relations, and on its real, derivative-free, parity-even, positive,
quadratic, uninserted single-trace domain derives the common kinetic
normalization.  Complex conjugation supplies \(\bar F\) but not a finite
occurrence theorem for \(F\oplus\bar F\); physical conjugate-sector completion
uses \(\mathsf R_5^{\rm PCT}\).  On
\(\mathsf R_5^{\rm SSQFT}\), the Fermi/CAR branch is derived chronologically
later than Paper~IV and local anomaly and one-loop statements are licensed.
Additional bulk Majorana targets, a Higgs stationary history, requested
low-energy coupling coordinates, thresholds, and renormalized continuum
dynamics retain their downstream domains. The fixed charged mass/incidence
packet and native weak-current readouts are generated descendants rather
than freely selected Yukawa entries.


% Paper V-A, after p5:prop:v2208-boundary-only-kernel.
\begin{proposition}[Boundary-only reduction cannot recover a kernel response]
\label{p5:prop:v246-boundary-elimination-no-kernel}
Let $B:\mathcal H_{12}\to\mathbb C^6$ be the fixed boundary
coisometry and $Q$ its kernel projector. Suppose a $C^2$ action on a
configuration chart factors as $S(f,y)=s(Bf,y)$. On any differentiable
stationary branch $y=y_*(Bf)$, its reduced action has
\[
 D S_{\rm eff}(f)[Qv]=0,\qquad
 D^2S_{\rm eff}(f)[Qv,w]=0.
\]
For a regular quadratic auxiliary block, the exact Hessian reduction is
\[
 H_{\rm eff}=B^\dagger(A-CD^{-1}C^\dagger)B,
 \qquad H_{\rm eff}Q=QH_{\rm eff}=0,
\]
where $A=A^\dagger$, $D=D^\dagger$ is invertible, and the full
Hessian blocks are $B^\dagger AB$, $B^\dagger C$, $D$.
\end{proposition}
\begin{proof}
The reduced function equals
$s(Bf,y_*(Bf))$ and is constant along every $f+tQv$ because $BQ=0$.
Differentiate that identity twice. In the quadratic case the stationary
auxiliary equation gives the stated Schur complement; multiplying it
by $Q$ on either side proves the claim.
\end{proof}
In particular, a boundary-only parent cannot produce the nonzero
$E_{163}$ direct or inverse collective role in the conditional neutral
rank word by regular elimination. Such a role requires an independently
proved full-owner kernel incidence before reduction. This condition does
not exclude the differently typed full-owner collective trace constructor.



% BEGIN FRAMEWORK v2.266: PROTECTED READOUT EXPORT FOR MATTER GEOMETRY
\subsection{Minimal protected observation interfaces and member-occurrence scope}
\label{p5:sec:v266-protected-observation}
\begin{corollary}[Mathematical chirality and population tests need no full tomography]
\label{p5:cor:v266-protected-observation}
On the intrinsic retained finite coordinate carrier, one can
mathematically measure either the native protected-block projector
$P_S$ or its chiral contrast $P_S\chi$ together with exact intrinsic
coordinate first moments with six outcomes at any classical label
radius $r\ge2\Rmom$. Both are directly and independently accessible
with seven outcomes at that radius; fewer outcomes cannot span the
corresponding six- and seven-dimensional real operator systems.
A five-outcome position readout, although its effects generate the
full native coordinate $C^*$-algebra, is blind to both in one setting.
\end{corollary}
\begin{proof}
The explicit effects, moment identities and dimension lower bounds
are Theorem~\ref{p2c:thm:v266-minimal-protected-readout}; the
algebra-versus-effect-span distinction is
Proposition~\ref{p2c:prop:v266-repeat-fingerprint}.
\end{proof}
The protected readout statement is about the original $11\oplus4$
carrier, not a derived physical Standard-Model member state. It
neither fixes an odd solder, matter embedding, charge normalization,
source occupation, interaction vertex, physical mass, nor a
same-action detector readout for charged matter.
% END FRAMEWORK v2.266: PROTECTED READOUT EXPORT FOR MATTER GEOMETRY


% BEGIN FRAMEWORK v2.321: PETERSEN INDUCED LOCAL MATTER-CARRIER COMPATIBILITY
\section{Petersen-induced local matter comparison and its finite status}
\label{p5:sec:v321-petersen-matter}

The Petersen--Hopf packet generated in Paper~I-A is a fifteen-dimensional
\emph{oriented edge} representation, not the independent fifteen-dimensional
level-one momentum carrier. We derive its exact restriction to one occurring
binary tetrahedral core before asking whether a physical matter operation
consumes it. This entire subsection has the explicit additional domain
\emph{an occurring embedded $2I$ support}; the local $2T$ theorem alone
does not force the icosahedral branch. Let $I=2I$, $H_i\cong2T$, $A_5=I/\{\pm1\}$ and
$A_{4,i}=H_i/\{\pm1\}$. Put $\mathscr W_P=\mathbb R^{E(P)}$, where a
chosen orientation of each Petersen edge determines its signed permutation
action. The edge orientation is a coordinate convention: reversing one edge
conjugates all operators by a diagonal sign matrix.

\begin{theorem}[Five local Clifford triples induce the Petersen edge module]
\label{p5:thm:v321-induced-local-edge-module}
For each $i\in\{1,\ldots,5\}$, let $V_i$ be the span of the three
Petersen edges whose two vertex-pairs omit the label $i$. Then
\[
 \boxed{\mathscr W_P=\bigoplus_{i=1}^5V_i,\qquad
 \dim_{\mathbb R}V_i=3,\qquad
 \mathscr W_P\cong\operatorname{Ind}_{A_4}^{A_5}(\mathbf3).}
\]
The summands are mutually orthogonal in the standard edge inner product;
$A_5$ permutes them, and the stabilizer $A_{4,i}$ acts on $V_i$ as its
irreducible real tetrahedral vector representation $\mathbf3$. Via the
oriented, overlap-signed cross-product directions of Paper~I-A, $V_i$ is
isometric and $A_{4,i}$-equivariant to the imaginary-quaternion three-axis
frame belonging to $H_i$. Under the full $A_5$ action,
\[
 \mathscr W_P\cong\mathbf4\oplus\mathbf5\oplus\mathbf3\oplus\mathbf3',
\]
recovering Paper~I-A's exact fifteen-edge Hodge resolution.
\end{theorem}
\begin{proof}
An edge of $KG(5,2)$ joins two disjoint two-subsets and therefore omits
exactly one label. For a fixed omitted label there are the three partitions
of the remaining four labels into two unordered pairs. This proves the
orthogonal five-by-three partition and the induced representation statement.
The fixed-label group $A_4$ permutes the three partitions with the signed
edge orientations. Its character is $3,-1,0$ on elements of order $1,2,3$,
respectively, hence it is the real irreducible tetrahedral $\mathbf3$.
The three normalized cross products form the orthogonal quaternionic frame
and transform with the same signed action. Frobenius reciprocity, or the
fifteen-edge projectors of Paper~I-A, gives the four distinct $A_5$
constituents displayed above.
\end{proof}

\begin{theorem}[Rooted $3+12$ reduction and component--family comparison]
\label{p5:thm:v321-rooted-3plus12-matter}
Fix the occurring $H_i$. Its stabilizer preserves
\[
 \mathscr W_P=V_i\oplus V_i^\perp,
 \qquad \dim(V_i,V_i^\perp)=(3,12),
\]
where $V_i^\perp$ is one free transitive $A_{4,i}$ edge orbit. Thus
\[
 \boxed{\mathscr W_P\!\downarrow_{A_{4,i}}
  \cong\mathbf3\oplus\operatorname{Reg}(A_4).}
\]
On the \emph{specified comparison action} of Paper~I-A's rooted
icosian-complement theorem, Paper~V-A's already generated modules obey
\[
 Q=\ell\oplus C\cong\mathbf1\oplus\mathbf3,
 \quad G\cong\mathbb C[C_3]
     \cong\mathbf1\oplus\mathbf1'\oplus\mathbf1'',
 \quad Q\otimes G\cong\operatorname{Reg}(A_4).
\]
Consequently there is a complex $A_{4,i}$-module isomorphism
\[
 \boxed{\mathbb C\otimes_{\mathbb R}\mathscr W_P
     \cong C\oplus(Q\otimes G).}
\]
The local $C$ comparison is fixed up to oriented quaternionic frame
convention. The regular twelve-channel comparison is not a canonical
intertwiner of the complete physical $SU(3)$ gauge representation, weak
representation, Dirac grading, charged source or family Yukawa operators.
\end{theorem}
\begin{proof}
For a fixed omitted root $i$, every other Petersen edge has one endpoint
containing $i$. It is indexed by an ordered pair $(j,k)$ of distinct labels
in the four-element complement of $i$. The natural $A_4$ action on these
twelve ordered pairs is free and transitive, hence the oriented edge
representation is a regular module after a compatible edge-sign choice.
The $A_4$ characters of $Q$ and $G$ are respectively $1+\mathbf3$ and
the sum of the three one-dimensional characters. Their product is
$\mathbf1\oplus\mathbf1'\oplus\mathbf1''\oplus3\mathbf3$, the regular
character. The physical proviso follows from the distinct source and
operator owners stipulated by the local-core factorization theorem.
\end{proof}

\begin{corollary}[No five-family, extra-color or new-Higgs inference]
\label{p5:cor:v321-petersen-nospurious-matter}
The five spaces $V_i$ form one $A_5$-induced geometric packet, not five
independent matter families. The constituted replication carrier remains
$G\cong\mathbb C^3$, the component carrier remains $\ell\oplus C$, and
the generated coordinate Higgs acts as the identity on $Q\otimes G$.
The geometric comparison alone produces no new $SU(3)$ color group, charge,
fermion species, Yukawa coefficient, pole, or mass. In particular it does
not contradict the established theorem that the actual local Standard--Model
core factors through a single $2T$ support.
\end{corollary}

\begin{remark}[Typed frontier for a physical cross-core operator]
\label{p5:rem:v321-crosscore-matter-frontier}
An action-owned realization must supply a specific intertwiner between the
rooted Petersen $12$-orbit and $Q\otimes G$, together with the applicable
weak/chiral, gauge and charged-source operations. $A_4$ equivalence of the
bare modules is necessary for this proposed route but not sufficient.
The mirror $\mathbf3'$ is a sector of the geometric cycle packet and is not
a newly generated physical matter multiplet. A physical source that couples
to it must be demonstrated in Paper~V-C's same-action calculus.
\end{remark}
% END FRAMEWORK v2.321




% BEGIN FRAMEWORK v2.338: NATIVE TWO-SOURCE PETERSEN--CLIFFORD BRIDGE
\section{Rooted oriented-pair and quadratic source capacities}
\label{p5:sec:v338-native-two-source}
All assertions below hold in the finite five-core permutation and
exterior-algebra carriers.  Neither source occurrence nor a physical
quark--antiquark embedding is inferred.

\begin{theorem}[One rooted scalar, no unrooted oriented-pair scalar]
\label{p5:thm:v338-oriented-pair-rooted-scalar}
For $W=\mathbb R^5$ and $U=\Lambda^2W$, the $A_5$-module is
$U\cong\mathbf4\oplus\mathbf3\oplus\mathbf3'$, so
$U^{A_5}=0$.  Fix one core $r$; the root stabilizer $A_4$
gives $U|_{A_4}\cong\mathbf1\oplus3\mathbf3$.
Its unique unit invariant vector is
\[
 w_r=e_r\wedge\frac12\sum_{j\ne r}e_j,
 \qquad \langle w_r,w_s\rangle=\begin{cases}1&r=s,\\-1/4&r\ne s.
 \end{cases}
\]
Consequently $\sum_rw_r=0$, and
$\sum_rw_rw_r^{\mathsf T}=\tfrac54P_4$, where $P_4$ is the
four-dimensional contrast projection of
Theorem~\ref{p1:thm:v338-signed-petersen-hodge}.
In particular $P_+w_r=P_-w_r=0$.
\end{theorem}
\begin{proof}
Decompose $W=\mathbf1\oplus\mathbf4$ under $A_5$ and
$\Lambda^2W$ as in the preceding geometry theorem.  Upon restriction
to the root stabilizer, the four-dimensional $A_5$ contrast restricts
to $\mathbf1\oplus\mathbf3$ and either icosahedral triplet restricts
to the tetrahedral triplet.  Direct wedge contractions give the
specified invariant, the Gram values and the tight-frame identity.
The unique $A_4$ fixed line lies in $P_4U$.
\end{proof}

\begin{theorem}[Quadratic cyclic phase channels and Pfaffian test]
\label{p5:thm:v338-quadratic-cyclic-source}
Fix root $0$ and write the four other cores as $1,2,3,4$.
For $x=\sum_{i<j}x_{ij}e_i\wedge e_j$ define
\[
 q_1=x_{12}x_{34},\qquad q_2=-x_{13}x_{24},\qquad
 q_3=x_{14}x_{23},\qquad \omega=e^{2\pi i/3}.
\]
The root-preserving $A_4$ permutes $(q_1,q_2,q_3)$ through
$A_4/V_4\cong C_3$.  Their normalized Fourier modes
\[
 q_0=(q_1+q_2+q_3)/\sqrt3,\quad
 q_\omega=(q_1+\omega^2q_2+\omega q_3)/\sqrt3,\quad
 q_{\bar\omega}=(q_1+\omega q_2+\omega^2q_3)/\sqrt3
\]
realize $\mathbf1\oplus\mathbf1'\oplus\mathbf1''$.
For real $x$, $q_{\bar\omega}=\overline{q_\omega}$ and
\[
 |q_\omega|\leq\frac{\|x\|^2}{2\sqrt3}.
\]
The global covariant Pfaffian vector
$\mathcal P(x)=*_5(x\wedge x)$ satisfies
\[
 \mathcal P(x)=0\ \Longleftrightarrow\ \operatorname{rank}(x)\leq2,
 \qquad \|\mathcal P(x)\|\leq\|x\|^2.
\]
For $x=(e_{12}+e_{34})/\sqrt2$ both bounds can be saturated
(with a root and orientation chosen consistently).
\end{theorem}
\begin{proof}
Every root-preserving even permutation induces one of the three
cyclic permutations of the three perfect matchings of four labels;
the three Fourier characters yield the irreducible lines.  Bound
the products $x_{ab}x_{cd}$ by half their squared coefficient
norms and use the Fourier triangle inequality.  In an orthonormal
skew normal form in five dimensions,
$x=s_1e_{12}+s_2e_{34}$, one has
$\mathcal P(x)=2s_1s_2e_5$ and
$\|x\|^2=s_1^2+s_2^2$.  The arithmetic--geometric mean inequality
proves the bound and the decomposability criterion.
\end{proof}

\begin{theorem}[Three equivalent fifteen-dimensional two-source carriers]
\label{p5:thm:v338-triple-clifford-edge-module}
For $U=\Lambda^2W$, the alternating two-source space splits
orthogonally and $A_5$-equivariantly as
\[
 \Lambda^2U=E_{15}\oplus O_{15}^+\oplus O_{15}^-,
 \qquad 45=15+15+15.
\]
The first summand consists of wedges of \emph{disjoint} pair labels,
which are exactly the fifteen Petersen edges.  The remaining thirty
basis directions have a shared root; the Hodge star on each rooted
four-dimensional complement splits their six-dimensional space into
three positive and three negative directions.  Each resulting
fifteen-dimensional $A_5$-module has character $15$ at the identity,
$-1$ on double transpositions, and zero on three- and five-cycles,
thus each is isomorphic to
$\mathbf3\oplus\mathbf3'\oplus\mathbf4\oplus\mathbf5$.
\end{theorem}
\begin{proof}
Disjoint unordered pair pairs number $\tfrac12\cdot10\cdot3=15$.
Each of the five possible common roots gives six overlapping
basis-pair wedges.  The rooted Hodge star is an involution with
three eigenvectors of each sign; even permutations carry each
rooted Hodge summand to the corresponding target root.  The resulting
three invariant summands all have the indicated character by the
signed rooted permutation action, so the irreducible decompositions
agree by the $A_5$ character table.
\end{proof}

\begin{proposition}[Bilinear Petersen--Clifford bridge, and its limitation]
\label{p5:prop:v338-bilinear-edge-map}
For two oriented-pair source amplitudes $x,y\in U$ define, on every
Petersen edge $\{A,B\}$,
\[
 \mathsf J(x,y)_{\{A,B\}}=x_Ay_B-x_By_A,
 \qquad A\cap B=\varnothing.
\]
With one consistently chosen oriented-edge sign convention this is
an alternating, surjective, $A_5$-equivariant map
$\Lambda^2U\twoheadrightarrow E_{15}$; the convention differs
from the existing rooted Clifford-edge convention only by a fixed
diagonal Levi-Civita sign gauge.  Its quadratic contact contraction is
\[
 \mathsf J^\dagger\mathsf J=P_{E_{15}},\qquad
 \operatorname{Spec}(P_{E_{15}})=\{1^{(15)},0^{(30)}\}.
\]
A chosen root separates a three-dimensional local Clifford triple
and its free, twelve-dimensional $A_4$ orbit; the latter is
representation-equivalent to the existing comparison module
$Q\otimes G\cong\operatorname{Reg}(A_4)$.  This equivalence does
\emph{not} fix a unique isometric map or a physical interaction.
\end{proposition}
\begin{proof}
Project the exterior product $x\wedge y\in\Lambda^2U$ onto the
fifteen disjoint basis wedges.  This is an orthogonal coordinate
projection, proving the contraction and spectrum.  Its sign action
is the standard induced orientation action and is gauge-equivalent
to the Clifford-edge module by orienting the omitted label.
The rooted edge orbit and its regular representation were established
in the preceding rooted Petersen construction.  The scalar contact
projector alone has one nonzero level and cannot distinguish two
positive meson rest masses.
\end{proof}
\begin{remark}[Representation deficit is a source possibility, not a particle]
The rooted linear carrier is $\mathbf1\oplus3\mathbf3$, whereas
$Q\otimes G=\mathbf1\oplus\mathbf1'\oplus\mathbf1''\oplus3\mathbf3$.
The missing two one-dimensional representation types appear in the
quadratic Fourier source above.  Symmetry nevertheless leaves
$\dim_{\mathbb C}\operatorname{Hom}_{A_4}(U,Q\otimes G)=10$
possible linear intertwiners.  No action-selected, color-singlet
physical map is supplied by this representation comparison.
\end{remark}
% END FRAMEWORK v2.338



% BEGIN FRAMEWORK v2.352: STATIC-SOURCE COMMUTANT INTERFACE
\clearpage
\subsection{Constant coordinate-compatible sources do not encode a flavor spectrum}
\label{p5a:sec:v352-static-source-interface}
\begin{corollary}[The momentum-side block-central source obstruction]
\label{p5a:cor:v352-static-source-obstruction}
On the canonical one-copy level-one momentum carrier, a constant
Hermitian source $M$ commuting with the \emph{entire} intrinsic
$M_{11}\oplus M_4$ algebra is necessarily
\[
 \boxed{M=\mu_G I_{11}\oplus\mu_S I_4},
\]
and therefore has at most two distinct singular-value magnitudes.
Such a source alone cannot produce a differentiated charged-family
nine-mass singular spectrum.  This restriction is distinct from the
root-involution no-go proved above, and it does \emph{not} apply
automatically to the $48$-dimensional fermion source, a noncommuting
flavor response, a frequency-dependent propagator, or a mass pole.
\end{corollary}
\begin{proof}
This is the commutant characterization of
Proposition~\ref{p2c:prop:v352-static-source-commutant}
restricted to its expressly stated one-copy and static-source domain.
No identification of the momentum and fermion field algebras is made.
\end{proof}
% END FRAMEWORK v2.352

\printbibliography
% BEGIN FRAMEWORK v2.241: ROOTED CONORMAL GRAVITY AUDIT
\section{The rooted source in the matched conormal boundary application}
\label{p5:sec:v229-rooted-boundary-application}
The rooted finite source $T_*=c(h_*)\otimes I_Q\otimes I_G$,
$h_*=-e_1$, retains the geometric origin and full rank already proved
here. For the explicit parent-diagonal matching $T_{\rm par}=c(u)$,
$u=(1,1,1,1)/2$, Proposition~\ref{p1b:prop:v229-parent-overlap}
evaluates its singlet overlap as $-1/2$ and complementary squared
weight as $3/4$. In the further frozen constant triplet realization
$g=2/a$, its squared background residual is $3/a^2$.
These are matching-dependent applications, not universal assertions for
an uncomputed left/right identification. A different actual matching
uses its own $\operatorname{Tr}(T_{\rm par}^\dagger T_*)/2$.

Theorem~\ref{p1b:thm:v229-rooted-gap} shows that leakage is compatible
with a full $m/2$ lower bound under its constant-mass and self-adjoint
domain hypotheses. Theorem~\ref{p1b:thm:v229-conormal-response}
evaluates the complete associated boundary source without deleting
the complementary modes. Neither statement identifies that normal
gap with a physical interacting fermion pole or selects a gravitational
scalar density. The charged hierarchy-origin status of this paper is
unchanged; the conormal application asks for no new primitive source.

% END FRAMEWORK v2.241: ROOTED CONORMAL GRAVITY AUDIT


% BEGIN FRAMEWORK v2.345: FIVEFOLD CORE-TRANSITION EXPORT
\clearpage
\section{Fivefold transitions do not preserve one rooted matter chart}
\label{p5:sec:v345-core-transition-boundary}
\begin{proposition}[A five-chart comparison is not an additional family]
\label{p5:prop:v345-core-fivefold-transport}
Let $H\cong2T$ be an embedded tetrahedral subgroup of
$G=2I$, and let $g$ belong to the twelve-element fivefold shell
$\mathcal S$ of Theorem~\ref{p1:thm:v345-f40-edge-regular}.
Then $g\notin H$ and on $\ell^2(G)$,
\[
 \boxed{P_HL_gP_H=0.}
\]
Moreover $g$ acts as a five-cycle on the five coset charts
$G/H$ and on the five conjugate tetrahedral-core labels.
The disjoint cosets $g_iH$ partition the $120$ group elements,
but the five conjugate subgroups $g_iHg_i^{-1}$ overlap.
The two five-element index sets serve different geometric roles;
neither is the three-dimensional generation carrier $G$ of
Paper~V-A.  A representation-theoretic identification is not
an action-derived matter-state embedding.
\end{proposition}
\begin{proof}
The scalar parts of elements of $H$ belong to
$\{0,\pm\tfrac12,\pm1\}$, while that of every $g\in\mathcal S$
is $(\sqrt5-1)/4$.  Left multiplication takes the support $H$
outside itself.  The quotient group action gives the five-cycle
statement, and the coset partition and subgroup conjugation
identities give the final distinctions.
\end{proof}
% END FRAMEWORK v2.345


% BEGIN FRAMEWORK v2.348: ROOTED SOURCE VS. TERMINAL MASS OPERATOR
\clearpage
\section{A unit Clifford solder does not by itself split fermion masses}
\label{p5a:sec:v348-involution-mass-owner}
\begin{theorem}[No spectral mass splitting from the root involution alone]
\label{p5a:thm:v348-root-involution-mass-nogo}
On the occurring rooted source branch of
Proposition~\ref{p5:prop:v392-rooted-member-full-rank}, retain the
self-adjoint odd operator $D_{F,*}^{\rm geom}$ with
$(D_{F,*}^{\rm geom})^2=I_F$.  Its modulus and its singular values are
exactly
\[
 \boxed{|D_{F,*}^{\rm geom}|=I_F,\qquad
        \operatorname{sv}(D_{F,*}^{\rm geom})=\{1^{(48)}\}.}
\]
For every scalar function $f$ defined on $\{1\}$,
\[
 \boxed{f((D_{F,*}^{\rm geom})^2)=f(1)I_F.}
\]
Consequently no terminal scheme depending \emph{only} on the squared
root source through scalar functional calculus can produce nine
nondegenerate charged-family masses.
\end{theorem}
\begin{proof}
Self-adjointness gives
$(D_{F,*}^{\rm geom})^\dagger D_{F,*}^{\rm geom}=I_F$.
The positive square root is $I_F$, and spectral calculus on the
singleton spectrum $\{1\}$ gives the second identity.  All singular
values are therefore equal.  The conclusion is limited to this functional
class; it does not forbid source-dependent operators using additional
noncommuting generated data.
\end{proof}
\begin{remark}[Owner and physical-pole firewall]
\label{p5a:rem:v348-source-mass-firewall}
The root involution constitutes a full-rank chiral solder, not a
nondegenerate observed mass spectrum.  The geometric terminal word of
Paper~V-B and a physical rest-mass/pole readout of the complete
same-action propagator remain different typed owners.  In particular,
applying a scalar function to the involution does not bypass the
source--propagation--detector or pole/residue occurrence conditions.
\end{remark}
% END FRAMEWORK v2.348


% BEGIN CHARGED-ACTION ORIGIN AND READOUT LEDGER v3.215
\begin{remark}[Charged rooted source, incidence and action-owner ledger]
\label{p5:rem:v3215-charged-origin-ledger}
The primitive rooted $D_4$ triality/Clifford odd member of
Theorem~\ref{p5:thm:v392-rooted-srcval-incidence} is an actual source
on its declared domain but has the family identity $I_G$ and therefore
does not itself supply the directed noncentral charged-family arrows.
The rooted $R_G,X,K_7,K_9$ data and the normalized polar lepton arrows
are established geometric incidence/operation data.  Their availability
in $M_3(\mathbb C)$ does not imply action-source occurrence.  For each
$s\in\{u,d,e\}$, Theorem~\ref{p5:thm:v397-charged-preterminal-srcval-realization}
consumes the independently generated, sector-owned directed arrows of
Definition~\ref{p5:def:v228-hierarchy-origin}; its charged occurrence
conclusion is conditional on $\mathsf{Cert}_{\rm hier}$.  The quark-arrow
certificates are not inherited from the charged-lepton polar chart.
The inherited term ``charged source closure'' means assembly and
co-realization \emph{on that certified-input domain}, not proof that a
missing action constructor has spontaneously selected the arrows.
The actual normalized mixed Hessian and any marked current response
must be calculated from one action.  The physical right-handed Hopf
orientation need not be selected for an observable factoring through
$RR^\dagger$, but right-sensitive readouts require their own occurring
mark.  Neither terminal mass coefficients nor scalar spectral traces
supply the independent source-origin certificate.  The action insertion,
physical charged pole/residue, G16 transport and hadronic decay readouts
remain separate proof obligations.
\end{remark}
% END CHARGED-ACTION ORIGIN AND READOUT LEDGER v3.215


% BEGIN CHARGED ORIGIN DECISION LEDGER (v3.216)

\begin{remark}[Two noncircular charged-source or pole closure routes]
\label{p5:rem:charged-origin-two-route-test}
The charged hierarchy origin certificate
\(\mathsf{Cert}_{\rm hier}\) of
Definition~\ref{p5:def:v228-hierarchy-origin} remains an independent
hypothesis of the charged source-assembly theorem. The mathematical
source-algebra no-go and signed-target image obstruction do not
upgrade a geometric matrix to an occurring source.
A direct source route must compute the canonically normalized actual
mixed action jet on the same source owner,
\[
 M_s^{\rm phys}=Z_{L,s}^{-1/2}
 \left.\frac{\delta^2S_*}{\delta\bar\psi_{L,s}\,\delta\psi_{R,s}}
 \right|_{\text{declared background}}Z_{R,s}^{-1/2},
\]
from a frozen, result-independent physical constructor and then
verify all three oriented contractions and the norm identity of
Proposition~\ref{p5:prop:v245-charged-four-invariant-source-test}.
Alternatively, an observable-minimal route may derive the complete
source-visible interacting inverse response and prove its physical
poles and residues without identifying a noncentral primitive source.
Neither route is completed by a candidate signed Gram, a fit to masses,
or regular reduction of a family-blind primitive source alone.
\end{remark}

% END CHARGED ORIGIN DECISION LEDGER

% BEGIN v3.058 ROOTED TWO-CONTACT RETURN MATHEMATICS
% PROPOSED, NOT MERGED: Paper5A-RootedTwoContactReturn-v1_1.tex
% Baseline: Paper5A-FiniteMatterGeometry-v3_217.tex from PapersMinimal v3_053.
% Typed status: fixed-chart exact finite mathematical return, NOT occurring charged source.
% Suggested insertion near the charged-source / noncentral-family source diagnostic section.
\subsection{Minimal rooted Hopf return correction and its source-jet boundary}
\label{p5:sec:v3218-rooted-two-contact-return}

\begin{theorem}[Exact two-contact Hurwitz return and positive chiral representative]
\label{p5:thm:v3218-two-contact-hurwitz-return}
Use the generated operators of
Theorem~\ref{p1:thm:v556-rooted-hurwitz-family-return-algebra}
with $K_9=I_G$, $K_7=2I_G+2J_G$ and
$R_G=\operatorname{diag}(1,-i,i)$, and the fixed rooted common family
chart of Theorem~\ref{p6:thm:v240-canonical-polar-hopf-triple}.
Put
\[
 P:=\frac{K_7-2K_9}{6}=cc^\dagger,
 \quad Q:=I_G-P,
 \quad W:=R_G^\dagger,
 \quad C:=QWP.
\]
Then $CC^\dagger=\frac89P_\mu$, $P_\tau=P$,
$P_e=I_G-P-\frac98CC^\dagger$.  For any positive finite charged
coefficients $e,m,t$ supplied on the stated role packet, put
$E=e^2$, $M=m^2$, $T=t^2$.  The *specified geometric target* left return
has the exact source-generator presentation
\begin{equation}\label{p5:eq:v3218-two-contact-hurwitz-return}
 \boxed{K_*=E I_G+(T-E)P+\frac98(M-E)CC^\dagger.}
\end{equation}
Its unique positive Hermitian chiral representative is
\begin{equation}\label{p5:eq:v3218-positive-candidate}
 \boxed{A_+=K_*^{1/2}=eI_G+(t-e)P+\frac98(m-e)CC^\dagger.}
\end{equation}
These are defined mathematical source/return \emph{candidates} in the
already generated finite family algebra.  Neither equality identifies
$A_+$ with the actual mixed chiral action jet or discharges
Definition~\ref{p5:def:v228-hierarchy-origin}.
\end{theorem}
\begin{proof}
By the rooted polar theorem, $\|QWc\|^2=8/9$ and
$P_\mu=(9/8)QWPW^\dagger Q$.  Therefore $P$ and
$(9/8)CC^\dagger$ are orthogonal rank-one projections, whose
orthogonal complement is $P_e$.  Express
$EP_e+MP_\mu+TP_\tau$ in the orthogonal resolution to obtain the
first identity.  The second follows by squaring the self-adjoint
linear combination with positive eigenvalues $e,m,t$; positive
square-root uniqueness gives the claim.
\end{proof}

\begin{theorem}[Mass-Vandermonde triangle and sharp commuting-baseline rank]
\label{p5:thm:v3218-minimal-return-rank}
Retain the fixed common $W=\operatorname{diag}(1,i,-i)$ eigenchart and
$K_*$ from \eqref{p5:eq:v3218-two-contact-hurwitz-return}.
Then
\begin{equation}\label{p5:eq:v3218-rooted-triangle}
 \boxed{\operatorname{Im}((K_*)_{12}(K_*)_{23}(K_*)_{31})
 =\frac{(E-M)(E-T)(M-T)}{12}.}
\end{equation}
Writing $H_W=(W-W^\dagger)/(2i)=\operatorname{diag}(0,1,-1)$,
\begin{equation}\label{p5:eq:v3218-rooted-commutator-cubic}
 \boxed{-i\operatorname{Tr}[H_W,K_*]^3
       =(E-M)(E-T)(M-T).}
\end{equation}
For the Hermitian commutant $\mathfrak A_W^{\rm sa}
=\{D=D^\dagger:[D,W]=0\}$,
\begin{equation}\label{p5:eq:v3218-rank-strata}
 \boxed{\min_{D\in\mathfrak A_W^{\rm sa}}
 \operatorname{rank}(K_*-D)
 =\#\{E,M,T\}_{\rm distinct}-1.}
\end{equation}
Thus three distinct finite squared mass coefficients necessitate a
rank-two noncommuting *return correction* relative to every commuting
Hermitian baseline, and the two-contact word of the preceding theorem
attains this minimum.
\end{theorem}
\begin{proof}
Set $x=E-M$ and $y=E+M-2T$.  In the printed $W$ eigenbasis,
\[
 (K_*)_{12}=\frac{-x-2y+3ix}{12},\quad
 (K_*)_{13}=\frac{-x-2y-3ix}{12},\quad
 (K_*)_{23}=\frac{2x-2y+3ix}{12}.
\]
Hence $(K_*)_{31}=(K_*)_{12}$ and their cyclic triple product
has imaginary part $x(y^2-x^2)/48$, equal to the displayed mass
Vandermonde divided by twelve.  For $H_W=\operatorname{diag}(0,1,-1)$,
a direct three-cycle expansion of $\operatorname{Tr}[H_W,K_*]^3$
gives twelve times $i$ times that imaginary part, proving the
commutator formula.

Because $W$ has distinct eigenvalues, every Hermitian $D$ commuting
with $W$ is diagonal in this chart.  If $K_*-D$ had rank one,
Hermiticity would make it $a vv^\dagger$ for some real $a$.
Its three offdiagonal entries would give the real triple product
$a^3|v_1v_2v_3|^2$, contradicting the nonzero displayed imaginary
part when $E,M,T$ are distinct.  The choice $D=E I_G$ yields the
rank-two correction $(M-E)P_\mu+(T-E)P_\tau$, attaining the bound.
When exactly two values coincide, the same target is a scalar
multiple of the identity plus a nonzero scalar multiple of a single
rooted $P_r$, hence has a rank-one commuting-baseline correction.
Each of the three rooted rank-one $P_r$ has a nonzero offdiagonal
entry in the $W$ eigenchart, excluding rank zero.  When all three
values coincide, the target is scalar.  These exhaust the cases.
\end{proof}

\begin{corollary}[Necessary same-chart noncommutation of an actual left-return source]
\label{p5:cor:v3218-source-commutator-threshold}
If an independently evaluated, canonically normalized square actual
source $N$ on this same common left/right $W$ chart satisfies
$NN^\dagger=K_*$, then, with $x=E-M$ and $y=E+M-2T$,
\[
 \boxed{\|[N,W]\|_{\rm HS}^2\ge
 \frac{23x^2-4xy+8y^2}{72\max(E,M,T)}.}
\]
This is a necessary inequality, not a sharp right-unitary orbit
minimization, and it supplies no physical occurrence result.
\end{corollary}
\begin{proof}
The direct finite calculation gives
$\|[K_*,W]\|_{\rm HS}^2=(23x^2-4xy+8y^2)/18$.
Use $[NN^\dagger,W]=N[N^\dagger,W]+[N,W]N^\dagger$,
$\|[N^\dagger,W]\|_{\rm HS}=\|[N,W]\|_{\rm HS}$
on this diagonal normal chart, and
$\|N\|_{\rm op}^2=\max(E,M,T)$.
The Hilbert--Schmidt product/triangle bounds give the claim.
\end{proof}

\begin{corollary}[Native-generator evaluation of the already established two-scalar return defect]
\label{p5:cor:v3218-native-action-defect}
For an \emph{independently evaluated actual} normalized same-action
return $B=M_{\rm act}M_{\rm act}^\dagger$, the Paper--III-F exact defect of
Theorem~\ref{p3f:thm:v267-two-scalar-return-defect} has the purely
native-generator presentation
\begin{align*}
 \mathscr F(B)={}&\operatorname{Tr}(B^2)-2E\operatorname{Tr}(B)
 -2(T-E)\operatorname{Tr}(PB)\\
 &-\frac94(M-E)\operatorname{Tr}(CC^\dagger B)
 +(E^2+M^2+T^2)\\
 ={}&\|B-K_*\|_{\rm HS}^2\ge0.
\end{align*}
Equality is equivalent to the full finite free left three-role return.
The displayed equation does not evaluate $B$ or provide a new
source-constituting variational functional.
\end{corollary}
\begin{proof}
Substitute \eqref{p5:eq:v3218-two-contact-hurwitz-return} into the
existing two-contraction squared-distance identity and use the
orthogonal role projector traces to obtain
$\operatorname{Tr}(K_*^2)=E^2+M^2+T^2$.
\end{proof}

\begin{remark}[Physical action, current and pole scope]
\label{p5:rem:v3218-two-contact-source-scope}
The rank result counts the least rank of a positive-return *correction*
relative to the negative-Hopf commutant; it does not count or require two
fundamental interaction vertices.  The triangle is a rephasing
invariant for a fixed rooted comparison, and the commutator cubic
involves the second comparator $H_W$.  Neither is by itself an
observed CP invariant.  No actual physical current has been proved to
consume that comparator.  The target $K_*$ is not inserted as the
physical action source.  The source-first acceptance criterion for an
independently evaluated actual $B=NN^\dagger$ is the *already existing*
Paper--III-F Theorem~\ref{p3f:thm:v267-two-scalar-return-defect},
which is stronger than a separate four-scalar reconstruction; the
hierarchy-arrow origin and interacting-pole gates retain their stated
independent statuses.
\end{remark}

% END v3.058 ROOTED TWO-CONTACT RETURN MATHEMATICS

\end{document}
